% !TEX program = xelatex
% ============================================================
%  高中化学系统讲义 —— 主文件
%  单文件版：由 main.tex 及 preamble.tex、chapters/*.tex 合并而成，内容完全相同
%  编译方法：xelatex 高中化学系统讲义.tex  两次（或 latexmk -xelatex）
% ============================================================
\documentclass[UTF8,a4paper,zihao=5,oneside,openany,fontset=none]{ctexbook}
% >>>>>>>> 开始：preamble.tex
% ============================================================
%  《高中化学系统讲义》导言区
%  编译：xelatex（两遍）  字体：Noto Serif/Sans CJK SC + TeX Gyre Termes
% ============================================================
\usepackage[a4paper,top=2.1cm,bottom=2.0cm,left=1.9cm,right=1.9cm,
  headheight=14pt,headsep=0.45cm,footskip=0.75cm]{geometry}

% ---------- 字体 ----------
\usepackage{amsmath,amssymb}
\usepackage{unicode-math}
\setmainfont{TeX Gyre Termes}
\setsansfont{TeX Gyre Heros}
\setmathfont{TeX Gyre Termes Math}
\setCJKmainfont{Noto Serif CJK SC}[BoldFont={Noto Sans CJK SC Bold},ItalicFont={AR PL KaitiM GB}]
\setCJKsansfont{Noto Sans CJK SC}[BoldFont={Noto Sans CJK SC Bold}]
\setCJKmonofont{Noto Sans Mono CJK SC}
\setCJKfamilyfont{zhsong}{Noto Serif CJK SC}[BoldFont={Noto Serif CJK SC Bold}]
\setCJKfamilyfont{zhhei}{Noto Sans CJK SC}[BoldFont={Noto Sans CJK SC Bold}]
\setCJKfamilyfont{zhkai}{AR PL KaitiM GB}[AutoFakeBold=2]
\setCJKfamilyfont{zhfs}{AR PL KaitiM GB}
\def\songti{\CJKfamily{zhsong}}
\def\heiti{\CJKfamily{zhhei}}
\def\kaishu{\CJKfamily{zhkai}}
\linespread{1.22}
\xeCJKsetcharclass{"2460}{"24FF}{1}   % ①②③ 等按中文字符处理
\xeCJKsetcharclass{"2160}{"217F}{1}   % Ⅰ Ⅱ Ⅲ 罗马数字
\xeCJKsetcharclass{"25A0}{"25FF}{1}   % △ ○ 等几何符号用中文字体
\xeCJKsetcharclass{"00B7}{"00B7}{0}   % 单位中的中点·按西文处理，避免 kJ·mol 断行
\xeCJKsetcharclass{"2013}{"2013}{0}   % en dash 用西文字体
\setlength{\parskip}{0.15em}

% ---------- 常用宏包 ----------
\usepackage[version=4]{mhchem}
\usepackage{chemfig}
\usepackage{tikz}
\usetikzlibrary{arrows.meta,calc,positioning,shapes.geometric,decorations.pathreplacing,patterns,3d,shadows}
\usepackage{pgfplots}
\pgfplotsset{compat=1.18}
\usepackage{booktabs,tabularx,multirow,array,makecell,colortbl}
\usepackage{enumitem}
\usepackage{multicol}
\usepackage[skins,breakable,most]{tcolorbox}
\usepackage{tasks}
\usepackage{fancyhdr}
\usepackage{xcolor}
\usepackage{caption}
\usepackage{wrapfig}
\usepackage{siunitx}
\sisetup{inter-unit-product=\ensuremath{\cdot},per-mode=symbol}
\usepackage[hidelinks,bookmarksnumbered]{hyperref}
\hypersetup{pdftitle={高中化学系统讲义},pdfsubject={以选择性必修1课件为主线 贯通高考化学知识体系（面向广东新高考）},pdfkeywords={高中化学,选择性必修1,广东高考,讲义}}

% ---------- 颜色 ----------
\definecolor{mainblue}{RGB}{0,84,147}
\definecolor{lightblue}{RGB}{232,242,251}
\definecolor{deepgreen}{RGB}{20,110,70}
\definecolor{lightgreen}{RGB}{233,246,238}
\definecolor{exorange}{RGB}{196,90,10}
\definecolor{lightorange}{RGB}{254,244,232}
\definecolor{warnred}{RGB}{180,30,40}
\definecolor{lightred}{RGB}{253,236,236}
\definecolor{gkpurple}{RGB}{100,45,140}
\definecolor{lightpurple}{RGB}{244,238,250}
\definecolor{graytext}{RGB}{90,90,90}

% ---------- 章节样式 ----------
\ctexset{
  chapter={
    format={\raggedright\heiti},
    nameformat={\color{mainblue}\Large},
    titleformat={\color{mainblue}\LARGE\bfseries},
    number={\arabic{chapter}},
    name={第,章},
    aftername={\quad},
    beforeskip={0pt},
    afterskip={1.0em},
    pagestyle=fancy
  },
  section={
    format={\Large\heiti\color{mainblue}\raggedright},
    number={\arabic{chapter}.\arabic{section}},
    beforeskip={1.0ex plus .3ex},
    afterskip={0.6ex}
  },
  subsection={
    format={\large\heiti\color{deepgreen}\raggedright},
    number={\arabic{chapter}.\arabic{section}.\arabic{subsection}},
    beforeskip={0.8ex plus .2ex},
    afterskip={0.4ex}
  },
  subsubsection={
    format={\normalsize\heiti\raggedright},
    number={},
    beforeskip={0.5ex},
    afterskip={0.2ex}
  },
  contentsname={目\quad 录}
}
\setcounter{secnumdepth}{2}
\setcounter{tocdepth}{1}

% ---------- 页眉页脚 ----------
\pagestyle{fancy}
\fancyhf{}
\fancyhead[L]{\small\heiti\color{mainblue}\leftmark}
\fancyhead[R]{\small\kaishu\color{graytext}高中化学系统讲义}
\fancyfoot[C]{\small\thepage}
\renewcommand{\headrulewidth}{0.5pt}
\renewcommand{\chaptermark}[1]{\markboth{第\,\thechapter\,章\quad #1}{}}
\fancypagestyle{plain}{\fancyhf{}\fancyfoot[C]{\small\thepage}\renewcommand{\headrulewidth}{0pt}}

% ---------- 列表 ----------
\setlist{nosep,leftmargin=2em}
\setlist[itemize]{label=\textcolor{mainblue}{\small$\blacktriangleright$}}
\setlist[enumerate,1]{label=(\arabic*)}
\settasks{label=\Alph*．,label-width=2.4em,item-indent=2.8em,before-skip=0.1em,after-skip=0.1em,column-sep=0.5em}
\setlength{\emergencystretch}{2.5em}

% ---------- 栏目环境 ----------
\tcbset{
  basestyle/.style={enhanced,breakable,boxrule=0.6pt,arc=1.5mm,left=2mm,right=2mm,top=1mm,bottom=1mm,
    before skip=0.55em,after skip=0.55em,fonttitle=\heiti\bfseries\small}
}
% 概念：核心概念与定义
\newtcolorbox{gainian}[1][]{basestyle,colback=lightblue,colframe=mainblue,
  title={\faicon 核心概念\ifx\relax#1\relax\else：#1\fi},attach boxed title to top left={xshift=3mm,yshift=-2mm},
  boxed title style={colback=mainblue,arc=1mm,boxrule=0pt},top=2.5mm}
% 本质：追本溯源（为什么）
\newtcolorbox{benzhi}[1][]{basestyle,colback=lightgreen,colframe=deepgreen,
  title={追本溯源\ifx\relax#1\relax\else：#1\fi},attach boxed title to top left={xshift=3mm,yshift=-2mm},
  boxed title style={colback=deepgreen,arc=1mm,boxrule=0pt},top=2.5mm}
% 方法：方法提炼
\newtcolorbox{fangfa}[1][]{basestyle,colback=lightorange,colframe=exorange,
  title={方法提炼\ifx\relax#1\relax\else：#1\fi},attach boxed title to top left={xshift=3mm,yshift=-2mm},
  boxed title style={colback=exorange,arc=1mm,boxrule=0pt},top=2.5mm}
% 易错：易错警示
\newtcolorbox{yicuo}[1][]{basestyle,colback=lightred,colframe=warnred,boxrule=0pt,leftrule=2.5pt,arc=0mm,
  title={易错警示\ifx\relax#1\relax\else：#1\fi},coltitle=warnred,colbacktitle=lightred,
  fonttitle=\heiti\bfseries\small,toptitle=0.5mm,bottomtitle=0mm}
% 勘误：原课件勘误
\newtcolorbox{kanwu}[1][]{basestyle,colback=white,colframe=warnred,boxrule=0.6pt,arc=0mm,
  borderline={0.3pt}{1.2pt}{warnred,dashed},
  title={课件勘误\ifx\relax#1\relax\else：#1\fi},coltitle=white,colbacktitle=warnred}
% 真题：真题链接（参数为来源）
\newtcolorbox{zhenti}[1]{basestyle,colback=lightpurple,colframe=gkpurple,
  title={真题链接\quad{\normalfont\small #1}},coltitle=white,colbacktitle=gkpurple}
% 拓展：视野拓展（高考不作要求或仅作了解）
\newtcolorbox{tuozhan}[1][]{basestyle,colback=white,colframe=graytext,boxrule=0.4pt,
  title={视野拓展\ifx\relax#1\relax\else：#1\fi},coltitle=graytext,colbacktitle=white,fonttitle=\heiti\small}
% 导学
\newtcolorbox{daoxue}{basestyle,colback=white,colframe=mainblue,boxrule=1pt,arc=2mm,
  title={本章导学},coltitle=white,colbacktitle=mainblue,fonttitle=\heiti\bfseries}

\def\faicon{}

% 例题：编号随章；上半题目、下半解析
\newcounter{liti}[chapter]
\renewcommand{\theliti}{\thechapter-\arabic{liti}}
\newtcolorbox{liti}[1][]{enhanced,breakable,colback=white,colframe=exorange,boxrule=0.5pt,arc=1mm,
  left=2mm,right=2mm,top=1mm,bottom=1mm,before skip=0.6em,after skip=0.6em,
  segmentation style={dashed,exorange},colbacklower=white,
  before upper={\stepcounter{liti}{\heiti\color{exorange}【例\,\theliti】}\ifx\relax#1\relax\else{\kaishu\small（#1）}\fi\ },
  before lower={\small}}
\newcommand{\silu}{{\heiti\color{exorange}【思路】}}
\newcommand{\jie}{{\heiti\color{exorange}【解析】}}
\newcommand{\daan}{{\heiti\color{exorange}【答案】}}
\newcommand{\xianjing}{{\heiti\color{warnred}【陷阱】}}
\newcommand{\dianbo}{{\heiti\color{deepgreen}【点拨】}}

% 习题与解析
\newlist{xiti}{enumerate}{1}
\setlist[xiti]{label={\textbf{\arabic*}.},leftmargin=1.6em,itemsep=0.35em,topsep=0.3em,resume=xiti}
\newcommand{\zu}[1]{\par\smallskip\noindent{\heiti\color{mainblue}\large #1}\par\smallskip}
\newcommand{\laiyuan}[1]{{\kaishu\small\color{graytext}（#1）}}
\newenvironment{jiexilist}{\small\begin{enumerate}[label={\textbf{\arabic*}.},leftmargin=1.6em,itemsep=0.3em,topsep=0.2em]}{\end{enumerate}}
\newcommand{\ans}[1]{{\heiti 答案：}#1\quad}

% 填空线、括号
\newcommand{\tk}[1][3em]{\underline{\hspace{#1}}}
\newcommand{\xz}{（\hspace{2em}）}

% 化学常用
\newcommand{\kjmol}{\,kJ·mol$^{-1}$}
\newcommand{\molL}{\,mol·L$^{-1}$}
\newcommand{\dH}{\Delta H}
\newcommand{\ksp}{K_{\mathrm{sp}}}
\newcommand{\ka}{K_{\mathrm{a}}}
\newcommand{\kb}{K_{\mathrm{b}}}
\newcommand{\kw}{K_{\mathrm{w}}}
\newcommand{\kh}{K_{\mathrm{h}}}
\newcommand{\kp}{K_{p}}
\newcommand{\NA}{N_{\mathrm{A}}}
\newcommand{\vz}{v_{\text{正}}}
\newcommand{\vn}{v_{\text{逆}}}
\newcommand{\kz}{k_{\text{正}}}
\newcommand{\kn}{k_{\text{逆}}}
\newcommand{\eqm}{\rightleftharpoons}

% 小节标题“章末”样式
\newcommand{\zhangmo}[1]{\section*{#1}\addcontentsline{toc}{section}{#1}}

% 表格行距
\renewcommand{\arraystretch}{1.25}
\setlength{\tabcolsep}{4pt}

% chemfig 设置
\setchemfig{atom sep=1.8em,bond offset=1pt,double bond sep=2.2pt,angle increment=30}

% 能量-历程图快捷命令：\energyprofile{起点}{峰}{终点}
\newcommand{\axisEP}[2]{%
  \draw[-Stealth] (0,0)--(#1,0) node[below left]{\small 反应进程};
  \draw[-Stealth] (0,0)--(0,#2) node[left]{\small 能量};}

% 热化学方程式中的焓变：\dHv{-285.8}
\newcommand{\dHv}[1]{\quad\ensuremath{\Delta H=#1\ \mathrm{kJ\cdot mol^{-1}}}}
% 题干中的选项（行内四选项，自动排版）
\newcommand{\xx}[4]{\begin{tasks}(4)\task #1 \task #2 \task #3 \task #4\end{tasks}}
\newcommand{\xxb}[4]{\begin{tasks}(2)\task #1 \task #2 \task #3 \task #4\end{tasks}}
\newcommand{\xxa}[4]{\begin{tasks}(1)\task #1 \task #2 \task #3 \task #4\end{tasks}}
% <<<<<<<< 结束：preamble.tex


\begin{document}
\frontmatter
% >>>>>>>> 开始：chapters/front.tex
% ============================================================
% 封面、前言、广东卷命题分析
% ============================================================
\begin{titlepage}
\begin{tikzpicture}[remember picture,overlay]
\fill[mainblue] (current page.north west) rectangle ([yshift=-9.5cm]current page.north east);
\fill[mainblue!80!black] ([yshift=-9.5cm]current page.north west) rectangle ([yshift=-10.1cm]current page.north east);
% 装饰：苯环与分子
\begin{scope}[shift={([xshift=-4.2cm,yshift=-4.6cm]current page.north east)},white,opacity=0.35,line width=1.2pt]
\foreach \a in {0,60,...,300} {\draw (\a:1.6)--(\a+60:1.6);}
\draw (0,0) circle (1.0);
\foreach \a in {30,150,270} {\draw (\a:1.6*0.866*1.15)--(\a:2.6);}
\end{scope}
\begin{scope}[shift={([xshift=3.2cm,yshift=-7.6cm]current page.north west)},white,opacity=0.25,line width=1pt]
\foreach \x/\y in {0/0,1.2/0.7,2.4/0,3.6/0.7,4.8/0}{\fill (\x,\y) circle (0.18);}
\draw (0,0)--(1.2,0.7)--(2.4,0)--(3.6,0.7)--(4.8,0);
\end{scope}
\node[anchor=west,text=white,font=\heiti\fontsize{40}{48}\selectfont] at ([xshift=2.2cm,yshift=-3.9cm]current page.north west) {高中化学系统讲义};
\node[anchor=west,text=white,font=\kaishu\Large] at ([xshift=2.3cm,yshift=-5.5cm]current page.north west) {以选择性必修1课件为主线\quad 贯通高考化学知识体系};
\node[anchor=west,text=white!85!mainblue,font=\heiti\large] at ([xshift=2.3cm,yshift=-6.6cm]current page.north west) {面向广东新高考\ ·\ 目标 90 分};
\end{tikzpicture}

\vspace*{10.2cm}
\begin{center}
\begin{minipage}{0.86\linewidth}
\begin{tcolorbox}[enhanced,colback=lightblue,colframe=mainblue,boxrule=0.8pt,arc=2mm,left=4mm,right=4mm,top=3mm,bottom=3mm]
\heiti\large 全书结构\normalfont\normalsize\par\medskip
\begin{tabular}{@{}ll@{}}
\textbf{预备篇} & 第0章\quad 高一必备基础速览 \\[2pt]
\textbf{原理篇（主线）} & 第1章 热效应\quad 第2章 速率与历程\quad 第3章 化学平衡 \\
& 第4章 水溶液中的离子平衡\quad 第5章 化学反应与电能 \\[2pt]
\textbf{结构篇} & 第6章 原子结构\quad 第7章 分子结构\quad 第8章 晶体结构 \\[2pt]
\textbf{有机篇} & 第9章 结构与烃\quad 第10章 烃的衍生物\quad 第11章 大分子与合成 \\[2pt]
\textbf{综合篇} & 第12章 化学实验综合\quad 第13章 元素化合物与化工流程 \\
& 第14章 高考综合大题训练（四道原创 14 分大题） \\
\end{tabular}
\end{tcolorbox}
\vspace{1.2cm}
{\kaishu\large 每章配有：本章导学\ ·\ 核心概念\ ·\ 追本溯源\ ·\ 例题精讲\ ·\ 真题链接\ ·\ 课件勘误\\[4pt]
\hspace*{5em}本章小结\ ·\ 方法提升\ ·\ 易错点归纳\\[4pt]
\hspace*{5em}分层练习（A 基础 / B 高考 / C 冲刺）\ ·\ 完整解析}
\end{minipage}
\end{center}
\vfill
\begin{center}
{\heiti\large 2026 年 10 月版}
\end{center}
\end{titlepage}

\chapter*{前\quad 言}
\markboth{前言}{}
\addcontentsline{toc}{chapter}{前言}

\textbf{一、这本讲义是什么。}本讲义以高二选择性必修1的七份课件（1.1 反应热、1.2 反应热的计算、2.1 化学反应速率、2.2 化学平衡、2.3 化学反应的方向、2.4 化学反应的调控，以及收录 2025 年广东卷原理题的演示文稿）为主线，逐份阅读、整合、去重后重写，并参考配套课后作业、必修与选择性必修各册思维导图、高考一轮复习知识清单 01～33，以及 2024—2026 年广东高考化学试卷，补齐课件尚未覆盖或讲解不充分的高考必备内容，形成从高一基础到高三综合的完整高考化学知识体系。

\textbf{二、怎样安排内容。}全书分五篇：预备篇只保留理解后续内容必需的高一基础；原理篇是课件主线，讲得最细，并补充了课件尚未讲到的第三、四章（离子平衡、电化学）；结构篇与有机篇对应选择性必修2、3；综合篇对应广东卷第17题（实验）和第18题（工艺流程），属于高三综合能力，最后一章用四道按广东卷第17～20题设计的原创大题做全书综合检测。知识深度以普通高中课程标准和高考要求为准：高中范围内能讲清楚的原理（如“为什么键能法是反应物减生成物”“为什么稀释时平衡正移而 $c(\ce{H+})$ 减小”）一律讲透；竞赛和大学层面的内容（如阿伦尼乌斯公式的推导、熵的统计意义）只在“视野拓展”中点到为止。

\textbf{三、栏目说明。}
\begin{itemize}
\item \textbf{核心概念}：必须准确理解和记忆的定义、规律与公式。
\item \textbf{追本溯源}：解释“为什么”，帮助理解而不是死记，是本书的特色。
\item \textbf{方法提炼}：可迁移的解题方法与答题模板。
\item \textbf{易错警示}：高频错误与概念陷阱。
\item \textbf{课件勘误}：对原课件中的笔误、不规范或有争议之处给出更正（汇总见附录A）。
\item \textbf{真题链接}：直接引用的高考真题，注明年份、地区、题号与来源，并给出完整解析。
\item \textbf{例题}：针对单个知识点的例题，分“思路—解析—陷阱/点拨”三部分。
\item \textbf{章末}：本章小结（知识总结、方法提升、易错点归纳）$+$ 分层练习（A组基础巩固、B组高考典型、C组冲刺高分）$+$ 练习解析。
\end{itemize}

\textbf{四、关于真题与题目来源。}本书直接引用的真题来自资料包中的《2024年广东省高考真题化学试卷》《2025年广东高考真题化学试题》《2026年高考广东卷化学高考真题（网络收集版）》，编写时逐题核对了题干与参考答案；其中 2026 年试卷为网络收集版，个别表述可能与官方版本略有出入。由于编写环境无法访问教育考试院等官方网页，本书未引用其他省份真题的原文；凡未标注真题来源的题目均为原创或依据课件、作业改编（标注“原创”），题目数据经过逐一验算。真题中涉及图片的，书中用文字或重新绘制的示意图准确转述其关键信息。

\textbf{五、使用建议。}①先读正文，遇到例题先自己做再看解析；②章末练习一定要“先做后对答案”，A 组应全对，B 组是高考常见难度，C 组用于冲击高分；③把“易错点归纳”和自己的错题对照，反复回看；④目标 90 分的同学，选择题力争只错一道，四道大题每题失分控制在 2～3 分以内，这要求概念零盲区、计算准确、表述规范——本书的“方法提炼”和“答题模板”正是为此而写。

\chapter*{广东卷命题分析}
\markboth{广东卷命题分析}{}
\addcontentsline{toc}{chapter}{广东卷命题分析}

\section*{一、试卷结构}
广东新高考化学满分 100 分，考试时间 75 分钟。第Ⅰ卷为 16 道单项选择题，共 44 分（第 1～10 题每题 2 分，第 11～16 题每题 4 分）；第Ⅱ卷为 4 道非选择题，每题 14 分，共 56 分，题型相对固定：第17题化学实验探究，第18题工艺流程，第19题化学反应原理综合，第20题有机化学综合。

\section*{二、2024—2026 年广东卷逐题考点}
\begin{center}
\small
\renewcommand{\arraystretch}{1.12}
\begin{tabularx}{\linewidth}{c>{\raggedright\arraybackslash}X>{\raggedright\arraybackslash}X>{\raggedright\arraybackslash}X}
\toprule
题号 & 2024 年 & 2025 年 & 2026 年（网络收集版） \\
\midrule
1 & 文物材料（有机高分子） & 传统工艺材料（合金） & 桥梁构件材质（有机物） \\
2 & 科技成就：变化、能量、电子式 & 多烯环配合物：键、大 $\pi$ 键 & 科技成就：放热、核素、周期 \\
3 & 计时器材料、石英、核素 & 摩尔质量、化合价、能量转化 & 竹：变化、\ce{SO2} 漂白、吸附 \\
4 & 饮食：糖类、油脂、极性分子 & 劳动项目：沉淀转化、盐析 & 成像：能量、\ce{SiO2}、共价晶体 \\
5 & 金属防腐（牺牲阳极） & 材料：单体、\ce{Fe3O4}、晶体 & 陈述关联：胶体、晶体、焰色 \\
6 & 苯甲酸重结晶 & 铁钉镀铜实验操作 & 海上风电防腐 \\
7 & 劳动项目与化学知识 & 物质转化链（S、Na、C、Al） & 阿伏加德罗常数 \\
8 & \ce{Cl2} 的制备与性质装置 & 陈述因果（浓硫酸、\ce{Fe^2+} 检验） & 结构解释性质（氢键、推电子） \\
9 & 陌生有机物结构与性质 & 第一电离能图与周期律 & \ce{Cl2} 制备与 Fe 反应实验 \\
10 & 阿伏加德罗常数 & 阿伏加德罗常数 & 岭南美食中的化学 \\
11 & 吸热反应、熵增、\ce{NH3} 实验 & \ce{H2O2} 分解与乙醇氧化实验 & 元素推断（3p 半充满） \\
12 & 元素推断（结构式） & CuCl 配位平衡、电荷守恒 & 卤代烃消去与水解实验 \\
13 & 陈述因果（酸性、冠醚） & 结构能否推测性质 & 价类二维图（C、N、Cl） \\
14 & 价类二维图（Mg、Al、Fe） & 水系可充电电池 & 方程式正误 \\
15 & 反应历程能量图（中间体、决速步） & Cu 与稀硝酸实验 & 平行反应能量图与 $c$–$t$ 图 \\
16 & 离子交换膜电解与计算 & 燃料电池催化机理 & 电解还原 \ce{CO2} 与海水脱盐 \\
17 & 硫的转化；滴定；速率与平衡建立 & 重结晶；pH 滴定求 $K_\text{a}$；方案设计 & 氨水；镀铜；铜氨配合物 pH 滴定 \\
18 & 炭渣提镓：配离子、$K_\text{sp}$、晶胞 & 多金属精矿：$K_\text{sp}$、萃取、晶胞、电解 & 金精矿提金：电解、配平、超分子、晶胞 \\
19 & 盖斯定律；催化；非水溶剂 $K_\text{a}$ & $\Delta G$；盖斯；平衡组成；$K$；选择性 & 轨道式；盖斯；速率数据；熵；$K$ \\
20 & 苯乙烯$\to$手性醇；同分异构；合成 & 光酶三组分反应；核磁；合成 & 羟基迁移；原子利用率；合成路线 \\
\bottomrule
\end{tabularx}
\end{center}

\section*{三、高频考点与命题趋势}
\begin{enumerate}
\item \textbf{每年必考的“固定题位”}：$N_\text{A}$、元素推断与周期律、价类二维图或物质转化、反应历程能量图（第15题）、两道电化学选择题、四道大题的题型。复习时应保证这些题位“零失误”。
\item \textbf{原理大题（第19题）三连考}：盖斯定律求 $\Delta H$（3/3）、平衡常数或吸收量计算（3/3）、结构小问（轨道表示式、价层电子排布 2/3）、选择性或速率数据分析（2/3）。计算要求“写出计算过程”，强调规范。
\item \textbf{流程题（第18题）高度综合}：$K_\text{sp}$ 与调 pH（3/3）、晶胞（3/3）、陌生氧化还原方程式（3/3）、配合物或超分子（3/3）——体现“必修元素化学$+$选必1原理$+$选必2结构”的融合。
\item \textbf{实验题（第17题）重定量与探究}：每年都有滴定（酸碱或 pH 滴定）计算，并以“提出假设—验证—修正”的完整探究链考查科学思维。
\item \textbf{有机题（第20题）重新反应的迁移}：以前沿成果（光酶催化、羟基迁移）为情境，考查官能团、同分异构、杂化与手性、合成路线设计。
\item \textbf{情境化与结构化}：选择题前几题以传统文化、科技成就、劳动项目、岭南生活为情境；“结构决定性质”的因果判断题逐年增多。
\end{enumerate}

\section*{四、冲刺 90 分的策略}
\begin{center}
\small
\begin{tabularx}{\linewidth}{p{3.2cm}cX}
\toprule
板块 & 目标得分 & 关键 \\
\midrule
选择题 1～10（20 分） & 20 & 概念准确：STSE、$N_\text{A}$、实验操作、结构与性质 \\
选择题 11～16（24 分） & $\ge20$ & 图像与电化学：本书第2、3、4、5章的读图方法 \\
第17题 实验（14 分） & $\ge12$ & 操作规范、滴定计算、假设与方案表述（第12章） \\
第18题 流程（14 分） & $\ge12$ & 预处理与条件控制、配平、$K_\text{sp}$、晶胞（第4、8、13章） \\
第19题 原理（14 分） & $\ge12$ & 盖斯定律、三段式、$K_p$、多重平衡与选择性（第1、3章） \\
第20题 有机（14 分） & $\ge12$ & 同分异构计数、合成路线（第9～11章） \\
\bottomrule
\end{tabularx}
\end{center}
% <<<<<<<< 结束：chapters/front.tex

\tableofcontents

\mainmatter
\setcounter{chapter}{-1}
% >>>>>>>> 开始：chapters/ch00.tex
% ============================================================
% 第0章 高一必备基础速览（必修一、必修二中理解后续内容所必需的部分）
% ============================================================
\chapter{高一必备基础速览}

\begin{daoxue}
\textbf{为什么先复习这些：}后面的反应热计算离不开\textbf{物质的量}；平衡、水解、电化学离不开\textbf{离子方程式}与\textbf{氧化还原}；物质结构与元素推断离不开\textbf{原子结构与周期律}。本章只保留这些“地基”，元素化合物的系统梳理放在第13章，实验基本操作放在第12章。

\textbf{高考考情：}阿伏加德罗常数（$N_\text{A}$）几乎每年一题（2024年第10题、2025年第10题、2026年第7题）；离子方程式正误、氧化还原方程式书写在大题中反复出现。
\end{daoxue}

\section{物质的量与阿伏加德罗常数}

\begin{gainian}[以物质的量为中心的换算]
\[n=\frac{m}{M}=\frac{N}{N_\text{A}}=\frac{V(\text{气体})}{V_\text{m}}=c\cdot V(\text{溶液})\]
$N_\text{A}\approx6.02\times10^{23}$ mol$^{-1}$；\textbf{标准状况}（0 ℃、101 kPa）下气体摩尔体积 $V_\text{m}=22.4$ L·mol$^{-1}$。阿伏加德罗定律：同温同压下，相同体积的任何气体含有相同数目的分子，即 $\dfrac{V_1}{V_2}=\dfrac{n_1}{n_2}$；同温同容下 $\dfrac{p_1}{p_2}=\dfrac{n_1}{n_2}$（这正是第3章用压强判断气体物质的量变化的依据）。
\end{gainian}

\begin{center}
\small
\begin{tabularx}{\linewidth}{p{2.4cm}X}
\toprule
陷阱类型 & 典型判断（$N_\text{A}$ 为阿伏加德罗常数的值） \\
\midrule
气体摩尔体积的适用 & 22.4 L·mol$^{-1}$ 只用于\textbf{标准状况下的气体}。标准状况下 \ce{H2O}、\ce{SO3}、\ce{CCl4}、HF、苯、乙醇不是气体；常温常压下的气体不能用 22.4 \\
微粒数与结构 & 1 mol \ce{^{18}O2} 比 1 mol \ce{^{16}O2} 多 $4N_\text{A}$ 个中子（2025年第10题A）；26 g \ce{C2H2} 含 $\sigma$ 键 $3N_\text{A}$（2024年第10题A）；标准状况下 22.4 L \ce{CH3Cl} 含极性键 $4N_\text{A}$（2026年第7题C） \\
电解质溶液 & 1 L 1 mol·L$^{-1}$ \ce{NH4NO3} 溶液中 \ce{NH4+} 少于 $N_\text{A}$（水解）；未给体积无法计算微粒数 \\
可逆反应、隐含反应 & 23 g \ce{NO2} 与 \ce{N2O4} 的混合气体含 N 原子 $0.5N_\text{A}$（用最简式 \ce{NO2} 计算，与比例无关；2026年第7题A 说成 $N_\text{A}$，错）；\ce{2NO2 <=> N2O4} 使分子数难以确定 \\
电子转移 & 1 mol Na 与 \ce{O2} 完全反应，无论生成 \ce{Na2O} 还是 \ce{Na2O2}，失电子都为 $N_\text{A}$（2026年第7题B，正确）；1 mol Fe 与水蒸气反应生成 $\frac43$ mol \ce{H2}（2025年第10题B）；1 mol \ce{Na2O2} 与水反应转移 $N_\text{A}$ 个电子 \\
\bottomrule
\end{tabularx}
\end{center}

\begin{fangfa}[化学计算的三种守恒]
\textbf{原子（质量）守恒}：反应前后某元素的原子总数不变，如第3章多重平衡中的“C 原子守恒”。\textbf{电荷守恒}：溶液呈电中性，阳离子所带正电荷总数等于阴离子所带负电荷总数（第4章）。\textbf{得失电子守恒}：氧化剂得电子总数等于还原剂失电子总数（氧化还原计算、电化学计算、滴定计算的核心）。
\end{fangfa}

\section{离子反应}

\subsection{电解质与离子方程式}
在水溶液或熔融状态下能导电的化合物是电解质（酸、碱、盐、金属氧化物、水）；\ce{CO2}、\ce{SO2}、\ce{NH3}、蔗糖、乙醇是非电解质（\ce{CO2} 的水溶液导电是因为生成了 \ce{H2CO3}）；单质和混合物既不是电解质也不是非电解质。

\textbf{书写离子方程式的“拆”与“不拆”}：易溶的强电解质（强酸、强碱、可溶性盐）拆成离子；弱电解质（弱酸、弱碱、水）、沉淀、气体、氧化物、单质写化学式。微溶物作反应物时，澄清溶液拆（澄清石灰水），浊液不拆（石灰乳）；浓硫酸不拆，浓盐酸、浓硝酸拆；\ce{HCO3-}、\ce{HSO3-} 等弱酸的酸式酸根不拆，而 \ce{HSO4-} 在水溶液中拆成 \ce{H+} 和 \ce{SO4^2-}。

\begin{fangfa}[离子方程式正误判断“五看”]
\textbf{一看}是否符合客观事实（Fe 与盐酸生成 \ce{Fe^2+}；Cu 不与稀硫酸反应；\ce{SO2} 通入 \ce{Ca(ClO)2} 溶液发生的是氧化还原反应）；\textbf{二看}拆分是否正确；\textbf{三看}原子、电荷、得失电子是否守恒（\ce{Fe^3+ + Cu = Fe^2+ + Cu^2+} 电荷不守恒）；\textbf{四看}量的关系（“少定多变”：以少量物质为 1 mol 来确定另一物质的用量，如少量 \ce{NaHCO3} 与足量 \ce{Ca(OH)2}：\ce{HCO3- + Ca^2+ + OH- = CaCO3 v + H2O}；足量 \ce{NaHCO3} 与少量 \ce{Ca(OH)2}：\ce{2HCO3- + Ca^2+ + 2OH- = CaCO3 v + CO3^2- + 2H2O}）；\textbf{五看}反应先后（\ce{FeBr2} 溶液中通入少量 \ce{Cl2}，还原性 \ce{Fe^2+}$>$\ce{Br-}，先氧化 \ce{Fe^2+}；\ce{FeI2} 中先氧化 \ce{I-}）。
\end{fangfa}

\begin{liti}[与量有关的离子方程式]
写出下列反应的离子方程式：(1) 向 NaOH 溶液中通入少量 \ce{CO2}；通入过量 \ce{CO2}。(2) 向 \ce{NH4HSO4} 溶液中滴加少量 NaOH 溶液。(3) 向明矾 [\ce{KAl(SO4)2*12H2O}] 溶液中滴加 \ce{Ba(OH)2} 溶液至 \ce{SO4^2-} 恰好沉淀完全。
\tcblower
\jie (1) 少量：\ce{CO2 + 2OH- = CO3^2- + H2O}；过量：\ce{CO2 + OH- = HCO3-}。(2) \ce{H+} 比 \ce{NH4+} 更易结合 \ce{OH-}，少量 NaOH 只中和 \ce{H+}：\ce{H+ + OH- = H2O}。(3) 设明矾 1 mol，含 2 mol \ce{SO4^2-}，需 2 mol \ce{Ba(OH)2}，带入 4 mol \ce{OH-}，恰好使 1 mol \ce{Al^3+} 转化为 \ce{[Al(OH)4]-}：\ce{Al^3+ + 2SO4^2- + 2Ba^2+ + 4OH- = 2BaSO4 v + [Al(OH)4]-}。\\
\dianbo “少定多变”：把量少的物质定为 1 mol，按它的需要确定另一物质的用量；有多种离子可与同一试剂反应时，按“谁先反应”逐步分析。
\end{liti}

\subsection{离子共存}
离子之间若发生以下反应则不能大量共存：①生成沉淀、气体或弱电解质（\ce{H+} 与 \ce{CO3^2-}、\ce{ClO-}，\ce{OH-} 与 \ce{NH4+}、\ce{HCO3-}）；②氧化还原反应（\ce{Fe^3+} 与 \ce{I-}、\ce{S^2-}；\ce{MnO4-}、\ce{ClO-}、\ce{NO3-}（\ce{H+}）与 \ce{Fe^2+}、\ce{I-}、\ce{SO3^2-}）；③相互促进的完全水解（\ce{Al^3+} 与 \ce{HCO3-}、\ce{CO3^2-}、\ce{[Al(OH)4]-}；\ce{Fe^3+} 与 \ce{HCO3-}）；④形成配合物（\ce{Fe^3+} 与 \ce{SCN-}）。注意隐含条件：“无色”排除 \ce{Cu^2+}（蓝）、\ce{Fe^2+}（浅绿）、\ce{Fe^3+}（黄）、\ce{MnO4-}（紫）；“由水电离出的 $c(\ce{H+})=10^{-13}$ mol·L$^{-1}$”的溶液可能呈强酸性也可能呈强碱性。

\section{氧化还原反应}

\begin{gainian}[基本概念]
本质是电子转移，特征是化合价变化。“\textbf{升失氧，还原剂；降得还，氧化剂}”——还原剂化合价升高、失电子、被氧化、发生氧化反应、得到氧化产物；氧化剂反之。
\end{gainian}

\textbf{三条规律}：①\textbf{强弱律}：氧化性：氧化剂$>$氧化产物；还原性：还原剂$>$还原产物。如 \ce{2Fe^3+ + 2I- = 2Fe^2+ + I2}，氧化性 \ce{Fe^3+}$>$\ce{I2}。②\textbf{先后律}：同一氧化剂遇多种还原剂，先与还原性最强的反应。③\textbf{价态律}：最高价只有氧化性，最低价只有还原性，中间价两者兼有；同种元素不同价态间反应“只靠拢、不交叉”（归中），如 \ce{KClO3 + 6HCl}（浓）\ce{= KCl + 3Cl2 ^ + 3H2O} 中，\ce{KClO3} 中的 Cl 降为 0 价而非 $-1$ 价，转移 5 个电子；同一物质中同一元素既升又降叫歧化，如 \ce{Cl2 + 2NaOH = NaCl + NaClO + H2O}。

\begin{liti}[得失电子守恒计算]
24 mL 0.05 mol·L$^{-1}$ 的 \ce{Na2SO3} 溶液恰好与 20 mL 0.02 mol·L$^{-1}$ 的 \ce{K2Cr2O7} 溶液完全反应，\ce{Na2SO3} 被氧化为 \ce{Na2SO4}，则 Cr 元素在还原产物中的化合价为多少？
\tcblower
\jie 失电子：$n(\ce{SO3^2-})\times2=0.024\times0.05\times2=2.4\times10^{-3}$ mol；得电子：$n(\ce{Cr})\times(6-x)=0.020\times0.02\times2\times(6-x)$。由 $2.4\times10^{-3}=8\times10^{-4}\times(6-x)$，得 $x=+3$。\\
\dianbo 守恒计算不需要写出完整方程式，只抓“谁得、谁失、各得失几个”。
\end{liti}

\begin{liti}[缺项配平]
高铁酸钾（\ce{K2FeO4}）是一种高效净水剂，可在碱性条件下制备：\\
\tk\ce{ClO-} $+$ \tk\ce{Fe(OH)3} $+$ \tk\tk $=$ \tk\ce{Cl-} $+$ \tk\ce{FeO4^2-} $+$ \tk\ce{H2O}
\tcblower
\jie ①找变价：Fe $+3\to+6$，每个失 3 个电子；Cl $+1\to-1$，每个得 2 个电子。最小公倍数 6：2 \ce{Fe(OH)3}、3 \ce{ClO-}。②配电荷：碱性介质缺项补 \ce{OH-}，左边 $-3+x(-1)$，右边 $3\times(-1)+2\times(-2)=-7$，得 $x=4$。③配 H、O：右边补 5 \ce{H2O}。\\
\ce{3ClO- + 2Fe(OH)3 + 4OH- = 3Cl- + 2FeO4^2- + 5H2O}。\\
\dianbo \ce{K2FeO4} 净水原理：$+6$ 价铁有强氧化性可消毒杀菌，还原产物 \ce{Fe^3+} 水解生成 \ce{Fe(OH)3} 胶体吸附杂质。
\end{liti}

\section{原子结构、化学键与元素周期律}
\textbf{原子结构}：${}^A_Z\mathrm{X}$ 中质量数 $A=$ 质子数 $Z+$ 中子数 $N$；质子数相同、中子数不同的核素互为同位素（化学性质几乎相同）。例：${}^3$He 与 ${}^4$He 中子数分别为 1、2（2026年第2题C）；${}^{135}_{55}$Cs 的质子数为 55（2024年第3题D）。

\textbf{化学键}：离子键（阴、阳离子间的静电作用，活泼金属与活泼非金属之间、铵盐中）；共价键（原子间共用电子对，同种原子间为非极性键，不同原子间为极性键）。含离子键的化合物一定是离子化合物，离子化合物中可以含共价键（NaOH、\ce{Na2O2}、\ce{NH4Cl}）。\textbf{电子式}书写要点：离子化合物中阴离子要加方括号并标电荷，如 NaCl：\ce{Na+}\,$\bigl[\,$\charge{0=\:,90=\:,180=\:,270=\:}{Cl}$\,\bigr]^-$；共价分子不加括号，每个原子达到稳定结构，如 \ce{H2O}：H\,\charge{0=\:,90=\:,180=\:,270=\:}{O}\,H（O 周围共 8 个电子，2024年第2题D 考查的正是漏写孤电子对的错误写法 H:O:H），\ce{HClO} 的结构式为 H—O—Cl（O 在中间）。

\textbf{元素周期律}（深入内容见第6章）：同周期从左到右金属性减弱、非金属性增强；同主族从上到下金属性增强、非金属性减弱。判断依据：\textbf{金属性}——单质与水或酸置换出氢的难易、最高价氧化物对应水化物的碱性强弱；\textbf{非金属性}——单质与 \ce{H2} 化合的难易、简单氢化物的稳定性、最高价氧化物对应水化物的酸性强弱（注意：不能用 HCl 与 \ce{H2S} 的酸性、HClO 与 \ce{H2SO4} 的酸性比较非金属性）。

% ---------------- 章末 ----------------
\zhangmo{本章小结}
\begin{enumerate}
\item \textbf{物质的量}：$n=\frac mM=\frac N{N_\text A}=\frac V{V_\text m}=cV$；$N_\text{A}$ 题从“状态、结构、溶液、可逆、电子转移”五个角度排查。
\item \textbf{离子反应}：拆分、守恒、量的关系、先后顺序；共存看沉淀气体弱电解质、氧化还原、双水解、配合四类反应及隐含条件。
\item \textbf{氧化还原}：升失氧还原剂；强弱、先后、价态三规律；计算抓得失电子守恒。
\item \textbf{结构与周期律}：$A=Z+N$；离子键、共价键与电子式；金属性、非金属性的判断依据。
\end{enumerate}
\textbf{易错点：}标准状况下 \ce{SO3}、HF 不是气体；\ce{Fe^3+} 与 Cu 反应的离子方程式电荷不守恒；用氢化物水溶液酸性比较非金属性；\ce{Na2O2} 电子式中 \ce{O2^2-} 写成两个 \ce{O^2-}。

\zhangmo{分层练习}
\zu{A组\quad 基础巩固}
\begin{xiti}[start=1]
\item \laiyuan{2025年广东卷第10题，来源：仓库所附真题文档}\ 设 $N_\text{A}$ 为阿伏加德罗常数的值。下列说法正确的是\xz
\xxa{1 mol \ce{^{18}O2} 的中子数，比 1 mol \ce{^{16}O2} 的多 $2N_\text{A}$}{1 mol Fe 与水蒸气完全反应，生成 \ce{H2} 的数目为 $2N_\text{A}$}{在 1 L 0.1 mol·L$^{-1}$ 的 \ce{NH4Cl} 溶液中，\ce{NH4+} 的数目为 $0.1N_\text{A}$}{标准状况下的 22.4 L \ce{Cl2} 与足量 \ce{H2} 反应，形成的共价键数目为 $2N_\text{A}$}

\item 下列离子方程式书写正确的是\xz
\xxa{铁与稀盐酸反应：\ce{2Fe + 6H+ = 2Fe^3+ + 3H2 ^}}{\ce{FeCl3} 溶液腐蚀铜板：\ce{Fe^3+ + Cu = Fe^2+ + Cu^2+}}{碳酸钙与醋酸反应：\ce{CaCO3 + 2CH3COOH = Ca^2+ + 2CH3COO- + CO2 ^ + H2O}}{向 \ce{FeBr2} 溶液中通入少量 \ce{Cl2}：\ce{2Br- + Cl2 = Br2 + 2Cl-}}

\item 常温下，下列各组离子在指定溶液中一定能大量共存的是\xz
\xxa{无色溶液中：\ce{K+}、\ce{Cu^2+}、\ce{SO4^2-}、\ce{Cl-}}{pH$=1$ 的溶液中：\ce{Fe^2+}、\ce{NO3-}、\ce{SO4^2-}、\ce{Na+}}{含大量 \ce{Fe^3+} 的溶液中：\ce{NH4+}、\ce{Na+}、\ce{Cl-}、\ce{SO4^2-}}{含大量 \ce{Al^3+} 的溶液中：\ce{Na+}、\ce{HCO3-}、\ce{Cl-}、\ce{SO4^2-}}

\item 下列物质中，既含离子键又含非极性共价键的是\xz
\xx{NaOH}{\ce{Na2O2}}{\ce{NH4Cl}}{\ce{H2O2}}
\end{xiti}

\zu{B组\quad 高考典型}
\begin{xiti}
\item \laiyuan{原创}\ 已知：①\ce{2Fe^3+ + 2I- = 2Fe^2+ + I2}；②\ce{2Fe^2+ + Br2 = 2Fe^3+ + 2Br-}。(1) 比较 \ce{Br2}、\ce{Fe^3+}、\ce{I2} 的氧化性强弱。(2) 向含 1 mol \ce{FeI2} 的溶液中通入 1.25 mol \ce{Cl2}，写出总反应的离子方程式。

\item \laiyuan{原创}\ 配平：\tk\ce{MnO4-} $+$ \tk\ce{H2C2O4} $+$ \tk\ce{H+} $=$ \tk\ce{Mn^2+} $+$ \tk\ce{CO2 ^} $+$ \tk\ce{H2O}；若反应中转移 0.5 mol 电子，生成标准状况下 \ce{CO2} 多少升？
\end{xiti}

\zu{C组\quad 冲刺高分}
\begin{xiti}

\item \laiyuan{原创}\ 设 $N_\text{A}$ 为阿伏加德罗常数的值。下列说法正确的是\xz
\xxa{7.8 g \ce{Na2O2} 中含有的阴离子数为 $0.1N_\text{A}$}{1 mol \ce{Cl2} 与足量 NaOH 溶液反应，转移电子数为 $2N_\text{A}$}{标准状况下，11.2 L \ce{SO3} 中含有的分子数为 $0.5N_\text{A}$}{1 L pH$=1$ 的 \ce{H2SO4} 溶液中含有 \ce{H+} 的数目为 $0.2N_\text{A}$}

\item \laiyuan{原创}\ 向含 0.1 mol \ce{FeBr2} 的溶液中通入 0.1 mol \ce{Cl2}。(1) 写出反应的离子方程式。(2) 被氧化的 \ce{Br-} 占 \ce{Br-} 总量的百分之几？(3) 若改为通入足量 \ce{Cl2}，写出离子方程式。
\end{xiti}

\zhangmo{练习解析}
\begin{jiexilist}
\item \ans{D}A：每个 \ce{^{18}O} 比 \ce{^{16}O} 多 2 个中子，每个 \ce{O2} 多 4 个，共多 $4N_\text{A}$；B：\ce{3Fe + 4H2O(g) ->[\text{高温}] Fe3O4 + 4H2}，1 mol Fe 生成 $\frac43$ mol \ce{H2}；C：\ce{NH4+} 水解，少于 $0.1N_\text{A}$；D：1 mol \ce{Cl2} 生成 2 mol HCl，形成 2 mol H—Cl 键，正确。
\item \ans{C}A：生成 \ce{Fe^2+}；B：电荷不守恒，应为 \ce{2Fe^3+ + Cu = 2Fe^2+ + Cu^2+}；C：碳酸钙、醋酸都不拆，正确；D：\ce{Fe^2+} 还原性强于 \ce{Br-}，少量 \ce{Cl2} 先氧化 \ce{Fe^2+}：\ce{2Fe^2+ + Cl2 = 2Fe^3+ + 2Cl-}。
\item \ans{C}A：\ce{Cu^2+} 呈蓝色；B：酸性条件下 \ce{NO3-} 氧化 \ce{Fe^2+}；C 可以共存；D：\ce{Al^3+} 与 \ce{HCO3-} 相互促进水解完全，生成 \ce{Al(OH)3} 和 \ce{CO2}。
\item \ans{B}NaOH 含离子键和极性键；\ce{Na2O2} 含离子键和 O—O 非极性键；\ce{NH4Cl} 含离子键、极性键和配位键；\ce{H2O2} 为共价化合物。
\item (1) \ce{Br2}$>$\ce{Fe^3+}$>$\ce{I2}（由①得 \ce{Fe^3+}$>$\ce{I2}，由②得 \ce{Br2}$>$\ce{Fe^3+}）。(2) 还原性 \ce{I-}$>$\ce{Fe^2+}，\ce{Cl2} 先氧化 2 mol \ce{I-}（消耗 1 mol \ce{Cl2}），剩余 0.25 mol \ce{Cl2} 氧化 0.5 mol \ce{Fe^2+}。按 $n(\ce{Fe^2+}):n(\ce{I-}):n(\ce{Cl2})=0.5:2:1.25=2:8:5$ 写总式：\ce{2Fe^2+ + 8I- + 5Cl2 = 2Fe^3+ + 4I2 + 10Cl-}。
\item 2、5、6、2、10、8；11.2 L。Mn：$+7\to+2$ 得 5 个电子；C：$+3\to+4$，每个 \ce{H2C2O4} 失 2 个电子。最小公倍数 10：2 \ce{MnO4-}、5 \ce{H2C2O4}，再由电荷守恒配 \ce{H+}（左边 $-2+x=+4$，$x=6$），最后配水。生成 10 mol \ce{CO2} 转移 10 mol 电子，转移 0.5 mol 电子生成 0.5 mol \ce{CO2}，标准状况下 11.2 L。
\item \ans{A}A：0.1 mol \ce{Na2O2} 含 0.1 mol \ce{O2^2-}，正确；B：\ce{Cl2 + 2NaOH = NaCl + NaClO + H2O}，1 mol \ce{Cl2} 转移 1 mol 电子；C：标准状况下 \ce{SO3} 不是气体；D：pH$=1$，$c(\ce{H+})=0.1$ mol·L$^{-1}$，1 L 中 \ce{H+} 为 $0.1N_\text{A}$。
\item (1) \ce{2Fe^2+ + 2Br- + 2Cl2 = 2Fe^3+ + Br2 + 4Cl-}。还原性 \ce{Fe^2+}$>$\ce{Br-}，0.1 mol \ce{Fe^2+} 先消耗 0.05 mol \ce{Cl2}，剩余 0.05 mol \ce{Cl2} 氧化 0.1 mol \ce{Br-}；参加反应的 \ce{Fe^2+}、\ce{Br-}、\ce{Cl2} 之比为 1:1:1。(2) 0.1 mol/0.2 mol$=50\%$。(3) \ce{2Fe^2+ + 4Br- + 3Cl2 = 2Fe^3+ + 2Br2 + 6Cl-}。
\end{jiexilist}
% <<<<<<<< 结束：chapters/ch00.tex

% >>>>>>>> 开始：chapters/ch01.tex
% ============================================================
% 第1章 化学反应的热效应（对应课件：选必一1.1 反应热、1.2 反应热的计算）
% ============================================================
\chapter{化学反应的热效应}

\begin{daoxue}
\textbf{知识主线：}反应热是什么（体系与环境、焓变）$\to$ 为什么有反应热（键能、物质能量）$\to$ 怎样测（中和反应反应热）$\to$ 怎样表示（热化学方程式、燃烧热）$\to$ 怎样算（盖斯定律等六种方法）。

\textbf{高考考情：}广东卷第19题（原理综合）几乎每年第一或第二小问考“盖斯定律求 $\Delta H$”（2024第19题(1)②、2025第19题(3)①、2026第19题(2)），约2分，必须稳拿；选择题常以“能量—反应进程图”考查放热吸热、活化能与催化剂（2024第15题、2026第15题）。

\textbf{学习目标：}①能从宏观（物质能量）和微观（化学键）两个角度解释反应热；②会规范书写热化学方程式并判断正误；③能熟练用盖斯定律、键能、燃烧热等计算 $\Delta H$，并会比较 $\Delta H$ 大小。
\end{daoxue}

\section{反应热与焓变}

\subsection{体系、环境与反应热}
研究化学反应的热效应，首先要划定研究对象。被研究的物质系统叫\textbf{体系}（如试管中的盐酸、NaOH溶液及发生的反应），与体系相互影响的其他部分叫\textbf{环境}（如试管、空气）。热量总是在体系和环境之间传递。

\begin{tuozhan}[三种热力学体系]
按体系与环境之间有无物质和能量交换，可分为：\textbf{敞开体系}（既有物质交换又有能量交换，如敞口烧杯中的反应）、\textbf{封闭体系}（只有能量交换、没有物质交换，如密闭容器中的反应）、\textbf{孤立体系}（既无物质交换也无能量交换，如理想的绝热密闭容器）。量热计的设计思路，就是尽量使反应体系接近“与外界无热交换”的状态，从而准确测出温度变化。
\end{tuozhan}

\begin{gainian}[反应热]
在\textbf{等温}条件下（反应前后体系温度相同），化学反应体系向环境释放或从环境吸收的热量，称为化学反应的热效应，简称\textbf{反应热}。
\end{gainian}

为什么强调“等温”？反应放出的热量会先使体系温度升高，只有让体系恢复到反应前的温度，传给环境的热量才完全由“化学变化”本身决定，否则一部分能量还“存”在体系的升温里。量热计测得的温度变化，正是用来反推这部分热量的。

\subsection{焓与焓变}
化学反应之所以有热效应，是因为反应前后体系的\textbf{内能}（$U$，体系内物质各种能量的总和）发生了变化。生产和科研中的反应大多在\textbf{等压}（敞口，压强等于大气压）下进行，为方便描述，引入与内能有关的物理量——\textbf{焓}（$H$）。

\begin{gainian}[焓变]
中学阶段规定：等压条件下进行的化学反应，其反应热等于反应的焓变，用 $\Delta H$ 表示，单位 kJ·mol$^{-1}$。
\[\Delta H=H(\text{生成物})-H(\text{反应物})\]
放热反应：体系能量降低，$\Delta H<0$；\quad 吸热反应：体系能量升高，$\Delta H>0$。
\end{gainian}

注意 $\Delta H$ 的正负号是“站在体系的立场”说的：体系放出热量，自身焓减少，故为负。“$\Delta H=-92.4$ kJ·mol$^{-1}$”中的“$-$”表示放热，不是“数值小”。

\subsection{从宏观角度认识：物质的能量}
\begin{center}
\begin{tikzpicture}[scale=0.72]
\begin{scope}
\draw[-Stealth] (0,0)--(4.4,0) node[below left]{\small 反应进程};
\draw[-Stealth] (0,0)--(0,3.3) node[left]{\small 焓 $H$};
\draw[thick] (0.4,2.4)--(1.6,2.4) node[midway,above]{\small 反应物};
\draw[thick] (2.6,0.8)--(3.8,0.8) node[midway,above]{\small 生成物};
\draw[-Stealth,warnred,thick] (2.2,2.4)--(2.2,0.8) node[midway,left]{\small$\Delta H<0$};
\node at (2.1,-1.0) {\small 放热反应};
\end{scope}
\begin{scope}[xshift=6.5cm]
\draw[-Stealth] (0,0)--(4.4,0) node[below left]{\small 反应进程};
\draw[-Stealth] (0,0)--(0,3.3) node[left]{\small 焓 $H$};
\draw[thick] (0.4,0.8)--(1.6,0.8) node[midway,above]{\small 反应物};
\draw[thick] (2.6,2.4)--(3.8,2.4) node[midway,above]{\small 生成物};
\draw[-Stealth,mainblue,thick] (2.2,0.8)--(2.2,2.4) node[midway,left]{\small$\Delta H>0$};
\node at (2.1,-1.0) {\small 吸热反应};
\end{scope}
\end{tikzpicture}
\end{center}

\begin{center}
\small
\begin{tabularx}{\linewidth}{>{\centering\arraybackslash}p{1.6cm}X}
\toprule
类型 & 常见实例 \\
\midrule
放热反应 & 燃烧；酸碱中和；活泼金属与水或酸；大多数化合反应；铝热反应；物质的缓慢氧化；\ce{2H2O2 = 2H2O + O2 ^}（分解但放热） \\
吸热反应 & 大多数分解反应；\ce{Ba(OH)2*8H2O} 与 \ce{NH4Cl} 的反应；\ce{C + H2O(g) ->[\text{高温}] CO + H2}；\ce{C + CO2 ->[\text{高温}] 2CO}（化合但吸热） \\
\bottomrule
\end{tabularx}
\end{center}

\begin{yicuo}
\begin{enumerate}
\item 反应吸热还是放热，只由反应物和生成物的能量相对高低决定，\textbf{与反应条件无关}。需要加热的反应不一定吸热（如木炭燃烧需点燃）；不需加热的反应也可能吸热（如 \ce{Ba(OH)2*8H2O} 与 \ce{NH4Cl} 常温搅拌即可反应）。“停止加热后仍能持续进行”的反应一般是放热反应。
\item “吸热反应”“放热反应”专指\textbf{化学反应}。浓硫酸稀释、NaOH 固体溶解放热，\ce{NH4NO3} 溶解吸热，水的汽化吸热，都是物理过程，不能叫“放热反应”或“吸热反应”。
\end{enumerate}
\end{yicuo}

\subsection{从微观角度认识：化学键的断裂与形成}
化学反应的本质是旧化学键断裂、新化学键形成。\textbf{断键要吸收能量，成键会释放能量}——两个原子靠近成键后能量降低才稳定，要把它们拉开就必须从外界补充同样多的能量。

\begin{gainian}[键能]
25 ℃、101 kPa 下，断开 1 mol 气态分子 AB(g) 中的共价键，使其生成气态原子 A(g) 和 B(g) 所吸收的能量。键能越大，化学键越牢固，含该键的分子越稳定。
\end{gainian}

以 \ce{H2(g) + Cl2(g) = 2HCl(g)} 为例（键能：H—H 436、Cl—Cl 243、H—Cl 431，单位 kJ·mol$^{-1}$）：断开 1 mol H—H 和 1 mol Cl—Cl 共吸收 $436+243=679$ kJ；形成 2 mol H—Cl 共放出 $2\times431=862$ kJ。放出多于吸收，净放热 183 kJ，即 $\Delta H\approx-183$ kJ·mol$^{-1}$，与实验值 $-184.6$ kJ·mol$^{-1}$ 很接近。这说明：\textbf{化学键断裂和形成时的能量变化，是化学反应中能量变化的主要原因}。

\begin{benzhi}[为什么 $\Delta H=\sum E(\text{反应物键能})-\sum E(\text{生成物键能})$]
设想反应分两步走：第一步把反应物的键全部拆成气态原子，吸热 $\sum E_\text{反}$（焓增加）；第二步原子重新组合成生成物，放热 $\sum E_\text{生}$（焓减少）。由于焓变只与始终态有关（见1.5节），$\Delta H=(+\sum E_\text{反})+(-\sum E_\text{生})$。所以“\textbf{键能算焓变：反应物减生成物}”，而“\textbf{物质能量算焓变：生成物减反应物}”，两者顺序恰好相反，原因是\textbf{键能越大，物质本身的能量反而越低}。
\end{benzhi}

用键能计算时，关键是数清 1 mol 物质中含多少摩尔化学键，特别是空间网状或特殊结构的物质：

\begin{center}
\small
\begin{tabular}{cccccccc}
\toprule
物质 & \ce{P4}（白磷） & 晶体硅 & \ce{SiO2} & 金刚石 & 石墨 & \ce{S8} & \ce{CO2} \\
\midrule
1 mol 中键的物质的量 & 6 mol P—P & 2 mol Si—Si & 4 mol Si—O & 2 mol C—C & 1.5 mol C—C & 8 mol S—S & 2 mol C=O \\
\bottomrule
\end{tabular}
\end{center}

\begin{liti}[键能计算]
\ce{SF6} 是一种优良的绝缘气体，分子中只存在 S—F 键。已知：1 mol S(s) 转化为气态硫原子吸收能量 280 kJ，断裂 1 mol F—F、1 mol S—F 分别需吸收能量 160 kJ、330 kJ。则 \ce{S(s) + 3F2(g) = SF6(g)} 的 $\Delta H$ 为
\xx{$-1780$ kJ·mol$^{-1}$}{$-1220$ kJ·mol$^{-1}$}{$-450$ kJ·mol$^{-1}$}{$+430$ kJ·mol$^{-1}$}
\tcblower
\silu “吸收的减去放出的”：使反应物变成气态原子所吸收的能量 $-$ 由气态原子形成生成物所放出的能量。\\
\jie 吸收：$280+3\times160=760$ kJ；放出：1 个 \ce{SF6} 含 6 个 S—F 键，$6\times330=1980$ kJ。$\Delta H=760-1980=-1220$ kJ·mol$^{-1}$。\daan B\\
\xianjing S(s) 变为气态 S 原子既要克服分子间作用力，又要断开 \ce{S8} 环中的 S—S 键，题目已直接给出“S(s)$\to$S(g)”的总能量，照用即可，不要再另算 S—S 键；不要漏乘 6。
\end{liti}

\section{中和反应反应热及其测定}

\begin{gainian}[中和反应反应热]
在稀溶液中，强酸和强碱发生中和反应生成 1 mol \ce{H2O(l)} 时所放出的热量。25 ℃、101 kPa 下：
\[\ce{H+(aq) + OH-(aq) = H2O(l)}\qquad \Delta H=-57.3\ \mathrm{kJ\cdot mol^{-1}}\]
（旧教材称“中和热”，新教材称“中和反应反应热”，含义相同。）
\end{gainian}

\textbf{实验原理：}用简易量热计（内筒、隔热层、外壳、杯盖、温度计、玻璃搅拌器）测出反应前后溶液的温度差，再由 $Q=cm\Delta t$ 求热量。

\textbf{实验步骤与数据处理：}用量筒量取 50 mL 0.50 mol·L$^{-1}$ 盐酸倒入内筒，测温；洗净并擦干温度计，再量取 50 mL 0.55 mol·L$^{-1}$ NaOH 溶液测温（取两者平均值为 $t_1$）；将 NaOH 溶液\textbf{一次迅速}倒入内筒，盖好杯盖，匀速搅拌，记录\textbf{最高温度} $t_2$；重复实验 2 次，取温差平均值。近似认为溶液密度为 1 g·mL$^{-1}$、比热容 $c=4.18$ J·g$^{-1}$·℃$^{-1}$：
\[Q=(50+50)\,\mathrm{g}\times4.18\times10^{-3}\,\mathrm{kJ\cdot g^{-1}\cdot ℃^{-1}}\times(t_2-t_1)=0.418(t_2-t_1)\ \mathrm{kJ}\]
生成 \ce{H2O} 为 0.025 mol，故 $\Delta H=-\dfrac{0.418(t_2-t_1)}{0.025}\ \mathrm{kJ\cdot mol^{-1}}$。

\begin{benzhi}[实验中的五个“为什么”]
\begin{enumerate}
\item \textbf{为何 NaOH 稍过量（0.55 vs 0.50）？}保证盐酸完全被中和，生成水的量由盐酸确定为 0.025 mol。
\item \textbf{为何要一次迅速倒入？}分多次或缓慢倒入，热量在操作中散失，$\Delta t$ 偏小。
\item \textbf{为何不能用铜（金属）搅拌器？}金属导热快，热量损失大。
\item \textbf{为何测完酸要洗净温度计再测碱？}否则残留的酸与碱在温度计上反应放热，使测得的碱初温偏高，$\Delta t$ 偏小。
\item \textbf{改用 100 mL 0.50 mol·L$^{-1}$ 盐酸与 100 mL 0.55 mol·L$^{-1}$ NaOH，结果会翻倍吗？}放出的总热量翻倍，但中和反应反应热是“生成 1 mol \ce{H2O}”的热量，数值基本不变。
\end{enumerate}
\end{benzhi}

\begin{center}
\small
\begin{tabularx}{\linewidth}{Xc}
\toprule
误差来源（与理论值 57.3 kJ·mol$^{-1}$ 比较测得的放热数值） & 结果 \\
\midrule
装置保温、隔热差；分多次加入碱液；用铜质搅拌器；温度计未洗直接测碱温 & 偏小 \\
用弱酸（如 \ce{CH3COOH}）或弱碱（如氨水）代替：弱电解质电离要吸热 & 偏小 \\
用浓硫酸或 NaOH 固体代替稀溶液：溶解（稀释）放热 & 偏大 \\
用稀硫酸与 \ce{Ba(OH)2} 溶液：生成 \ce{BaSO4} 沉淀也放热 & 偏大 \\
量取盐酸时仰视读数：实际 HCl 多于 0.025 mol，生成水多，但仍按 0.025 mol 计算 & 偏大 \\
\bottomrule
\end{tabularx}
\end{center}

\begin{liti}[误差分析]
某同学测得稀盐酸与稀 NaOH 溶液反应生成 1 mol \ce{H2O(l)} 放出热量 52.3 kJ，造成这一结果的原因\textbf{不可能}是
\xxb{实验装置保温、隔热效果差}{用量筒量取盐酸时仰视读数}{分多次将 NaOH 溶液倒入量热计内筒}{用测量盐酸的温度计直接测定 NaOH 溶液的温度}
\tcblower
\jie 52.3 < 57.3，测得值偏小。A、C 热量散失，偏小；D 碱初温读数偏高，$\Delta t$ 偏小；B 仰视读数使量取的盐酸偏多，实际生成的水多于 0.025 mol，放热更多，而计算时仍除以 0.025 mol，测得值偏大，不可能导致偏小。\daan B
\end{liti}

\section{热化学方程式}

化学方程式只表明物质变化；\textbf{热化学方程式}同时表明物质变化和能量变化。例如 \[\ce{H2(g) + Cl2(g) = 2HCl(g)}\dHv{-184.6}\]
表示：25 ℃、101 kPa 下，1 mol \ce{H2(g)} 与 1 mol \ce{Cl2(g)} 完全反应生成 2 mol \ce{HCl(g)}，放出 184.6 kJ 热量。

\begin{benzhi}[$\Delta H$ 单位中的 “mol$^{-1}$” 指什么？]
指“每摩尔反应”，即\textbf{按方程式所写的计量数完成一次反应}。所以 $\Delta H$ 的数值必须与计量数对应：计量数加倍，$\Delta H$ 加倍；反应逆向书写，$\Delta H$ 变号。对可逆反应，$\Delta H$ 指按方程式\textbf{完全}转化时的热效应，例如 \ce{N2(g) + 3H2(g) <=> 2NH3(g)} $\Delta H=-92.4$ kJ·mol$^{-1}$，在密闭容器中充入 1 mol \ce{N2} 和 3 mol \ce{H2}，因不能完全转化，实际放热\textbf{小于} 92.4 kJ。
\end{benzhi}

\textbf{书写规则：}
\begin{enumerate}
\item \textbf{注明聚集状态}：g、l、s、aq。因为同一物质状态不同，能量不同：等质量水的焓 $H(\mathrm{s})<H(\mathrm{l})<H(\mathrm{g})$。如 \ce{H2(g) + 1/2O2(g) = H2O(g)} $\Delta H=-241.8$ kJ·mol$^{-1}$，而生成 \ce{H2O(l)} 时 $\Delta H=-285.8$ kJ·mol$^{-1}$，差值 44 kJ 正是 1 mol 水蒸气液化放出的热。同素异形体要注明名称，如 C（石墨）。
\item \textbf{计量数只表示物质的量}，可以是分数；$\Delta H$ 与计量数成正比。
\item \textbf{$\Delta H$ 写在方程式右边}，放热为“$-$”，吸热为“$+$”，单位 kJ·mol$^{-1}$。
\item 25 ℃、101 kPa 下可不注明温度压强；一般不写反应条件；标明状态后不再用“↑”“↓”。
\end{enumerate}

\begin{fangfa}[热化学方程式正误“五审”]
\textbf{一审符号}（放热“$-$”、吸热“$+$”）；\textbf{二审单位}（kJ·mol$^{-1}$，易误写为 kJ）；\textbf{三审状态}（溶液中反应常漏写 aq、水的状态）；\textbf{四审数值与计量数对应}；\textbf{五审是否符合概念}（表示燃烧热时可燃物计量数为 1 且生成指定产物；表示中和反应反应热时生成 1 mol \ce{H2O(l)}）。
\end{fangfa}

\begin{kanwu}[课件 1.1 第47页例1]
课件写作 \ce{N2(g) + 3H2(aq) = 2NH3(g)}，\ce{H2} 的状态“aq”错误，应为 \ce{N2(g) + 3H2(g) = 2NH3(g)} $\Delta H=-92.2$ kJ·mol$^{-1}$（教材合成氨数据常用 $-92.4$ kJ·mol$^{-1}$，二者为不同来源的测定值）。
\end{kanwu}

\begin{liti}[由实验事实写热化学方程式]
常温下：(1) 8 g \ce{CH4} 完全燃烧生成 \ce{CO2(g)} 和 \ce{H2O(l)}，放出 445.0 kJ 热量；(2) 14 g CO 完全燃烧放出 141.5 kJ 热量。分别写出热化学方程式。
\tcblower
\silu 先把质量换算成物质的量，再按所选计量数“放大”热量。\\
\jie (1) 8 g \ce{CH4} 为 0.5 mol，1 mol 放热 890.0 kJ：\ce{CH4(g) + 2O2(g) = CO2(g) + 2H2O(l)}\dHv{-890.0}。\\
(2) 14 g CO 为 0.5 mol，若写 2 mol CO 则放热 $141.5\times4=566$ kJ：\ce{2CO(g) + O2(g) = 2CO2(g)}\dHv{-566}（写成 1 mol CO 时 $\Delta H=-283$ kJ·mol$^{-1}$ 亦正确）。
\end{liti}

\section{燃烧热与能源}

\begin{gainian}[燃烧热]
101 kPa 时，\textbf{1 mol} 纯物质\textbf{完全燃烧}生成\textbf{指定产物}时所放出的热量。指定产物：C 元素$\to$\ce{CO2(g)}，H 元素$\to$\ce{H2O(l)}，S 元素$\to$\ce{SO2(g)}，N 元素$\to$\ce{N2(g)}。
\end{gainian}

例如 \ce{CH4(g) + 2O2(g) = CO2(g) + 2H2O(l)} $\Delta H=-890.3$ kJ·mol$^{-1}$，可表述为“\ce{CH4} 的燃烧热为 890.3 kJ·mol$^{-1}$”或“\ce{CH4} 的燃烧热 $\Delta H=-890.3$ kJ·mol$^{-1}$”。注意：文字叙述“燃烧热”用正值；一旦写成“$\Delta H$”，必须带负号。

\begin{center}
\small
\begin{tabular}{lll}
\toprule
& 燃烧热 & 中和反应反应热 \\
\midrule
定量标准 & 可燃物为 1 mol & 生成 \ce{H2O(l)} 为 1 mol \\
产物要求 & 指定产物（如 \ce{H2O(l)}，不能是 \ce{H2O(g)}；C 生成 \ce{CO2} 而非 CO） & 稀溶液、强酸强碱 \\
数值 & 随可燃物不同而不同 & 强酸强碱稀溶液均约为 57.3 kJ·mol$^{-1}$ \\
\bottomrule
\end{tabular}
\end{center}

\begin{kanwu}[课件 1.1 第56页]
课件把“Cl$\to$\ce{Cl2(g)}”列为燃烧的指定产物，教材无此规定，含氯有机物燃烧时氯一般转化为 HCl。高考只要求掌握 C、H、S、N 四种元素的指定产物，不必记忆其他规定。
\end{kanwu}

\textbf{燃烧热的测定：}把待测物质放在充满 \ce{O2} 的密封金属容器（氧弹）中，将氧弹置于盛有一定量水的量热计内筒里，点火使物质完全燃烧，测出水温的变化即可求出燃烧热（弹式量热计）。

\textbf{燃料的选择：}比较燃料优劣常用“单位质量放热”（热值）。\ce{H2} 的燃烧热 285.8 kJ·mol$^{-1}$，换算为 $285.8/2\approx143$ kJ·g$^{-1}$，远高于 \ce{CH4}（约 55.6 kJ·g$^{-1}$）和碳（约 32.8 kJ·g$^{-1}$），且产物只有水，因此液氢—液氧是高能、清洁的火箭推进剂；缺点是液氢需低温储存、密度小、储罐体积大。\textbf{能源}分为化石燃料（煤、石油、天然气，不可再生）和新能源（太阳能、氢能、风能、地热能、海洋能、生物质能，多可再生）；节能的关键是提高燃料燃烧效率（如使燃料与空气充分接触、适当过量空气）和充分利用余热。

\begin{liti}[燃烧热的计算]
已知：\ce{2H2(g) + O2(g) = 2H2O(l)} $\Delta H=-571.6$ kJ·mol$^{-1}$；\ce{C3H8(g) + 5O2(g) = 3CO2(g) + 4H2O(l)} $\Delta H=-2220$ kJ·mol$^{-1}$。(1) \ce{H2} 的燃烧热为\tk；(2) \ce{H2} 和 \ce{C3H8} 的混合气体共 5 mol，完全燃烧放热 3846.75 kJ，则 $V(\ce{H2}):V(\ce{C3H8})=$\tk。
\tcblower
\jie (1) 燃烧热以 1 mol 可燃物计：$571.6\div2=285.8$ kJ·mol$^{-1}$（不是 571.6）。\\
(2) 设 \ce{H2} 为 $x$ mol、\ce{C3H8} 为 $y$ mol：$x+y=5$，$285.8x+2220y=3846.75$，解得 $x=3.75$，$y=1.25$，体积比 $3:1$。\\
\dianbo 也可用“平均值法”：平均每摩尔放热 $3846.75/5=769.35$ kJ，十字交叉 $(2220-769.35):(769.35-285.8)=1450.65:483.55=3:1$。
\end{liti}

\section{盖斯定律}

有些反应的反应热无法直接测定。例如 \ce{C(s) + 1/2O2(g) = CO(g)}：碳燃烧时总有部分生成 \ce{CO2}，无法控制只生成 CO。怎么得到它的 $\Delta H$？

\begin{gainian}[盖斯定律]
一个化学反应，不管是一步完成还是分几步完成，其反应热是相同的。即：在一定条件下，化学反应的反应热只与反应体系的\textbf{始态和终态}有关，而与反应的\textbf{途径}无关。
\end{gainian}

\begin{benzhi}[为什么与途径无关]
焓像“海拔高度”一样是\textbf{状态的属性}：从山脚到山顶，无论走哪条路，高度差都一样。若两条途径的 $\Delta H$ 不同，就可以沿放热多的途径正向、沿吸热少的途径逆向构成循环，凭空得到能量——这违背能量守恒定律。
\end{benzhi}

\begin{center}
\begin{tikzpicture}[scale=0.8,every node/.style={font=\small}]
\node (A) at (0,0) {\ce{C(s) + O2(g)}};
\node (B) at (6,0) {\ce{CO2(g)}};
\node (C) at (3,-1.6) {\ce{CO(g) + 1/2O2(g)}};
\draw[-Stealth,thick] (A)--(B) node[midway,above]{$\Delta H_1=-393.5$ kJ·mol$^{-1}$};
\draw[-Stealth,thick] (A)--(C) node[midway,left=2pt]{$\Delta H_3=\ ?$};
\draw[-Stealth,thick] (C)--(B) node[midway,right=2pt]{$\Delta H_2=-283.0$ kJ·mol$^{-1}$};
\end{tikzpicture}
\end{center}
由盖斯定律 $\Delta H_1=\Delta H_3+\Delta H_2$，得 $\Delta H_3=-393.5-(-283.0)=-110.5$ kJ·mol$^{-1}$。

\begin{fangfa}[加和法“四步走”]
\textbf{①找}：在已知方程式中找出目标方程式各物质（优先找\textbf{只在一个已知式中出现}的物质，它决定该式的系数和方向）；\textbf{②调}：按目标式调整已知式的方向（逆写则 $\Delta H$ 变号）和系数（乘以几，$\Delta H$ 也乘以几）；\textbf{③加}：将调整后的式子相加，$\Delta H$ 同步相加；\textbf{④查}：检查中间物质是否全部消去、状态是否一致。
\end{fangfa}

\begin{liti}[盖斯定律的计算]
用 CuO 与 Cu 高温烧结可制取 \ce{Cu2O}。已知：\\
①\ce{2Cu(s) + O2(g) = 2CuO(s)} $\Delta H_1=-314$ kJ·mol$^{-1}$；②\ce{2Cu2O(s) + O2(g) = 4CuO(s)} $\Delta H_2=-292$ kJ·mol$^{-1}$。\\
则 \ce{CuO(s) + Cu(s) = Cu2O(s)} 的 $\Delta H=$\tk。
\tcblower
\silu 目标式中 \ce{Cu} 只出现在①（反应物侧，与目标同侧），\ce{Cu2O} 只出现在②（反应物侧，与目标异侧，需逆写）。\\
\jie 目标式$=\frac12\times$(①$-$②)：$\frac12[2\ce{Cu} + \ce{O2} - 2\ce{Cu2O} - \ce{O2} \to 2\ce{CuO}-4\ce{CuO}]$，整理得 \ce{Cu + CuO = Cu2O}。$\Delta H=\frac12(\Delta H_1-\Delta H_2)=\frac12(-314+292)=-11$ kJ·mol$^{-1}$。\\
\xianjing \ce{O2}、\ce{CuO} 在两式中都出现，它们只能“被动消去”，不要拿它们来定系数。
\end{liti}

\begin{zhenti}{2026年广东卷第19题(2)\quad 来源：仓库所附真题文档（网络收集版）}
硫化物的氧化可为生物活动提供能量：\\
\ce{2H2S(aq) + O2(aq) = 2S(s) + 2H2O(l)} \quad $\Delta H_1$\\
\ce{2S(s) + 3O2(aq) + 2H2O(l) = 2SO4^2-(aq) + 4H+(aq)} \quad $\Delta H_2$\\
\ce{SO4^2-(aq) + 3H2S(aq) + 2H+(aq) = 4S(s) + 4H2O(l)} \quad $\Delta H_3$\\
则 $\Delta H_3=$\tk（用 $\Delta H_1$ 和 $\Delta H_2$ 表示）。

\textbf{解析：}\ce{SO4^2-} 只出现在第2式的生成物侧，而目标式中它在反应物侧，且系数为 1，故取 $-\frac12\times$第2式；\ce{H2S} 只出现在第1式，目标式需 3 个，故取 $\frac32\times$第1式。验证：$\frac32$×①$-\frac12$×② 中 \ce{O2} 为 $\frac32-\frac32=0$，恰好消去。答案：$\Delta H_3=\frac32\Delta H_1-\frac12\Delta H_2$。
\end{zhenti}

\begin{zhenti}{2025年广东卷第19题(3)①\quad 来源：仓库所附真题文档}
以 \ce{TiCl4} 为原料制备 \ce{TiCl3} 的体系中存在：(a)\ce{Ti(s) + TiCl4(g) <=> 2TiCl2(s)} $\Delta H_1$；(b)\ce{TiCl2(s) + TiCl4(g) <=> 2TiCl3(g)} $\Delta H=+200.1$ kJ·mol$^{-1}$；(c)\ce{Ti(s) + 3TiCl4(g) <=> 4TiCl3(g)} $\Delta H=+132.4$ kJ·mol$^{-1}$。求 $\Delta H_1$。

\textbf{解析：}(a)$=$(c)$-2\times$(b)：$\ce{Ti} + 3\ce{TiCl4} - 2\ce{TiCl2} - 2\ce{TiCl4} \to 4\ce{TiCl3} - 4\ce{TiCl3}$，即 \ce{Ti + TiCl4 = 2TiCl2}。$\Delta H_1=132.4-2\times200.1=-267.8$ kJ·mol$^{-1}$。
\end{zhenti}

\begin{liti}[盖斯定律在火箭推进剂中的应用]
肼（\ce{N2H4}）与 \ce{N2O4} 可作火箭推进剂。已知：①\ce{N2(g) + 2O2(g) = N2O4(l)} $\Delta H_1=-19.5$ kJ·mol$^{-1}$；②\ce{N2H4(l) + O2(g) = N2(g) + 2H2O(g)} $\Delta H_2=-534.2$ kJ·mol$^{-1}$。写出 \ce{N2H4(l)} 与 \ce{N2O4(l)} 反应生成 \ce{N2(g)} 和 \ce{H2O(g)} 的热化学方程式，并说明该推进剂的优点。
\tcblower
\jie 目标：\ce{2N2H4(l) + N2O4(l) = 3N2(g) + 4H2O(g)}。\ce{N2H4} 只在②中，需 2 个，取 2×②；\ce{N2O4} 只在①中且在生成物侧，目标中在反应物侧，取 $-$①。$\Delta H=2\Delta H_2-\Delta H_1=-1068.4+19.5=-1048.9$ kJ·mol$^{-1}$。优点：放热多；产物为 \ce{N2} 和 \ce{H2O}，无污染；生成大量气体，推力大。
\end{liti}

\subsection{反应热大小的比较}
比较 $\Delta H$ 必须\textbf{带符号}比较。放热反应放热越多，$\Delta H$ 越\textbf{小}（越负）。最可靠的方法是画能量示意图：

\begin{center}
\small
\begin{tabularx}{\linewidth}{lXl}
\toprule
情形 & 分析 & 结论 \\
\midrule
生成物状态不同 & \ce{2H2(g) + O2(g) = 2H2O(l)} $\Delta H_1$；…$=$\ce{2H2O(g)} $\Delta H_2$。生成液态水多放出液化热 & $\Delta H_1<\Delta H_2$ \\
反应物状态不同 & \ce{S(g) + O2(g) = SO2(g)} $\Delta H_1$；\ce{S(s) + O2(g) = SO2(g)} $\Delta H_2$。气态硫能量高，放热多 & $\Delta H_1<\Delta H_2$ \\
燃烧程度不同 & \ce{C(s) + O2(g) = CO2(g)} $\Delta H_1$；\ce{C(s) + 1/2O2(g) = CO(g)} $\Delta H_2$。完全燃烧放热多 & $\Delta H_1<\Delta H_2$ \\
同素异形体 & C（石墨）$\to$ C（金刚石）$\Delta H=+1.9$ kJ·mol$^{-1}$：石墨能量低、更稳定，等量燃烧金刚石放热多 & $\Delta H_\text{金刚石}<\Delta H_\text{石墨}$ \\
\bottomrule
\end{tabularx}
\end{center}

\begin{yicuo}
“能量越低越稳定”：由 C（石墨）$\to$ C（金刚石）吸热可知石墨更稳定；由 Sn（灰）$\to$ Sn（白）$\Delta H=+2.1$ kJ·mol$^{-1}$ 可知灰锡的能量（焓）较低。但锡的稳定形态还与温度有关：高于 13.2 ℃ 时白锡稳定（常温下锡以白锡形式存在），低于 13.2 ℃ 时白锡才会缓慢转化为灰锡（锡器在寒冷环境中会“自行毁坏”）。
\end{yicuo}

\section{反应热计算的六种方法}

\begin{center}
\small
\begin{tabularx}{\linewidth}{p{2.4cm}Xp{4.6cm}}
\toprule
依据 & 计算式 & 注意 \\
\midrule
①物质能量（焓） & $\Delta H=\sum H_\text{生成物}-\sum H_\text{反应物}$ & 生成物减反应物 \\
②键能 & $\Delta H=\sum E_\text{反应物键能}-\sum E_\text{生成物键能}$ & 反应物减生成物；数准键数 \\
③活化能 & $\Delta H=E_\text{a(正)}-E_\text{a(逆)}$ & 催化剂同等降低两者，$\Delta H$ 不变 \\
④燃烧热 & $\Delta H=\sum\nu\Delta_\text{c}H_\text{反应物}-\sum\nu\Delta_\text{c}H_\text{生成物}$ & 代入带负号的 $\Delta_\text{c}H$；\ce{O2}、\ce{CO2}、\ce{H2O(l)} 的燃烧热为 0 \\
⑤盖斯定律 & 已知式加减组合，$\Delta H$ 同步加减 & 方向、系数同步调整 \\
⑥热量比例 & $Q=n\cdot|\Delta H|$ & 可逆反应实际放热 $<|\Delta H|$ \\
\bottomrule
\end{tabularx}
\end{center}

\begin{benzhi}[为什么燃烧热法也是“反应物减生成物”]
把反应物和生成物都“燃烧”到同一组最终产物（\ce{CO2}、\ce{H2O(l)} 等）：反应物 $\xrightarrow{\Delta H}$ 生成物 $\xrightarrow{\sum\Delta_\text{c}H_\text{生}}$ 最终产物，与 反应物 $\xrightarrow{\sum\Delta_\text{c}H_\text{反}}$ 最终产物 始终态相同，故 $\Delta H+\sum\Delta_\text{c}H_\text{生}=\sum\Delta_\text{c}H_\text{反}$。
\end{benzhi}

\begin{kanwu}[课件 1.2 第19页]
课件公式“$\Delta H=\sum$反应物的燃烧热$-\sum$生成物的燃烧热”只在“燃烧热”取\textbf{带负号的 $\Delta_\text{c}H$}时成立；若按“燃烧热为正值”的文字定义代入，所得 $\Delta H$ 符号恰好相反。安全写法：$\Delta H=\sum\Delta_\text{c}H_\text{反}-\sum\Delta_\text{c}H_\text{生}$，或用正值 $Q$ 写作 $\Delta H=\sum Q_\text{生}-\sum Q_\text{反}$。
\end{kanwu}

\begin{liti}[燃烧热法]
已知 \ce{H2(g)}、\ce{C2H4(g)} 和 \ce{C2H5OH(l)} 的燃烧热 $\Delta H$ 分别为 $-285.8$、$-1411.0$、$-1366.8$ kJ·mol$^{-1}$。则 \ce{C2H4(g) + H2O(l) = C2H5OH(l)} 的 $\Delta H$ 为\tk。
\tcblower
\jie \ce{H2O(l)} 是燃烧的最终产物，燃烧热为 0；\ce{H2} 不参与目标反应，是干扰数据。$\Delta H=(-1411.0+0)-(-1366.8)=-44.2$ kJ·mol$^{-1}$。\\
\xianjing 若直接用正值“$1411.0-1366.8=+44.2$”就错了符号。
\end{liti}

\begin{liti}[活化能与能量图]
\ce{N2O} 和 CO 在 \ce{Pt2O+} 表面转化为无害气体：\ce{N2O(g) + CO(g) = N2(g) + CO2(g)}。其能量变化如图，正反应活化能为 134 kJ·mol$^{-1}$，反应物与生成物能量差为 226 kJ·mol$^{-1}$（生成物更低）。(1) 该反应的 $\Delta H=$\tk；(2) 逆反应活化能为\tk；(3) 若不用 \ce{Pt2O+}，正反应活化能\tk，$\Delta H$\tk（填“增大”“减小”或“不变”）。
\begin{center}
\begin{tikzpicture}[scale=0.62]
\axisEP{7}{4.6}
\draw[thick,mainblue] (0.3,2.2)--(1.2,2.2) .. controls (2.2,2.2) and (2.4,3.9) .. (3.2,3.9) .. controls (4,3.9) and (4.3,0.6) .. (5.4,0.6)--(6.5,0.6);
\draw[dashed] (1.2,2.2)--(3.2,2.2); \draw[dashed] (3.2,3.9)--(6.3,3.9); \draw[dashed] (1.2,2.2)--(6.3,2.2);
\draw[<->] (3.2,2.2)--(3.2,3.9) node[midway,left]{\scriptsize 134};
\draw[<->] (5.9,0.6)--(5.9,2.2) node[midway,right]{\scriptsize 226};
\draw[<->] (6.4,0.6)--(6.4,3.9) node[midway,right]{\scriptsize ?};
\end{tikzpicture}
\end{center}
\tcblower
\jie (1) 生成物能量低，放热：$\Delta H=-226$ kJ·mol$^{-1}$。(2) $\Delta H=E_\text{a正}-E_\text{a逆}$，$E_\text{a逆}=134+226=360$ kJ·mol$^{-1}$。(3) \ce{Pt2O+} 是催化剂，去掉后活化能增大；$\Delta H$ 只与始终态有关，不变。
\end{liti}

\begin{liti}[两步反应的能量图]
反应 \ce{A -> C} 分两步进行：①\ce{A -> B}，②\ce{B -> C}。反应过程的能量变化如图（$E_1$、$E_2$、$E_3$、$E_4$ 均为活化能，单位 kJ·mol$^{-1}$）。下列说法\textbf{不正确}的是
\begin{center}
\begin{tikzpicture}[scale=0.6,font=\scriptsize]
\axisEP{9}{5}
\draw[thick,mainblue] (0.3,1.8)--(0.8,1.8) .. controls (1.6,1.8) and (1.8,4.0) .. (2.4,4.0) .. controls (3.0,4.0) and (3.2,2.8) .. (3.9,2.8) .. controls (4.6,2.8) and (4.8,4.4) .. (5.4,4.4) .. controls (6.0,4.4) and (6.3,0.7) .. (7.2,0.7)--(8.3,0.7);
\node[below] at (0.6,1.8) {A}; \node[below] at (3.9,2.8) {B}; \node[below] at (7.8,0.7) {C};
\draw[<->] (1.2,1.8)--(1.2,4.0) node[midway,left]{$E_1$}; \draw[dashed] (0.8,1.8)--(1.5,1.8);
\draw[<->] (3.1,2.8)--(3.1,4.0) node[midway,left]{$E_2$}; \draw[dashed] (2.4,4.0)--(3.3,4.0);
\draw[<->] (4.6,2.8)--(4.6,4.4) node[midway,left]{$E_3$}; \draw[dashed] (3.9,2.8)--(4.8,2.8);
\draw[<->] (6.6,0.7)--(6.6,4.4) node[midway,right]{$E_4$}; \draw[dashed] (5.4,4.4)--(6.8,4.4);
\end{tikzpicture}
\end{center}
\xxb{三种物质中 C 最稳定}{\ce{A -> B} 的活化能为 $E_1$}{\ce{B -> C} 的 $\Delta H=E_4-E_3$}{加入催化剂不能改变反应的焓变}
\tcblower
\jie 能量越低越稳定，C 最稳定，A 对；A$\to$B 的正反应活化能是 $E_1$，B 对；B$\to$C 的 $\Delta H=E_3-E_4$（正反应活化能减逆反应活化能），且 $E_4>E_3$，$\Delta H<0$，C 错；催化剂不改变焓变，D 对。答案：C。\\
\dianbo 整个反应 \ce{A -> C} 的 $\Delta H=(E_1-E_2)+(E_3-E_4)$，体现盖斯定律。
\end{liti}

\begin{tuozhan}[标准摩尔生成焓]
由最稳定单质生成 1 mol 化合物的焓变称为该物质的标准摩尔生成焓 $\Delta_\text{f}H_\text{m}^\ominus$，最稳定单质的 $\Delta_\text{f}H_\text{m}^\ominus=0$。$\Delta H=\sum\nu\Delta_\text{f}H_\text{生成物}-\sum\nu\Delta_\text{f}H_\text{反应物}$（“生成物减反应物”，与焓的计算一致）。高考若出现，题目会给出定义，会用即可。
\end{tuozhan}

% ------------------------------------------------------------
\zhangmo{本章小结}
\subsection*{一、知识总结}
\begin{center}
\small
\begin{tabularx}{\linewidth}{p{2.2cm}X}
\toprule
反应热本质 & 宏观：反应物与生成物焓不同；微观：断键吸热 $<$ 成键放热则放热 \\
$\Delta H$ 符号 & 放热 $\Delta H<0$，吸热 $\Delta H>0$；与条件无关 \\
测定 & 中和反应反应热 $-57.3$ kJ·mol$^{-1}$；$\Delta H=-cm\Delta t/n(\ce{H2O})$ \\
表示 & 热化学方程式：标状态、$\Delta H$ 与计量数对应、可逆反应指完全转化 \\
燃烧热 & 1 mol 可燃物、完全燃烧、指定产物（\ce{CO2(g)}、\ce{H2O(l)}、\ce{SO2(g)}、\ce{N2(g)}） \\
计算 & 能量（生$-$反）、键能（反$-$生）、活化能（正$-$逆）、燃烧热（反$-$生）、盖斯定律、比例 \\
\bottomrule
\end{tabularx}
\end{center}

\subsection*{二、方法提升}
\begin{enumerate}
\item \textbf{盖斯定律“定系数”技巧}：先找只在一个已知式中出现的物质，由它在目标式中的位置和系数确定该已知式的方向和倍数；最后用出现在多个式中的物质验算。
\item \textbf{比较 $\Delta H$ 大小}：画能量图，带符号比较；放热越多 $\Delta H$ 越小，吸热越多 $\Delta H$ 越大。
\item \textbf{能量图读图}：峰顶对应过渡态，起点到峰顶为正反应活化能，终点到峰顶为逆反应活化能，起终点差为 $\Delta H$。
\end{enumerate}

\subsection*{三、易错点归纳}
\begin{enumerate}
\item 把“需要加热”等同于吸热；把物理过程（溶解、稀释、汽化）的热效应称为“吸热反应”“放热反应”。
\item 写热化学方程式漏写状态、$\Delta H$ 单位写成 kJ、放热漏“$-$”、计量数加倍 $\Delta H$ 不变。
\item 燃烧热生成 \ce{H2O(g)} 或 CO；燃烧热的可燃物系数不为 1。
\item 可逆反应中误认为实际放热等于 $|\Delta H|$。
\item 键能法与能量法的减法顺序弄反；数错 \ce{P4}、\ce{SiO2}、晶体硅中的键数。
\item 催化剂能改变活化能，但不改变 $\Delta H$。
\end{enumerate}

% ------------------------------------------------------------
\zhangmo{分层练习}
\zu{A组\quad 基础巩固}
\begin{xiti}[start=1]
\item 下列说法正确的是\xz
\xxa{吸热反应一定需要加热才能发生}{化学反应中，反应物的总能量大于生成物的总能量时，反应放热，$\Delta H<0$}{NaOH 固体溶于水放出热量，该过程属于放热反应}{同温同压下，\ce{H2(g) + Cl2(g) = 2HCl(g)} 在光照和点燃条件下的 $\Delta H$ 不同}

\item 下列有关中和反应反应热测定实验的说法正确的是\xz
\xxa{可用温度计代替玻璃搅拌器搅拌溶液}{强酸与强碱的稀溶液反应生成 1 mol \ce{H2O(l)} 放出的热量都约为 57.3 kJ，因此稀硫酸与 \ce{Ba(OH)2} 稀溶液反应也可用于测定}{应读取混合溶液的最高温度作为终止温度}{用 \ce{CH3COOH} 溶液代替盐酸，测得的放热数值偏大}

\item 已知 1 g \ce{H2} 完全燃烧生成液态水放出 142.9 kJ 热量。下列热化学方程式正确的是\xz
\xxa{\ce{2H2(g) + O2(g) = 2H2O(l)} \quad $\Delta H=-142.9$ kJ·mol$^{-1}$}{\ce{H2(g) + 1/2O2(g) = H2O(l)} \quad $\Delta H=-285.8$ kJ·mol$^{-1}$}{\ce{2H2 + O2 = 2H2O} \quad $\Delta H=-571.6$ kJ·mol$^{-1}$}{\ce{H2(g) + 1/2O2(g) = H2O(g)} \quad $\Delta H=-285.8$ kJ·mol$^{-1}$}

\item 下列关于燃烧热的说法正确的是\xz
\xxa{1 mol 碳燃烧生成 CO 时放出的热量是碳的燃烧热}{表示 \ce{H2} 燃烧热的热化学方程式中，生成物应为 \ce{H2O(g)}}{\ce{CH4} 的燃烧热为 890.3 kJ·mol$^{-1}$，则 \ce{CH4(g) + 2O2(g) = CO2(g) + 2H2O(l)} $\Delta H=-890.3$ kJ·mol$^{-1}$}{燃烧热的数值与可燃物的物质的量成正比}

\item 已知部分键能（kJ·mol$^{-1}$）：H—H 436，O=O 498，H—O 463。则 \ce{2H2(g) + O2(g) = 2H2O(g)} 的 $\Delta H=$\tk；若已知 \ce{H2O(l) = H2O(g)} $\Delta H=+44$ kJ·mol$^{-1}$，则 \ce{H2(g) + 1/2O2(g) = H2O(l)} 的 $\Delta H=$\tk。

\item 化学反应 \ce{A2(g) + B2(g) = 2AB(g)} 的能量变化如图（$a$、$b$ 均为正值，单位 kJ·mol$^{-1}$）。下列叙述正确的是\xz
\begin{center}
\begin{tikzpicture}[scale=0.55]
\axisEP{6.5}{4.4}
\draw[thick] (0.3,1)--(1.3,1) node[above]{\scriptsize 1 mol A$_2$+1 mol B$_2$};
\draw[thick] (1.6,3.8)--(3.4,3.8) node[above]{\scriptsize 2 mol A+2 mol B};
\draw[thick] (4.4,2.2)--(5.9,2.2) node[above]{\scriptsize 2 mol AB};
\draw[<->] (2,1)--(2,3.8) node[midway,right]{\scriptsize $a$};
\draw[<->] (4.6,2.2)--(4.6,3.8) node[midway,right]{\scriptsize $b$};
\draw[dashed] (1.3,1)--(2,1); \draw[dashed] (3.4,3.8)--(4.6,3.8);
\end{tikzpicture}
\end{center}
\xxb{每生成 2 mol AB(g) 吸收 $b$ kJ 热量}{$\Delta H=+(a-b)$ kJ·mol$^{-1}$}{反应物总能量高于生成物总能量}{断裂 1 mol A—A 和 1 mol B—B 放出 $a$ kJ 能量}

\item 已知 \ce{H2(g) + I2(g) <=> 2HI(g)} $\Delta H=-9.48$ kJ·mol$^{-1}$。下列说法正确的是\xz
\xxa{1 mol \ce{H2(g)} 与 1 mol \ce{I2(g)} 在密闭容器中充分反应，放出 9.48 kJ 热量}{生成 2 mol HI(g) 时放出 9.48 kJ 热量}{\ce{2HI(g) <=> H2(g) + I2(s)} 的 $\Delta H=+9.48$ kJ·mol$^{-1}$}{1 mol \ce{H2(g)} 与 1 mol \ce{I2(g)} 的总能量低于 2 mol HI(g) 的能量}
\end{xiti}

\zu{B组\quad 高考典型}
\begin{xiti}
\item \laiyuan{原创}\ \ce{CO2} 加氢制甲醇是实现“碳中和”的重要途径。已知：\\
①\ce{CO2(g) + H2(g) = CO(g) + H2O(g)} \quad $\Delta H_1=+41.2$ kJ·mol$^{-1}$\\
②\ce{CO(g) + 2H2(g) = CH3OH(g)} \quad $\Delta H_2=-90.7$ kJ·mol$^{-1}$\\
则 \ce{CO2(g) + 3H2(g) = CH3OH(g) + H2O(g)} 的 $\Delta H_3=$\tk kJ·mol$^{-1}$；该反应在\tk（填“高温”或“低温”）下更有利于自发进行（提示：$\Delta S<0$，可学完第3章再做）。

\item 下列各组热化学方程式中，$\Delta H_1>\Delta H_2$ 的是\xz
\xxa{\ce{C(s) + O2(g) = CO2(g)} $\Delta H_1$；\ce{C(s) + 1/2O2(g) = CO(g)} $\Delta H_2$}{\ce{S(s) + O2(g) = SO2(g)} $\Delta H_1$；\ce{S(g) + O2(g) = SO2(g)} $\Delta H_2$}{\ce{H2(g) + 1/2O2(g) = H2O(l)} $\Delta H_1$；\ce{H2(g) + 1/2O2(g) = H2O(g)} $\Delta H_2$}{\ce{CaO(s) + H2O(l) = Ca(OH)2(s)} $\Delta H_1$；\ce{CaCO3(s) = CaO(s) + CO2(g)} $\Delta H_2$}

\item 发射火箭可用液氢和液氧作推进剂。已知：①\ce{H2(g) + 1/2O2(g) = H2O(l)} $\Delta H_1=-285.8$ kJ·mol$^{-1}$；②\ce{H2(g) = H2(l)} $\Delta H_2=-0.92$ kJ·mol$^{-1}$；③\ce{O2(g) = O2(l)} $\Delta H_3=-6.84$ kJ·mol$^{-1}$；④\ce{H2O(l) = H2O(g)} $\Delta H_4=+44.0$ kJ·mol$^{-1}$。则 \ce{H2(l) + 1/2O2(l) = H2O(g)} 的 $\Delta H$ 为\xz
\xx{$+237.46$ kJ·mol$^{-1}$}{$-474.92$ kJ·mol$^{-1}$}{$-118.73$ kJ·mol$^{-1}$}{$-237.46$ kJ·mol$^{-1}$}

\item \laiyuan{2024年广东卷第19题(1)②，来源：仓库所附真题文档}\ 酸催化下 \ce{NaNO2} 与 \ce{NH4Cl} 混合溶液的反应（反应a）可用于石油开采中油路解堵。反应a：\ce{NO2^-(aq) + NH4^+(aq) = N2(g) + 2H2O(l)}。已知：\ce{NaNO2(s) + NH4Cl(s) = N2(g) + NaCl(s) + 2H2O(l)} $\Delta H_1$；\ce{NaNO2(s)} 溶解为 \ce{Na+(aq) + NO2^-(aq)} 的焓变为 $\Delta H_2$；\ce{NH4Cl(s)} 溶解为 \ce{Cl-(aq) + NH4^+(aq)} 的焓变为 $\Delta H_3$；\ce{NaCl(s)} 溶解为 \ce{Na+(aq) + Cl-(aq)} 的焓变为 $\Delta H_4$。则反应a的 $\Delta H=$\tk。

\item 已知石墨和金刚石的燃烧热 $\Delta H$ 分别为 $-393.5$ kJ·mol$^{-1}$、$-395.4$ kJ·mol$^{-1}$。下列说法正确的是\xz
\xxa{由石墨制备金刚石是放热过程}{等质量时，石墨的能量比金刚石的高}{石墨比金刚石稳定，石墨中 C—C 键的平均键能更大}{C（石墨，s）$=$ C（金刚石，s）$\Delta H=-1.9$ kJ·mol$^{-1}$}

\item \laiyuan{原创}\ 用 \ce{CH4} 催化还原 \ce{NO_x} 可消除氮氧化物污染。已知：①\ce{CH4(g) + 4NO2(g) = 4NO(g) + CO2(g) + 2H2O(g)} $\Delta H_1=-574$ kJ·mol$^{-1}$；②\ce{CH4(g) + 4NO(g) = 2N2(g) + CO2(g) + 2H2O(g)} $\Delta H_2=-1160$ kJ·mol$^{-1}$。(1) 写出 \ce{CH4} 将 \ce{NO2} 还原为 \ce{N2} 的热化学方程式（\ce{CH4} 的计量数为 1）。(2) 若用标准状况下 4.48 L \ce{CH4} 将 \ce{NO2} 还原为 \ce{N2}，放出的热量为多少？转移电子多少摩尔？

\item \laiyuan{原创}\ 已知键能（kJ·mol$^{-1}$）：\ce{N#N} 946，H—H 436，N—H 391。(1) 估算 \ce{N2(g) + 3H2(g) = 2NH3(g)} 的 $\Delta H$。(2) 从键能角度解释为什么 \ce{N2} 化学性质稳定，而合成氨反应却是放热反应。
\end{xiti}

\zu{C组\quad 冲刺高分}
\begin{xiti}
\item 白磷（\ce{P4}）、\ce{P4O6}、\ce{P4O10} 的分子结构：\ce{P4} 为正四面体，6 条棱为 P—P 键；\ce{P4O6} 是在 \ce{P4} 的每条 P—P 键中插入 1 个 O 原子；\ce{P4O10} 是在 \ce{P4O6} 基础上每个 P 原子再以 P=O 连接 1 个 O 原子。部分键能（kJ·mol$^{-1}$）：P—P 198，P—O 360，O=O 498，P=O $x$。
(1) 已知 1 mol 白磷完全燃烧生成 \ce{P4O10(s)} 的 $\Delta H=-2982$ kJ·mol$^{-1}$，求 $x$（忽略固态与气态的差别）。
(2) 0.5 mol 白磷与 \ce{O2} 反应完全生成 \ce{P4O6} 放出多少热量？
(3) 已知 \ce{P4}（白磷，s）$=$ 4P（红磷，s）$\Delta H=-29.2$ kJ·mol$^{-1}$，判断白磷、红磷哪种更稳定，并比较二者燃烧热的大小。

\item \laiyuan{原创}\ 甲烷干重整（\ce{CH4 + CO2 -> 2CO + 2H2}）可同时利用两种温室气体。已知 25 ℃ 时 \ce{CH4(g)}、\ce{CO(g)}、\ce{H2(g)} 的燃烧热 $\Delta H$ 分别为 $-890.3$、$-283.0$、$-285.8$ kJ·mol$^{-1}$；\ce{H2O(l) = H2O(g)} $\Delta H=+44.0$ kJ·mol$^{-1}$。
(1) 求 \ce{CH4(g) + CO2(g) = 2CO(g) + 2H2(g)} 的 $\Delta H$。
(2) 若将产物中的 \ce{H2} 全部燃烧生成水蒸气，求 \ce{2H2(g) + O2(g) = 2H2O(g)} 的 $\Delta H$。
(3) 有人认为“干重整是吸热反应，所以没有利用价值”。结合 (1)(2) 及能量守恒，评价该观点。

\item \laiyuan{原创}\ \ce{CO2} 加氢制甲醇可经两条途径：途径Ⅰ直接反应 \ce{CO2(g) + 3H2(g) = CH3OH(g) + H2O(g)} $\Delta H_1$；途径Ⅱ先发生 \ce{CO2(g) + H2(g) = CO(g) + H2O(g)} $\Delta H_2=+41.2$ kJ·mol$^{-1}$，再发生 \ce{CO(g) + 2H2(g) = CH3OH(g)} $\Delta H_3=-90.7$ kJ·mol$^{-1}$。已知 \ce{H2O(l) = H2O(g)} $\Delta H=+44.0$ kJ·mol$^{-1}$。(1) 求 $\Delta H_1$。(2) 求 \ce{CO2(g) + 3H2(g) = CH3OH(g) + H2O(l)} 的 $\Delta H$。(3) 画出途径Ⅱ的能量示意图（标出中间状态），并判断哪一步需要从外界获得能量。(4) 若在途径Ⅱ中加入只加快第二步的催化剂，总反应的 $\Delta H$ 是否改变？
\end{xiti}

% ------------------------------------------------------------
\zhangmo{练习解析}
\begin{jiexilist}
\item \ans{B}A 错：吸热反应不一定需加热，如 \ce{Ba(OH)2*8H2O} 与 \ce{NH4Cl}。B 对：反应物总能量高，多余能量放出，$\Delta H<0$。C 错：溶解是物理过程，不是“放热反应”。D 错：$\Delta H$ 只与始终态有关，与条件无关。
\item \ans{C}A 错：温度计不能搅拌，易损坏且测温不准。B 错：生成 \ce{BaSO4} 沉淀也放热，测得值偏大，不能用于测定中和反应反应热。C 对：反应放热使温度升高，最高温度对应反应恰好完成、散热最少的时刻。D 错：醋酸电离吸热，测得放热数值偏小。
\item \ans{B}1 g \ce{H2} 为 0.5 mol，故 1 mol \ce{H2} 生成液态水放热 285.8 kJ。A 中 2 mol \ce{H2} 应为 $-571.6$；C 未标状态；D 中生成的是 \ce{H2O(g)}，放热应少于 285.8 kJ。
\item \ans{C}A 错：碳的指定产物是 \ce{CO2(g)}。B 错：H 元素的指定产物是 \ce{H2O(l)}。C 对：燃烧热用 $\Delta H$ 表示时带负号。D 错：燃烧热以 1 mol 可燃物为标准，是定值；放出的总热量才与物质的量成正比。
\item \ans{$-482$ kJ·mol$^{-1}$；$-285$ kJ·mol$^{-1}$}$\Delta H=2\times436+498-4\times463=872+498-1852=-482$ kJ·mol$^{-1}$（1 个 \ce{H2O} 含 2 个 H—O）。生成 1 mol \ce{H2O(g)}：$-241$ kJ·mol$^{-1}$；再液化放热 44 kJ：$-241-44=-285$ kJ·mol$^{-1}$。
\item \ans{B}由图：断开 1 mol \ce{A2} 和 1 mol \ce{B2} 生成气态原子吸收 $a$ kJ，原子形成 2 mol AB 放出 $b$ kJ，且 $a>b$（2 mol AB 能量高于反应物），故 $\Delta H=+(a-b)$ kJ·mol$^{-1}$，反应吸热，A、C 错；断键吸热，D 错。
\item \ans{B}A：可逆反应不能进行完全，放热小于 9.48 kJ；B：$\Delta H$ 对应按方程式完全生成 2 mol HI(g)，正确；C：生成物为 \ce{I2(s)}，与 \ce{I2(g)} 能量不同，$\Delta H$ 不等于 $+9.48$；D：放热反应，反应物总能量高于生成物。
\item \ans{$-49.5$；低温}①$+$②得目标式，$\Delta H_3=41.2-90.7=-49.5$ kJ·mol$^{-1}$。$\Delta H<0$、$\Delta S<0$，由 $\Delta G=\Delta H-T\Delta S<0$ 可知低温有利于自发。
\item \ans{B}A：完全燃烧放热多，$\Delta H_1<\Delta H_2$。B：S(g) 能量高于 S(s)，燃烧放热更多，$\Delta H_2<\Delta H_1$，即 $\Delta H_1>\Delta H_2$，符合题意。C：生成液态水多放出液化热，$\Delta H_1<\Delta H_2$。D：\ce{CaO} 与水化合放热 $\Delta H_1<0$，\ce{CaCO3} 分解吸热 $\Delta H_2>0$，$\Delta H_1<\Delta H_2$。
\item \ans{D}目标式$=$①$-$②$-\frac12$③$+$④：\ce{H2(l)} 只在②中且需在反应物侧，故 $-$②；\ce{O2(l)} 只在③中，取 $-\frac12$③；\ce{H2O(g)} 只在④中，取 $+$④。$\Delta H=-285.8+0.92+\frac12\times6.84+44.0=-237.46$ kJ·mol$^{-1}$。
\item \ans{$\Delta H_1-\Delta H_2-\Delta H_3+\Delta H_4$}设计途径：\ce{NO2^-(aq) + NH4^+(aq)}（与 \ce{Na+}、\ce{Cl-} 一起）$\xrightarrow{-\Delta H_2-\Delta H_3}$ \ce{NaNO2(s) + NH4Cl(s)} $\xrightarrow{\Delta H_1}$ \ce{N2(g) + NaCl(s) + 2H2O(l)} $\xrightarrow{\Delta H_4}$ \ce{N2(g) + Na+(aq) + Cl-(aq) + 2H2O(l)}。\ce{Na+}、\ce{Cl-} 为旁观离子，两端相同，故反应a的 $\Delta H=\Delta H_1-\Delta H_2-\Delta H_3+\Delta H_4$。
\item \ans{C}由燃烧热求得 C（石墨）$\to$ C（金刚石）$\Delta H=-393.5-(-395.4)=+1.9$ kJ·mol$^{-1}$，吸热，A、D 错；石墨能量低、更稳定，B 错；石墨更稳定，其中化学键更牢固，C 对。
\item (1) (①$+$②)$\div2$：\ce{CH4(g) + 2NO2(g) = N2(g) + CO2(g) + 2H2O(g)} $\Delta H=\frac{-574-1160}{2}=-867$ kJ·mol$^{-1}$。(2) $n(\ce{CH4})=0.2$ mol，放热 $0.2\times867=173.4$ kJ；C 由 $-4\to+4$，每个 \ce{CH4} 失 8 个电子，转移 1.6 mol 电子。
\item (1) $\Delta H=946+3\times436-6\times391=-92$ kJ·mol$^{-1}$。(2) \ce{N#N} 键能很大，断开它需要很高的能量（活化能大），所以 \ce{N2} 很稳定、常温难反应；但生成物中 6 个 N—H 键放出的能量（2346 kJ）大于断开 1 个 \ce{N#N} 和 3 个 H—H 吸收的能量（2254 kJ），所以总反应放热。“稳定”说的是反应难易（与活化能有关），“放热”说的是能量高低，两者不矛盾。
\item \textbf{(1)} $x=585$。\ce{P4 + 5O2 = P4O10}：断键 $6\times198+5\times498=3678$ kJ；\ce{P4O10} 含 12 个 P—O 和 4 个 P=O，成键 $12\times360+4x$。$3678-(4320+4x)=-2982$，解得 $x=585$。
\textbf{(2)} 819 kJ。\ce{P4 + 3O2 = P4O6}：$\Delta H=6\times198+3\times498-12\times360=-1638$ kJ·mol$^{-1}$；0.5 mol 放热 819 kJ。
\textbf{(3)} 白磷转化为红磷放热，红磷能量低，红磷更稳定；等质量（同为 4 mol P）燃烧时，白磷放热更多，白磷的燃烧热（以 \ce{P4} 计）比 4 mol 红磷燃烧放出的热量多 29.2 kJ，即单位质量白磷燃烧放热更多。
\item \textbf{(1)} $\Delta H=+247.3$ kJ·mol$^{-1}$。燃烧热法：$\Delta H=\Delta_\text{c}H(\ce{CH4})+\Delta_\text{c}H(\ce{CO2})-2\Delta_\text{c}H(\ce{CO})-2\Delta_\text{c}H(\ce{H2})=-890.3+0+566.0+571.6=+247.3$ kJ·mol$^{-1}$（\ce{CO2} 是最终产物，燃烧热为 0）。
\textbf{(2)} \ce{2H2(g) + O2(g) = 2H2O(l)} $\Delta H=-571.6$ kJ·mol$^{-1}$，再加 \ce{2H2O(l) = 2H2O(g)} $+88.0$，得 $\Delta H=-483.6$ kJ·mol$^{-1}$。
\textbf{(3)} 观点片面。干重整吸收的能量（247.3 kJ）储存在产物 CO 和 \ce{H2} 中，产物燃烧时可以释放出来：1 mol \ce{CH4} 直接燃烧放热 890.3 kJ，而干重整产物 2 mol CO 和 2 mol \ce{H2} 燃烧（生成液态水）放热 $566.0+571.6=1137.6$ kJ，恰好比 890.3 kJ 多 247.3 kJ（能量守恒）。若用太阳能、工业余热等廉价能量驱动干重整，就能把低品位热能转化为化学能储存，同时消耗 \ce{CO2}、得到合成气（CO、\ce{H2}，化工原料），具有利用价值。
\item (1) $\Delta H_1=\Delta H_2+\Delta H_3=41.2-90.7=-49.5$ kJ·mol$^{-1}$。(2) 再加上 \ce{H2O(g) = H2O(l)}（$-44.0$）：$-49.5-44.0=-93.5$ kJ·mol$^{-1}$。(3) 示意图：起点“\ce{CO2(g) + 3H2(g)}”，第一步升高 41.2 kJ 到“\ce{CO(g) + H2O(g) + 2H2(g)}”，第二步降低 90.7 kJ 到终点“\ce{CH3OH(g) + H2O(g)}”，终点比起点低 49.5 kJ；第一步吸热，需要从外界获得能量。(4) 不改变，$\Delta H$ 只与始终态有关。
\end{jiexilist}
% <<<<<<<< 结束：chapters/ch01.tex

% >>>>>>>> 开始：chapters/ch02.tex
% ============================================================
% 第2章 化学反应速率与反应历程（对应课件：选必一2.1 化学反应速率）
% ============================================================
\chapter{化学反应速率与反应历程}

\begin{daoxue}
\textbf{知识主线：}怎样描述快慢（速率的表示、计算、测定）$\to$ 什么因素影响快慢（浓度、压强、温度、催化剂）$\to$ 为什么会影响（碰撞理论、活化能）$\to$ 反应是怎样一步步发生的（基元反应、反应历程、催化机理）$\to$ 定量描述（速率方程、速率常数）。

\textbf{高考考情：}广东卷选择题第15题连续考查“反应历程能量图”（2024年考中间体、决速步与催化剂，2026年考平行反应的能量图与浓度—时间图）；原理大题常考速率计算和实验数据分析（2026第19题(3)）；实验题考控制变量与速率测定（2024第17题(3)）。

\textbf{学习目标：}①熟练进行速率计算与比较；②能用碰撞理论解释各因素的影响；③会读能量—历程图和催化循环图，找出决速步、中间体和催化剂；④理解速率方程与速率常数。
\end{daoxue}

\section{化学反应速率的表示与测定}

\begin{gainian}[化学反应速率]
用单位时间内反应物浓度的减少或生成物浓度的增加来表示（均取正值）：
\[v(\mathrm{B})=\frac{\Delta c(\mathrm{B})}{\Delta t}=\frac{\Delta n(\mathrm{B})}{V\cdot\Delta t}\]
常用单位 mol·L$^{-1}$·s$^{-1}$ 或 mol·L$^{-1}$·min$^{-1}$（$1\ \mathrm{mol\cdot L^{-1}\cdot s^{-1}}=60\ \mathrm{mol\cdot L^{-1}\cdot min^{-1}}$）。
\end{gainian}

理解这个定义要抓住三点：
\begin{enumerate}
\item 这样算出的是一段时间内的\textbf{平均速率}，不是某一时刻的瞬时速率（瞬时速率是 $c$–$t$ 曲线上该点切线斜率的绝对值）。
\item \textbf{固体和纯液体}的浓度视为常数，不能用它们的浓度变化表示速率（可用单位时间内质量或物质的量的变化表示）。
\item 同一反应 $m\mathrm{A}+n\mathrm{B}=p\mathrm{C}+q\mathrm{D}$，用不同物质表示的速率\textbf{数值可能不同，但意义相同}，且
\[v(\mathrm{A}):v(\mathrm{B}):v(\mathrm{C}):v(\mathrm{D})=m:n:p:q,\quad\text{即}\ \frac{v(\mathrm{A})}{m}=\frac{v(\mathrm{B})}{n}=\frac{v(\mathrm{C})}{p}=\frac{v(\mathrm{D})}{q}\]
因为在同一时间、同一体积内，各物质浓度变化量之比必然等于计量数之比。所以表示速率时必须指明是哪种物质。
\end{enumerate}

\begin{kanwu}[课件 2.1 第4页]
课件称“纯液体表示速率可用 L/s 或 L/min，纯固体可用 g/s 或 g/min”。规范表述为：纯液体、纯固体的浓度视为常数，不能用浓度变化表示速率；需要时可用单位时间内物质的量或质量的变化表示其反应快慢。
\end{kanwu}

\begin{fangfa}[比较速率快慢的“归一法”]
①统一单位；②都换算成同一种物质表示的速率，或计算 $\dfrac{v(\mathrm{X})}{\text{X的计量数}}$；③比较数值大小。
\end{fangfa}

\begin{liti}[速率比较]
对于 \ce{3A(g) + B(g) = 2C(g) + 2D(g)}，在不同条件下测得：①$v(\mathrm{A})=0.6$ mol·L$^{-1}$·min$^{-1}$；②$v(\mathrm{B})=0.45$ mol·L$^{-1}$·min$^{-1}$；③$v(\mathrm{C})=0.015$ mol·L$^{-1}$·s$^{-1}$；④$v(\mathrm{D})=0.45$ mol·L$^{-1}$·min$^{-1}$。反应最快的是\tk。
\tcblower
\jie 统一为 mol·L$^{-1}$·min$^{-1}$ 并除以计量数：①$0.6/3=0.2$；②$0.45/1=0.45$；③$0.015\times60/2=0.45$；④$0.45/2=0.225$。②与③相同且最快。\\
\xianjing ③若忘记把“s”换成“min”，会算出 0.0075 而误判为最慢。
\end{liti}

\begin{liti}[“三段式”计算]
容积为 2 L 的密闭容器中，将 2 mol L 气体和 3 mol M 气体混合，发生 \ce{2L(g) + 3M(g) = $x$Q(g) + 3R(g)}，10 s 末生成 2.4 mol R，并测得 Q 的浓度为 0.4 mol·L$^{-1}$。求：(1) 10 s 末 L 的浓度；(2) 前 10 s 内 $v(\mathrm{M})$；(3) $x$。
\tcblower
\silu 列“起始—转化—某时刻”三段，\textbf{只有转化量之比等于计量数之比}。\\
\jie \begin{tabular}[t]{lcccc}
 & 2L & 3M & $x$Q & 3R \\
起始/mol & 2 & 3 & 0 & 0 \\
转化/mol & 1.6 & 2.4 & 0.8 & 2.4 \\
10 s 末/mol & 0.4 & 0.6 & 0.8 & 2.4
\end{tabular}\\
(1) $c(\mathrm{L})=0.4/2=0.2$ mol·L$^{-1}$；(2) $v(\mathrm{M})=\dfrac{2.4}{2\times10}=0.12$ mol·L$^{-1}$·s$^{-1}$；(3) $n(\mathrm{Q})=0.4\times2=0.8$ mol，$x:3=0.8:2.4$，$x=1$。
\end{liti}

\textbf{速率的测定：}原则上，任何一种\textbf{随反应进行而明显变化、又便于测量}的量都可以用来测速率。例如锌与稀硫酸反应，可测一定时间内收集的 \ce{H2} 体积、收集一定体积 \ce{H2} 所需的时间、溶液 pH 的变化、锌粒质量的减少等；有颜色变化的反应可用比色法；离子浓度变化明显的可用电导法；生成沉淀的反应可测出现浑浊所需的时间；也可以每隔一段时间取样，用\textbf{滴定法}测某物质的浓度。

\begin{liti}[用量气法测定速率]
用锌粒与 40 mL 1 mol·L$^{-1}$ 稀硫酸反应，用注射器收集 \ce{H2}，测得 0～10 s 内收集到 \ce{H2} 44.8 mL（已换算为标准状况，溶液体积变化忽略不计）。(1) 用 \ce{H2} 的体积表示该段时间的速率；(2) 用 $c(\ce{H+})$ 的变化表示的速率为多少 mol·L$^{-1}$·s$^{-1}$？(3) 用 \ce{Zn^2+} 表示呢？
\tcblower
\jie (1) $v(\ce{H2})=\dfrac{44.8\ \mathrm{mL}}{10\ \mathrm{s}}=4.48$ mL·s$^{-1}$。(2) $n(\ce{H2})=2\times10^{-3}$ mol，消耗 \ce{H+} $4\times10^{-3}$ mol，$\Delta c(\ce{H+})=\dfrac{4\times10^{-3}}{0.040}=0.1$ mol·L$^{-1}$，$v(\ce{H+})=0.01$ mol·L$^{-1}$·s$^{-1}$。(3) $v(\ce{Zn^2+})=0.005$ mol·L$^{-1}$·s$^{-1}$（为 $v(\ce{H+})$ 的一半）。\\
\dianbo 用不同的量表示速率，单位不同、数值不同，但描述的是同一反应的快慢。
\end{liti}

\section{影响化学反应速率的因素}

\subsection{内因与外因}
决定反应速率的主要因素是\textbf{反应物本身的性质}（内因），如钠与水剧烈反应、镁与冷水反应很慢。对同一反应，外界条件（外因）也会显著改变速率：

\begin{center}
\small
\begin{tabularx}{\linewidth}{p{1.8cm}Xp{4.6cm}}
\toprule
条件 & 规律（其他条件不变时） & 说明 \\
\midrule
浓度 & 增大反应物浓度，速率增大 & 固体、纯液体的量不影响速率 \\
压强 & 对有气体参加的反应，压缩体积增大压强，速率增大 & 压强通过改变\textbf{浓度}起作用 \\
温度 & 升高温度，速率增大（每升高 10 ℃，速率通常增大到原来的 2～4 倍） & 正、逆反应速率都增大；与吸热放热无关 \\
催化剂 & 使用（正）催化剂，速率显著增大 & 同等程度改变正、逆反应速率 \\
其他 & 增大固体表面积（粉碎）、形成原电池、光照、超声波、搅拌等 & 原电池使锌与酸反应加快 \\
\bottomrule
\end{tabularx}
\end{center}

\begin{benzhi}[压强为什么只能“通过浓度”起作用]
对气体反应，速率取决于单位体积内的反应物分子数（浓度）。“增大压强”如果没有使反应物浓度变化，速率就不变。三种情况：
\begin{enumerate}
\item \textbf{恒温压缩体积}：各气体浓度都增大，正、逆反应速率都增大。
\item \textbf{恒温恒容充入惰性气体}（不参与反应的 He、Ar、\ce{N2} 等）：总压增大，但各反应物浓度不变，\textbf{速率不变}。
\item \textbf{恒温恒压充入惰性气体}：为保持压强不变，容器体积增大，各反应物浓度减小，\textbf{速率减小}。
\end{enumerate}
\end{benzhi}

\begin{liti}[惰性气体与速率]
在带活塞的密闭容器中发生 \ce{3Fe(s) + 4H2O(g) ->[\text{高温}] Fe3O4(s) + 4H2(g)}。判断下列操作对反应速率的影响：(1) 增加铁片的质量（表面积不变）；(2) 容积不变，充入更多水蒸气；(3) 压强不变，充入 Ar；(4) 容积不变，充入 Ar。
\tcblower
\jie (1) 不变：固体浓度视为常数。(2) 增大：反应物 \ce{H2O(g)} 浓度增大。(3) 减小：体积增大，\ce{H2O(g)}、\ce{H2(g)} 浓度减小。(4) 不变：各物质浓度均不变。
\end{liti}

\subsection{控制变量法探究影响因素}
研究某一因素对速率的影响时，必须保证\textbf{其他条件完全相同}，只改变这一个因素。溶液实验中，常通过\textbf{补加蒸馏水使总体积相同}，从而只改变某一反应物的浓度。

常用实验：①\ce{Na2S2O3 + H2SO4 = Na2SO4 + SO2 ^ + S v + H2O}，以“出现浑浊所需时间”衡量速率，探究浓度、温度的影响；②\ce{2MnO4^- + 5H2C2O4 + 6H+ = 2Mn^2+ + 10CO2 ^ + 8H2O}，以“褪色时间”衡量速率，需保证 \ce{KMnO4} 少量、\ce{H2C2O4} 过量，否则不褪色；③\ce{2H2O2 = 2H2O + O2 ^}，比较 \ce{FeCl3}、\ce{CuSO4} 的催化效果（为排除阴离子干扰，可改用 \ce{Fe2(SO4)3} 与 \ce{CuSO4} 比较）。

\begin{benzhi}[速率—时间曲线为什么会“先增大后减小”]
锌与足量盐酸反应：开始一段时间速率逐渐增大，是因为反应放热使溶液温度升高（温度的影响占主导）；后来速率减小，是因为 $c(\ce{H+})$ 减小（浓度的影响占主导）。类似地，\ce{KMnO4} 与 \ce{H2C2O4} 反应开始很慢、后来突然加快，是因为生成的 \ce{Mn^2+} 对反应有催化作用（\textbf{自催化}）。分析这类曲线时要找出“两种相反的影响”，再判断哪种占主导。
\end{benzhi}

\begin{zhenti}{2024年广东卷第17题(3)②\quad 来源：仓库所附真题文档}
（节选改述）探究 \ce{Na2CO3} 溶液浓度对反应 \ce{CaSO4(s) + CO3^2-(aq) <=> CaCO3(s) + SO4^2-(aq)} 速率的影响：已知 \ce{Na2CO3} 溶液浓度为 $\frac{c_1V_1}{V_0}$ mol·L$^{-1}$（$V_0$ mL 该溶液滴定消耗 $c_1$ mol·L$^{-1}$ 盐酸 $V_1$ mL，终点时 \ce{CO3^2-} 转化为 \ce{HCO3^-}）。按下表配制总体积相同的溶液，分别加入 $m_1$ g \ce{CaSO4} 固体，反应 $t_1$ min 后过滤，取 $V_0$ mL 滤液，用同一盐酸滴定，记录消耗盐酸的体积。

\begin{center}\small
\begin{tabular}{ccccc}
\toprule
序号 & $V(\ce{Na2CO3})$/mL & $V(\ce{H2O})$/mL & $V$(滤液)/mL & $V_\text{消耗}$(盐酸)/mL \\
\midrule
a & 100.0 & 0 & $V_0$ & $2V_1/5$ \\
b & 80.0 & $x$ & $V_0$ & $3V_1/10$ \\
\bottomrule
\end{tabular}
\end{center}
则 $x=$\tk，测得的平均反应速率之比 $v_\text{a}:v_\text{b}=$\tk。

\textbf{解析：}总体积相同，$x=20.0$。速率之比等于相同时间内 \ce{CO3^2-} 浓度减少量之比，以 $V_0$ mL 溶液中“消耗盐酸量”表示 \ce{CO3^2-}：a 中起始相当于 $c_1V_1$，剩余 $\frac25c_1V_1$，减少 $\frac35c_1V_1$；b 中起始为 $0.8c_1V_1$，剩余 $\frac{3}{10}c_1V_1$，减少 $\frac12c_1V_1$。$v_\text{a}:v_\text{b}=\frac35:\frac12=6:5$。
\end{zhenti}

\section{碰撞理论与活化能}

\subsection{有效碰撞与活化分子}
气体分子间的碰撞极其频繁（例如 773 K 时，浓度为 $1\times10^{-3}$ mol·L$^{-1}$ 的 HI 气体中，每升气体内分子每秒碰撞约 $3.5\times10^{28}$ 次），若每次碰撞都发生反应，所有气体反应都将瞬间完成。事实并非如此，说明\textbf{只有少数碰撞能引发反应}。

\begin{gainian}[有效碰撞、活化分子与活化能]
能够发生化学反应的碰撞叫\textbf{有效碰撞}，需同时满足：①分子具有足够高的能量；②碰撞时取向合适。能量足够、能发生有效碰撞的分子叫\textbf{活化分子}。活化分子具有的平均能量与反应物分子平均能量之差叫\textbf{活化能}（$E_\text{a}$）。活化能越小，活化分子所占比例越大，反应越快。
\end{gainian}

\begin{center}
\begin{tikzpicture}[scale=0.85]
\begin{scope}
\draw[-Stealth] (0,0)--(5.4,0) node[below left]{\small 分子能量};
\draw[-Stealth] (0,0)--(0,2.8) node[left]{\small 分子比例};
\fill[mainblue!30] (3.025,0) -- plot[domain=5.5:9,samples=40] ({0.55*\x},{4.0*\x*\x*exp(-\x)}) -- (4.95,0) -- cycle;
\fill[warnred!30] (2.2,0) -- plot[domain=4:5.5,samples=30] ({0.55*\x},{4.0*\x*\x*exp(-\x)}) -- (3.025,0) -- cycle;
\draw[thick] plot[domain=0:9,samples=90] ({0.55*\x},{4.0*\x*\x*exp(-\x)});
\draw[dashed] (3.025,0)--(3.025,1.0) node[above]{\scriptsize $E_\text{a}$};
\draw[dashed] (2.2,0)--(2.2,1.55) node[above]{\scriptsize $E_\text{a}'$};
\node[font=\scriptsize,align=left,anchor=west] at (3.0,2.3) {\textcolor{mainblue}{蓝}：无催化剂时的活化分子\\\textcolor{warnred}{红}$+$\textcolor{mainblue}{蓝}：有催化剂时的活化分子};
\end{scope}
\begin{scope}[xshift=8.8cm]
\axisEP{5.2}{3.2}
\draw[thick] (0.3,1.4)--(0.9,1.4) .. controls (1.7,1.4) and (1.9,2.8) .. (2.5,2.8) .. controls (3.1,2.8) and (3.3,0.6) .. (4.1,0.6)--(4.8,0.6);
\draw[thick,dashed,warnred] (0.9,1.4) .. controls (1.7,1.4) and (1.9,2.1) .. (2.5,2.1) .. controls (3.1,2.1) and (3.3,0.6) .. (4.1,0.6);
\draw[<->] (1.0,1.4)--(1.0,2.8) node[midway,left]{\scriptsize $E_1$};
\draw[<->] (4.5,0.6)--(4.5,2.8) node[midway,right]{\scriptsize $E_2$};
\draw[dotted] (1.0,2.8)--(4.5,2.8);
\node[font=\scriptsize] at (2.5,3.05) {过渡态};
\node[font=\scriptsize,warnred] at (3.2,1.9) {催化剂};
\end{scope}
\end{tikzpicture}
\end{center}
右图中 $E_1$ 为正反应活化能，$E_2$ 为逆反应活化能，$\Delta H=E_1-E_2$。催化剂改变反应途径，使 $E_1$、$E_2$ 降低相同的数值，所以 $\Delta H$ 不变。

\subsection{用碰撞理论解释外因的影响}
反应速率由“单位体积内单位时间的有效碰撞次数”决定，而它正比于\textbf{单位体积内的活化分子数}$=$单位体积内分子总数 $\times$ 活化分子百分数。

\begin{center}
\small
\begin{tabular}{lcccc}
\toprule
条件改变 & 单位体积分子总数 & 活化分子百分数 & 单位体积活化分子数 & 速率 \\
\midrule
增大反应物浓度（或压缩气体体积） & 增多 & 不变 & 增多 & 增大 \\
升高温度 & 不变 & 增大 & 增多 & 增大 \\
使用催化剂 & 不变 & 增大 & 增多 & 增大 \\
\bottomrule
\end{tabular}
\end{center}

\begin{yicuo}
\begin{enumerate}
\item 浓度（压强）只增加活化分子\textbf{数目}，不改变活化分子\textbf{百分数}；温度、催化剂才改变百分数。
\item 升高温度既有分子能量升高（更多分子成为活化分子，主要原因），也有碰撞频率增大（次要原因）。
\item 升温对吸热、放热反应都加快速率，对正、逆反应都加快；只是加快的程度不同，平衡才会移动（第3章）。
\end{enumerate}
\end{yicuo}

\section{基元反应、反应历程与催化机理}

\subsection{基元反应与反应历程}
大多数化学反应并非“一步到位”，而是经过若干个简单步骤才完成。每一步都叫一个\textbf{基元反应}；这些基元反应的总和叫\textbf{反应历程}（反应机理）。例如 \ce{2HI = H2 + I2} 的历程是：
\[\text{①}\ \ce{2HI -> H2 + 2I*}\qquad \text{②}\ \ce{2I* -> I2}\]
\ce{I*} 是带单电子的原子，称为\textbf{自由基}；它在①中生成、在②中消耗，是\textbf{中间产物}。

\begin{gainian}[决速步]
多步反应的总速率由其中\textbf{最慢}的一步决定，这一步称为决速步（速控步）。在能量—历程图上，决速步是\textbf{自身能垒（该步过渡态能量 $-$ 该步起始物质能量）最大}的那一步。
\end{gainian}

\begin{tuozhan}[飞秒化学——“看见”过渡态]
过渡态存在的时间极短（约 $10^{-13}$ s 量级）。化学家利用飞秒（1 fs $=10^{-15}$ s）激光脉冲，像“高速摄影”一样观测反应过程中化学键的断裂与形成，直接“看到”了反应物经过渡态转化为产物的过程（泽维尔因此获 1999 年诺贝尔化学奖）。这为“基元反应”“过渡态”等概念提供了实验依据。
\end{tuozhan}

\begin{yicuo}
判断决速步看的是“\textbf{每一步自己的爬坡高度}”，而不是哪个峰的绝对位置最高。例如某步起点很低、峰也较高，其能垒可能最大；另一步峰虽最高，但起点也高，能垒未必最大。数据常以 eV（单个分子）给出，读图时逐步做减法。
\end{yicuo}

\begin{liti}[多步反应的能量图]
某反应历程的相对能量如下（单位 eV）：反应物 0 $\to$ 过渡态 TS1 $+0.82$ $\to$ 中间体 M1 $-0.35$ $\to$ TS2 $+0.60$ $\to$ M2 $-0.20$ $\to$ TS3 $+0.25$ $\to$ 生成物 $-0.90$。(1) 决速步是第几步？其能垒为多少？(2) 该总反应吸热还是放热？(3) 能量最低的中间体是哪个？
\tcblower
\jie (1) 三步的能垒：第一步 $0.82-0=0.82$ eV；第二步 $0.60-(-0.35)=0.95$ eV；第三步 $0.25-(-0.20)=0.45$ eV。决速步为第二步（\ce{M1 -> M2}），能垒 0.95 eV——虽然 TS1 的绝对能量最高，但它不是决速步。(2) 生成物能量 $-0.90$ eV 低于反应物，放热。(3) M1（$-0.35$ eV）。
\end{liti}

\subsection{催化剂的作用机理}
催化剂参与反应、改变反应历程，但反应前后其质量和化学性质不变。例如 \ce{Fe+} 催化 \ce{N2O} 与 CO 的反应：
\[\text{①}\ \ce{Fe+ + N2O -> FeO+ + N2}\qquad\text{②}\ \ce{FeO+ + CO -> Fe+ + CO2}\qquad\text{总：}\ \ce{N2O + CO ->[\ce{Fe+}] N2 + CO2}\]

\begin{fangfa}[区分催化剂与中间产物]
\textbf{催化剂}：先参与反应（出现在前面某步的反应物中），后又再生（出现在后面某步的生成物中），总方程式中不出现，如上例的 \ce{Fe+}。\textbf{中间产物}：先生成、后消耗，总方程式中也不出现，如上例的 \ce{FeO+}。在“循环图”中，从外部进入循环的是\textbf{反应物}，离开循环的是\textbf{生成物}，循环圈上的物种中，作为循环起点并最终回到原状的是\textbf{催化剂}，其余是中间体。
\end{fangfa}

\begin{liti}[读催化循环图]
某金属催化剂 M 催化 \ce{X + Y -> Z} 的循环机理为：①\ce{M + X -> M-X}；②\ce{M-X + Y -> M-XY}；③\ce{M-XY -> M + Z}。
\begin{center}
\begin{tikzpicture}[font=\small,scale=0.9]
\node (M) at (90:1.5) {M}; \node (MX) at (-30:1.5) {M—X}; \node (MXY) at (210:1.5) {M—XY};
\draw[-Stealth,thick,mainblue] (M) to[bend left=30] node[right=3pt]{\scriptsize ① $+$X} (MX);
\draw[-Stealth,thick,mainblue] (MX) to[bend left=30] node[below]{\scriptsize ② $+$Y} (MXY);
\draw[-Stealth,thick,mainblue] (MXY) to[bend left=30] node[left=3pt]{\scriptsize ③ $-$Z} (M);
\end{tikzpicture}
\end{center}
(1) 指出反应物、生成物、催化剂和中间体；(2) 若第②步最慢，增大 $c(\mathrm{X})$ 对总速率影响大还是增大 $c(\mathrm{Y})$ 影响大？
\tcblower
\jie (1) 从外部进入循环的 X、Y 是反应物，离开循环的 Z 是生成物；M 先消耗后再生，是催化剂；M—X、M—XY 先生成后消耗，是中间体。(2) 总速率由第②步决定，其速率与 $c(\ce{M-X})$、$c(\mathrm{Y})$ 有关；直接增大 $c(\mathrm{Y})$ 能加快决速步，影响更大（增大 $c(\mathrm{X})$ 只有在第①步建立的 M—X 浓度随之增大时才有作用）。
\end{liti}

\begin{benzhi}[催化剂“能做什么、不能做什么”]
\textbf{能}：降低活化能、加快正逆反应速率（同等倍数）、缩短达到平衡的时间；有\textbf{选择性}——在多个可能的反应中只加快某一个，从而提高目标产物的选择性。\textbf{不能}：改变 $\Delta H$、改变平衡常数和平衡转化率、使不能发生（非自发）的反应发生。此外，催化剂需要合适的温度（活性温度），温度过高可能失活；某些杂质会使催化剂“中毒”，所以工业原料气要净化。
\end{benzhi}

\begin{zhenti}{2024年广东卷第15题\quad 来源：仓库所附真题文档}
对反应 \ce{S(g) <=> T(g)}（I 为中间产物），相同条件下：①加入催化剂，反应达到平衡所需时间大幅缩短；②提高反应温度，$c_\text{平}(\mathrm{S})/c_\text{平}(\mathrm{T})$ 增大，$c_\text{平}(\mathrm{S})/c_\text{平}(\mathrm{I})$ 减小。基于以上事实，可能的反应历程示意图（实线为无催化剂，虚线为有催化剂）为
\begin{center}
\begin{tikzpicture}[xscale=0.42,yscale=0.42,font=\scriptsize]
\foreach \lab/\xs/\Ilev/\cat in {A/0/1.6/1, B/7.5/1.6/2, C/15/0.4/1, D/22.5/0.4/2}{
\begin{scope}[xshift=\xs cm]
\draw[-Stealth] (0,0)--(6.2,0); \draw[-Stealth] (0,0)--(0,5.2);
\node[left] at (-0.2,4.6) {\lab.};
\draw[thick] (0.3,1.1)--(0.6,1.1) .. controls (1.2,1.1) and (1.3,4.6) .. (1.9,4.6) .. controls (2.5,4.6) and (2.7,\Ilev) .. (3.1,\Ilev) .. controls (3.5,\Ilev) and (3.7,3.6) .. (4.1,3.6) .. controls (4.5,3.6) and (4.8,0.3) .. (5.3,0.3)--(5.7,0.3);
\ifnum\cat=1 \draw[dashed,thick] (0.6,1.1) .. controls (1.0,1.1) and (1.1,2.9) .. (1.5,2.9) .. controls (1.8,2.9) and (2.0,2.5) .. (2.2,2.5) .. controls (2.4,2.5) and (2.5,3.4) .. (2.7,3.4) .. controls (2.9,3.4) and (2.95,\Ilev) .. (3.1,\Ilev);
\else \draw[dashed,thick] (3.1,\Ilev) .. controls (3.4,\Ilev) and (3.5,2.6) .. (3.7,2.6) .. controls (3.85,2.6) and (3.9,2.2) .. (4.05,2.2) .. controls (4.2,2.2) and (4.3,2.7) .. (4.45,2.7) .. controls (4.7,2.7) and (4.85,0.3) .. (5.3,0.3);
\fi
\node[below] at (0.45,1.1) {S}; \node[below] at (3.1,\Ilev) {I}; \node[below] at (5.5,0.3) {T};
\end{scope}}
\end{tikzpicture}
\end{center}
\textbf{解析：}②中升温 $c(\mathrm{S})/c(\mathrm{T})$ 增大，说明 \ce{S <=> T} 平衡逆移，正反应放热，T 的能量低于 S；升温 $c(\mathrm{S})/c(\mathrm{I})$ 减小，说明 \ce{S <=> I} 正向吸热，I 的能量高于 S，排除 C、D（I 低于 S）。①中加入催化剂使达到平衡的时间\textbf{大幅}缩短，说明催化剂降低的是决速步（能垒大的第一步）的活化能，A 正确；B 中催化剂只降低了能垒较小的第二步，对总速率影响不大。答案：A。
\end{zhenti}

\begin{zhenti}{2026年广东卷第15题\quad 来源：仓库所附真题文档（网络收集版）}
物质 A(g) 可同时转化为它的同分异构体 B(g) 和 C(g)。图a（能量—反应过程）显示：A 生成 C 的活化能较小，C 的能量略低于 A，B 的能量最低；图b 为某温度下恒容密闭容器中各组分浓度随时间的变化：$t_1$ 时 B、C 浓度相等（均为 $x$），A 的浓度为 $y$；$t_2$ 时 A、B 浓度相等（均为 $z$）；C 的浓度先增大后减小。则
\xxb{$t_1$ 时，$m(\mathrm{A})=2m(\mathrm{B})$}{$2z>x+y$}{体系中，反应 \ce{A(g) <=> C(g)} 先达到平衡}{体系平衡后，升高温度，B 的体积分数增加}
\textbf{解析：}设 A 起始浓度为 $c_0$，三者互为同分异构体，由原子守恒，任意时刻 $c(\mathrm{A})+c(\mathrm{B})+c(\mathrm{C})=c_0$。$t_1$：$y+2x=c_0$，A 与 B 的质量关系无法确定，A 错。$t_2$：设 $c(\mathrm{C})=a$，$2z+a=c_0=2x+y$，得 $2z=x+y+(x-a)$；由图b可知 $t_2$ 时 C 的浓度小于 $t_1$ 时（C 先增后减，$t_2$ 时已降到 $x$ 以下），即 $a<x$，故 $2z>x+y$，B 正确。两个反应共用反应物 A，A 的浓度不变时两个平衡同时建立，不存在谁“先达到平衡”，C 错。B 的能量低于 A，\ce{A -> B} 放热，升温平衡逆移，B 的体积分数减小，D 错。答案：B。

\dianbo C 生成快（活化能小）但最终不如 B 多（B 更稳定）：这体现了\textbf{动力学（速率）控制与热力学（平衡）控制}的区别——短时间内产物以“生成快”的为主，足够长时间后以“更稳定”的为主。
\end{zhenti}

\section{速率方程与速率常数}

\begin{gainian}[速率方程]
一定温度下，反应 $a\mathrm{A}+b\mathrm{B}\to$ 产物 的瞬时速率常可表示为
\[v=k\cdot c^{\,m}(\mathrm{A})\cdot c^{\,n}(\mathrm{B})\]
$k$ 为\textbf{速率常数}，与温度、催化剂、固体表面性质等有关，\textbf{与浓度无关}；$m$、$n$ 由实验测定，对\textbf{基元反应}才等于计量数（即 $m=a$，$n=b$）。$k$ 越大，反应越快。
\end{gainian}

\textbf{确定速率方程的方法：}控制变量——保持其他物质浓度不变，某物质浓度变为原来的 2 倍时，若速率变为 2 倍，则其指数为 1；变为 4 倍，指数为 2；不变，指数为 0。

\begin{benzhi}[温度如何影响速率常数——阿伦尼乌斯公式]
$k=A\mathrm{e}^{-E_\text{a}/(RT)}$，取对数得 $\ln k=\ln A-\dfrac{E_\text{a}}{R}\cdot\dfrac1T$。以 $\ln k$（或 $R\ln k$）对 $\dfrac1T$ 作图得直线，\textbf{斜率 $=-\dfrac{E_\text{a}}{R}$}（或 $-E_\text{a}$）。由此可知：①升高温度，$k$ 增大；②$E_\text{a}$ 越大，直线越陡，速率对温度越敏感；③催化剂降低 $E_\text{a}$，直线变平缓。高考中常给出 $R\ln k$–$\frac1T$ 图，要求由斜率求活化能或比较催化剂优劣，只需会读图，不需记忆公式。
\end{benzhi}

\begin{zhenti}{2026年广东卷第19题(3)\quad 来源：仓库所附真题文档（网络收集版）}
\ce{H2S} 的氧化可在卟啉铁配合物 [\ce{FeL(H2O)}]（用 Z 表示）催化下进行，其中一步为 \ce{[FeL(H2O)] + HS- <=> [FeL(HS)]- + H2O}。为研究不同因素对该反应的影响，进行 4 组实验：Ⅰ～Ⅲ 组仅初始 $c_0(\ce{HS-})$ 不同（$c_\text{Ⅰ}<c_\text{Ⅱ}<c_\text{Ⅲ}$）；Ⅳ 与 Ⅱ 组仅温度不同。$0\sim t$ 时间内 Z 的平均反应速率 $-\frac{\Delta c(\mathrm{Z})}{\Delta t}$（单位 $10^{-6}$ mol·L$^{-1}$·s$^{-1}$）：Ⅰ 为 1.0，Ⅱ 为 2.0，Ⅲ 为 $y$，Ⅳ 为 1.0。\\
①下列说法正确的有\tk。A．$0\sim t$ 内反应速率：Ⅰ$>$Ⅱ$>$Ⅲ\quad B．Z 的平衡转化率：Ⅰ$<$Ⅲ\quad C．实验温度：Ⅱ$>$Ⅳ\quad D．平衡时 Z 的浓度：Ⅱ 和 Ⅳ 相等\\
②Ⅰ、Ⅱ 和 Ⅲ 的数据点分布在一条直线上，且 $(c_\text{Ⅲ}-c_\text{Ⅱ})/(c_\text{Ⅱ}-c_\text{Ⅰ})=1.2$，则 $y=$\tk。

\textbf{解析：}①随 $c_0(\ce{HS-})$ 增大，Z 的反应速率增大，A 错；另一反应物 \ce{HS-} 浓度越大，Z 的平衡转化率越大，B 对；Ⅱ、Ⅳ 浓度相同而 Ⅱ 速率大，温度更高，C 对；温度不同，平衡常数不同，平衡时 Z 的浓度不相等，D 错。选 BC。②三点共线，斜率相同：$\dfrac{y-2.0}{c_\text{Ⅲ}-c_\text{Ⅱ}}=\dfrac{2.0-1.0}{c_\text{Ⅱ}-c_\text{Ⅰ}}$，$y-2.0=1.2$，$y=3.2$。
\end{zhenti}

\textbf{速率常数与平衡常数：}对基元可逆反应 \ce{$a$A + $b$B <=> $c$C + $d$D}，$\vz=\kz c^a(\mathrm{A})c^b(\mathrm{B})$，$\vn=\kn c^c(\mathrm{C})c^d(\mathrm{D})$。平衡时 $\vz=\vn$，得 $K=\dfrac{\kz}{\kn}$。升高温度，$\kz$、$\kn$ 都增大；若正反应放热，则 $\kn$ 增大得更多，$K$ 减小（详见第3章）。

% ------------------------------------------------------------
\zhangmo{本章小结}
\subsection*{一、知识总结}
\begin{center}
\small
\begin{tabularx}{\linewidth}{p{2.4cm}X}
\toprule
速率表示 & $v=\Delta c/\Delta t$（平均速率）；速率之比$=$计量数之比；固体、纯液体不用浓度表示 \\
影响因素 & 内因决定；浓度、压强（通过浓度）、温度、催化剂、表面积等外因 \\
碰撞理论 & 浓度改变活化分子数；温度、催化剂改变活化分子百分数 \\
反应历程 & 基元反应；决速步$=$自身能垒最大的一步；催化剂先消耗后再生，中间体先生成后消耗 \\
速率方程 & $v=kc^m(\mathrm{A})c^n(\mathrm{B})$，指数实验测定；$k$ 只与温度、催化剂等有关；基元反应 $K=\kz/\kn$ \\
\bottomrule
\end{tabularx}
\end{center}

\subsection*{二、方法提升}
\begin{enumerate}
\item \textbf{速率计算}：一律先列三段式，只用“转化量”与计量数成比例。
\item \textbf{惰性气体问题}：先判断容器是恒容还是恒压，再看反应物浓度是否改变。
\item \textbf{历程图三问}：吸放热看起终点；决速步看每步能垒；催化剂、中间体看“先进后出”还是“先出后进”。
\item \textbf{数据分析题}：用控制变量找“只有一个条件不同”的两组实验对比；涉及多组数据时，注意线性关系、倍数关系。
\end{enumerate}

\subsection*{三、易错点归纳}
\begin{enumerate}
\item 误用固体或纯液体的浓度变化计算速率；比较速率时忘记统一单位。
\item 认为“增大压强速率一定增大”，忽略恒容充入惰性气体的情况。
\item 认为升温只加快吸热反应；认为催化剂只加快正反应。
\item 把能量最高的过渡态所在步骤当作决速步；把中间体当作催化剂。
\item 认为速率方程中的指数一定等于方程式计量数（只有基元反应才是）。
\end{enumerate}

% ------------------------------------------------------------
\zhangmo{分层练习}
\zu{A组\quad 基础巩固}
\begin{xiti}[start=1]
\item 下列关于化学反应速率的说法错误的是\xz
\xxa{化学反应速率是衡量化学反应进行快慢的物理量}{决定化学反应速率的主要因素是反应物本身的性质}{某反应在一段时间内的速率为 0.5 mol·L$^{-1}$·s$^{-1}$，是指该段时间内反应物和生成物的浓度变化都为 0.5 mol·L$^{-1}$}{可以用单位时间内固体质量的变化来表示有固体参加的反应的快慢}

\item 合成氨 \ce{N2 + 3H2 <=> 2NH3} 在不同条件下测得：①$v(\ce{N2})=0.6$ mol·L$^{-1}$·min$^{-1}$；②$v(\ce{H2})=1.2$ mol·L$^{-1}$·min$^{-1}$；③$v(\ce{NH3})=0.03$ mol·L$^{-1}$·s$^{-1}$。反应速率由大到小的顺序为\tk。

\item 一定温度下，在 2 L 密闭容器中充入 1 mol \ce{H2(g)} 和 1 mol \ce{I2(g)}，发生 \ce{H2(g) + I2(g) <=> 2HI(g)} $\Delta H<0$。下列说法正确的是\xz
\xxa{保持容积不变，再充入 1 mol \ce{H2(g)}，反应速率加快}{保持容积不变，充入 1 mol \ce{N2}（不参与反应），反应速率加快}{保持压强不变，充入 1 mol \ce{N2}（不参与反应），反应速率不变}{保持压强不变，再充入 1 mol \ce{H2(g)} 和 1 mol \ce{I2(g)}，反应速率加快}

\item 下列有关有效碰撞理论的说法正确的是\xz
\xxa{增大反应物浓度，可增大活化分子百分数}{升高温度，活化分子百分数增大，有效碰撞频率增大}{催化剂能降低反应的活化能，因此能使 $\Delta H$ 减小}{活化分子之间的碰撞一定是有效碰撞}

\item 下列事实能说明影响化学反应速率的决定性因素是反应物本身性质的是\xz
\xxa{大小、形状相同的镁片、铁片分别与同浓度盐酸反应，前者产生 \ce{H2} 更快}{Cu 与浓硝酸反应比与稀硝酸反应快}{\ce{N2} 与 \ce{O2} 在常温常压下不反应，放电时可反应}{Cu 能与浓硫酸反应，而不与稀硫酸反应}

\item 下列关于催化剂的说法正确的是\xz
\xxa{催化剂不参加化学反应，反应前后质量不变}{催化剂可以提高反应物的平衡转化率}{同一可逆反应中，催化剂对正、逆反应速率的改变程度相同}{任何温度下催化剂的催化效果都相同}

\item 已知温度每升高 10 ℃，某反应速率增大为原来的 2 倍。若将温度由 20 ℃ 升高到 50 ℃，反应速率增大为原来的\tk{}倍；若原来反应需 16 min 完成，50 ℃ 时约需\tk min。
\end{xiti}

\zu{B组\quad 高考典型}
\begin{xiti}
\item \laiyuan{原创}\ 用 0.10 mol·L$^{-1}$ \ce{Na2S2O3} 溶液与 0.10 mol·L$^{-1}$ \ce{H2SO4} 溶液探究影响速率的因素，记录出现浑浊所需时间：
\begin{center}\small
\begin{tabular}{cccccc}
\toprule
实验 & 温度/℃ & $V(\ce{Na2S2O3})$/mL & $V(\ce{H2SO4})$/mL & $V(\ce{H2O})$/mL & 出现浑浊时间/s \\
\midrule
① & 20 & 10.0 & 10.0 & 0 & $t_1$ \\
② & 20 & 5.0 & 10.0 & $a$ & $t_2$ \\
③ & 40 & 10.0 & 10.0 & $b$ & $t_3$ \\
\bottomrule
\end{tabular}
\end{center}
(1) 写出反应的离子方程式。(2) $a=$\tk，$b=$\tk。(3) 实验①②探究的是\tk 对速率的影响，预测 $t_1$\tk$t_2$（填“$>$”“$<$”或“$=$”）。(4) 实验中能否用“生成 \ce{SO2} 气体的体积”来衡量速率？说明理由。

\item \laiyuan{原创}\ \ce{Fe+} 可催化 \ce{N2O(g) + CO(g) = N2(g) + CO2(g)}，反应分两步：①\ce{Fe+ + N2O -> FeO+ + N2}；②\ce{FeO+ + CO -> Fe+ + CO2}。其能量变化示意图中，第①步活化能为 $E_1$，第②步活化能为 $E_2$，且 $E_1>E_2$，两步均放热。下列说法错误的是\xz
\xxa{\ce{Fe+} 是催化剂，\ce{FeO+} 是中间产物}{总反应速率主要由第①步决定}{使用 \ce{Fe+} 后总反应的 $\Delta H$ 减小}{第①步反应的速率比第②步慢}

\item 已知 \ce{2NO(g) + 2H2(g) = N2(g) + 2H2O(g)}，800 ℃ 时测得如下数据：
\begin{center}\small
\begin{tabular}{cccc}
\toprule
实验 & $c(\ce{NO})$/(mol·L$^{-1}$) & $c(\ce{H2})$/(mol·L$^{-1}$) & 初始速率/(mol·L$^{-1}$·s$^{-1}$) \\
\midrule
1 & 0.10 & 0.10 & $2.0\times10^{-3}$ \\
2 & 0.20 & 0.10 & $8.0\times10^{-3}$ \\
3 & 0.10 & 0.20 & $4.0\times10^{-3}$ \\
\bottomrule
\end{tabular}
\end{center}
(1) 写出速率方程，并求速率常数 $k$（注明单位）。(2) 求 $c(\ce{NO})=0.30$ mol·L$^{-1}$、$c(\ce{H2})=0.05$ mol·L$^{-1}$ 时的速率。(3) 有人提出该反应的历程：Ⅰ．\ce{2NO <=> N2O2}（快）；Ⅱ．\ce{N2O2 + H2 -> N2O + H2O}（慢）；Ⅲ．\ce{N2O + H2 -> N2 + H2O}（快）。指出决速步和中间产物。

\item \laiyuan{原创}\ 某反应在两种催化剂 a、b 作用下，$\ln k$ 与 $\frac1T$ 的关系为直线：催化剂 a 的直线过点 $(2.0\times10^{-3},\,2.0)$ 和 $(3.0\times10^{-3},\,-4.0)$；催化剂 b 的直线过点 $(2.0\times10^{-3},\,3.0)$ 和 $(3.0\times10^{-3},\,0)$（$\frac1T$ 单位 K$^{-1}$）。已知 $\ln k=\ln A-\frac{E_\text{a}}{R}\cdot\frac1T$，$R=8.314$ J·mol$^{-1}$·K$^{-1}$。(1) 分别求两种催化剂下的活化能（kJ·mol$^{-1}$，保留三位有效数字）。(2) 哪种催化剂效果更好？哪种情况下反应速率对温度更敏感？

\item \laiyuan{原创}\ 用酸性 \ce{KMnO4} 溶液与 \ce{H2C2O4} 溶液反应探究影响速率的因素（\ce{KMnO4} 均少量，溶液总体积相同）：
\begin{center}\small
\begin{tabular}{ccccc}
\toprule
实验 & 温度/℃ & 加入的其他物质 & $c(\ce{H2C2O4})$ & 褪色时间/s \\
\midrule
① & 25 & 无 & 0.10 mol·L$^{-1}$ & 40 \\
② & 25 & 少量 \ce{MnSO4} 固体 & 0.10 mol·L$^{-1}$ & 4 \\
③ & 25 & 无 & 0.20 mol·L$^{-1}$ & 25 \\
④ & 50 & 无 & 0.10 mol·L$^{-1}$ & 10 \\
\bottomrule
\end{tabular}
\end{center}
(1) 写出反应的离子方程式。(2) 对比实验\tk{}可探究浓度的影响，对比实验\tk{}可探究温度的影响。(3) 实验②说明了什么？据此解释实验①中“开始褪色很慢，一段时间后突然加快”的现象。

\item \laiyuan{原创}\ 某反应在催化剂 A、B 作用下的历程如下（相对能量，eV）：催化剂 A：反应物 0 $\to$ TS1 $+1.10$ $\to$ 中间体 $+0.20$ $\to$ TS2 $+0.85$ $\to$ 产物 $-0.50$；催化剂 B：反应物 0 $\to$ TS1$'$ $+0.70$ $\to$ 中间体 $-0.40$ $\to$ TS2$'$ $+0.60$ $\to$ 产物 $-0.50$。(1) 分别指出两种催化剂下的决速步及其能垒。(2) 哪种催化剂效果更好？(3) 两种催化剂下反应的 $\Delta H$ 是否相同？
\end{xiti}

\zu{C组\quad 冲刺高分}
\begin{xiti}
\item \laiyuan{原创}\ 某温度下，向恒容密闭容器中充入 A(g)，同时发生两个反应：\ce{A(g) <=> B(g)}（反应Ⅰ，活化能大，$\Delta H_1<0$）和 \ce{A(g) <=> C(g)}（反应Ⅱ，活化能小，$\Delta H_2<0$，且 $|\Delta H_2|<|\Delta H_1|$）。
(1) 反应开始一段时间内，B、C 中哪一种产物的浓度增大得快？为什么？
(2) 足够长时间后，平衡体系中 B、C 哪一种含量高？为什么？
(3) 若要在较短时间内获得尽可能多的 C，应如何控制反应时间？若要获得 B，应如何操作？
(4) 若选用一种只加快反应Ⅰ的催化剂，平衡时 B 的体积分数是否改变？说明理由。

\item \laiyuan{原创}\ 在恒容密闭容器中发生基元反应 \ce{2NO2(g) <=> N2O4(g)} $\Delta H<0$，$\vz=\kz c^2(\ce{NO2})$，$\vn=\kn c(\ce{N2O4})$。$T_1$ 时，$\kz=2\kn$（浓度单位 mol·L$^{-1}$）。
(1) $T_1$ 时该反应的平衡常数 $K=$\tk。
(2) $T_1$ 时，若某时刻 $c(\ce{NO2})=0.5$ mol·L$^{-1}$，$c(\ce{N2O4})=1.0$ mol·L$^{-1}$，判断此时 $\vz$ 与 $\vn$ 的大小。
(3) 升高温度至 $T_2$，$\kz$、$\kn$ 如何变化？$\kz/\kn$ 增大还是减小？

\item \laiyuan{原创}\ 反应 \ce{2NO(g) + O2(g) = 2NO2(g)} 的速率方程为 $v=kc^2(\ce{NO})c(\ce{O2})$。有人提出两步机理：Ⅰ．\ce{2NO <=> N2O2}（快速平衡，$\Delta H<0$）；Ⅱ．\ce{N2O2 + O2 -> 2NO2}（慢）。(1) 指出决速步和中间产物。(2) 用第Ⅰ步的平衡常数 $K_1$ 和第Ⅱ步的速率常数 $k_2$ 表示总反应的 $k$。(3) 实验发现：升高温度，该反应的速率反而略有减小。结合机理解释。
\end{xiti}

% ------------------------------------------------------------
\zhangmo{练习解析}
\begin{jiexilist}
\item \ans{C}速率是单位时间内浓度的变化量，单位是 mol·L$^{-1}$·s$^{-1}$；而且用不同物质表示时数值不一定相同，C 错。D 对：可用单位时间内固体质量（或物质的量）的变化表示快慢，只是不能用其“浓度”表示。
\item \ans{③$>$①$>$②}都换算为 $v(\ce{N2})$（mol·L$^{-1}$·min$^{-1}$）：①0.6；②$1.2/3=0.4$；③$0.03\times60/2=0.9$。注意③的单位是“s”，必须先乘 60 换算为“min”。
\item \ans{A}A：反应物 \ce{H2} 浓度增大，速率加快。B：恒容充 \ce{N2}，各物质浓度不变，速率不变。C：恒压充 \ce{N2}，体积增大，浓度减小，速率减小。D：恒压下按原比例再充入 \ce{H2}、\ce{I2}，体积同比例增大，各物质浓度与原来相同，速率不变。
\item \ans{B}A 错：浓度增大的是单位体积活化分子数，百分数不变。C 错：催化剂不改变 $\Delta H$。D 错：活化分子还需取向合适才能发生有效碰撞。
\item \ans{A}A 中两种金属本身性质不同，其他条件相同，体现内因。B 是浓度的影响；C 是放电（外界能量）的影响；D 是浓度不同导致性质不同（浓硫酸有强氧化性），不属于“不同反应物本身性质”的比较。
\item \ans{C}A 错：催化剂参与反应，只是反应前后质量和化学性质不变。B 错：催化剂不能使平衡移动。D 错：催化剂有活性温度范围。
\item 8；2。升高 30 ℃，即 3 个 10 ℃，速率变为 $2^3=8$ 倍，时间变为原来的 $\frac18$。
\item (1) \ce{S2O3^2- + 2H+ = S v + SO2 ^ + H2O}。(2) $a=5.0$，$b=0$：①总体积 20.0 mL，②中补水 5.0 mL 使总体积相同；③只改变温度，不加水。(3) \ce{Na2S2O3} 浓度；$t_1<t_2$（①浓度大，速率快，出现浑浊所需时间短）。(4) 不宜。\ce{SO2} 易溶于水，生成的 \ce{SO2} 有相当一部分溶于溶液，测得的气体体积不能准确反映反应进程。
\item \ans{C}\ce{Fe+} 先消耗后再生，是催化剂；\ce{FeO+} 先生成后消耗，是中间产物，A 对。$E_1>E_2$，第①步慢，决定总速率，B、D 对。催化剂不改变 $\Delta H$，C 错。
\item (1) 由实验1、2：$c(\ce{NO})$ 加倍，速率变为 4 倍，NO 的指数为 2；由实验1、3：$c(\ce{H2})$ 加倍，速率加倍，\ce{H2} 的指数为 1。$v=kc^2(\ce{NO})c(\ce{H2})$；$k=\dfrac{2.0\times10^{-3}}{0.10^2\times0.10}=2.0$ L$^2$·mol$^{-2}$·s$^{-1}$。(2) $v=2.0\times0.30^2\times0.05=9.0\times10^{-3}$ mol·L$^{-1}$·s$^{-1}$。(3) 决速步为Ⅱ（慢反应）；中间产物为 \ce{N2O2} 和 \ce{N2O}。（说明：第Ⅰ步为快速平衡，\ce{N2O2} 浓度正比于 $c^2(\ce{NO})$，所以慢步Ⅱ的速率 $\propto c^2(\ce{NO})c(\ce{H2})$，与实验测得的速率方程一致，这正是用实验检验机理的思路。）
\item (1) 斜率 $=-\dfrac{E_\text{a}}{R}$。a：斜率 $=\dfrac{-4.0-2.0}{1.0\times10^{-3}}=-6.0\times10^3$ K，$E_\text{a}=6.0\times10^3\times8.314$ J·mol$^{-1}\approx49.9$ kJ·mol$^{-1}$；b：斜率 $=\dfrac{0-3.0}{1.0\times10^{-3}}=-3.0\times10^3$ K，$E_\text{a}\approx24.9$ kJ·mol$^{-1}$。(2) b 的活化能小，且在所给温度范围内 $\ln k$ 都更大，催化效果更好；a 的直线更陡（$E_\text{a}$ 更大），使用 a 时速率对温度更敏感。
\item (1) \ce{2MnO4^- + 5H2C2O4 + 6H+ = 2Mn^2+ + 10CO2 ^ + 8H2O}。(2) ①③；①④。(3) \ce{Mn^2+} 对该反应有催化作用；实验①开始时 \ce{Mn^2+} 很少，反应很慢，随着反应生成 \ce{Mn^2+}（自催化），速率突然加快。
\item (1) A：第一步能垒 1.10 eV，第二步 $0.85-0.20=0.65$ eV，决速步为第一步（1.10 eV）；B：第一步 0.70 eV，第二步 $0.60-(-0.40)=1.00$ eV，决速步为第二步（1.00 eV）。(2) B 的最大能垒（1.00 eV）小于 A（1.10 eV），B 效果更好。(3) 相同，都为 $-0.50$ eV（每个分子），$\Delta H$ 与催化剂无关。注意：催化剂 B 下 TS2$'$（$+0.60$ eV）的绝对能量虽低于 TS1$'$（$+0.70$ eV），但决速步却是第二步——必须逐步计算能垒。
\item (1) C。反应Ⅱ活化能小，速率快，开始阶段 C 生成得快。(2) B。两反应都放热，$|\Delta H_1|>|\Delta H_2|$，B 的能量最低、最稳定，平衡（足够长时间后）B 的含量高；此时生成的 C 会经 A 逐步转化为 B。(3) 获得 C：控制较短的反应时间，在 C 浓度达到最大值附近时及时终止反应（分离产物）；获得 B：延长反应时间使体系接近平衡（也可选用能加快反应Ⅰ的催化剂）。(4) 不改变。催化剂不改变平衡常数，平衡状态只由温度和起始组成决定，B 的平衡体积分数不变；催化剂只缩短 B 达到平衡含量所需的时间。
\item (1) 2。平衡时 $\kz c^2(\ce{NO2})=\kn c(\ce{N2O4})$，$K=\dfrac{c(\ce{N2O4})}{c^2(\ce{NO2})}=\dfrac{\kz}{\kn}=2$。(2) $Q=\dfrac{1.0}{0.5^2}=4>K=2$，反应逆向进行，$\vz<\vn$（也可直接算：$\vz=2\kn\times0.25=0.5\kn$，$\vn=\kn\times1.0$，$\vz<\vn$）。(3) 升温，$\kz$、$\kn$ 都增大；正反应放热，升温平衡逆移，$K$ 减小，故 $\kz/\kn$ 减小（$\kn$ 增大的倍数更大）。
\item (1) 第Ⅱ步为决速步；中间产物为 \ce{N2O2}。(2) 第Ⅰ步快速平衡：$c(\ce{N2O2})=K_1c^2(\ce{NO})$；总速率由慢步决定：$v=k_2c(\ce{N2O2})c(\ce{O2})=k_2K_1c^2(\ce{NO})c(\ce{O2})$，所以 $k=k_2K_1$。(3) 升高温度，$k_2$ 增大；但第Ⅰ步放热，升温使 $K_1$ 减小，$c(\ce{N2O2})$ 降低。若 $K_1$ 减小的程度超过 $k_2$ 增大的程度，$k=k_2K_1$ 减小，总反应速率反而减小。这说明“升温速率一定增大”只对基元反应成立，对复杂反应要具体分析。
\end{jiexilist}
% <<<<<<<< 结束：chapters/ch02.tex

% >>>>>>>> 开始：chapters/ch03.tex
% ============================================================
% 第3章 化学平衡（对应课件：选必一2.2 化学平衡、2.3 化学反应方向、2.4 化学反应调控）
% ============================================================
\chapter{化学平衡}

\begin{daoxue}
\textbf{知识主线：}可逆反应能否进行到底（平衡状态及其判断）$\to$ 进行到什么程度（平衡常数 $K$、$K_p$、转化率）$\to$ 条件改变后怎样变化（平衡移动、勒夏特列原理、图像）$\to$ 能不能自发发生（$\Delta G=\Delta H-T\Delta S$）$\to$ 工业上怎样选择条件（化学反应的调控）$\to$ 多个反应同时发生怎么办（多重平衡与选择性）。

\textbf{高考考情：}本章是广东卷第19题的“主战场”：2025年考 $\Delta G$ 判断、平衡组成图、$K$ 的计算和催化剂选择性；2026年考吸收平衡的 $K$ 计算；2024年第17题考“如何证明平衡已建立”。选择题常考平衡标志、平衡移动和图像。分值约 10～14 分，是冲击 90 分必须攻克的板块。

\textbf{学习目标：}①会判断平衡状态；②熟练书写 $K$、$K_p$ 并用三段式计算；③能用 $Q$ 与 $K$ 的关系解释平衡移动；④掌握图像分析的通用方法；⑤会用 $\Delta G$ 判断自发性；⑥能处理多重平衡中的转化率、选择性和 $K_p$ 计算。
\end{daoxue}

\section{化学平衡状态}

\subsection{可逆反应与平衡的建立}
在\textbf{同一条件}下，既能向正反应方向进行、又能向逆反应方向进行的反应叫可逆反应，用“\ce{<=>}”表示。可逆反应不能进行到底，任一时刻反应物与生成物共存。注意 \ce{2H2 + O2 ->[\text{点燃}] 2H2O} 与 \ce{2H2O ->[\text{电解}] 2H2 ^ + O2 ^} 条件不同，不是可逆反应。

以 \ce{N2 + 3H2 <=> 2NH3} 为例：开始时反应物浓度最大，$\vz$ 最大，$\vn=0$；随着反应进行，$\vz$ 减小、$\vn$ 增大，直到 $\vz=\vn$，各物质浓度不再改变——这就是\textbf{化学平衡状态}。

\begin{center}
\begin{tikzpicture}[scale=0.75,font=\small]
\begin{scope}
\draw[-Stealth] (0,0)--(5,0) node[below left]{$t$}; \draw[-Stealth] (0,0)--(0,3.2) node[left]{$c$};
\draw[thick,mainblue] (0,2.8) .. controls (1,1.8) and (1.8,1.6) .. (2.6,1.6)--(4.8,1.6) node[right]{\ce{H2}};
\draw[thick,deepgreen] (0,1.0) .. controls (1,0.7) and (1.8,0.62) .. (2.6,0.62)--(4.8,0.62) node[right]{\ce{N2}};
\draw[thick,warnred] (0,0) .. controls (1,0.6) and (1.8,0.75) .. (2.6,0.76)--(4.8,0.76) node[above right]{\ce{NH3}};
\draw[dashed] (2.6,0)--(2.6,3) node[above]{$t_1$};
\end{scope}
\begin{scope}[xshift=7cm]
\draw[-Stealth] (0,0)--(5,0) node[below left]{$t$}; \draw[-Stealth] (0,0)--(0,3.2) node[left]{$v$};
\draw[thick,mainblue] (0.02,2.8) .. controls (1,1.7) and (1.8,1.4) .. (2.6,1.4)--(4.8,1.4);
\draw[thick,warnred] (0.02,0.02) .. controls (1,1.0) and (1.8,1.38) .. (2.6,1.4);
\node[mainblue] at (0.9,2.4) {$\vz$}; \node[warnred] at (1.1,0.5) {$\vn$};
\node at (3.8,1.75) {$\vz=\vn\neq0$};
\draw[dashed] (2.6,0)--(2.6,3) node[above]{$t_1$};
\end{scope}
\end{tikzpicture}
\end{center}

\begin{gainian}[化学平衡状态的特征：逆、等、动、定、变]
\textbf{逆}：研究对象是可逆反应；\textbf{等}：$\vz=\vn\neq0$（本质）；\textbf{动}：动态平衡，反应没有停止；\textbf{定}：各组分的浓度（含量）保持不变（现象）；\textbf{变}：条件改变，平衡会被破坏而移动。化学平衡是该条件下反应的\textbf{最大限度}，决定了反应物的最大（平衡）转化率。
\end{gainian}

\subsection{平衡状态的判断}
判断依据只有两条：一是“正逆相等”，二是“变量不变”。

\textbf{(1) 正逆相等：}同一物质的生成速率等于消耗速率；用不同物质表示时，必须\textbf{一个表示正反应、一个表示逆反应}，且速率之比等于计量数之比。以 \ce{N2 + 3H2 <=> 2NH3} 为例：“$2\vz(\ce{N2})=\vn(\ce{NH3})$”是平衡标志；“$3v(\ce{N2})=v(\ce{H2})$”任何时刻都成立，不能作标志；“断裂 3 mol H—H 键的同时断裂 6 mol N—H 键”表示正逆速率相等，是平衡标志；“断裂 3 mol H—H 键的同时形成 6 mol N—H 键”都表示正反应，不是标志。

\textbf{(2) 变量不变：}一个量若随反应进行而变化（变量），则它不变时反应达到平衡；若它本来就不随反应变化（定量），则不能作为标志。

\begin{center}
\small
\begin{tabularx}{\linewidth}{p{3.4cm}XXX}
\toprule
物理量（不变时） & \ce{N2(g) + 3H2(g) <=> 2NH3(g)}（气体分子数变化） & \ce{H2(g) + I2(g) <=> 2HI(g)}（气体分子数不变） & \ce{C(s) + H2O(g) <=> CO(g) + H2(g)}（有固体） \\
\midrule
恒温恒容：总压强 $p$ & 是 & 否（$n$ 总不变） & 是 \\
恒温恒压：总体积 $V$ & 是 & 否 & 是 \\
气体总物质的量 $n$ & 是 & 否 & 是 \\
恒容：气体密度 $\rho=m_\text{气}/V$ & 否（$m$、$V$ 都不变） & 否 & 是（气体质量变化） \\
平均摩尔质量 $\bar M=m_\text{气}/n_\text{气}$ & 是 & 否 & 是 \\
有色物质的颜色 & — & 是（\ce{I2} 浓度变化） & — \\
\bottomrule
\end{tabularx}
\end{center}

\begin{yicuo}
\begin{enumerate}
\item “各物质浓度之比等于计量数之比”“各物质浓度相等”都不一定是平衡状态——平衡时浓度\textbf{不变}，而不是\textbf{相等或成比例}。
\item 反应物为固体、生成的气体组成比例固定的分解反应要特别小心：如 \ce{NH4Cl(s) <=> NH3(g) + HCl(g)}，两种气体始终按 1:1 生成，$\bar M$ 和 \ce{NH3} 的体积分数始终不变，不能作标志；但恒容时压强、气体密度是变量，可以作标志。
\item \textbf{绝热}容器中，反应热使体系温度变化，“温度不变”可作为平衡标志；同时平衡常数也会随温度变化。
\end{enumerate}
\end{yicuo}

\begin{liti}[平衡标志综合判断]
可逆反应 \ce{2NO2(g) <=> 2NO(g) + O2(g)} 在恒容密闭容器中进行：①单位时间内生成 $n$ mol \ce{O2} 的同时生成 $2n$ mol \ce{NO2}；②单位时间内生成 $n$ mol \ce{O2} 的同时生成 $2n$ mol NO；③用 \ce{NO2}、NO、\ce{O2} 表示的速率之比为 2:2:1；④混合气体颜色不再改变；⑤混合气体密度不再改变；⑥混合气体压强不再改变；⑦混合气体平均相对分子质量不再改变。能说明达到平衡状态的是\tk。
\tcblower
\jie ①生成 \ce{O2} 为正反应、生成 \ce{NO2} 为逆反应，且比例 1:2 正确，是。②都是正反应，否。③任何时刻都成立，否。④\ce{NO2} 红棕色，颜色不变即浓度不变，是。⑤恒容且全为气体，密度始终不变，否。⑥气体分子数变化，压强是变量，是。⑦$m$ 不变、$n$ 变化，$\bar M$ 是变量，是。答案：①④⑥⑦。
\end{liti}

\begin{liti}[恒温恒压下的体积计算]
某体积可变的密闭容器中盛有适量 A 和 B 的混合气体，在一定条件下发生 \ce{A(g) + 3B(g) <=> 2C(g)}。维持温度和压强不变，达到平衡时容器体积为 $V$ L，其中 C 的体积分数为 10\%。则原混合气体的体积为\tk L，平衡时 A 消耗了\tk L。
\tcblower
\silu 恒温恒压下，气体体积之比等于物质的量之比，可直接用“体积”列三段式。\\
\jie 平衡时 $V(\mathrm{C})=0.1V$。每生成 2 体积 C，气体总体积减少 $1+3-2=2$ 体积，所以减少 $0.1V$；原混合气体为 $V+0.1V=1.1V$ L。消耗 A 为 $0.1V\times\frac12=0.05V$ L，消耗 B 为 $0.15V$ L。\\
\dianbo “差量法”：总体积的变化量与某物质变化量成比例，可省去设未知数。
\end{liti}

\section{化学平衡常数}

\subsection{平衡常数的表达式}
800 ℃ 时，对 \ce{CO(g) + H2O(g) <=> CO2(g) + H2(g)}，无论起始投料如何，平衡时 $\dfrac{c(\ce{CO2})\cdot c(\ce{H2})}{c(\ce{CO})\cdot c(\ce{H2O})}$ 都约等于 1.0。大量实验表明这一规律具有普遍性。

\begin{gainian}[化学平衡常数]
一定温度下，可逆反应 \ce{$m$A(g) + $n$B(g) <=> $p$C(g) + $q$D(g)} 达到平衡时，生成物浓度幂之积与反应物浓度幂之积的比值是一个常数：
\[K=\frac{c^{\,p}(\mathrm{C})\cdot c^{\,q}(\mathrm{D})}{c^{\,m}(\mathrm{A})\cdot c^{\,n}(\mathrm{B})}\qquad(\text{各浓度均为平衡浓度})\]
对指定反应，$K$ \textbf{只与温度有关}，与浓度、压强、催化剂无关。
\end{gainian}

\textbf{书写要点：}
\begin{enumerate}
\item \textbf{固体、纯液体不写入表达式}（其“浓度”视为常数）。如 \ce{CaCO3(s) <=> CaO(s) + CO2(g)}，$K=c(\ce{CO2})$。
\item 稀水溶液中的反应，水不写入；\textbf{非水体系中有水参与}时（如酯化反应 \ce{CH3COOH + C2H5OH <=> CH3COOC2H5 + H2O}），水要写入。
\item 离子反应先写成离子方程式再写 $K$。如 \ce{Cl2 + H2O <=> H+ + Cl- + HClO}，$K=\dfrac{c(\ce{H+})c(\ce{Cl-})c(\ce{HClO})}{c(\ce{Cl2})}$。
\item $K$ 与方程式的写法一一对应：
\end{enumerate}
\begin{center}
\small
\begin{tabular}{ll}
\toprule
方程式的变化 & $K$ 的变化 \\
\midrule
计量数扩大 $n$ 倍 & $K'=K^n$ \\
正、逆反应 & $K_\text{逆}=1/K_\text{正}$ \\
两式相加（相减） & $K=K_1\cdot K_2$（$K=K_1/K_2$） \\
\bottomrule
\end{tabular}
\end{center}
这与盖斯定律“$\Delta H$ 相加”正好对应：\textbf{方程式相加，$\Delta H$ 相加，$K$ 相乘}。

\subsection{平衡常数的意义与应用}
\textbf{①衡量反应进行的程度：}$K$ 越大，正反应进行得越完全，反应物平衡转化率越大。一般 $K>10^5$ 认为反应基本完全。如 25 ℃ 时 \ce{H2 + F2 <=> 2HF} 的 $K\approx6.5\times10^{95}$，几乎完全；\ce{H2 + I2 <=> 2HI} 的 $K\approx867$，明显不完全。但 $K$ 只反映“限度”（热力学），不反映快慢（动力学）：$K$ 很大的反应也可能很慢。

\textbf{②判断反应方向：}任意时刻的浓度幂之积比值叫\textbf{浓度商} $Q$（形式同 $K$，但用任意时刻浓度）。
\[Q<K\ \text{反应正向进行};\qquad Q=K\ \text{平衡状态};\qquad Q>K\ \text{反应逆向进行}\]

\textbf{③判断反应的热效应：}升高温度，$K$ 增大，则正反应吸热；$K$ 减小，则正反应放热。

\textbf{④计算平衡组成和转化率。}

\begin{gainian}[转化率、产率与选择性]
\[\alpha(\text{反应物})=\frac{\text{转化的量}}{\text{起始的量}}\times100\%,\qquad \text{产率}=\frac{\text{实际产量}}{\text{理论产量}}\times100\%\]
\[\text{目标产物的选择性}=\frac{\text{转化为目标产物的反应物的量}}{\text{反应物转化的总量}}\times100\%,\qquad \text{目标产物产率}=\alpha\times\text{选择性}\]
\end{gainian}

\begin{liti}[三段式与 $Q$、$K$]
将 0.1 mol CO 与 0.1 mol \ce{H2O(g)} 充入 10 L 密闭容器，800 ℃ 达平衡时 $c(\ce{CO})=0.005$ mol·L$^{-1}$。(1) 求 $K$ 和 CO 的平衡转化率。(2) 同温下，CO 仍为 0.1 mol，\ce{H2O(g)} 改为 0.3 mol，求 CO 的平衡转化率。(3) 同温下，CO 为 0.1 mol、\ce{H2O(g)} 为 0.375 mol，某时刻测得 $c(\ce{CO})=0.0025$ mol·L$^{-1}$，此时是否平衡？反应向哪个方向进行？
\tcblower
\jie (1) 起始 $c(\ce{CO})=c(\ce{H2O})=0.01$，转化 0.005，平衡时四种物质均为 0.005 mol·L$^{-1}$，$K=\dfrac{0.005\times0.005}{0.005\times0.005}=1$，$\alpha(\ce{CO})=50\%$。\\
(2) 设 CO 转化 $x$ mol·L$^{-1}$：$\dfrac{x^2}{(0.01-x)(0.03-x)}=1$，解得 $x=0.0075$，$\alpha(\ce{CO})=75\%$。增大 \ce{H2O} 的浓度，CO 的转化率提高（而 \ce{H2O} 自身的转化率由 50\% 降为 25\%）。\\
(3) 此时 CO 转化了 0.0075，$c(\ce{H2O})=0.0375-0.0075=0.03$，$Q=\dfrac{0.0075^2}{0.0025\times0.03}=0.75<K=1$，未平衡，反应正向进行。
\end{liti}

\begin{liti}[平衡常数的组合与反应方向]
已知：①\ce{2H2(g) + CO(g) <=> CH3OH(g)} 的 $K_1$ 在 500 ℃、800 ℃ 时分别为 2.5、0.15；②\ce{CO2(g) + H2(g) <=> CO(g) + H2O(g)} 的 $K_2$ 分别为 1.0、2.5。(1) 写出 ③\ce{CO2(g) + 3H2(g) <=> CH3OH(g) + H2O(g)} 的 $K_3$ 与 $K_1$、$K_2$ 的关系，并判断反应③的 $\Delta H$ 正负。(2) 500 ℃ 时，测得反应③某时刻 \ce{H2}、\ce{CO2}、\ce{CH3OH}、\ce{H2O} 的浓度（mol·L$^{-1}$）分别为 0.8、0.1、0.3、0.15，判断此时 $\vz$ 与 $\vn$ 的大小。
\tcblower
\jie (1) ③$=$①$+$②，$K_3=K_1\cdot K_2$。500 ℃：$K_3=2.5$；800 ℃：$K_3=0.375$。升温 $K_3$ 减小，正反应放热，$\Delta H<0$。\\
(2) $Q=\dfrac{0.3\times0.15}{0.1\times0.8^3}=\dfrac{0.045}{0.0512}\approx0.88<K_3=2.5$，反应正向进行，$\vz>\vn$。
\end{liti}

\subsection{压强平衡常数 $K_p$}
对气相反应，常用平衡分压代替平衡浓度，得到\textbf{压强平衡常数} $K_p$。分压 $=$ 总压 $\times$ 物质的量分数，即 $p(\mathrm{B})=p_\text{总}\cdot\dfrac{n(\mathrm{B})}{n_\text{总}}$。如 \ce{N2(g) + 3H2(g) <=> 2NH3(g)}：
\[K_p=\frac{p^2(\ce{NH3})}{p(\ce{N2})\cdot p^3(\ce{H2})}\]
$K_p$ 同样只与温度有关。有些题目把分压除以标准压强 $p^\ominus=100$ kPa 使其无单位（如 2026 年广东卷第19题），按题给定义计算即可。

\begin{kanwu}[课件 2.2 第38～39页]
课件中压强平衡常数、物质的量分数平衡常数对应的方程式写作“\ce{$m$A(g) + $n$B(g) <=> $e$C(g) + $f$d(g)}”，其中“d”应为“D”。课件给出的关系 $K_p=K_c(RT)^{(e+f)-(m+n)}$ 与 $K_x=K_p\cdot p_\text{总}^{(m+n)-(e+f)}$ 正确，了解即可，高考一般直接用分压计算 $K_p$。
\end{kanwu}

\begin{fangfa}[求 $K_p$ 的步骤]
①列三段式求平衡时各气体的物质的量和\textbf{气体总物质的量}（别漏掉不反应的稀释气体如 Ar、水蒸气）；②求各物质的量分数；③乘以总压得分压；④代入表达式。若反应前后气体分子数相等，总压和 $n_\text{总}$ 可约去，直接用物质的量代入。
\end{fangfa}

\begin{liti}[$K_p$ 的计算]
一定温度、恒定总压 $p=100$ kPa 下，\ce{N2O4(g) <=> 2NO2(g)} 达到平衡时 \ce{N2O4} 的转化率为 50\%。求 $K_p$；若保持温度不变，将总压改为 200 kPa，平衡转化率变大还是变小？
\tcblower
\jie 设起始 \ce{N2O4} 1 mol，平衡时 \ce{N2O4} 0.5 mol、\ce{NO2} 1 mol，$n_\text{总}=1.5$ mol。$p(\ce{N2O4})=\frac13p$，$p(\ce{NO2})=\frac23p$，$K_p=\dfrac{(\frac23p)^2}{\frac13p}=\frac43p=133$ kPa（一般式：$K_p=\dfrac{4\alpha^2}{1-\alpha^2}p$）。总压增大，平衡向气体分子数减少的方向（逆向）移动，转化率变小；由 $\dfrac{4\alpha^2}{1-\alpha^2}\times200=133$ 可解得 $\alpha\approx0.38$。
\end{liti}

\section{化学平衡的移动}

\subsection{平衡移动的本质}
条件改变使 $\vz\neq\vn$，原平衡被破坏，在新条件下建立新平衡的过程叫\textbf{平衡移动}。$\vz>\vn$ 时正向移动，$\vz<\vn$ 时逆向移动。从热力学看，就是条件改变使 $Q\neq K$：$Q<K$ 正向移动，$Q>K$ 逆向移动。用“$Q$ 与 $K$”分析，往往比“速率”更清楚。

\subsection{浓度、压强、温度、催化剂的影响}
\textbf{(1) 浓度：}增大反应物浓度或减小生成物浓度，$Q$ 减小，$Q<K$，平衡正向移动；反之逆向移动。实验：\ce{Fe^3+ + 3SCN- <=> Fe(SCN)3}（红色），加 KSCN 溶液红色加深；加铁粉（\ce{2Fe^3+ + Fe = 3Fe^2+}，消耗 \ce{Fe^3+}）红色变浅。注意：加 KCl 固体，\ce{K+}、\ce{Cl-} 不参与反应，平衡\textbf{不移动}；改变固体或纯液体的量，平衡也不移动。

\textbf{(2) 压强：}只影响有气体参与、且反应前后气体分子数不等的反应。对 \ce{2NO2(g) <=> N2O4(g)}，压缩体积到原来的一半，各浓度都变为 2 倍：
\[Q=\frac{2c(\ce{N2O4})}{[2c(\ce{NO2})]^2}=\frac{K}{2}<K\]
平衡正向（气体分子数减少的方向）移动。现象：颜色\textbf{先变深}（\ce{NO2} 浓度瞬间增大）\textbf{后变浅}（平衡正移），但最终仍比原来深（只能“减弱”不能“抵消”）。对 \ce{H2(g) + I2(g) <=> 2HI(g)}，压缩时 $Q=K$ 不变，平衡不移动，但颜色因浓度增大而加深。

\textbf{(3) 温度：}升高温度，平衡向\textbf{吸热}方向移动；降低温度，向放热方向移动。温度是\textbf{唯一能改变 $K$} 的因素。实验：\ce{2NO2(g) <=> N2O4(g)} $\Delta H<0$，热水中红棕色加深，冰水中变浅。

\textbf{(4) 催化剂：}同等程度改变 $\vz$、$\vn$，平衡不移动，但缩短达到平衡的时间。

\begin{gainian}[勒夏特列原理（平衡移动原理）]
如果改变影响平衡的一个因素（浓度、压强、温度），平衡就向着能够\textbf{减弱}这种改变的方向移动。适用于已达平衡的体系、只改变一个条件的情况；移动的结果只能减弱、不能抵消这种改变。
\end{gainian}

\begin{benzhi}[“充入惰性气体”平衡怎样移动]
\textbf{恒温恒容}：总压增大，但各反应物、生成物的浓度（分压）不变，$Q=K$，平衡\textbf{不移动}。\textbf{恒温恒压}：容器体积增大，各气体浓度（分压）同比例减小，相当于\textbf{减压}，平衡向气体分子数增大的方向移动。工业上常在恒压下通入水蒸气（如乙苯脱氢制苯乙烯）或 Ar 稀释，就是为了提高气体分子数增大的反应的平衡转化率。
\end{benzhi}

\subsection{平衡移动与转化率的变化规律}
\begin{center}
\small
\begin{tabularx}{\linewidth}{p{5.6cm}X}
\toprule
操作 & 结果 \\
\midrule
增大反应物 A 的浓度（反应 \ce{A + B <=> C}） & 平衡正移，B 的转化率增大，A 自身的转化率减小 \\
恒温恒容，按原比例同时增加所有反应物 & 相当于加压：若正反应气体分子数减小，各反应物转化率增大；分子数不变，转化率不变 \\
恒温恒容，再加入一定量的唯一反应物（如 \ce{2NO2 <=> N2O4} 中再加 \ce{NO2}） & 相当于加压：\ce{NO2} 的转化率增大（新平衡中 \ce{N2O4} 的体积分数增大） \\
恒温恒压，按原比例投料 & 新平衡与原平衡等效，各组分的百分含量和转化率不变 \\
\bottomrule
\end{tabularx}
\end{center}

\begin{fangfa}[“虚拟容器法”（放缩法）]
恒容容器中按比例增加投料，可设想先把新增物质放在另一个相同容器中，在相同条件下建立与原来\textbf{完全相同}的平衡，再把两个容器“压”到一个容器中——这就把“加料”转化为“加压”问题。同理，一个容器中扩大体积，可看作“减压”。这一思想也是理解\textbf{等效平衡}的钥匙：恒温恒压下，投料比例相同即等效；恒温恒容下，气体分子数不变的反应按比例投料即等效（百分含量相同），分子数改变的反应需“一边倒”后投料量完全相同才等效。
\end{fangfa}

\begin{liti}[等效平衡]
恒温恒容的甲、乙、丙三个相同容器中发生 \ce{2SO2(g) + O2(g) <=> 2SO3(g)}：甲中投入 2 mol \ce{SO2} 和 1 mol \ce{O2}，乙中投入 2 mol \ce{SO3}，丙中投入 4 mol \ce{SO2} 和 2 mol \ce{O2}。比较：(1) 甲、乙平衡时 \ce{SO3} 的体积分数；(2) 甲中 \ce{SO2} 的转化率 $\alpha_\text{甲}$ 与乙中 \ce{SO3} 的转化率 $\alpha_\text{乙}$ 之和；(3) 甲、丙中 \ce{SO2} 的转化率。
\tcblower
\jie (1) 把乙中 2 mol \ce{SO3} “一边倒”为 2 mol \ce{SO2} 和 1 mol \ce{O2}，与甲完全相同，恒温恒容下两者为等效平衡，平衡组成完全相同，\ce{SO3} 体积分数相等。(2) 设甲平衡时 \ce{SO3} 为 $x$ mol，则 $\alpha_\text{甲}=\frac x2$；乙平衡时 \ce{SO3} 也为 $x$ mol，$\alpha_\text{乙}=\frac{2-x}{2}$，二者之和为 1。(3) 丙相当于把两个甲“压缩”到一个容器中，加压平衡正移，$\alpha_\text{丙}>\alpha_\text{甲}$。
\end{liti}

\begin{liti}[浓度变化倍数的分析]
某温度下，密闭容器中发生 \ce{$a$X(g) + $b$Y(g) <=> $c$Z(g) + $d$W(g)}，达平衡后将容器体积扩大为原来的 2 倍，再次达到平衡时 W 的浓度为原平衡的 0.65 倍。则平衡向\tk{}方向移动，$a+b$\tk$c+d$，X 的转化率\tk。
\tcblower
\silu 先假设“平衡不移动”，再与实际比较。\\
\jie 体积扩大为 2 倍的瞬间，W 浓度变为原来的 0.5 倍；最终为 0.65 倍 $>0.5$ 倍，说明平衡\textbf{正向}移动生成了更多 W。减压平衡向气体分子数增大的方向移动，故 $a+b<c+d$；X 的转化率增大。\\
\xianjing 若只看“0.65 倍 $<1$ 倍”就判断“逆移”，就犯了忽视体积变化的错误。
\end{liti}

\begin{yicuo}
下列事实\textbf{不能}用勒夏特列原理解释：使用催化剂加快合成氨；\ce{H2 + I2 <=> 2HI} 平衡体系加压颜色变深；工业合成氨选择 500 ℃左右（从平衡看低温才有利，选高温是为了速率和催化剂活性）。能解释的如：光照新制氯水 pH 减小（\ce{HClO} 分解，\ce{Cl2 + H2O <=> HCl + HClO} 正移）；浓氨水加 NaOH 固体制氨；增大压强有利于 \ce{SO2} 转化为 \ce{SO3}。
\end{yicuo}

\section{化学平衡图像}

\begin{fangfa}[图像题通用四步]
\textbf{一看面}：纵、横坐标代表什么；\textbf{二看线}：曲线的变化趋势，是否为平衡线；\textbf{三看点}：起点、拐点、交点、最高（低）点、终点；\textbf{四看辅助线}：作等温线、等压线。两个基本原则：\textbf{①“先拐先平，数值大”}——先出现拐点（先达平衡）的曲线对应的温度高或压强大或使用了催化剂；\textbf{②“定一议二”}——图中有三个变量时，先固定一个，再讨论另外两个的关系。
\end{fangfa}

\subsection{速率—时间（$v$–$t$）图像}
\begin{center}
\begin{tikzpicture}[scale=0.55,font=\scriptsize]
\newcommand{\vtaxes}{\draw[-Stealth] (0,0)--(4.2,0) node[below]{$t$}; \draw[-Stealth] (0,0)--(0,3.2) node[left]{$v$}; \draw[thick] (0,1.3)--(1.2,1.3);}
\begin{scope}
\vtaxes \draw[thick,mainblue] (1.2,2.5) .. controls (2,1.9) .. (3.8,1.75); \draw[thick,warnred] (1.2,1.3) .. controls (2,1.6) .. (3.8,1.75); \draw[densely dotted] (1.2,1.3)--(1.2,2.5);
\node at (2.1,-0.6) {①增大反应物浓度};\node[mainblue] at (2.0,2.5){$\vz$};\node[warnred] at (2.2,1.25){$\vn$};
\end{scope}
\begin{scope}[xshift=5.2cm]
\vtaxes \draw[thick,mainblue] (1.2,2.0) .. controls (2,2.25) .. (3.8,2.35); \draw[thick,warnred] (1.2,2.7) .. controls (2,2.45) .. (3.8,2.35); \draw[densely dotted] (1.2,1.3)--(1.2,2.7);
\node[align=center] at (2.1,-0.8) {②升温\\（正反应放热）};\node[warnred] at (2.2,2.9){$\vn$};\node[mainblue] at (2.0,1.75){$\vz$};
\end{scope}
\begin{scope}[xshift=10.4cm]
\vtaxes \draw[thick,mainblue] (1.2,2.7) .. controls (2,2.45) .. (3.8,2.35); \draw[thick,warnred] (1.2,2.0) .. controls (2,2.25) .. (3.8,2.35); \draw[densely dotted] (1.2,1.3)--(1.2,2.7);
\node[align=center] at (2.1,-0.8) {③加压（正反应\\气体分子数减小）};\node[mainblue] at (2.2,2.9){$\vz$};\node[warnred] at (2.0,1.75){$\vn$};
\end{scope}
\begin{scope}[xshift=15.6cm]
\vtaxes \draw[thick,deepgreen] (1.2,2.4)--(3.8,2.4); \draw[densely dotted] (1.2,1.3)--(1.2,2.4);
\node[align=center] at (2.1,-0.8) {④催化剂\\（或等体积反应加压）};\node[deepgreen] at (2.5,2.75){$\vz'=\vn'$};
\end{scope}
\end{tikzpicture}
\end{center}
读图要点：①\textbf{连续}（某一速率从原值开始变化）——改变的是浓度；\textbf{断点}（两条线都跳变）——改变的是温度或压强；\textbf{平台}（两线同步跳到同一值）——催化剂或等体积反应改变压强。②跳到上方的那条速率线，决定平衡移动的方向。

\begin{liti}[多阶段 $v$-$t$ 图像]
合成氨反应 \ce{N2(g) + 3H2(g) <=> 2NH3(g)} $\Delta H<0$ 在某时段内的 $v$–$t$ 图如下（$t_1$、$t_3$、$t_4$ 时刻各改变一个条件）。(1) 判断各时刻改变的条件；(2) 哪段时间内 \ce{NH3} 的百分含量最高？
\begin{center}
\begin{tikzpicture}[xscale=0.95,yscale=0.8,font=\scriptsize]
\draw[-Stealth] (0,0)--(8.4,0) node[below]{$t$}; \draw[-Stealth] (0,0)--(0,3.6) node[left]{$v$};
\draw[thick] (0,1.2)--(1.2,1.2);
\draw[thick,warnred] (1.2,2.6) .. controls (1.8,2.2) .. (2.6,2.05); \draw[thick,mainblue] (1.2,1.8) .. controls (1.8,1.95) .. (2.6,2.05);
\draw[thick] (2.6,2.05)--(3.6,2.05);
\draw[thick,deepgreen] (3.6,2.9)--(5.2,2.9);
\draw[thick,mainblue] (5.2,1.4) .. controls (5.9,1.55) .. (6.8,1.62); \draw[thick,warnred] (5.2,1.9) .. controls (5.9,1.7) .. (6.8,1.62);
\draw[thick] (6.8,1.62)--(8.0,1.62);
\foreach \x/\l in {1.2/t_1,2.6/t_2,3.6/t_3,5.2/t_4,6.8/t_5}{\draw[densely dotted] (\x,0)--(\x,3.1); \node[below] at (\x,0) {$\l$};}
\node[warnred] at (1.65,2.8) {$\vn$}; \node[mainblue] at (1.65,1.55) {$\vz$};
\node[deepgreen] at (4.4,3.15) {$\vz=\vn$};
\node[warnred] at (5.65,2.05) {$\vn$}; \node[mainblue] at (5.65,1.25) {$\vz$};
\end{tikzpicture}
\end{center}
\tcblower
\jie (1) $t_1$：$\vz$、$\vn$ 都突然增大（断点），且 $\vn>\vz$，平衡逆移。若是加压，平衡应正移，所以是\textbf{升高温度}（放热反应升温逆移）。$t_3$：两速率同等程度增大（平台），平衡不移动，是\textbf{使用催化剂}（该反应气体分子数变化，不能是加压）。$t_4$：两速率都突然减小且 $\vn>\vz$，平衡逆移，是\textbf{减小压强}（若是降温，平衡应正移）。\\
(2) $t_1$ 后平衡逆移，$t_4$ 后又逆移，\ce{NH3} 含量逐次降低，所以 $0\sim t_1$ 段 \ce{NH3} 百分含量最高。\\
\xianjing 不要把“$t_3\sim t_4$ 速率最大”误认为“\ce{NH3} 含量最高”——速率大小与平衡组成是两回事。
\end{liti}

\subsection{含量（转化率）—时间—温度（压强）图像}
\begin{center}
\begin{tikzpicture}[scale=0.6,font=\scriptsize]
\begin{scope}
\draw[-Stealth] (0,0)--(5,0) node[below]{$t$}; \draw[-Stealth] (0,0)--(0,3.4) node[left]{C\%};
\draw[thick,warnred] (0,0) .. controls (0.6,1.6) .. (1.1,1.7)--(4.6,1.7) node[right]{$T_2$};
\draw[thick,mainblue] (0,0) .. controls (1.5,2.5) .. (2.4,2.7)--(4.6,2.7) node[right]{$T_1$};
\node[align=center] at (2.4,-0.9) {$T_2$ 先拐，$T_2>T_1$；高温 C\% 低\\$\Rightarrow$ 正反应放热};
\end{scope}
\begin{scope}[xshift=7.8cm]
\draw[-Stealth] (0,0)--(5,0) node[below]{$T$}; \draw[-Stealth] (0,0)--(0,3.4) node[left]{$\alpha$(A)};
\draw[thick,mainblue] (0.2,2.9) .. controls (2,2.6) .. (4.6,1.6) node[right]{$p_1$};
\draw[thick,warnred] (0.2,2.2) .. controls (2,1.9) .. (4.6,0.9) node[right]{$p_2$};
\draw[dashed] (2.5,0)--(2.5,3.0);
\node[align=center] at (2.4,-0.9) {升温 $\alpha$ 减小$\Rightarrow$放热；作等温线，\\若 $p_1>p_2$ 则正反应气体分子数减小};
\end{scope}
\end{tikzpicture}
\end{center}

\subsection{特殊点的分析}
\textbf{(1) 平衡曲线外的点：}平衡曲线上每一点都是平衡点。曲线外的点，看它要“走向”曲线需要正移还是逆移：若某点的生成物含量低于同条件下的平衡值，则反应还要正向进行，$\vz>\vn$。

\textbf{(2) 最高点（最低点）：}在“反应进行相同时间后测产物含量（或转化率）—温度”曲线中，若出现先升后降，则最高点之前多为“未达平衡”，升温加快速率，含量增大；最高点之后已达平衡，升温使平衡逆向移动（正反应放热），含量减小。也可能是催化剂在高温下活性降低所致——答题时要结合题给信息判断。

\begin{liti}[最高点图像]
在不同温度下，向相同的恒容密闭容器中充入等量 CO 和 \ce{H2}，发生 \ce{CO(g) + 2H2(g) <=> CH3OH(g)}，反应相同时间后测得 \ce{CH3OH} 的体积分数 $\varphi$ 与温度 $T$ 的关系如图（在 $T_3$ 处最高）。(1) $T_1$、$T_5$ 中，哪个点一定处于平衡状态？(2) 该反应 $\Delta H$\tk0。(3) $T_3$ 之前 $\varphi$ 随温度升高而增大的原因是\tk。
\begin{center}
\begin{tikzpicture}[scale=0.55,font=\scriptsize]
\draw[-Stealth] (0,0)--(6.5,0) node[below]{$T$}; \draw[-Stealth] (0,0)--(0,3.2) node[left]{$\varphi(\ce{CH3OH})$};
\draw[thick,mainblue] (0.4,0.4) .. controls (1.8,1.2) and (2.4,2.7) .. (3,2.7) .. controls (3.6,2.7) and (4.6,1.6) .. (6,1.0);
\foreach \x/\l in {0.9/1,2.0/2,3.0/3,4.4/4,5.6/5} {\draw[dashed] (\x,0)--(\x,0.12) node[below=2pt]{$T_\l$};}
\end{tikzpicture}
\end{center}
\tcblower
\jie (1) $T_5$。最高点之后的点（$T_4$、$T_5$）已达平衡；$T_1$ 时速率慢，相同时间内尚未达到平衡。(2) $<$：最高点之后升温 $\varphi$ 减小，平衡逆移，正反应放热。(3) $T_3$ 之前反应未达平衡，升高温度反应速率加快，相同时间内生成的 \ce{CH3OH} 增多。
\end{liti}

\section{化学反应的方向}

\subsection{自发过程与焓判据、熵判据}
在一定条件下，不借助外力就能自动进行的过程叫\textbf{自发过程}，如水往低处流、冰在常温下融化、墨水扩散。能自发进行到显著程度的反应叫\textbf{自发反应}。“自发”只说明有发生的\textbf{趋势}，不涉及快慢：\ce{H2} 与 \ce{O2} 常温下自发，但不点燃几乎不反应。

\textbf{焓判据：}放热反应体系能量降低，多数能自发。但有些吸热反应也自发，如 \ce{Ba(OH)2*8H2O} 与 \ce{NH4Cl} 的反应、\ce{NaHCO3} 与盐酸的反应。

\textbf{熵判据：}熵（$S$）是衡量体系混乱度的物理量。同一物质 $S(\mathrm{g})>S(\mathrm{l})>S(\mathrm{s})$；\textbf{气体分子数增多的反应 $\Delta S>0$}。熵增有利于自发，但有些熵减反应也自发，如 \ce{4Fe(OH)2(s) + O2(g) + 2H2O(l) = 4Fe(OH)3(s)}。

\begin{kanwu}[课件 2.3 第16页]
课件中“\ce{Fe(OH)2(s) + O2(g) + H2O(g) = Fe(OH)3(s)}”未配平且水的状态不当，应为 \ce{4Fe(OH)2(s) + O2(g) + 2H2O(l) = 4Fe(OH)3(s)}（这是氢氧化亚铁被氧化，而不是“制备氢氧化亚铁”）。
\end{kanwu}

\subsection{自由能判据}
\begin{gainian}[$\Delta G=\Delta H-T\Delta S$]
恒温恒压下：$\Delta G<0$，反应能自发进行；$\Delta G=0$，平衡状态；$\Delta G>0$，不能自发进行（逆反应自发）。
\end{gainian}

\begin{center}
\small
\begin{tabular}{cccl}
\toprule
$\Delta H$ & $\Delta S$ & $\Delta G=\Delta H-T\Delta S$ & 自发性 \\
\midrule
$<0$ & $>0$ & 恒 $<0$ & 任何温度都自发 \\
$>0$ & $<0$ & 恒 $>0$ & 任何温度都不自发 \\
$<0$ & $<0$ & 低温 $<0$ & 低温自发（如合成氨） \\
$>0$ & $>0$ & 高温 $<0$ & 高温自发（如 \ce{CaCO3} 分解） \\
\bottomrule
\end{tabular}
\end{center}
转折温度 $T=\dfrac{\Delta H}{\Delta S}$。计算时注意单位：$\Delta H$ 常用 kJ·mol$^{-1}$，$\Delta S$ 常用 J·mol$^{-1}$·K$^{-1}$，需统一。例如 \ce{CaCO3(s) = CaO(s) + CO2(g)}，$\Delta H=+178.5$ kJ·mol$^{-1}$，$\Delta S=+160.5$ J·mol$^{-1}$·K$^{-1}$，自发分解所需温度 $T>\dfrac{178.5\times10^3}{160.5}\approx1112$ K。

\begin{zhenti}{2025年广东卷第19题(2)\quad 来源：仓库所附真题文档}
298 K 下，反应 \ce{TiO2(s) + 2C(s) + 2Cl2(g) = TiCl4(g) + 2CO(g)} 的 $\Delta H<0$、$\Delta S>0$，则 298 K 下该反应\tk（填“能”或“不能”）自发进行。

\textbf{解析：}$\Delta H<0$、$\Delta S>0$，$\Delta G=\Delta H-T\Delta S$ 在任何温度下都小于 0，能自发进行。
\end{zhenti}

\begin{liti}[由自发性推断]
常温下，反应 \ce{A(s) + B(g) = 2X(g)} 不能自发进行。则该反应的 $\Delta H$\tk0。\ce{CH3OH(l) + NH3(g) = CH3NH2(g) + H2O(g)} 在某温度下自发向右进行，若 $|\Delta H|=17$ kJ·mol$^{-1}$，$|\Delta H-T\Delta S|=17$ kJ·mol$^{-1}$，则 $\Delta H$\tk0。
\tcblower
\jie 第一个反应气体分子数增大，$\Delta S>0$；若 $\Delta H\le0$，则 $\Delta G<0$ 一定自发，与事实矛盾，故 $\Delta H>0$。第二个反应自发，$\Delta G=-17$；该反应气体分子数增加，$\Delta S>0$，$-T\Delta S<0$。若 $\Delta H=-17$，则 $\Delta G=-17-T\Delta S<-17$，矛盾；故 $\Delta H=+17$ kJ·mol$^{-1}$，$\Delta H>0$。
\end{liti}

\section{化学反应的调控}

\subsection{工业合成氨条件的选择}
\ce{N2(g) + 3H2(g) <=> 2NH3(g)} $\Delta H=-92.4$ kJ·mol$^{-1}$。特点：可逆、放热、气体分子数减小（$\Delta S<0$，低温自发）。

\begin{center}
\small
\begin{tabularx}{\linewidth}{p{1.6cm}XXp{4.6cm}}
\toprule
条件 & 从速率看 & 从平衡看 & 实际选择及理由 \\
\midrule
压强 & 越大越快 & 越大越有利 & 10～30 MPa：压强越大，对设备材料和动力要求越高，成本增加 \\
温度 & 越高越快 & 越低越有利 & 400～500 ℃：兼顾速率；且铁触媒在 500 ℃ 左右活性最大 \\
催化剂 & 加快 & 无影响 & 铁触媒；原料气需净化，防止催化剂“中毒” \\
浓度 & 增大反应物浓度 & 增大反应物浓度、减小 \ce{NH3} 浓度 & 冷却使 \ce{NH3} 液化分离；未反应的 \ce{N2}、\ce{H2} 循环使用并及时补充；\ce{N2} 价廉，适当提高其比例以提高 \ce{H2} 的转化率 \\
\bottomrule
\end{tabularx}
\end{center}

\begin{center}
\small
\begin{tabular}{ccccccc}
\toprule
\multirow{2}{*}{温度/℃} & \multicolumn{6}{c}{平衡混合物中 \ce{NH3} 的体积分数/\%（$n(\ce{N2}):n(\ce{H2})=1:3$）} \\
\cmidrule(lr){2-7}
& 0.1 MPa & 10 MPa & 20 MPa & 30 MPa & 60 MPa & 100 MPa \\
\midrule
200 & 15.3 & 81.5 & 86.4 & 89.9 & 95.4 & 98.8 \\
300 & 2.20 & 52.0 & 64.2 & 71.0 & 84.2 & 92.6 \\
400 & 0.40 & 25.1 & 38.2 & 47.0 & 65.2 & 79.8 \\
500 & 0.10 & 10.6 & 19.1 & 26.4 & 42.2 & 57.5 \\
600 & 0.05 & 4.50 & 9.10 & 13.8 & 23.1 & 31.4 \\
\bottomrule
\end{tabular}
\end{center}
\noindent\small 数据分析：同一温度下压强越大、同一压强下温度越低，\ce{NH3} 的平衡含量越高，与勒夏特列原理的预测一致；但 200 ℃ 时速率太慢，实际只能折中选择 400～500 ℃，靠循环利用原料弥补单程转化率的不足。\normalsize

工艺流程中：\textbf{净化}防催化剂中毒；\textbf{压缩机}加压；\textbf{热交换器}用反应放出的热预热原料气，节约能源；\textbf{冷却}使氨液化分离，促使平衡正移；\textbf{循环}提高原料利用率。

\textbf{类比——硫酸工业中 \ce{SO2} 的催化氧化}：\ce{2SO2(g) + O2(g) <=> 2SO3(g)} $\Delta H=-196.6$ kJ·mol$^{-1}$。常压下 450 ℃ 时 \ce{SO2} 平衡转化率已达 97.5\%，加压收益很小而成本大，故采用\textbf{常压}；温度 400～500 ℃（兼顾速率与 \ce{V2O5} 催化剂活性）；通入过量空气提高 \ce{SO2} 转化率；尾气中的 \ce{SO2} 需回收处理。

\begin{fangfa}[“条件选择原因”的答题模板]
从三个角度作答，缺一不可：①\textbf{速率}：该条件使反应速率较快；②\textbf{平衡}：该条件下平衡转化率（产率）较高或可接受；③\textbf{实际因素}：催化剂活性温度、设备承受能力、能耗与成本、安全与环保。例如“选择 500 ℃左右”：温度较高反应速率快，且铁触媒在 500 ℃左右活性最大；温度过高则平衡转化率降低、能耗增大。
\end{fangfa}

\section{平衡计算的常见题型}

\begin{fangfa}[五类常见平衡计算]
\textbf{①求转化率、平衡浓度}：三段式，转化量成比例。\\
\textbf{②已知 $K$ 求组成}：设未知数代入 $K$ 表达式，常解一元二次方程；若 $K$ 很大或很小，可先假设“几乎完全反应”或“几乎不反应”再近似。\\
\textbf{③压强、密度、平均摩尔质量的变化}：恒温恒容 $\dfrac{p_\text{后}}{p_\text{前}}=\dfrac{n_\text{后}}{n_\text{前}}$；恒温恒压 $\dfrac{V_\text{后}}{V_\text{前}}=\dfrac{n_\text{后}}{n_\text{前}}$；$\bar M=\dfrac{m_\text{总}}{n_\text{总}}$。\\
\textbf{④平衡后再加料}：先判断是恒容还是恒压，用“虚拟容器法”转化为等效平衡或加压问题。\\
\textbf{⑤范围（极值）问题}：可逆反应不能进行到底，各物质的量介于“完全正向”与“完全逆向”两个极端之间。
\end{fangfa}

\begin{liti}[极值法判断范围]
在密闭容器中进行 \ce{X2(g) + Y2(g) <=> 2Z(g)}，已知 \ce{X2}、\ce{Y2}、Z 的起始浓度分别为 0.1、0.3、0.2 mol·L$^{-1}$，一定条件下达到平衡。下列浓度可能存在的是
\xx{Z 为 0.4 mol·L$^{-1}$}{\ce{X2} 为 0.2 mol·L$^{-1}$}{\ce{Y2} 为 0.4 mol·L$^{-1}$}{Z 为 0.3 mol·L$^{-1}$}
\tcblower
\jie 假设完全正向：\ce{X2} 0，\ce{Y2} 0.2，Z 0.4；假设完全逆向：\ce{X2} 0.2，\ce{Y2} 0.4，Z 0。可逆反应不能进行到底，所以 $0<c(\ce{X2})<0.2$，$0.2<c(\ce{Y2})<0.4$，$0<c(\mathrm{Z})<0.4$。A、B、C 都取到了端点值，不可能；D 可能。答案：D。
\end{liti}

\begin{liti}[由压强变化求转化率]
恒温恒容下，向密闭容器中充入 2 mol \ce{SO2} 和 1 mol \ce{O2}，发生 \ce{2SO2(g) + O2(g) <=> 2SO3(g)}，平衡时容器内压强为起始时的 $\frac56$。求 \ce{SO2} 的平衡转化率及平衡时 \ce{SO3} 的体积分数。
\tcblower
\jie 设转化 \ce{O2} $x$ mol，则气体总物质的量减少 $x$ mol：$\dfrac{3-x}{3}=\dfrac56$，$x=0.5$。转化 \ce{SO2} 1 mol，$\alpha(\ce{SO2})=50\%$；平衡时 \ce{SO3} 1 mol，总 2.5 mol，体积分数 40\%。
\end{liti}

\section{多重平衡与选择性}

实际体系中往往同时存在几个反应，它们共用某些物质（\textbf{竞争反应}或\textbf{连续反应}）。处理多重平衡的要点：
\begin{enumerate}
\item \textbf{每个反应各设一个变量}（转化量 $x$、$y$），列出各物质的平衡量；共用物质的量要把各反应的贡献相加。
\item 用\textbf{原子守恒}代替繁琐的三段式往往更快：如由 C 原子守恒得 $n(\ce{CO2})+n(\ce{CO})+n(\ce{CH3OH})=n_0(\ce{CO2})$。
\item 每个反应的 $K$ 都要用\textbf{同一平衡体系中}各物质的平衡浓度（分压）计算，共用物质的浓度只有一个。
\item 分析温度对转化率的影响时，要看哪个反应占主导：放热反应在低温段占主导，吸热反应在高温段占主导。
\end{enumerate}

\begin{liti}[多重平衡的计算]
\ce{CO2} 加氢体系：反应Ⅰ \ce{CO2(g) + 3H2(g) <=> CH3OH(g) + H2O(g)} $\Delta H_1=-49.5$ kJ·mol$^{-1}$；反应Ⅱ \ce{CO2(g) + H2(g) <=> CO(g) + H2O(g)} $\Delta H_2=+41.2$ kJ·mol$^{-1}$。恒压 $p$ 下充入 1 mol \ce{CO2} 和 3 mol \ce{H2}，平衡时 \ce{CO2} 转化率为 30\%，\ce{CH3OH} 的选择性为 80\%。(1) 求平衡时各物质的物质的量；(2) 求反应Ⅱ的 $K_p$；(3) 实验发现：温度升高，\ce{CO2} 的平衡转化率先减小后增大，解释原因。
\tcblower
\jie (1) 转化的 \ce{CO2} 为 0.3 mol，其中 $0.3\times80\%=0.24$ mol 生成 \ce{CH3OH}，0.06 mol 生成 CO。\ce{H2} 消耗 $3\times0.24+0.06=0.78$ mol，剩 2.22 mol；\ce{H2O} 由两反应共生成 $0.24+0.06=0.30$ mol。平衡时：\ce{CO2} 0.70，\ce{H2} 2.22，\ce{CH3OH} 0.24，\ce{H2O} 0.30，CO 0.06，$n_\text{总}=3.52$ mol。\\
(2) 反应Ⅱ前后气体分子数相等，$K_p=\dfrac{n(\ce{CO})n(\ce{H2O})}{n(\ce{CO2})n(\ce{H2})}=\dfrac{0.06\times0.30}{0.70\times2.22}\approx0.0116$。反应Ⅰ前后气体分子数不等，需代入分压（$p$ 为平衡总压）：\[K_p(\text{Ⅰ})=\dfrac{\dfrac{0.24}{3.52}p\cdot\dfrac{0.30}{3.52}p}{\dfrac{0.70}{3.52}p\cdot\left(\dfrac{2.22}{3.52}p\right)^3}\approx\dfrac{0.116}{p^2}\]
(3) 反应Ⅰ放热、反应Ⅱ吸热。低温段以反应Ⅰ为主，升温反应Ⅰ逆移，\ce{CO2} 转化率减小；高温段以反应Ⅱ为主，升温反应Ⅱ正移，\ce{CO2} 转化率增大。同时 \ce{CH3OH} 的选择性随温度升高持续下降。
\end{liti}

\begin{zhenti}{2025年广东卷第19题(3)②③\quad 来源：仓库所附真题文档}
将 5.0 mol \ce{TiCl4} 与 10.0 mol Ti 放入容积为 $V_0$ L 的恒容密闭容器中，体系存在：(a)\ce{Ti(s) + TiCl4(g) <=> 2TiCl2(s)} $\Delta H_1=-267.8$ kJ·mol$^{-1}$；(b)\ce{TiCl2(s) + TiCl4(g) <=> 2TiCl3(g)} $\Delta H=+200.1$ kJ·mol$^{-1}$；(c)\ce{Ti(s) + 3TiCl4(g) <=> 4TiCl3(g)} $\Delta H=+132.4$ kJ·mol$^{-1}$；(d)\ce{TiCl3(g) <=> TiCl3(s)} $\Delta H_2$（凝华，放热）。不同温度下平衡体系组成如图（1073～1273 K）：Ti(s) 的量随温度升高而增大，$T_0$ 时为 8.5 mol；曲线Ⅱ随温度升高而增大，$T_0$ 时为 6.0 mol；曲线Ⅰ随温度升高由约 9 mol 减小至接近 0；\ce{TiCl3(s)}、\ce{TiCl4(g)} 的量均很少。②曲线Ⅰ对应的物质为\tk。③温度 $T_0$ K 下（此时 \ce{TiCl2(s)}、\ce{TiCl3(s)} 均为 0），$n[\ce{TiCl4(g)}]=$\tk mol，反应(c)的平衡常数 $K=$\tk（列出算式）。

\textbf{解析：}②(a)放热、(d)放热，升温逆移；(b)(c)吸热，升温正移。所以 \ce{TiCl3(g)} 随温度升高增多（曲线Ⅱ），\ce{TiCl2(s)} 随温度升高减少（曲线Ⅰ）。③由 Ti 原子守恒：$n(\ce{TiCl4})=10.0+5.0-8.5-6.0=0.5$ mol（可用 Cl 守恒验证：$4\times0.5+3\times6.0=20=4\times5.0$）。Ti 为固体不写入，$K=\dfrac{c^4(\ce{TiCl3})}{c^3(\ce{TiCl4})}=\dfrac{(6.0/V_0)^4}{(0.5/V_0)^3}$。
\end{zhenti}

% >>>>>>>> 开始：chapters/ch03ex.tex
% ---------------- 第3章 章末 ----------------
\zhangmo{本章小结}
\subsection*{一、知识总结}
\begin{center}
\small
\begin{tabularx}{\linewidth}{p{2.4cm}X}
\toprule
平衡状态 & 逆、等、动、定、变；判断：正逆相等（一正一逆、成比例）或变量不变 \\
平衡常数 & 只与温度有关；固体纯液体不写；方程式相加 $K$ 相乘；$Q$ 与 $K$ 比较定方向；$K_p$ 用分压 \\
平衡移动 & 浓度、压强（看气体分子数）、温度（唯一改变 $K$）、催化剂（不移动）；勒夏特列原理“减弱不抵消” \\
图像 & 四看；先拐先平数值大；定一议二；曲线外的点看趋向；最高点前未平衡、后平衡 \\
方向 & $\Delta G=\Delta H-T\Delta S<0$ 自发；四种组合；转折温度 $\Delta H/\Delta S$ \\
调控 & 合成氨 10～30 MPa、400～500 ℃、铁触媒、液化分离与循环；兼顾速率、平衡、成本 \\
多重平衡 & 每个反应设变量；原子守恒；同一体系共用浓度；选择性$=$生成目标产物所耗反应物/反应物总转化量 \\
\bottomrule
\end{tabularx}
\end{center}

\subsection*{二、方法提升}
\begin{enumerate}
\item \textbf{三段式是平衡计算的“万能钥匙”}：起始、转化、平衡三行，转化量与计量数成比例；有 $K_p$ 时加一行“分压”。
\item \textbf{判断移动方向先用 $Q$ 与 $K$}：浓度、压强变化看 $Q$ 怎么变，温度变化看 $K$ 怎么变。
\item \textbf{“假设不移动再比较”}：处理“浓度变为原来的几倍”类问题，先算假设不移动时的值，再与实际值比较。
\item \textbf{多重平衡温度曲线}：先判断各反应的吸放热，再分温度段讨论谁占主导。
\end{enumerate}

\subsection*{三、易错点归纳}
\begin{enumerate}
\item 把“浓度相等”“浓度比等于计量数比”当成平衡标志；忘记检查一个量是否为“变量”。
\item $K$ 表达式中写入固体、纯液体；用起始浓度或某时刻浓度代替平衡浓度计算 $K$。
\item 误认为增大压强平衡一定移动；恒容充惰性气体判为“加压，平衡移动”。
\item 误认为平衡正移，反应物转化率一定增大（增大某反应物浓度时其自身转化率减小）。
\item 误认为“自发反应一定能迅速发生”；计算 $T=\Delta H/\Delta S$ 时忘统一 kJ 与 J。
\item 选择合成氨温度时答“温度高有利于平衡正移”（错，低温才有利于平衡）。
\end{enumerate}

\zhangmo{分层练习}
\zu{A组\quad 基础巩固}
\begin{xiti}[start=1]
\item 一定温度下的恒容容器中发生 \ce{2A(g) + B(s) <=> C(g) + D(g)}，下列描述能表明反应已达平衡的是\xz\\
①混合气体的压强不变\quad ②混合气体的密度不变\quad ③C(g) 的浓度不变\quad ④A、C、D 的浓度之比为 2:1:1\quad ⑤单位时间内生成 $n$ mol C，同时生成 $n$ mol D\quad ⑥单位时间内生成 $n$ mol D，同时生成 $2n$ mol A
\xx{①②⑤}{②③⑤}{②③⑥}{①③⑥}

\item 下列关于化学平衡常数的说法正确的是\xz
\xxa{化学平衡常数表示反应进行的程度，与温度无关}{两种物质反应，不管怎样书写化学方程式，平衡常数均不变}{化学平衡常数等于某一时刻生成物浓度幂之积与反应物浓度幂之积的比值}{温度一定时，对于给定的反应，正、逆反应的平衡常数互为倒数}

\item 某温度下 \ce{N2(g) + 3H2(g) <=> 2NH3(g)} 的平衡常数为 $K$。同温下，\ce{1/2N2(g) + 3/2H2(g) <=> NH3(g)} 的平衡常数为\tk，\ce{2NH3(g) <=> N2(g) + 3H2(g)} 的平衡常数为\tk。

\item 下列事实\textbf{不能}用勒夏特列原理解释的是\xz
\xxa{光照新制的氯水时，溶液的 pH 逐渐减小}{加催化剂，使 \ce{N2} 和 \ce{H2} 在一定条件下转化为 \ce{NH3}}{可用浓氨水和氢氧化钠固体快速制氨}{增大压强，有利于 \ce{SO2} 与 \ce{O2} 反应生成 \ce{SO3}}

\item 对于 \ce{2SO2(g) + O2(g) <=> 2SO3(g)} $\Delta H<0$，在恒温恒容条件下达到平衡后，下列措施能提高 \ce{SO2} 平衡转化率的是\xz\\
①升高温度\quad ②再充入 \ce{O2}\quad ③充入 He\quad ④再按原比例充入 \ce{SO2} 和 \ce{O2}\quad ⑤使用催化剂\quad ⑥移走部分 \ce{SO3}
\xx{①②⑥}{②④⑥}{②③④}{④⑤⑥}

\item 下列反应在任何温度下均能自发进行的是\xz
\xxa{\ce{2N2(g) + O2(g) = 2N2O(g)} \quad $\Delta H=+163$ kJ·mol$^{-1}$}{\ce{Ag(s) + 1/2Cl2(g) = AgCl(s)} \quad $\Delta H=-127$ kJ·mol$^{-1}$}{\ce{HgO(s) = Hg(l) + 1/2O2(g)} \quad $\Delta H=+91$ kJ·mol$^{-1}$}{\ce{H2O2(l) = 1/2O2(g) + H2O(l)} \quad $\Delta H=-98$ kJ·mol$^{-1}$}

\item 在 \ce{2NO2(g) <=> N2O4(g)} $\Delta H<0$ 的平衡体系中，下列操作能使混合气体颜色变深、且平衡逆向移动的是\xz
\xxa{恒温时压缩容器体积}{恒容时升高温度}{恒容时充入 Ar}{恒压时充入 Ar}
\end{xiti}

\zu{B组\quad 高考典型}
\begin{xiti}
\item 某温度下，在 2 L 密闭容器中充入气体 A、B，其浓度分别为 2 mol·L$^{-1}$、1 mol·L$^{-1}$，发生 \ce{3A(g) + 2B(g) <=> 4C(\text{?}) + 2D(\text{?})}。达到平衡时生成 1.6 mol C，且反应前后的压强之比为 5:4（相同温度）。(1) 判断 C、D 的状态，写出平衡常数表达式；(2) 求 B 的平衡转化率和该温度下的 $K$（保留两位有效数字）；(3) 增加 C 的量，A、B 的转化率如何变化？

\item \laiyuan{原创}\ 对于 \ce{$m$A(g) + $n$B(g) <=> $p$C(g)}，在不同条件下 C 的体积分数随时间变化如下：曲线 a（$T_2$、$p_2$）最先出现拐点，平台高度 0.30；曲线 b（$T_1$、$p_2$）其次出现拐点，平台高度 0.50；曲线 c（$T_1$、$p_1$）最后出现拐点，平台高度 0.35。(1) 比较 $T_1$、$T_2$ 和 $p_1$、$p_2$ 的大小；(2) 判断 $\Delta H$ 的正负及 $m+n$ 与 $p$ 的关系；(3) 在恒压下，\ce{2SO2(g) + O2(g) <=> 2SO3(g)} 的 \ce{SO2} 平衡转化率随温度升高而降低，投料比 $m=n(\ce{SO2})/n(\ce{O2})$ 分别为 $m_1$、$m_2$、$m_3$ 的三条曲线自上而下排列，比较 $m_1$、$m_2$、$m_3$ 的大小。

\item 硫酸工业中 \ce{2SO2(g) + O2(g) <=> 2SO3(g)} $\Delta H=-196.6$ kJ·mol$^{-1}$。不同条件下 \ce{SO2} 的平衡转化率：450 ℃ 时 0.1 MPa 为 97.5\%，10 MPa 为 99.7\%；550 ℃ 时 0.1 MPa 为 85.6\%，10 MPa 为 98.3\%。(1) 实际生产为何选择常压？(2) 为何选择 400～500 ℃ 而不是更低温度？(3) 为何通入过量空气？(4) 尾气为何要处理？

\item \laiyuan{2024年广东卷第17题(4)，来源：仓库所附真题文档}\ 兴趣小组探究反应 \ce{CaSO4(s) + CO3^2-(aq) <=> CaCO3(s) + SO4^2-(aq)}（反应Ⅰ）平衡的建立：将 1.00 g 硫酸钙（$M=136$ g·mol$^{-1}$）加入 100.0 mL 0.100 mol·L$^{-1}$ \ce{Na2CO3} 溶液中，在 25 ℃ 和搅拌条件下，用 pH 计测得体系 pH 随时间先减小，$t_2$ min 后不变。甲同学认为反应Ⅰ已达平衡；乙同学认为证据不足，提出：假设1 硫酸钙固体已完全消耗；假设2 硫酸钙有剩余，但被碳酸钙沉淀包裹。(1) 乙设计实验：将反应混合物固液分离，取少量沉淀置于试管中，滴加\tk，现象为\tk，沉淀完全溶解；继续向试管中滴加\tk，无白色沉淀生成。由此得出假设1成立。(2) 从数据角度说明假设1为何成立。(3) 写出优化的实验方案，并给出反应Ⅰ平衡已建立的判断依据。

\item \laiyuan{原创}\ 在恒定总压 20 MPa、一定温度下，向反应器中充入 1 mol \ce{N2} 和 3 mol \ce{H2}，发生 \ce{N2(g) + 3H2(g) <=> 2NH3(g)}，平衡时 \ce{NH3} 的物质的量分数为 25\%。(1) 求 \ce{H2} 的平衡转化率。(2) 求该温度下的 $K_p$（保留三位有效数字，单位 MPa$^{-2}$）。(3) 若改为在恒温恒容条件下进行（起始压强仍为 20 MPa），\ce{NH3} 的平衡物质的量分数将\tk（填“大于”“小于”或“等于”）25\%。

\item \laiyuan{原创}\ \ce{CO2(g) + 3H2(g) <=> CH3OH(g) + H2O(g)} $\Delta H<0$。在压强一定时，按不同投料比 $m=\dfrac{n(\ce{H2})}{n(\ce{CO2})}$（$m_1$、$m_2$、$m_3$）进行实验，\ce{CO2} 的平衡转化率随温度变化的三条曲线均随温度升高而下降，且同一温度下 $m_1$ 对应的转化率最高、$m_3$ 最低。(1) 比较 $m_1$、$m_2$、$m_3$ 的大小。(2) 解释曲线随温度升高而下降的原因。(3) 若同一投料比下，曲线上 a 点（$T_a$）与 b 点（$T_b$，$T_b>T_a$）的平衡常数分别为 $K_a$、$K_b$，比较大小。
\end{xiti}

\zu{C组\quad 冲刺高分}
\begin{xiti}
\item \laiyuan{原创}\ 乙苯脱氢制苯乙烯：\ce{C6H5CH2CH3(g) <=> C6H5CH=CH2(g) + H2(g)} $\Delta H>0$。某温度、总压恒为 100 kPa 时，以纯乙苯为原料，乙苯的平衡转化率为 40\%。(1) 求该温度下的 $K_p$。(2) 保持温度和总压不变，改为通入 $n(\ce{H2O}):n(\text{乙苯})=9:1$ 的混合气（水蒸气不参与反应），求乙苯的平衡转化率（保留两位有效数字）。(3) 解释工业上掺入大量水蒸气的原因（至少答两点）。

\item \laiyuan{2025年广东卷第19题(4)，来源：仓库所附真题文档}\ 钛基催化剂可催化储氢物质肼（\ce{N2H4}）的分解：(e)\ce{N2H4 = N2 + 2H2}；(f)\ce{3N2H4 = N2 + 4NH3}。以肼的水溶液为原料（含 \ce{N2H4} 的物质的量为 $n_0$）进行实验，得到 $n_1/n_0$、$n_2/n_0$ 随时间 $t$ 的变化曲线，其中 $n_1$ 为 \ce{H2} 与 \ce{N2} 的物质的量之和，$n_2$ 为剩余 \ce{N2H4} 的物质的量；$t_0$ min 时 $n_1/n_0=y_1$，$n_2/n_0=y_2$。设 $n_\text{e}$ 为 $0\sim t$ 时间段内反应(e)消耗 \ce{N2H4} 的物质的量，该时间段内催化剂的选择性用 $\dfrac{n_\text{e}}{n_0-n_2}\times100\%$ 表示。①$0\sim t_0$ min 内，\ce{N2H4} 的转化率为\tk（用含 $y_2$ 的代数式表示）。②$0\sim t_0$ min 内，催化剂的选择性为\tk（用含 $y_1$ 与 $y_2$ 的代数式表示，写出推导过程）。

\item \laiyuan{2026年广东卷第19题(5)，来源：仓库所附真题文档（网络收集版）}\ 某低共熔溶剂 D(l) 可用于吸收 HCl(g)，吸收过程可表示为 \ce{HCl(g) <=> HCl(sol)}，\ce{$x$HCl(sol) + D(l) <=> D*$x$HCl(sol)}。在一定 $T$ 和 HCl(g) 压强 $p$ 下，吸收至饱和时，每摩尔 D 吸收 HCl(g) 的量为 $\alpha+\beta$，其中 $\alpha=\dfrac{n(\ce{HCl})}{n_0(\mathrm{D})}$，$\beta=\dfrac{x\cdot n(\mathrm{DH})}{n_0(\mathrm{D})}$（$x$ 为仅与温度有关的系数；$n_0(\mathrm{D})$ 为 D(l) 的初始物质的量；$n(\ce{HCl})$ 为平衡时 HCl(sol) 的物质的量；$n(\mathrm{DH})$ 为平衡时 \ce{D*$x$HCl(sol)} 的物质的量）。由图读得：303 K、$p=100$ kPa 时 $\alpha=0.15$、$\beta=0.86$；350 K 时 $\alpha$ 始终为 0.00，$p=40$ kPa 时 $\beta=0.19$，$p=100$ kPa 时 $\beta=0.25$。
①1 mol D 在 303 K、$p=100$ kPa 下吸收 HCl(g) 至饱和，再升温至 350 K、降压至 $p=40$ kPa，达平衡后可释放\tk mol HCl(g)。
②形成 \ce{D*$x$HCl(sol)} 的总过程为 \ce{$x$HCl(g) + D(l) <=> D*$x$HCl(sol)}，其平衡常数 $K=\dfrac{n(\mathrm{DH})}{n(\mathrm{D})\cdot(p/100\,\mathrm{kPa})^x}$。已知 350 K 时 $K=1.0$，则在 350 K、$p=64$ kPa 下，1 mol D 最多可吸收\tk mol HCl(g)（写出计算过程，结果保留两位有效数字）。

\item \laiyuan{原创}\ 一定温度下，在恒容密闭容器中充入 \ce{NO2}，起始压强为 100 kPa，发生 \ce{2NO2(g) <=> N2O4(g)}，该反应可视为基元反应：$\vz=\kz p^2(\ce{NO2})$，$\vn=\kn p(\ce{N2O4})$。达到平衡时总压为 70 kPa。(1) 求平衡时 \ce{NO2}、\ce{N2O4} 的分压及 $K_p$。(2) 用 $\kz$、$\kn$ 表示 $K_p$。(3) 当 \ce{NO2} 的分压为 60 kPa 时，$\dfrac{\vz}{\vn}=$\tk。(4) 升高温度，平衡体系中 \ce{NO2} 的体积分数增大，判断该反应的 $\Delta H$ 正负。

\item \laiyuan{改编自选必1课件“演示文稿1”中的练习，数据为本书设定}\ 羰基法提纯镍：\ce{Ni(s) + 4CO(g) <=> Ni(CO)4(g)}。在不同总压下，平衡时 \ce{Ni(CO)4} 的体积分数 $x$ 随温度升高而减小，且同一温度下总压越大 $x$ 越大。(1) 判断该反应的 $\Delta H$\tk0，从热力学角度考虑，\tk{}有利于 \ce{Ni(CO)4} 的生成（写出两点）。(2) 设起始只充入 CO（Ni 足量），若总压为 1.0 MPa、100 ℃ 时 $x=0.90$，求 CO 的平衡转化率（保留三位有效数字）和该温度下的 $K_p$（单位 MPa$^{-3}$）。(3) 工业上先在约 50 ℃ 使粗镍与 CO 反应生成 \ce{Ni(CO)4}，再将气体送到约 230 ℃ 的分解器得到纯镍。结合 (1) 说明其原理。
\end{xiti}

\zhangmo{练习解析}
\begin{jiexilist}
\item \ans{C}B 为固体，气体分子数 $2\to2$ 不变，恒容时压强始终不变，①不能作标志；气体质量随反应变化，密度是变量，②能；③浓度不变，能；④浓度比不能说明不变；⑤都表示正反应；⑥生成 D 为正、生成 A 为逆，$n:2n=1:2$ 与计量数比相符，能。选②③⑥。
\item \ans{D}$K$ 与温度有关，A 错；方程式写法不同 $K$ 不同，B 错；$K$ 用平衡浓度，C 错；正逆反应 $K$ 互为倒数，D 对。
\item \ans{$K^{1/2}$（即 $\sqrt K$）；$1/K$}计量数变为 $\frac12$，$K$ 开平方；逆反应取倒数。
\item \ans{B}催化剂不使平衡移动，不能用勒夏特列原理解释。A：\ce{HClO} 见光分解，\ce{Cl2 + H2O <=> H+ + Cl- + HClO} 正移，$c(\ce{H+})$ 增大；C：NaOH 溶解放热且 $c(\ce{OH-})$ 增大，使 \ce{NH3 + H2O <=> NH3*H2O <=> NH4+ + OH-} 逆移；D：加压平衡向气体分子数减小的方向移动。
\item \ans{B}①升温平衡逆移，转化率降低；②增大 \ce{O2} 浓度，\ce{SO2} 转化率提高；③恒容充 He，平衡不移动；④按比例加料相当于加压，平衡正移，转化率提高；⑤催化剂不影响平衡；⑥移走生成物，平衡正移。选②④⑥。
\item \ans{D}任何温度自发需 $\Delta H<0$ 且 $\Delta S>0$。A：吸热、气体分子数减少（熵减），任何温度不自发；B：放热、熵减，低温自发；C：吸热、熵增，高温自发；D：放热、生成气体熵增，任何温度自发。
\item \ans{B}A：压缩体积，\ce{NO2} 浓度瞬间增大，颜色变深，但平衡正向（气体分子数减小的方向）移动；B：升温，平衡向吸热的逆反应方向移动，\ce{NO2} 增多，颜色变深，正确；C：恒容充 Ar，各物质浓度不变，平衡不移动，颜色不变；D：恒压充 Ar，体积增大，相当于减压，平衡逆向移动，但 \ce{NO2} 浓度减小，颜色变浅。
\item (1) 起始 A 4 mol、B 2 mol（共 6 mol 气体）；转化 A 1.2 mol、B 0.8 mol、C 1.6 mol、D 0.8 mol。平衡气体总量 $6\times\frac45=4.8$ mol。若 C、D 均为气体，平衡时 $6-2.0+2.4=6.4$ mol，不符；若只有 D 为气体，$6-2.0+0.8=4.8$ mol，符合。故 C 为非气态（固体或液体），D 为气体，$K=\dfrac{c^2(\mathrm{D})}{c^3(\mathrm{A})\cdot c^2(\mathrm{B})}$。(2) $\alpha(\mathrm{B})=0.8/2=40\%$；平衡浓度 $c(\mathrm{A})=1.4$、$c(\mathrm{B})=0.6$、$c(\mathrm{D})=0.4$ mol·L$^{-1}$，$K=\dfrac{0.4^2}{1.4^3\times0.6^2}\approx0.16$。(3) C 非气态，增加 C 的量平衡不移动，A、B 转化率不变。
\item (1) 曲线 a、b 压强相同，a 先拐，$T_2>T_1$；曲线 b、c 温度相同，b 先拐，$p_2>p_1$。(2) 同为 $p_2$，高温（a）C\% 低，升温平衡逆移，$\Delta H<0$；同为 $T_1$，高压（b）C\% 高，加压平衡正移，$m+n>p$。(3) 投料比 $m$ 越小，\ce{O2} 相对越多，\ce{SO2} 转化率越高，最上方曲线 $m$ 最小：$m_1<m_2<m_3$。
\item (1) 常压下 \ce{SO2} 平衡转化率已很高（450 ℃ 时 97.5\%），加压提高有限，而对设备要求和能耗大幅增加，不经济。(2) 温度过低反应速率太慢，且催化剂（\ce{V2O5}）在 400～500 ℃ 活性最高；温度过高平衡转化率下降。(3) 增大 \ce{O2} 浓度，使平衡正移，提高 \ce{SO2} 的转化率（空气廉价易得）。(4) \ce{SO2} 有毒、会形成酸雨污染环境，回收还可提高原料利用率。
\item (1) 过量稀盐酸；有气体产生；\ce{BaCl2} 溶液。（\ce{CaCO3} 溶于盐酸放出 \ce{CO2}；若沉淀中有 \ce{CaSO4}，溶解后加 \ce{BaCl2} 会生成 \ce{BaSO4} 白色沉淀，无沉淀说明无 \ce{CaSO4}。）(2) $n(\ce{CaSO4})=1.00/136\approx7.4\times10^{-3}$ mol，$n(\ce{CO3^2-})=0.100\times0.100=1.00\times10^{-2}$ mol，\ce{CaSO4} 不足；反应Ⅰ的平衡常数很大，\ce{CaSO4} 可完全转化。pH 不变只是因为一种反应物耗尽，并非平衡。(3) 要保证 \ce{CaSO4} 有剩余：按完全反应计算需 $0.0100\times136=1.36$ g，故将明显多于 1.36 g（如 2.00 g）的硫酸钙加入 100.0 mL 0.100 mol·L$^{-1}$ \ce{Na2CO3} 溶液中，在 25 ℃、搅拌条件下用 pH 计测 pH；当 pH 不变时，加入 \ce{Na2SO4} 固体，若 pH 增大（\ce{SO4^2-} 浓度增大使反应Ⅰ逆移，$c(\ce{CO3^2-})$ 增大，碱性增强），说明反应Ⅰ的平衡已建立（体现了“平衡可以被打破并移动”这一特征）。
\item (1) 设 \ce{N2} 转化 $x$ mol：平衡时 \ce{N2} $(1-x)$、\ce{H2} $(3-3x)$、\ce{NH3} $2x$，总 $(4-2x)$ mol。$\dfrac{2x}{4-2x}=0.25$，$x=0.4$；$\alpha(\ce{H2})=\dfrac{1.2}{3}=40\%$。(2) $n_\text{总}=3.2$ mol，物质的量分数：\ce{N2} 0.1875、\ce{H2} 0.5625、\ce{NH3} 0.25。$K_p=\dfrac{(0.25\times20)^2}{(0.1875\times20)\times(0.5625\times20)^3}\approx4.68\times10^{-3}$ MPa$^{-2}$。(3) 小于。恒容时随反应进行气体分子数减少、压强降低，相当于在比恒压时更低的压强下建立平衡，平衡转化率较低。
\item (1) $m_1>m_2>m_3$：\ce{H2} 相对越多，\ce{CO2} 的转化率越高。(2) 正反应放热，升温平衡逆向移动，\ce{CO2} 平衡转化率下降。(3) $K_a>K_b$（放热反应，温度越高 $K$ 越小；$K$ 只与温度有关，与投料比无关）。
\item (1) 设乙苯 1 mol，平衡时乙苯 0.6 mol、苯乙烯 0.4 mol、\ce{H2} 0.4 mol，$n_\text{总}=1.4$ mol。$K_p=\dfrac{(\frac{0.4}{1.4}p)^2}{\frac{0.6}{1.4}p}=\dfrac{0.16}{0.84}\times100\approx19$ kPa。(2) 设转化率为 $\alpha$，$n_\text{总}=10+\alpha$：$K_p=\dfrac{\alpha^2p}{(1-\alpha)(10+\alpha)}=19.05$，整理得 $119.05\alpha^2+171.45\alpha-190.5=0$，$\alpha\approx0.74$，即约 74\%（远高于 40\%）。(3) ①恒压下水蒸气稀释，降低各反应物、产物的分压，相当于减压，平衡向气体分子数增大的正方向移动，提高乙苯转化率；②水蒸气热容大，可为吸热反应供热、维持反应温度；③水蒸气可与催化剂表面积炭反应（\ce{C + H2O -> CO + H2}），减少积炭，延长催化剂寿命。
\item ①$(1-y_2)\times100\%$。起始 $n_2/n_0=1$，$t_0$ 时为 $y_2$，转化了 $n_0(1-y_2)$。②设反应(e)消耗 \ce{N2H4} $n_\text{e}$，反应(f)消耗 $n_\text{f}$：$n_\text{e}+n_\text{f}=n_0(1-y_2)$；(e)生成 $n_\text{e}$ mol \ce{N2} 和 $2n_\text{e}$ mol \ce{H2}，(f)生成 $\frac13n_\text{f}$ mol \ce{N2}：$3n_\text{e}+\frac13n_\text{f}=y_1n_0$。将第一式乘 $\frac13$ 后与第二式相减：$\frac83n_\text{e}=n_0\left(y_1-\frac{1-y_2}{3}\right)$，$n_\text{e}=\dfrac{(3y_1+y_2-1)n_0}{8}$。选择性 $=\dfrac{n_\text{e}}{n_0(1-y_2)}\times100\%=\dfrac{3y_1+y_2-1}{8-8y_2}\times100\%$。
\item ①0.82。303 K、100 kPa 饱和时每摩尔 D 吸收 $\alpha+\beta=0.15+0.86=1.01$ mol；350 K、40 kPa 平衡时为 $0+0.19=0.19$ mol；释放 $1.01-0.19=0.82$ mol。②0.22。350 K 时 $\alpha=0$，HCl 只以 \ce{D*$x$HCl} 形式被吸收。$p=100$ kPa 时 $\beta=0.25$：$n(\mathrm{DH})=\dfrac{0.25}{x}$ mol，$n(\mathrm{D})=1-\dfrac{0.25}{x}$，$K=\dfrac{0.25/x}{(1-0.25/x)\times1^x}=1.0$，解得 $\dfrac{0.25}{x}=0.5$，$x=0.5$。$p=64$ kPa 时：$\dfrac{n(\mathrm{DH})}{[1-n(\mathrm{DH})]\times0.64^{0.5}}=1.0$，即 $n(\mathrm{DH})=0.8[1-n(\mathrm{DH})]$，$n(\mathrm{DH})=\dfrac49$ mol；吸收 HCl 的量 $=\beta=x\cdot n(\mathrm{DH})=0.5\times\dfrac49\approx0.22$ mol。
\item (1) 设生成 \ce{N2O4} 的分压为 $x$：$p(\ce{NO2})=100-2x$，总压 $100-x=70$，$x=30$ kPa；$p(\ce{NO2})=40$ kPa，$p(\ce{N2O4})=30$ kPa，$K_p=\dfrac{30}{40^2}=0.01875$ kPa$^{-1}$。(2) 平衡时 $\vz=\vn$，$K_p=\dfrac{\kz}{\kn}$。(3) 此时 $p(\ce{N2O4})=\dfrac{100-60}{2}=20$ kPa，$\dfrac{\vz}{\vn}=\dfrac{\kz\times60^2}{\kn\times20}=0.01875\times180=3.375$（$>1$，反应正向进行）。(4) 升温后 \ce{NO2} 的体积分数增大，说明平衡逆向移动，逆反应吸热，正反应放热，$\Delta H<0$。
\item (1) $<$；低温、高压。升温 $x$ 减小，平衡逆移，正反应放热；加压 $x$ 增大，有利于气体分子数减小的正反应。(2) 设起始 CO 为 1 mol，生成 \ce{Ni(CO)4} $y$ mol：剩余 CO $(1-4y)$ mol，气体总量 $(1-3y)$ mol。$\dfrac{y}{1-3y}=0.90$，$y=\dfrac{0.90}{3.7}\approx0.243$；$\alpha(\ce{CO})=4y\approx97.3\%$。$p(\ce{Ni(CO)4})=0.90$ MPa，$p(\ce{CO})=0.10$ MPa，$K_p=\dfrac{0.90}{0.10^4}=9.0\times10^3$ MPa$^{-3}$。(3) 低温时平衡正移，粗镍中的 Ni 转化为气态 \ce{Ni(CO)4} 与杂质分离；高温时平衡逆移，\ce{Ni(CO)4} 分解析出纯镍，CO 可循环使用——利用温度对平衡的影响实现“搬运”提纯。
\end{jiexilist}
% <<<<<<<< 结束：chapters/ch03ex.tex
% <<<<<<<< 结束：chapters/ch03.tex

% >>>>>>>> 开始：chapters/ch04.tex
% ============================================================
% 第4章 水溶液中的离子反应与平衡（选必1第三章，课件未覆盖，补充）
% ============================================================
\chapter{水溶液中的离子反应与平衡}

\begin{daoxue}
\textbf{知识主线：}把第3章的平衡思想用到水溶液中——弱电解质的\textbf{电离平衡}（$K_\text{a}$、$K_\text{b}$）$\to$ 水的\textbf{电离平衡}（$K_\text{w}$、pH、中和滴定）$\to$ 盐的\textbf{水解平衡}（$K_\text{h}$）$\to$ 难溶电解质的\textbf{沉淀溶解平衡}（$K_\text{sp}$）。四种平衡常数都只与温度有关，都可用“$Q$ 与 $K$”判断方向。

\textbf{高考考情：}广东卷每年都在大题中考查：2025年第17题用 pH 滴定曲线求苯甲酸的 $K_\text{a}$、设计测 HClO 的 $K_\text{a}$ 的方案；2026年第17题用 pH–$V$ 突跃求反应比例；2024年第19题求非水溶剂中弱酸的 $K_\text{a}$；工艺流程题（第18题）每年都有 $K_\text{sp}$ 计算与调 pH。选择题考电荷守恒、配离子平衡（2025第12题）、盐的水解（2026第7题D）。

\textbf{学习目标：}①会分析电离、水解平衡的移动；②熟练进行 pH、$K_\text{a}$、$K_\text{h}$、$K_\text{sp}$ 计算；③掌握三大守恒并会比较粒子浓度；④会读滴定曲线、分布系数图、沉淀溶解平衡图。
\end{daoxue}

\section{弱电解质的电离平衡}

\subsection{强电解质与弱电解质}
在水溶液中\textbf{全部电离}的电解质是强电解质（强酸、强碱、绝大多数盐）；\textbf{部分电离}的是弱电解质（弱酸、弱碱、水等）。强电解质电离用“$=$”，弱电解质用“\ce{<=>}”，多元弱酸分步电离（以第一步为主），多元弱碱一步写出：
\[\ce{H2CO3 <=> H+ + HCO3-},\quad \ce{HCO3- <=> H+ + CO3^2-};\qquad \ce{Fe(OH)3 <=> Fe^3+ + 3OH-}\]
\ce{NaHSO4} 在水中：\ce{NaHSO4 = Na+ + H+ + SO4^2-}；熔融时：\ce{NaHSO4 = Na+ + HSO4-}。\ce{NaHCO3 = Na+ + HCO3-}（\ce{HCO3-} 再部分电离）。

\begin{yicuo}
溶液导电能力取决于\textbf{离子浓度和离子所带电荷}，与电解质强弱无直接关系：很稀的盐酸导电性可能不如较浓的醋酸。\ce{BaSO4} 难溶，但溶解的部分全部电离，是强电解质。
\end{yicuo}

\subsection{电离平衡及其移动}
以 \ce{CH3COOH <=> CH3COO- + H+}（$\Delta H>0$）为例，平衡移动遵循勒夏特列原理：

\begin{center}
\small
\begin{tabular}{lccccc}
\toprule
改变条件 & 平衡移动 & $n(\ce{H+})$ & $c(\ce{H+})$ & 电离度 $\alpha$ & $K_\text{a}$ \\
\midrule
加水稀释 & 正向 & 增大 & 减小 & 增大 & 不变 \\
升高温度 & 正向 & 增大 & 增大 & 增大 & 增大 \\
加少量 \ce{CH3COONa} 固体 & 逆向 & 减小 & 减小 & 减小 & 不变 \\
通入少量 HCl & 逆向 & 增大 & 增大 & 减小 & 不变 \\
加少量 NaOH 固体 & 正向 & 减小 & 减小 & 增大 & 不变 \\
加冰醋酸 & 正向 & 增大 & 增大 & 减小 & 不变 \\
\bottomrule
\end{tabular}
\end{center}

\begin{benzhi}[稀释时为什么平衡正移，$c(\ce{H+})$ 却减小]
稀释到原体积的 2 倍，各粒子浓度瞬间都减半：$Q=\dfrac{\frac12c(\ce{CH3COO-})\cdot\frac12c(\ce{H+})}{\frac12c(\ce{CH3COOH})}=\dfrac{K_\text{a}}{2}<K_\text{a}$，平衡正向移动（“越稀越电离”）。但移动只能部分补偿，$c(\ce{H+})$ 仍比原来小。
\end{benzhi}

\subsection{电离平衡常数}
\begin{gainian}[电离常数]
\[\ce{HA <=> H+ + A-}\qquad K_\text{a}=\frac{c(\ce{H+})\cdot c(\ce{A-})}{c(\ce{HA})}\]
$K_\text{a}$（$K_\text{b}$）只与温度有关，升温增大。$K_\text{a}$ 越大，酸性越强。多元弱酸 $K_{\text{a}1}\gg K_{\text{a}2}$，酸性主要由第一步决定。
\end{gainian}

\begin{fangfa}[$K_\text{a}$ 的三个用途]
\textbf{①计算：}$c$ mol·L$^{-1}$ 的弱酸 HA，若电离程度很小，$c(\ce{H+})\approx\sqrt{K_\text{a}\cdot c}$。如 0.1 mol·L$^{-1}$ \ce{CH3COOH}（$K_\text{a}=1.75\times10^{-5}$），$c(\ce{H+})\approx1.3\times10^{-3}$ mol·L$^{-1}$，pH$\approx2.9$。\\
\textbf{②判断反应能否发生（强酸制弱酸）：}只有 $K_\text{a}$ 大的酸能制 $K_\text{a}$ 小的酸。已知 $K_\text{a}$：\ce{H2CO3}（$K_{\text{a}1}=4.5\times10^{-7}$，$K_{\text{a}2}=4.7\times10^{-11}$）$>$ \ce{HClO}（$4.0\times10^{-8}$）$>$ \ce{HCO3-}，所以向 NaClO 溶液中通入 \ce{CO2}，无论多少，都只生成 \ce{HCO3-}：\ce{ClO- + CO2 + H2O = HClO + HCO3-}；苯酚钠溶液中通 \ce{CO2} 同理只生成 \ce{HCO3-}。\\
\textbf{③判断比值变化：}$\dfrac{c(\ce{A-})}{c(\ce{HA})}=\dfrac{K_\text{a}}{c(\ce{H+})}$。如稀释醋酸，$c(\ce{H+})$ 减小，$\dfrac{c(\ce{CH3COO-})}{c(\ce{CH3COOH})}$ 增大。遇到“比值如何变化”，先把比值改写成“常数/某一浓度”。
\end{fangfa}

\begin{liti}[电离常数与电离度的计算]
25 ℃ 时，0.10 mol·L$^{-1}$ 一元弱酸 HA 溶液的 pH$=3.0$。(1) 求 HA 的电离度 $\alpha$ 和 $K_\text{a}$；(2) 将该溶液稀释到 0.0010 mol·L$^{-1}$，估算此时的电离度，说明“越稀越电离”。
\tcblower
\jie (1) $c(\ce{H+})=c(\ce{A-})=10^{-3}$ mol·L$^{-1}$，$\alpha=\dfrac{10^{-3}}{0.10}=1.0\%$；$K_\text{a}=\dfrac{10^{-3}\times10^{-3}}{0.10-10^{-3}}\approx1.0\times10^{-5}$（电离程度很小，分母近似为 0.10）。\\
(2) $K_\text{a}$ 不变，$c(\ce{H+})\approx\sqrt{1.0\times10^{-5}\times0.0010}=1.0\times10^{-4}$ mol·L$^{-1}$，$\alpha\approx\dfrac{10^{-4}}{0.0010}=10\%$。浓度变为原来的 $\frac1{100}$，电离度增大到 10 倍（$\alpha\approx\sqrt{K_\text{a}/c}$），但 $c(\ce{H+})$ 只有原来的 $\frac1{10}$。（此时 $\alpha$ 已超过 5\%，精确解一元二次方程得 $c(\ce{H+})\approx9.5\times10^{-5}$ mol·L$^{-1}$，$\alpha\approx9.5\%$，结论不变。）\\
\dianbo 近似公式 $c(\ce{H+})\approx\sqrt{K_\text{a}c}$ 适用于 $\alpha<5\%$ 左右；电离度较大时要解一元二次方程。
\end{liti}

\begin{center}
\small
\begin{tabularx}{\linewidth}{lXX}
\toprule
比较项目 & 等物质的量浓度的盐酸与醋酸 & 等 pH 的盐酸与醋酸 \\
\midrule
$c(\ce{H+})$ & 盐酸大 & 相等 \\
中和 NaOH 的能力 & 相等 & 醋酸大（醋酸浓度大得多） \\
与足量 Zn 反应 & 开始盐酸快，产生 \ce{H2} 总量相等 & 开始速率相等，后来醋酸快；醋酸产生 \ce{H2} 多 \\
稀释相同倍数后 pH & 盐酸 pH 小 & 盐酸 pH 变化大，醋酸 pH 变化小 \\
\bottomrule
\end{tabularx}
\end{center}

\begin{zhenti}{2025年广东卷第17题(3)\quad 来源：仓库所附真题文档}
测定常温下苯甲酸饱和溶液的浓度 $c_0$ 和苯甲酸的 $K_\text{a}$：取 50.00 mL 苯甲酸饱和溶液，用 0.1000 mol·L$^{-1}$ NaOH 溶液滴定，用 pH 计测得 pH 随 $V$(NaOH) 的变化曲线：起点 pH$=2.89$，加入 7.00 mL 时 pH$=4.20$，加入 14.00 mL 时出现突跃。①$c_0=$\tk mol·L$^{-1}$；②苯甲酸的 $K_\text{a}=$\tk（列出算式，水的电离可忽略）。

\textbf{解析：}①突跃处恰好中和：$c_0=\dfrac{0.1000\times14.00}{50.00}=0.028$ mol·L$^{-1}$。②起点即饱和苯甲酸溶液，$c(\ce{H+})=c(\text{苯甲酸根})=10^{-2.89}$ mol·L$^{-1}$，$K_\text{a}=\dfrac{10^{-2.89}\times10^{-2.89}}{0.028-10^{-2.89}}$。\\
\dianbo 加入 7.00 mL（半中和点）时，溶液中苯甲酸与苯甲酸根浓度近似相等，$K_\text{a}\approx c(\ce{H+})=10^{-4.20}$，这是读滴定曲线求 $K_\text{a}$ 的另一常用方法。两种方法结果略有差异，源于近似处理，按题目要求作答。
\end{zhenti}

\section{水的电离与溶液的酸碱性}

\subsection{水的离子积}
水是极弱的电解质：\ce{H2O <=> H+ + OH-}（$\Delta H>0$）。\textbf{水的离子积} $K_\text{w}=c(\ce{H+})\cdot c(\ce{OH-})$，25 ℃ 时 $K_\text{w}=1.0\times10^{-14}$，100 ℃ 时约为 $5.5\times10^{-13}$（习题中常近似取 $1.0\times10^{-12}$，本节图像题即按此取值）。$K_\text{w}$ 只与温度有关，适用于任何稀水溶液。

\begin{center}
\small
\begin{tabular}{lll}
\toprule
条件 & 对水的电离 & 实例（25 ℃，计算水电离出的 $c(\ce{H+})_\text{水}$） \\
\midrule
加酸或碱 & 抑制 & 0.01 mol·L$^{-1}$ HCl：$c(\ce{OH-})_\text{水}=c(\ce{H+})_\text{水}=10^{-12}$ mol·L$^{-1}$ \\
加能水解的盐 & 促进 & pH$=5$ 的 \ce{NH4Cl} 溶液：$c(\ce{H+})_\text{水}=10^{-5}$ mol·L$^{-1}$ \\
升高温度 & 促进 & $K_\text{w}$ 增大，但纯水仍呈中性 \\
加活泼金属（如 Na） & 促进 & 消耗 \ce{H+} \\
\bottomrule
\end{tabular}
\end{center}

\begin{yicuo}
判断溶液酸碱性的根本依据是 $c(\ce{H+})$ 与 $c(\ce{OH-})$ 的相对大小，而不是“pH 是否等于 7”。100 ℃ 时纯水的 pH 约为 6（$<7$），但仍呈中性。“由水电离出的 $c(\ce{H+})=10^{-12}$ mol·L$^{-1}$”的溶液可能是 pH$=2$ 的酸，也可能是 pH$=12$ 的碱。
\end{yicuo}

\begin{liti}[水的电离平衡曲线]
水的电离平衡曲线如图，A 点（$10^{-7}$，$10^{-7}$）在 25 ℃ 曲线上，B 点（$10^{-6}$，$10^{-6}$）在 100 ℃ 曲线上。(1) 100 ℃ 时 $K_\text{w}=$\tk。(2) 100 ℃ 时，pH$=6$ 的溶液呈\tk{}性。(3) 25 ℃ 时，从 A 点变到曲线上 C 点（$c(\ce{H+})=10^{-6}$，$c(\ce{OH-})=10^{-8}$）可采用的方法是\tk。(4) 能否通过升温使溶液由 A 点变到 C 点？
\begin{center}
\begin{tikzpicture}
\begin{axis}[width=5.6cm,height=4.4cm,xlabel={$c(\ce{H+})$/(mol·L$^{-1}$)},ylabel={$c(\ce{OH-})$/(mol·L$^{-1}$)},xmode=log,ymode=log,xmin=1e-9,xmax=1e-4,ymin=1e-9,ymax=1e-4,tick label style={font=\scriptsize},label style={font=\scriptsize},legend style={font=\scriptsize,at={(1.03,0.5)},anchor=west}]
\addplot[thick,mainblue,domain=1e-9:1e-5,samples=40]{1e-14/x};
\addplot[thick,warnred,domain=1e-8:1e-4,samples=40]{1e-12/x};
\addplot[only marks,mark=*,mark size=1.5pt] coordinates {(1e-7,1e-7) (1e-6,1e-6) (1e-6,1e-8)};
\node[font=\scriptsize,above right] at (axis cs:1e-7,1e-7) {A};
\node[font=\scriptsize,above right] at (axis cs:1e-6,1e-6) {B};
\node[font=\scriptsize,above right] at (axis cs:1e-6,1e-8) {C};
\legend{25 ℃,100 ℃}
\end{axis}
\end{tikzpicture}
\end{center}
\tcblower
\jie (1) $10^{-6}\times10^{-6}=1.0\times10^{-12}$。(2) 中性（100 ℃ 时中性溶液 pH$=6$）。(3) C 点仍在 25 ℃ 曲线上，$c(\ce{H+})>c(\ce{OH-})$，呈酸性，可加入少量酸（或强酸弱碱盐如 \ce{NH4Cl}）。(4) 不能。升温使 $K_\text{w}$ 增大，纯水中 $c(\ce{H+})$、$c(\ce{OH-})$ 同时增大（沿对角线移向 B 点），而 C 点 $c(\ce{OH-})$ 减小。
\end{liti}

\subsection{pH 的计算}
pH$=-\lg c(\ce{H+})$。计算思路：\textbf{酸性溶液先求 $c(\ce{H+})$，碱性溶液先求 $c(\ce{OH-})$ 再由 $K_\text{w}$ 换算}。

\begin{enumerate}
\item \textbf{稀释：}pH$=a$ 的强酸稀释 $10^n$ 倍，pH$=a+n$（但无限稀释只能接近 7，不能超过 7）；弱酸稀释 $10^n$ 倍，$a<$pH$<a+n$。
\item \textbf{强酸与强碱混合：}先判断谁过量。如 25 ℃ 时 pH$=2$ 的盐酸与 pH$=12$ 的 NaOH 等体积混合，pH$=7$；pH$=2$ 的醋酸与 pH$=12$ 的 NaOH 等体积混合，醋酸大量过量，溶液呈酸性。
\item \textbf{两强碱混合：}必须先算 $c(\ce{OH-})$ 的平均值，不能直接平均 $c(\ce{H+})$。
\end{enumerate}

\begin{liti}[pH 的计算]
25 ℃ 时：(1) 0.005 mol·L$^{-1}$ \ce{H2SO4} 溶液的 pH；(2) pH$=10$ 与 pH$=12$ 的 NaOH 溶液等体积混合后的 pH（$\lg5\approx0.7$）；(3) pH$=13$ 的 NaOH 溶液与 pH$=2$ 的盐酸混合后溶液呈中性，求二者体积比。
\tcblower
\jie (1) $c(\ce{H+})=0.01$ mol·L$^{-1}$，pH$=2$。\\
(2) 必须先算 $c(\ce{OH-})$：$\dfrac{10^{-4}+10^{-2}}{2}\approx5\times10^{-3}$ mol·L$^{-1}$，pOH$\approx2.3$，pH$\approx11.7$。（若错误地平均 $c(\ce{H+})$，会得到 pH$\approx10.3$。）\\
(3) $10^{-1}V_\text{碱}=10^{-2}V_\text{酸}$，$V_\text{碱}:V_\text{酸}=1:10$。
\end{liti}

\subsection{酸碱中和滴定}
\textbf{原理：}$c_\text{待}=\dfrac{c_\text{标}\cdot V_\text{标}}{V_\text{待}}$（一元酸碱）。\textbf{关键}：准确测量体积（滴定管精确到 0.01 mL）、准确判断终点（指示剂）。

\textbf{操作要点：}检漏 $\to$ 洗涤 $\to$ \textbf{用待装液润洗 2～3 次} $\to$ 装液 $\to$ 赶气泡 $\to$ 调液面、读数 $\to$ 滴定（左手控制活塞，右手摇锥形瓶，眼睛注视锥形瓶中溶液颜色变化）$\to$ 读数 $\to$ 重复 2～3 次取平均值。锥形瓶\textbf{不能}用待测液润洗。

\textbf{指示剂选择：}指示剂变色范围应落在滴定突跃范围内。强碱滴定强酸：酚酞或甲基橙均可；强碱滴定弱酸（终点溶液呈碱性）：酚酞；强酸滴定弱碱（终点呈酸性）：甲基橙。

\begin{center}
\begin{tikzpicture}
\begin{axis}[width=7.8cm,height=5.2cm,xlabel={$V(\ce{NaOH})$/mL},ylabel={pH},xmin=0,xmax=40,ymin=0,ymax=14,
 ytick={1,4.76,7,8.73,13},xtick={0,10,20,30,40},grid=major,grid style={gray!20},legend pos=south east,legend style={font=\scriptsize},tick label style={font=\scriptsize},label style={font=\small}]
\addplot[thick,mainblue,smooth] coordinates {(0,1.00) (2,1.09) (5,1.22) (10,1.48) (15,1.85) (18,2.28) (19,2.59) (19.5,2.90) (19.8,3.30) (19.96,4.00) (20,7.00) (20.04,10.00) (20.2,10.70) (20.5,11.09) (21,11.39) (22,11.68) (25,12.05) (30,12.30) (35,12.44) (40,12.52)};
\addplot[thick,warnred,smooth,dashed] coordinates {(0,2.88) (2,3.81) (5,4.28) (10,4.76) (15,5.23) (18,5.71) (19,6.04) (19.5,6.35) (19.8,6.75) (19.96,7.45) (20,8.73) (20.04,10.00) (20.2,10.70) (20.5,11.09) (21,11.39) (22,11.68) (25,12.05) (30,12.30) (35,12.44) (40,12.52)};
\legend{滴定 HCl,滴定 \ce{CH3COOH}}
\end{axis}
\end{tikzpicture}
\end{center}
\noindent\small 图：0.1000 mol·L$^{-1}$ NaOH 滴定 20.00 mL 0.1000 mol·L$^{-1}$ 盐酸与醋酸（计算值）。滴定醋酸时：起点 pH$\approx2.9$；半中和点（10.00 mL）pH$=\mathrm{p}K_\text{a}=4.76$；终点 pH$\approx8.7$（\ce{CH3COONa} 水解呈碱性），突跃范围窄，必须用酚酞。\normalsize

\begin{fangfa}[终点判断与误差分析]
\textbf{终点描述模板：}“当滴入最后半滴（一滴）××标准溶液时，溶液由\tk 色变为\tk 色，且半分钟内不恢复原色”。\\
\textbf{误差分析：}一切误差都归结为对 $V_\text{标}$ 的影响（$c_\text{待}\propto V_\text{标}$）：滴定管未用标准液润洗（标准液被稀释，$V_\text{标}$ 偏大，结果偏高）；锥形瓶用待测液润洗（偏高）；滴定前尖嘴有气泡、滴定后消失（偏高）；滴定后俯视读数（读数偏小，偏低）；锥形瓶中有少量蒸馏水（无影响）。
\end{fangfa}

\textbf{氧化还原滴定}同理：\ce{KMnO4} 自身作指示剂（终点溶液变为浅紫红色）；碘量法用淀粉作指示剂（\ce{I2 + 2S2O3^2- = 2I- + S4O6^2-}，终点蓝色褪去）。计算抓住\textbf{得失电子守恒}或关系式。

\section{盐类的水解}

\subsection{水解的本质与规律}
盐电离出的弱酸根离子或弱碱阳离子，与水电离出的 \ce{H+} 或 \ce{OH-} 结合生成弱电解质，\textbf{破坏了水的电离平衡}，使溶液中 $c(\ce{H+})\neq c(\ce{OH-})$。如 \ce{CH3COO- + H2O <=> CH3COOH + OH-}，\ce{NH4+ + H2O <=> NH3*H2O + H+}。

\begin{gainian}[水解规律与水解常数]
\textbf{有弱才水解，无弱不水解；越弱越水解，都弱都水解；谁强显谁性，同强显中性。}\\
水解常数：$K_\text{h}=\dfrac{K_\text{w}}{K_\text{a}}$（弱酸根）或 $\dfrac{K_\text{w}}{K_\text{b}}$（弱碱阳离子）。对应的酸越弱（$K_\text{a}$ 越小），$K_\text{h}$ 越大，水解程度越大。如 $K_\text{h}(\ce{CO3^2-})=\dfrac{K_\text{w}}{K_{\text{a}2}}=\dfrac{10^{-14}}{4.7\times10^{-11}}\approx2.1\times10^{-4}$。
\end{gainian}

\textbf{书写：}水解程度一般很小，用“\ce{<=>}”，不标“↑”“↓”；多元弱酸根分步写（\ce{CO3^2- + H2O <=> HCO3- + OH-} 为主）；多元弱碱阳离子一步写（\ce{Fe^3+ + 3H2O <=> Fe(OH)3 + 3H+}）。\textbf{相互促进且彻底的水解}用“$=$”并标沉淀、气体，如泡沫灭火器原理：\ce{Al^3+ + 3HCO3- = Al(OH)3 v + 3CO2 ^}。

\begin{benzhi}[酸式盐溶液为什么有酸有碱]
\ce{HA-} 既能电离（\ce{HA- <=> H+ + A^2-}，$K_{\text{a}2}$）又能水解（\ce{HA- + H2O <=> H2A + OH-}，$K_{\text{h}2}=K_\text{w}/K_{\text{a}1}$），比较二者即可：
\ce{HCO3-}：$K_{\text{h}2}=\dfrac{10^{-14}}{4.5\times10^{-7}}\approx2.2\times10^{-8}>K_{\text{a}2}=4.7\times10^{-11}$，水解为主，\ce{NaHCO3} 溶液呈碱性；
\ce{HSO3-}：$K_{\text{h}2}=\dfrac{10^{-14}}{1.4\times10^{-2}}\approx7\times10^{-13}<K_{\text{a}2}\approx6.0\times10^{-8}$，电离为主，\ce{NaHSO3} 溶液呈酸性。同理 \ce{NaH2PO4} 呈酸性，\ce{Na2HPO4} 呈碱性，\ce{NaHC2O4} 呈酸性。
\end{benzhi}

\begin{liti}[水解常数与盐溶液的 pH]
已知 25 ℃ 时 \ce{CH3COOH} 的 $K_\text{a}=1.75\times10^{-5}$。(1) 求 0.10 mol·L$^{-1}$ \ce{CH3COONa} 溶液的 pH（$\sqrt{5.7}\approx2.4$，$\lg2.4\approx0.38$）；(2) 说明 \ce{CH3COONa} 溶液中 $\dfrac{c(\ce{CH3COO-})}{c(\ce{CH3COOH})\cdot c(\ce{OH-})}$ 在加水稀释时如何变化。
\tcblower
\jie (1) $K_\text{h}=\dfrac{K_\text{w}}{K_\text{a}}=\dfrac{10^{-14}}{1.75\times10^{-5}}\approx5.7\times10^{-10}$；$c(\ce{OH-})\approx\sqrt{5.7\times10^{-10}\times0.10}=\sqrt{5.7}\times10^{-5.5}\approx7.6\times10^{-6}$ mol·L$^{-1}$，pOH$\approx5.12$，pH$\approx8.88$。\\
(2) 该比值 $=\dfrac{1}{K_\text{h}}$，温度不变则不变。\\
\dianbo 0.10 mol·L$^{-1}$ 的醋酸 pH$\approx2.9$，而其钠盐溶液 pH$\approx8.9$——“弱酸的电离”与“弱酸根的水解”程度都很小，这是比较粒子浓度大小的基础。
\end{liti}

\subsection{影响水解的因素与水解的应用}
水解是吸热过程：\textbf{升温}促进水解；\textbf{稀释}促进水解（水解程度增大，但 $c(\ce{OH-})$ 或 $c(\ce{H+})$ 减小）；加酸抑制弱碱阳离子水解，加碱抑制弱酸根水解。

\begin{center}
\small
\begin{tabularx}{\linewidth}{p{3.4cm}X}
\toprule
应用 & 原理与实例 \\
\midrule
配制易水解的盐溶液 & 配 \ce{FeCl3} 溶液先溶于浓盐酸再稀释；配 \ce{FeSO4} 溶液加少量稀硫酸（抑制水解）和铁粉（防氧化） \\
盐溶液蒸干产物 & \ce{FeCl3}、\ce{AlCl3}、\ce{MgCl2} 溶液蒸干得氢氧化物（HCl 挥发促进水解），灼烧得氧化物；\ce{Fe2(SO4)3}（\ce{H2SO4} 难挥发）蒸干仍得原物质；\ce{Na2SO3} 蒸干被氧化为 \ce{Na2SO4}；\ce{NaHCO3} 受热分解为 \ce{Na2CO3}。在\textbf{ HCl 气流中}加热 \ce{MgCl2*6H2O} 可得无水 \ce{MgCl2} \\
净水 & 明矾、\ce{FeCl3}：\ce{Al^3+ + 3H2O <=> Al(OH)3}（胶体）\ce{+ 3H+}，胶体吸附杂质 \\
热纯碱去油污 & 升温促进 \ce{CO3^2-} 水解，碱性增强 \\
除杂（调 pH） & 除去 \ce{MgCl2} 溶液中的 \ce{Fe^3+}：加 MgO、\ce{Mg(OH)2} 或 \ce{MgCO3} 消耗 \ce{H+}，促进 \ce{Fe^3+} 水解成 \ce{Fe(OH)3} 沉淀，且不引入新杂质 \\
化肥使用 & 草木灰（\ce{K2CO3}）不能与铵态氮肥混用（相互促进水解，放出 \ce{NH3}） \\
\bottomrule
\end{tabularx}
\end{center}

\begin{zhenti}{2025年广东卷第17题(5)\quad 来源：仓库所附真题文档}
小组讨论后，选用 0.100 mol·L$^{-1}$ NaClO 溶液（含少量 NaCl）测定弱酸 HClO 的 $K_\text{a}$（不用滴定法的原因之一是 HClO 见光易分解：\ce{2HClO ->[\text{光照}] 2H+ + 2Cl- + O2 ^}）。简述该方案（包括所用仪器及数据处理思路）。

\textbf{参考答案：}取适量 0.100 mol·L$^{-1}$ NaClO 溶液于烧杯中，用 pH 计测定 25 ℃ 时的 pH；由于 \ce{ClO-} 水解（\ce{ClO- + H2O <=> HClO + OH-}），由 pH 求出 $c(\ce{OH-})$，进而求水解常数 $K_\text{h}=\dfrac{c^2(\ce{OH-})}{0.100-c(\ce{OH-})}$，再由 $K_\text{a}=\dfrac{K_\text{w}}{K_\text{h}}$ 求 HClO 的 $K_\text{a}$。\\
\dianbo 体现“$K_\text{a}\cdot K_\text{h}=K_\text{w}$”的灵活运用：酸不稳定时，可通过测其强碱盐溶液的 pH 间接求 $K_\text{a}$。
\end{zhenti}

\section{溶液中粒子浓度的关系}

\subsection{三大守恒}
以 $c$ mol·L$^{-1}$ \ce{Na2CO3} 溶液为例：
\begin{enumerate}
\item \textbf{电荷守恒}（溶液呈电中性）：$c(\ce{Na+})+c(\ce{H+})=2c(\ce{CO3^2-})+c(\ce{HCO3-})+c(\ce{OH-})$。注意 \ce{CO3^2-} 前的系数 2 是\textbf{电荷数}。
\item \textbf{元素质量守恒（物料守恒）}：$n(\mathrm{Na}):n(\mathrm{C})=2:1$，$c(\ce{Na+})=2[c(\ce{CO3^2-})+c(\ce{HCO3-})+c(\ce{H2CO3})]$。
\item \textbf{质子守恒}：水电离出的 \ce{H+} 与 \ce{OH-} 相等，\ce{H+} 被 \ce{CO3^2-} 结合后的去向要“追踪”：$c(\ce{OH-})=c(\ce{H+})+c(\ce{HCO3-})+2c(\ce{H2CO3})$（可由电荷守恒减去物料守恒得到）。
\end{enumerate}

\begin{fangfa}[质子守恒的“基准法”]
选溶液中大量存在的原始粒子（\ce{CO3^2-}、\ce{H2O}）为基准：得质子的产物（\ce{HCO3-} 得 1 个，\ce{H2CO3} 得 2 个，\ce{H+}（即 \ce{H3O+}）得 1 个）与失质子的产物（\ce{OH-} 失 1 个）——\textbf{得质子总数 $=$ 失质子总数}。\ce{NaHCO3} 溶液以 \ce{HCO3-}、\ce{H2O} 为基准：$c(\ce{H+})+c(\ce{H2CO3})=c(\ce{OH-})+c(\ce{CO3^2-})$。
\end{fangfa}

\subsection{粒子浓度大小比较}
\textbf{单一溶液}：先看电离和水解的主次，“\textbf{电离、水解都是微弱的}”。
\begin{itemize}
\item \ce{CH3COONa}：$c(\ce{Na+})>c(\ce{CH3COO-})>c(\ce{OH-})>c(\ce{H+})$。
\item \ce{Na2CO3}：$c(\ce{Na+})>c(\ce{CO3^2-})>c(\ce{OH-})>c(\ce{HCO3-})>c(\ce{H+})$（\ce{OH-} 来自两步水解和水的电离，故多于 \ce{HCO3-}）。
\item \ce{NaHCO3}：$c(\ce{Na+})>c(\ce{HCO3-})>c(\ce{OH-})>c(\ce{H+})$，且水解大于电离，$c(\ce{H2CO3})>c(\ce{CO3^2-})$。
\end{itemize}

\textbf{混合溶液}：先看反应，再看“剩什么、剩多少”，最后比较电离与水解谁大。
\begin{center}
\small
\begin{tabularx}{\linewidth}{lp{3.8cm}X}
\toprule
等浓度混合体系 & 比较依据 & 结论 \\
\midrule
\ce{CH3COOH}与\ce{CH3COONa} & $K_\text{a}=1.75\times10^{-5}>K_\text{h}=5.7\times10^{-10}$，电离为主 & 酸性；$c(\ce{CH3COO-})>c(\ce{Na+})>c(\ce{CH3COOH})>c(\ce{H+})>c(\ce{OH-})$ \\
\ce{NH3*H2O}与\ce{NH4Cl} & $K_\text{b}>K_\text{h}$，电离为主 & 碱性；$c(\ce{NH4+})>c(\ce{Cl-})>c(\ce{NH3*H2O})>c(\ce{OH-})>c(\ce{H+})$ \\
HCN 与 NaCN & $K_\text{a}=6.2\times10^{-10}<K_\text{h}=1.6\times10^{-5}$，水解为主 & 碱性；$c(\ce{HCN})>c(\ce{Na+})>c(\ce{CN-})>c(\ce{OH-})>c(\ce{H+})$ \\
\bottomrule
\end{tabularx}
\end{center}

\begin{fangfa}[滴定曲线上的“五个关键点”（NaOH 滴定 \ce{CH3COOH}）]
\textbf{起点}：只有 \ce{CH3COOH}，看其电离；\textbf{半中和点}：\ce{CH3COOH} 与 \ce{CH3COONa} 等量，pH$\approx$p$K_\text{a}$；\textbf{中性点}：$c(\ce{H+})=c(\ce{OH-})$，由电荷守恒 $c(\ce{Na+})=c(\ce{CH3COO-})$，此时酸略过量；\textbf{恰好中和点}：溶质为 \ce{CH3COONa}，呈碱性，水的电离程度最大；\textbf{过量点}：\ce{NaOH} 与 \ce{CH3COONa} 混合，强碱抑制水解。
\end{fangfa}

\begin{liti}[粒子浓度关系的正误判断]
常温下，向 20 mL 0.1 mol·L$^{-1}$ \ce{CH3COOH} 溶液中逐滴加入 0.1 mol·L$^{-1}$ NaOH 溶液。下列说法正确的是
\xxa{加入 10 mL NaOH 时：$c(\ce{CH3COO-})+c(\ce{CH3COOH})=2c(\ce{Na+})$}{pH$=7$ 时，加入 NaOH 溶液的体积为 20 mL}{加入 20 mL NaOH 时：$c(\ce{Na+})>c(\ce{CH3COO-})>c(\ce{H+})>c(\ce{OH-})$}{滴定过程中，水的电离程度一直增大}
\tcblower
\jie A：此时 $n(\ce{CH3COOH})_\text{总}=2n(\ce{Na+})$，物料守恒成立，正确。B：恰好中和时溶液呈碱性，pH$=7$ 时 NaOH 不足 20 mL。C：恰好中和时溶液呈碱性，应为 $c(\ce{OH-})>c(\ce{H+})$。D：恰好中和前，水的电离程度逐渐增大；之后 NaOH 过量，抑制水的电离，程度又减小。答案：A。
\end{liti}

\section{沉淀溶解平衡}

\subsection{溶度积}
难溶电解质在水中存在溶解与沉淀的平衡，如 \ce{AgCl(s) <=> Ag+(aq) + Cl-(aq)}。
\begin{gainian}[溶度积 $K_\text{sp}$]
\ce{A_mB_n(s) <=> $m$A^{n+}(aq) + $n$B^{m-}(aq)}，$K_\text{sp}=c^{\,m}(\ce{A^{n+}})\cdot c^{\,n}(\ce{B^{m-}})$，只与难溶电解质本身和温度有关。离子积 $Q>K_\text{sp}$，析出沉淀；$Q=K_\text{sp}$，饱和（平衡）；$Q<K_\text{sp}$，沉淀溶解。
\end{gainian}

\begin{yicuo}
\begin{enumerate}
\item 只有\textbf{同类型}（阴阳离子个数比相同）的难溶物才能直接用 $K_\text{sp}$ 比较溶解度。如 $K_\text{sp}(\ce{AgCl})=1.8\times10^{-10}>K_\text{sp}(\ce{Ag2CrO4})\approx1.1\times10^{-12}$，但饱和溶液中 $c(\ce{Ag+})$：AgCl 约 $1.3\times10^{-5}$ mol·L$^{-1}$，\ce{Ag2CrO4} 约 $1.3\times10^{-4}$ mol·L$^{-1}$，后者溶解度反而大。
\item 一般认为离子浓度 $\le1.0\times10^{-5}$ mol·L$^{-1}$ 时，该离子已“沉淀完全”。
\end{enumerate}
\end{yicuo}

\subsection{沉淀的生成、溶解与转化}
\textbf{(1) 沉淀的生成——调 pH 除杂：}以 0.1 mol·L$^{-1}$ \ce{MgCl2} 溶液中除去少量 \ce{Fe^3+} 为例（$K_\text{sp}[\ce{Fe(OH)3}]=2.8\times10^{-39}$，$K_\text{sp}[\ce{Mg(OH)2}]=5.6\times10^{-12}$）：
\ce{Fe^3+} 沉淀完全时 $c(\ce{OH-})=\sqrt[3]{\dfrac{2.8\times10^{-39}}{1.0\times10^{-5}}}\approx6.5\times10^{-12}$ mol·L$^{-1}$，pH$\approx2.8$；
\ce{Mg^2+} 开始沉淀时 $c(\ce{OH-})=\sqrt{\dfrac{5.6\times10^{-12}}{0.1}}\approx7.5\times10^{-6}$ mol·L$^{-1}$，pH$\approx8.9$。
所以应调节 $2.8\le$pH$<8.9$（实际取 3～4 左右即可），用 MgO 等调节。

\textbf{(2) 沉淀的溶解：}加入能与沉淀电离出的离子反应的试剂，使 $Q<K_\text{sp}$。如 \ce{CaCO3} 溶于盐酸；\ce{Mg(OH)2} 溶于 \ce{NH4Cl} 溶液（\ce{NH4+} 结合 \ce{OH-}）；AgCl 溶于氨水（生成 \ce{[Ag(NH3)2]+}）。

\textbf{(3) 沉淀的转化：}一般由溶解度小的向更小的转化，转化反应的平衡常数 $K=\dfrac{K_\text{sp}(\text{原沉淀})}{K_\text{sp}(\text{新沉淀})}$（同类型）。
\begin{itemize}
\item AgCl（白）$\to$ AgBr（浅黄）$\to$ AgI（黄）：\ce{AgCl(s) + I-(aq) <=> AgI(s) + Cl-(aq)}，$K=\dfrac{1.8\times10^{-10}}{8.5\times10^{-17}}\approx2\times10^6$，很彻底。
\item 锅炉除水垢：\ce{CaSO4(s) + CO3^2-(aq) <=> CaCO3(s) + SO4^2-(aq)}，$K=\dfrac{4.9\times10^{-5}}{3.4\times10^{-9}}\approx1.4\times10^4$，再用酸除去 \ce{CaCO3}。
\item 溶解度小的也能转化为溶解度稍大的：\ce{BaSO4(s) + CO3^2-(aq) <=> BaCO3(s) + SO4^2-(aq)}，$K=\dfrac{1.1\times10^{-10}}{2.6\times10^{-9}}\approx0.04$，用饱和 \ce{Na2CO3} 溶液\textbf{多次}处理即可转化（只要 $c(\ce{CO3^2-})$ 大于 $c(\ce{SO4^2-})$ 的约 24 倍）。
\item 用 FeS、MnS 等除去废水中的 \ce{Cu^2+}、\ce{Hg^2+}：\ce{FeS(s) + Cu^2+(aq) <=> CuS(s) + Fe^2+(aq)}。
\end{itemize}

\begin{zhenti}{2025年广东卷第18题(3)\quad 来源：仓库所附真题文档}
从多金属精矿中提取金属的流程中，“沉铝”时需保证 \ce{Cu^2+}、\ce{Ni^2+} 不沉淀。已知“滤液1”中 $c(\ce{Cu^2+})=0.022$ mol·L$^{-1}$、$c(\ce{Ni^2+})=0.042$ mol·L$^{-1}$，$K_\text{sp}[\ce{Cu(OH)2}]=2.2\times10^{-20}$，$K_\text{sp}[\ce{Ni(OH)2}]=5.5\times10^{-16}$。“沉铝”时 pH 最高可调至\tk（溶液体积变化可忽略）。

\textbf{解析：}比较哪种离子先沉淀：\ce{Cu^2+} 开始沉淀时 $c^2(\ce{OH-})=\dfrac{2.2\times10^{-20}}{0.022}=1.0\times10^{-18}$，$c(\ce{OH-})=10^{-9}$ mol·L$^{-1}$，pH$=5$；\ce{Ni^2+} 开始沉淀时 $c^2(\ce{OH-})=\dfrac{5.5\times10^{-16}}{0.042}\approx1.3\times10^{-14}$，pH$\approx7.1$。\ce{Cu^2+} 先沉淀，pH 最高调至 5。
\end{zhenti}

\section{离子平衡图像与配位平衡}

\subsection{分布系数图}
弱酸各型体占总浓度的分数叫分布系数 $\delta$。以二元弱酸 \ce{H2A} 为例，三条曲线两两相交：\ce{H2A} 与 \ce{HA-} 交点处 $c(\ce{H2A})=c(\ce{HA-})$，代入 $K_{\text{a}1}=\dfrac{c(\ce{H+})c(\ce{HA-})}{c(\ce{H2A})}$ 得 \textbf{pH$=$p$K_{\text{a}1}$}；\ce{HA-} 与 \ce{A^2-} 交点处 pH$=$p$K_{\text{a}2}$。

\begin{center}
\begin{tikzpicture}
\begin{axis}[width=7.6cm,height=4.6cm,xlabel={pH},ylabel={$\delta$},xmin=0,xmax=7,ymin=0,ymax=1.05,
 xtick={0,1.25,2,4.27,6,7},xticklabels={0,1.25,2,4.27,6,7},ytick={0,0.5,1},grid=major,grid style={gray!20},tick label style={font=\scriptsize},label style={font=\small},legend style={font=\scriptsize,at={(1.02,0.5)},anchor=west}]
\addplot[thick,mainblue,domain=0:7,samples=120] {1/(1+10^(x-1.25)+10^(2*x-5.52))};
\addplot[thick,deepgreen,domain=0:7,samples=120] {10^(x-1.25)/(1+10^(x-1.25)+10^(2*x-5.52))};
\addplot[thick,warnred,domain=0:7,samples=120] {10^(2*x-5.52)/(1+10^(x-1.25)+10^(2*x-5.52))};
\legend{\ce{H2A},\ce{HA-},\ce{A^2-}}
\end{axis}
\end{tikzpicture}
\end{center}
\noindent\small 图：某二元弱酸 \ce{H2A}（p$K_{\text{a}1}=1.25$，p$K_{\text{a}2}=4.27$）的分布系数图。\ce{HA-} 的分布系数在 pH$=\frac{\mathrm pK_{\text a1}+\mathrm pK_{\text a2}}2\approx2.76$ 时最大。由 $K_{\text{h}2}=10^{-14}/10^{-1.25}=10^{-12.75}<K_{\text{a}2}=10^{-4.27}$ 可知 \ce{NaHA} 溶液呈酸性。\normalsize

\subsection{对数图与沉淀溶解平衡曲线}
\textbf{对数图}（纵轴为 $\lg c$ 或 $-\lg c$）：把幂运算变成线性关系，读图时把交点、特殊点的坐标代入 $K$ 的表达式即可。如难溶物 \ce{MX2}，$\lg c(\ce{M^2+})=\lg K_\text{sp}-2\lg c(\ce{X-})$，以 $\lg c(\ce{M^2+})$ 对 $\lg c(\ce{X-})$ 作图是斜率为 $-2$ 的直线。\textbf{沉淀溶解平衡曲线}上的点都是饱和溶液；曲线上方（离子浓度大）的点 $Q>K_\text{sp}$，有沉淀析出；下方的点为不饱和溶液。温度不变，曲线不变；同一曲线上各点 $K_\text{sp}$ 相等。

\begin{liti}[沉淀溶解平衡曲线]
某温度下，\ce{BaSO4} 在水中的沉淀溶解平衡曲线如图（$K_\text{sp}=1.0\times10^{-10}$，图中 a、c 在曲线上，b 在曲线上方，d 在曲线下方）。下列说法正确的是
\begin{center}
\begin{tikzpicture}
\begin{axis}[width=6.2cm,height=4.6cm,xlabel={$c(\ce{SO4^2-})/(10^{-5}\,\mathrm{mol\cdot L^{-1}})$},ylabel={$c(\ce{Ba^2+})/(10^{-5}\,\mathrm{mol\cdot L^{-1}})$},xmin=0,xmax=4.4,ymin=0,ymax=4.4,tick label style={font=\scriptsize},label style={font=\scriptsize},grid=major,grid style={gray!20}]
\addplot[thick,mainblue,domain=0.24:4.2,samples=80]{1/x};
\addplot[only marks,mark=*,mark size=1.6pt] coordinates {(0.5,2) (2,0.5) (2,2) (1,0.5)};
\node[font=\scriptsize,above right] at (axis cs:0.5,2) {a};
\node[font=\scriptsize,above right] at (axis cs:2,0.5) {c};
\node[font=\scriptsize,above right] at (axis cs:2,2) {b};
\node[font=\scriptsize,below left] at (axis cs:1,0.5) {d};
\end{axis}
\end{tikzpicture}
\end{center}
\xxa{加入 \ce{Na2SO4} 固体可使溶液由 c 点变到 a 点}{b 点对应的溶液中无 \ce{BaSO4} 沉淀生成}{蒸发溶剂可使溶液由 d 点变到曲线上 a、c 之间的某点}{升高温度可使溶液由 c 点变到 a 点}
\tcblower
\jie 曲线上的点为饱和溶液，$c(\ce{Ba^2+})\cdot c(\ce{SO4^2-})=K_\text{sp}$。A：加 \ce{Na2SO4}，$c(\ce{SO4^2-})$ 增大，平衡逆移，$c(\ce{Ba^2+})$ 减小，应沿曲线向右下方移动（a$\to$c），而不是 c$\to$a，错误。B：b 点 $Q=2\times10^{-5}\times2\times10^{-5}=4\times10^{-10}>K_\text{sp}$，过饱和，有沉淀生成，错误。C：蒸发时 $c(\ce{Ba^2+})$、$c(\ce{SO4^2-})$ 同比例增大，d 点（$c(\ce{SO4^2-})=1$，$c(\ce{Ba^2+})=0.5$，单位 $10^{-5}$ mol·L$^{-1}$）沿直线 $c(\ce{Ba^2+})=0.5c(\ce{SO4^2-})$ 移动，与曲线交于 $c(\ce{SO4^2-})=\sqrt2\approx1.41$ 处，位于 a（0.5）与 c（2）之间，正确。D：升温 $K_\text{sp}$ 改变，曲线整体移动，c 点不会沿原曲线变到 a 点，错误。答案：C。\\
\dianbo 曲线上的移动只能靠改变某一离子的浓度（温度不变）；“蒸发浓缩”使两种离子浓度同比增大，“加入某离子”使另一离子浓度减小。
\end{liti}

\subsection{溶液中的配位平衡}
金属离子与配体形成配离子也是可逆的，如 \ce{Cu^2+ + 4NH3 <=> [Cu(NH3)4]^2+}，\ce{Ag+ + 2NH3 <=> [Ag(NH3)2]+}，\ce{Fe^3+ + 3SCN- <=> Fe(SCN)3}。配体浓度越大，越有利于形成配位数高的配离子；加入能结合配体或中心离子的试剂（如酸结合 \ce{NH3}），配离子会解离。

\begin{zhenti}{2025年广东卷第12题\quad 来源：仓库所附真题文档}
CuCl 微溶于水，但在 \ce{Cl-} 浓度较高的溶液中因形成 \ce{[CuCl2]-} 和 \ce{[CuCl3]^2-} 而溶解。将适量 CuCl 完全溶于盐酸，得到含 \ce{[CuCl2]-} 和 \ce{[CuCl3]^2-} 的溶液。下列叙述正确的是
\xxa{加水稀释，\ce{[CuCl3]^2-} 浓度一定下降}{向溶液中加入少量 NaCl 固体，\ce{[CuCl2]-} 浓度一定上升}{\ce{H[CuCl2]} 的电离方程式为 \ce{H[CuCl2] = H+ + Cu+ + 2Cl-}}{体系中，$c(\ce{Cu+})+c(\ce{H+})=c(\ce{Cl-})+c(\ce{[CuCl2]-})+c(\ce{[CuCl3]^2-})+c(\ce{OH-})$}
\textbf{解析：}A：加水稀释，溶液体积增大，且 \ce{[CuCl2]- + Cl- <=> [CuCl3]^2-} 因 $c(\ce{Cl-})$ 减小而逆移，\ce{[CuCl3]^2-} 浓度一定下降，正确。B：加 NaCl，$c(\ce{Cl-})$ 增大，可能使 \ce{[CuCl2]-} 转化为 \ce{[CuCl3]^2-}，浓度不一定上升。C：配离子是一个整体，应为 \ce{H[CuCl2] = H+ + [CuCl2]-}。D：电荷守恒中 \ce{[CuCl3]^2-} 的系数应为 2。答案：A。
\end{zhenti}

% >>>>>>>> 开始：chapters/ch04ex.tex
% ---------------- 第4章 章末 ----------------
\zhangmo{本章小结}
\subsection*{一、知识总结}
\begin{center}
\small
\begin{tabularx}{\linewidth}{p{2.0cm}p{3.6cm}X}
\toprule
平衡 & 常数 & 要点 \\
\midrule
电离平衡 & $K_\text{a}=\frac{c(\ce{H+})c(\ce{A-})}{c(\ce{HA})}$ & 越稀越电离、越热越电离；$K_\text{a}$ 大酸性强；强酸制弱酸 \\
水的电离 & $K_\text{w}=c(\ce{H+})c(\ce{OH-})$ & 酸碱抑制、水解促进、升温促进；酸碱性看 $c(\ce{H+})$ 与 $c(\ce{OH-})$ 的相对大小 \\
盐类水解 & $K_\text{h}=K_\text{w}/K_\text{a}$ & 越弱越水解；酸式盐比较 $K_{\text{a}2}$ 与 $K_{\text{h}2}$；水解应用 \\
沉淀溶解 & $K_\text{sp}=c^m(\ce{A^{n+}})c^n(\ce{B^{m-}})$ & $Q$ 与 $K_\text{sp}$ 比较；沉淀完全 $\le10^{-5}$ mol·L$^{-1}$；转化 $K=K_\text{sp1}/K_\text{sp2}$ \\
\bottomrule
\end{tabularx}
\end{center}

\subsection*{二、方法提升}
\begin{enumerate}
\item \textbf{比值变化题}：把比值改写成“$K$/某浓度”或“$K_1/K_2$”，再看那个浓度怎么变。
\item \textbf{粒子浓度题}：先写出三大守恒；比较大小时先找主要过程（电离还是水解），再定次序；等式题优先用守恒组合检验。
\item \textbf{图像题}：找交点、特殊点代入 $K$ 表达式；对数图注意斜率的意义。
\item \textbf{$K_\text{sp}$ 计算}：开始沉淀用原浓度，沉淀完全用 $10^{-5}$；比较不同离子谁先沉淀要分别计算所需 $c(\ce{OH-})$（或 pH）。
\end{enumerate}

\subsection*{三、易错点归纳}
\begin{enumerate}
\item 把导电性强弱等同于电解质强弱；把难溶等同于弱电解质。
\item 认为稀释弱酸时所有离子浓度都减小（$c(\ce{OH-})$ 增大）。
\item 电荷守恒漏乘电荷数；物料守恒漏写某种型体。
\item 用 pH$=7$ 判断中性而忽略温度；把“水电离的 $c(\ce{H+})$”与溶液的 $c(\ce{H+})$ 混淆。
\item 不同类型难溶物直接用 $K_\text{sp}$ 比较溶解度；认为 $K_\text{sp}$ 小的物质不能转化为 $K_\text{sp}$ 大的物质。
\item 写多元弱酸根水解时一步写成 \ce{CO3^2- + 2H2O <=> H2CO3 + 2OH-}（错）。
\end{enumerate}

\zhangmo{分层练习}
\zu{A组\quad 基础巩固}
\begin{xiti}[start=1]
\item 常温下，将 0.1 mol·L$^{-1}$ \ce{CH3COOH} 溶液加水稀释，下列各量增大的是\xz
\xx{$c(\ce{H+})$}{$c(\ce{CH3COOH})$}{$\dfrac{c(\ce{CH3COO-})}{c(\ce{CH3COOH})}$}{$K_\text{a}$}

\item 下列说法正确的是\xz
\xxa{25 ℃ 时，某溶液中由水电离出的 $c(\ce{H+})=10^{-12}$ mol·L$^{-1}$，该溶液一定呈酸性}{100 ℃ 时纯水的 pH$\approx6$，此时纯水呈酸性}{向水中加入少量 \ce{NH4Cl} 固体，水的电离平衡正向移动，$K_\text{w}$ 不变}{向水中加入少量 NaOH 固体，水的电离平衡逆向移动，$c(\ce{H+})$ 增大}

\item 常温下，有 pH 均为 3 的盐酸和醋酸各 10 mL，下列说法正确的是\xz
\xxa{两溶液的物质的量浓度相等}{分别与足量锌反应，醋酸产生的 \ce{H2} 多}{分别加水稀释 100 倍，盐酸的 pH 比醋酸小}{分别与等浓度 NaOH 溶液反应至中性，消耗 NaOH 一样多}

\item 下列事实与盐类水解\textbf{无关}的是\xz
\xxa{\ce{NH4Cl} 溶液可用于除铁锈}{实验室盛放 \ce{Na2CO3} 溶液的试剂瓶不能用磨口玻璃塞}{用 \ce{Na2S} 溶液除去废水中的 \ce{Hg^2+}}{加热蒸干 \ce{AlCl3} 溶液得不到 \ce{AlCl3} 固体}

\item 对 0.1 mol·L$^{-1}$ \ce{Na2S} 溶液，下列关系正确的是\xz
\xxa{$c(\ce{Na+})=2c(\ce{S^2-})$}{$c(\ce{Na+})+c(\ce{H+})=c(\ce{S^2-})+c(\ce{HS-})+c(\ce{OH-})$}{$c(\ce{OH-})=c(\ce{H+})+c(\ce{HS-})+2c(\ce{H2S})$}{$c(\ce{Na+})>c(\ce{S^2-})>c(\ce{HS-})>c(\ce{OH-})$}

\item 已知 25 ℃ 时 $K_\text{sp}(\ce{AgCl})=1.8\times10^{-10}$。下列说法正确的是\xz
\xxa{AgCl 在水中和在 0.1 mol·L$^{-1}$ NaCl 溶液中的溶解度相同}{向 AgCl 饱和溶液中加入少量 NaCl 固体，$K_\text{sp}(\ce{AgCl})$ 减小}{将 $2\times10^{-4}$ mol·L$^{-1}$ \ce{AgNO3} 溶液与 $2\times10^{-4}$ mol·L$^{-1}$ NaCl 溶液等体积混合，有沉淀生成}{AgCl 饱和溶液中 $c(\ce{Ag+})=1.8\times10^{-10}$ mol·L$^{-1}$}

\item 下列水解方程式书写正确的是\xz
\xxa{\ce{CO3^2- + 2H2O <=> H2CO3 + 2OH-}}{\ce{Fe^3+ + 3H2O = Fe(OH)3 v + 3H+}}{\ce{NH4+ + H2O <=> NH3*H2O + H+}}{\ce{HCO3- + H2O <=> CO3^2- + H3O+}（表示 \ce{HCO3-} 的水解）}
\end{xiti}

\zu{B组\quad 高考典型}
\begin{xiti}
\item \laiyuan{原创}\ 25 ℃ 时：\ce{H2CO3} $K_{\text{a}1}=4.5\times10^{-7}$、$K_{\text{a}2}=4.7\times10^{-11}$；HClO $K_\text{a}=4.0\times10^{-8}$；HCN $K_\text{a}=6.2\times10^{-10}$；\ce{CH3COOH} $K_\text{a}=1.75\times10^{-5}$。(1) 等浓度的 \ce{CH3COONa}、NaClO、NaCN、\ce{Na2CO3} 溶液，pH 由大到小的顺序为\tk。(2) 写出向 NaCN 溶液中通入少量 \ce{CO2} 的离子方程式。(3) 0.1 mol·L$^{-1}$ NaClO 溶液的 pH 约为\tk（$\sqrt{2.5}\approx1.6$，$\lg1.6\approx0.2$）。

\item 用 0.1000 mol·L$^{-1}$ NaOH 标准溶液滴定 25.00 mL 未知浓度的醋酸。(1) 应选用的指示剂是\tk，终点现象为\tk。(2) 三次滴定消耗 NaOH 溶液体积分别为 20.02 mL、19.98 mL、20.80 mL，醋酸的浓度为\tk。(3) 下列操作使测定结果偏高的是\tk（填序号）：①碱式滴定管未用标准液润洗；②锥形瓶用蒸馏水洗后未干燥；③滴定前仰视读数、滴定后平视读数；④滴定前滴定管尖嘴有气泡，滴定后气泡消失。

\item \laiyuan{2024年广东卷第19题(3)，来源：仓库所附真题文档}\ 一定温度下，某研究组用分光光度法测定两种一元弱酸 HX（X 为 A 或 B）在某非水溶剂中的 $K_\text{a}$。a．选择合适的指示剂 HIn，$K_\text{a}(\ce{HIn})=3.6\times10^{-20}$；其钾盐为 KIn。b．向 KIn 溶液中加入 HX，发生反应 \ce{In- + HX <=> X- + HIn}。KIn 起始物质的量为 $n_0(\ce{KIn})$，加入 HX 的物质的量为 $n(\ce{HX})$，平衡时测得 $c_\text{平}(\ce{In-})/c_\text{平}(\ce{HIn})$ 随 $n(\ce{HX})/n_0(\ce{KIn})$ 的变化：当 $n(\ce{HX})/n_0(\ce{KIn})=1.0$ 时，HA 体系中该比值为 3.0，HB 体系中该比值小于 1.0。已知该溶剂本身不电离，钾盐在该溶剂中完全电离。①计算 $K_\text{a}(\ce{HA})$（写出计算过程，结果保留两位有效数字）。②在该溶剂中，$K_\text{a}(\ce{HB})$\tk$K_\text{a}(\ce{HA})$；$K_\text{a}(\ce{HB})$\tk$K_\text{a}(\ce{HIn})$（填“$>$”“$<$”或“$=$”）。

\item \laiyuan{2026年广东卷第17题(4)节选，来源：仓库所附真题文档（网络收集版）}\ 探究铜氨配合物 $Q$（\ce{[Cu(NH3)4]SO4*H2O}）与酸的反应。资料：向铜氨溶液中滴加硫酸，依次发生 (a)\ce{[Cu(NH3)4]^2+ + 2H+ + 2H2O = Cu(OH)2 v + 4NH4+}；(b)\ce{Cu(OH)2 + 2H+ = Cu^2+ + 2H2O}。实验：(ii) 称取 $w$ g Q 于烧杯中，移取 $V_0$ mL 稀氨水使 Q 完全溶解，得待测液；(iii) 搅拌下向待测液中滴加 0.1 mol·L$^{-1}$ \ce{H2SO4} 溶液，记录 pH–$V$ 数据；(iv) 移取 $V_0$ mL 上述稀氨水，用\tk{}溶液滴至终点，记录体积为 $V_1$ mL。①(iii) 所得 pH–$V$ 曲线呈现两次 pH 突跃，对应终点体积依次为 $V_2$ mL 和 $V_3$ mL，则反应(a)、(b)消耗酸的体积比 $V_\text{a}:V_\text{b}=$\tk（用 $V_1$、$V_2$、$V_3$ 表示）。②多次实验均得 $V_\text{a}:V_\text{b}>1$。有人提出假设：(iii) 中生成的沉淀除 \ce{Cu(OH)2} 外还含有 \ce{Cu2(OH)2SO4}。某次实验 $V_\text{a}:V_\text{b}=21:19$，则第一个终点时 $\dfrac{n[\ce{Cu(OH)2}]}{n[\ce{Cu2(OH)2SO4}]}=$\tk。

\item \laiyuan{原创}\ 某废液中含 0.10 mol·L$^{-1}$ \ce{Zn^2+} 和少量 \ce{Fe^3+}、\ce{Fe^2+}。已知 $K_\text{sp}[\ce{Fe(OH)3}]=2.8\times10^{-39}$，$K_\text{sp}[\ce{Fe(OH)2}]=4.9\times10^{-17}$，$K_\text{sp}[\ce{Zn(OH)2}]=1.0\times10^{-17}$，$\lg2\approx0.3$，$\lg6.5\approx0.8$。(1) 为什么要先加 \ce{H2O2} 再调 pH？写出相关离子方程式。(2) 计算使 \ce{Fe^3+} 沉淀完全而 \ce{Zn^2+} 不沉淀的 pH 范围。(3) 可选用什么物质调节 pH？

\item \laiyuan{原创}\ 常温下，将 0.2 mol·L$^{-1}$ 氨水与 0.1 mol·L$^{-1}$ 盐酸等体积混合（$K_\text{b}(\ce{NH3*H2O})=1.8\times10^{-5}$）。(1) 混合后溶液中的溶质及其浓度；(2) 判断溶液的酸碱性并说明理由；(3) 将下列粒子浓度由大到小排序：\ce{NH4+}、\ce{Cl-}、\ce{NH3*H2O}、\ce{OH-}、\ce{H+}；(4) 写出该溶液的物料守恒式。

\item \laiyuan{原创}\ 已知 $K_\text{sp}(\ce{AgCl})=1.8\times10^{-10}$，$K_\text{sp}(\ce{AgI})=8.5\times10^{-17}$。(1) 向含 AgCl 沉淀的饱和溶液中加入 KI 溶液，发生的反应为\tk，其平衡常数 $K=$\tk。(2) 若要使 0.01 mol AgCl 完全转化为 AgI（溶液体积 1 L），所加 KI 的物质的量至少为多少（结果保留两位有效数字）？
\end{xiti}

\zu{C组\quad 冲刺高分}
\begin{xiti}
\item \laiyuan{原创}\ 常温下，向 \ce{H2A} 溶液中滴加 NaOH 溶液，\ce{H2A}、\ce{HA-}、\ce{A^2-} 的分布系数 $\delta$ 随 pH 变化：\ce{H2A} 与 \ce{HA-} 曲线交于 pH$=1.2$，\ce{HA-} 与 \ce{A^2-} 曲线交于 pH$=4.2$。(1) 求 $K_{\text{a}1}$、$K_{\text{a}2}$。(2) NaHA 溶液呈酸性还是碱性？通过计算说明。(3) 当 $c(\ce{H2A})=c(\ce{A^2-})$ 时，溶液 pH 为多少？(4) 求 $\dfrac{c^2(\ce{HA-})}{c(\ce{H2A})\cdot c(\ce{A^2-})}$ 的值，它是否随 pH 变化？(5) 溶液 pH$=7$ 时，写出 $c(\ce{Na+})$ 与含 A 粒子浓度的关系式。

\item \laiyuan{2026年广东卷第18题(3)，来源：仓库所附真题文档（网络收集版）}\ 金精矿提金流程中，“沉锑”用铁粉还原 \ce{[SbCl4]-} 得到 Sb，“沉锑”后的滤液含 \ce{Fe^2+}、\ce{Mg^2+} 等。“电解”在用阴离子交换膜隔开的电解池中进行，一定条件下实现“沉锑”试剂（铁粉）的回收利用。忽略生成气体的反应。已知 $K_\text{sp}$（298 K）：\ce{Fe(OH)3} $2.8\times10^{-39}$，\ce{Fe(OH)2} $4.9\times10^{-17}$，\ce{Mg(OH)2} $5.6\times10^{-12}$。①阴极反应式为\tk。②初始电解液中 $c(\ce{Fe^2+})=1.4$ mol·L$^{-1}$，$c(\ce{Mg^2+})=0.056$ mol·L$^{-1}$。为防止阳极室电解过程中产生沉淀，理论上最终阳极室 $c(\ce{OH-})$ 应不大于\tk mol·L$^{-1}$（列算式）。

\item \laiyuan{原创}\ 用 0.1000 mol·L$^{-1}$ 盐酸滴定 20.00 mL 0.1000 mol·L$^{-1}$ \ce{Na2CO3} 溶液，pH 计记录的滴定曲线有两个突跃：第一个终点在 20.00 mL 附近（pH 约 8.3），第二个终点在 40.00 mL 附近（pH 约 3.9）；加入 10.00 mL 时 pH$=10.33$，加入 30.00 mL 时 pH$=6.35$。(1) 写出两个阶段的离子方程式。(2) 求 \ce{H2CO3} 的 $K_{\text{a}1}$、$K_{\text{a}2}$（用 10 的幂表示）。(3) 若用指示剂代替 pH 计，两个终点应分别选用什么指示剂？(4) 加入 20.00 mL 盐酸时，溶液中的主要溶质是什么？写出此时的质子守恒式。
\end{xiti}

\zhangmo{练习解析}
\begin{jiexilist}
\item \ans{C}稀释时 $c(\ce{H+})$、$c(\ce{CH3COOH})$ 减小；$\dfrac{c(\ce{CH3COO-})}{c(\ce{CH3COOH})}=\dfrac{K_\text{a}}{c(\ce{H+})}$，$c(\ce{H+})$ 减小，比值增大；$K_\text{a}$ 只与温度有关，不变。
\item \ans{C}A：水的电离被抑制，可能是酸也可能是碱。B：纯水中 $c(\ce{H+})=c(\ce{OH-})$，始终呈中性。C：\ce{NH4+} 水解促进水的电离，温度不变 $K_\text{w}$ 不变，正确。D：加 NaOH，$c(\ce{OH-})$ 增大，$c(\ce{H+})$ 减小。
\item \ans{B}pH 相同，醋酸浓度远大于盐酸，A 错；醋酸总物质的量大，产生 \ce{H2} 多，B 对；稀释 100 倍，盐酸 pH 变为 5，醋酸继续电离，pH 小于 5，C 错；醋酸消耗 NaOH 多，D 错。
\item \ans{C}A：\ce{NH4+} 水解呈酸性，可溶解铁锈；B：\ce{CO3^2-} 水解呈碱性，会与玻璃中 \ce{SiO2} 反应使瓶塞粘住；C：\ce{Hg^2+ + S^2- = HgS v}，是沉淀反应，与水解无关；D：\ce{AlCl3} 水解，HCl 挥发，蒸干得 \ce{Al(OH)3}。
\item \ans{C}A：\ce{S^2-} 水解，$c(\ce{Na+})>2c(\ce{S^2-})$，物料守恒应为 $c(\ce{Na+})=2[c(\ce{S^2-})+c(\ce{HS-})+c(\ce{H2S})]$。B：漏乘电荷数，应为 $2c(\ce{S^2-})$。C：质子守恒，正确。D：\ce{OH-} 来自两步水解和水的电离，$c(\ce{OH-})>c(\ce{HS-})$。
\item \ans{C}A：同离子效应，AgCl 在 NaCl 溶液中溶解度更小；B：$K_\text{sp}$ 只与温度有关，不变；C：混合后 $c(\ce{Ag+})=c(\ce{Cl-})=1\times10^{-4}$ mol·L$^{-1}$，$Q=1\times10^{-8}>K_\text{sp}=1.8\times10^{-10}$，有沉淀生成，正确；D：$c(\ce{Ag+})=\sqrt{1.8\times10^{-10}}\approx1.3\times10^{-5}$ mol·L$^{-1}$。
\item \ans{C}A：多元弱酸根应分步水解；B：水解程度小，用“\ce{<=>}”且不标沉淀；C 正确；D 表示的是 \ce{HCO3-} 的电离，水解应为 \ce{HCO3- + H2O <=> H2CO3 + OH-}。
\item (1) \ce{Na2CO3}$>$NaCN$>$NaClO$>$\ce{CH3COONa}。对应酸的强弱：\ce{HCO3-}（$4.7\times10^{-11}$）$<$HCN（$6.2\times10^{-10}$）$<$HClO（$4.0\times10^{-8}$）$<$\ce{CH3COOH}（$1.75\times10^{-5}$），越弱越水解，碱性越强。（\ce{NaHCO3} 既水解又电离，其溶液 pH 约为 8.3 且几乎不随浓度变化，不能只凭 $K_{\text{a}1}$ 与其他盐排序，故本题不列入。）(2) \ce{CN- + CO2 + H2O = HCN + HCO3-}（HCN 的 $K_\text{a}$ 大于 \ce{HCO3-}，只能生成 \ce{HCO3-}）。(3) $K_\text{h}=\dfrac{10^{-14}}{4.0\times10^{-8}}=2.5\times10^{-7}$，$c(\ce{OH-})\approx\sqrt{2.5\times10^{-7}\times0.1}=\sqrt{2.5}\times10^{-4}\approx1.6\times10^{-4}$ mol·L$^{-1}$，pOH$\approx3.8$，pH$\approx10.2$。
\item (1) 酚酞；当滴入最后半滴 NaOH 溶液时，溶液由无色变为浅红色，且半分钟内不褪色。(2) 第三次数据（20.80 mL）偏差过大，应舍去，取前两次平均值 20.00 mL；$c=\dfrac{0.1000\times20.00}{25.00}=0.08000$ mol·L$^{-1}$。(3) ①④。①标准液被稀释，消耗体积偏大，偏高；②无影响；③滴定前仰视读数偏大，$V_\text{标}$ 测得值偏小，偏低；④气泡体积被计入，偏高。
\item ①设 KIn 与 HA 起始浓度均为 $c_0$，转化 $x$：平衡时 $c(\ce{In-})=c(\ce{HA})=c_0-x$，$c(\ce{A-})=c(\ce{HIn})=x$。由 $\dfrac{c_0-x}{x}=3.0$ 得 $x=0.25c_0$。该反应 $K_1=\dfrac{c(\ce{A-})c(\ce{HIn})}{c(\ce{In-})c(\ce{HA})}=\dfrac{c(\ce{A-})c(\ce{H+})}{c(\ce{HA})}\cdot\dfrac{c(\ce{HIn})}{c(\ce{In-})c(\ce{H+})}=\dfrac{K_\text{a}(\ce{HA})}{K_\text{a}(\ce{HIn})}$，而 $K_1=\dfrac{0.25c_0\times0.25c_0}{0.75c_0\times0.75c_0}=\dfrac19$，故 $K_\text{a}(\ce{HA})=\dfrac{3.6\times10^{-20}}{9}=4.0\times10^{-21}$。②$>$；$>$。HB 体系中 $\dfrac{c_0-y}{y}<1.0$，$y>0.5c_0$，$K_2=\dfrac{y^2}{(c_0-y)^2}>1>K_1$，所以 $K_\text{a}(\ce{HB})>K_\text{a}(\ce{HA})$；又 $K_2=\dfrac{K_\text{a}(\ce{HB})}{K_\text{a}(\ce{HIn})}>1$，$K_\text{a}(\ce{HB})>K_\text{a}(\ce{HIn})$。
\item (iv) 0.1 mol·L$^{-1}$ \ce{H2SO4}（与(iii)中滴定所用溶液相同，作空白对照，测出 $V_0$ mL 稀氨水自身消耗的酸）。①$(V_2-V_1):(V_3-V_2)$。第一个突跃对应反应(a)完成，此前消耗的酸包括稀氨水消耗的 $V_1$ 和反应(a)消耗的部分，$V_\text{a}=V_2-V_1$；第二个突跃对应反应(b)完成，$V_\text{b}=V_3-V_2$。②18。设 \ce{Cu(OH)2} $x$ mol、\ce{Cu2(OH)2SO4} $y$ mol：生成沉淀时 \ce{[Cu(NH3)4]^2+ + 2H+ + 2H2O = Cu(OH)2 v + 4NH4+}，\ce{2[Cu(NH3)4]^2+ + 6H+ + SO4^2- + 2H2O = Cu2(OH)2SO4 v + 8NH4+}，消耗 \ce{H+} $(2x+6y)$ mol；溶解沉淀时 \ce{Cu(OH)2 + 2H+ = Cu^2+ + 2H2O}，\ce{Cu2(OH)2SO4 + 2H+ = 2Cu^2+ + SO4^2- + 2H2O}，消耗 $(2x+2y)$ mol。$\dfrac{2x+6y}{2x+2y}=\dfrac{21}{19}$，解得 $x=18y$。
\item (1) $K_\text{sp}[\ce{Fe(OH)2}]$ 与 $K_\text{sp}[\ce{Zn(OH)2}]$ 接近，\ce{Fe^2+} 沉淀时 \ce{Zn^2+} 也会沉淀，无法分离；先用 \ce{H2O2} 将 \ce{Fe^2+} 氧化为 \ce{Fe^3+}，\ce{Fe^3+} 在很低 pH 下即可沉淀完全：\ce{2Fe^2+ + H2O2 + 2H+ = 2Fe^3+ + 2H2O}。(2) \ce{Fe^3+} 沉淀完全：$c(\ce{OH-})=\sqrt[3]{2.8\times10^{-39}/10^{-5}}\approx6.5\times10^{-12}$，pOH$\approx11.2$，pH$\approx2.8$；\ce{Zn^2+} 开始沉淀：$c(\ce{OH-})=\sqrt{1.0\times10^{-17}/0.10}=10^{-8}$，pH$=6$。范围：$2.8\le$pH$<6$。(3) ZnO、\ce{Zn(OH)2} 或 \ce{ZnCO3}（消耗 \ce{H+} 且不引入新杂质）。
\item (1) \ce{NH4Cl} 0.05 mol·L$^{-1}$ 和 \ce{NH3*H2O} 0.05 mol·L$^{-1}$。(2) 碱性：\ce{NH3*H2O} 的电离常数 $K_\text{b}=1.8\times10^{-5}$ 大于 \ce{NH4+} 的水解常数 $K_\text{h}=\frac{10^{-14}}{1.8\times10^{-5}}\approx5.6\times10^{-10}$，电离程度大于水解程度。(3) $c(\ce{NH4+})>c(\ce{Cl-})>c(\ce{NH3*H2O})>c(\ce{OH-})>c(\ce{H+})$。(4) $c(\ce{NH4+})+c(\ce{NH3*H2O})=2c(\ce{Cl-})$。
\item (1) \ce{AgCl(s) + I-(aq) <=> AgI(s) + Cl-(aq)}；$K=\dfrac{c(\ce{Cl-})}{c(\ce{I-})}=\dfrac{1.8\times10^{-10}}{8.5\times10^{-17}}\approx2.1\times10^6$。(2) 完全转化后溶液中 $c(\ce{Cl-})=0.01$ mol·L$^{-1}$，平衡时 $c(\ce{I-})=\dfrac{0.01}{2.1\times10^6}\approx4.7\times10^{-9}$ mol·L$^{-1}$，可忽略；消耗 \ce{I-} 0.01 mol，所以至少加 KI 约 0.010 mol（理论上略多于 0.01 mol）。
\item (1) 交点处两种型体浓度相等：$K_{\text{a}1}=c(\ce{H+})=10^{-1.2}$，$K_{\text{a}2}=10^{-4.2}$。(2) 酸性。\ce{HA-} 的水解常数 $K_{\text{h}2}=\dfrac{10^{-14}}{10^{-1.2}}=10^{-12.8}<K_{\text{a}2}=10^{-4.2}$，电离程度大于水解程度。(3) $K_{\text{a}1}K_{\text{a}2}=\dfrac{c^2(\ce{H+})c(\ce{A^2-})}{c(\ce{H2A})}$，$c(\ce{H2A})=c(\ce{A^2-})$ 时 $c^2(\ce{H+})=10^{-5.4}$，pH$=2.7$。(4) $\dfrac{c^2(\ce{HA-})}{c(\ce{H2A})c(\ce{A^2-})}=\dfrac{K_{\text{a}1}}{K_{\text{a}2}}=10^{3}$，是常数，不随 pH 变化。(5) pH$=7$ 时 $c(\ce{H+})=c(\ce{OH-})$，由电荷守恒 $c(\ce{Na+})=c(\ce{HA-})+2c(\ce{A^2-})$。
\item ①\ce{Fe^2+ + 2e- = Fe}（回收铁粉，供“沉锑”再用）。②$\sqrt[3]{\dfrac{2.8\times10^{-39}}{1.4}}$。阳极上 \ce{Fe^2+} 失电子生成 \ce{Fe^3+}，最终阳极室 $c(\ce{Fe^3+})=1.4$ mol·L$^{-1}$。比较各离子开始沉淀所需 $c(\ce{OH-})$：\ce{Fe(OH)3} 为 $\sqrt[3]{2.8\times10^{-39}/1.4}\approx1.3\times10^{-13}$；\ce{Fe(OH)2}（若有剩余 \ce{Fe^2+}）约 $5.9\times10^{-9}$；\ce{Mg(OH)2} 为 $\sqrt{5.6\times10^{-12}/0.056}=10^{-5}$。\ce{Fe(OH)3} 所需最小，故 $c(\ce{OH-})\le\sqrt[3]{2.8\times10^{-39}/1.4}$ mol·L$^{-1}$。
\item (1) 第一阶段：\ce{CO3^2- + H+ = HCO3-}；第二阶段：\ce{HCO3- + H+ = H2CO3}（\ce{H2CO3 = CO2 ^ + H2O}）。(2) 加入 10.00 mL 时 \ce{CO3^2-} 恰好转化一半，$c(\ce{CO3^2-})=c(\ce{HCO3-})$，$K_{\text{a}2}=c(\ce{H+})=10^{-10.33}$；加入 30.00 mL 时 $c(\ce{HCO3-})=c(\ce{H2CO3})$，$K_{\text{a}1}=10^{-6.35}$。(3) 第一终点（pH$\approx$8.3）用酚酞；第二终点（pH$\approx$3.9）用甲基橙。(4) \ce{NaHCO3}（和 NaCl）；质子守恒：$c(\ce{H+})+c(\ce{H2CO3})=c(\ce{OH-})+c(\ce{CO3^2-})$。
\end{jiexilist}
% <<<<<<<< 结束：chapters/ch04ex.tex
% <<<<<<<< 结束：chapters/ch04.tex

% >>>>>>>> 开始：chapters/ch05.tex
% ============================================================
% 第5章 化学反应与电能（选必1第四章，课件未覆盖，补充）
% ============================================================
\chapter{化学反应与电能}

\begin{daoxue}
\textbf{知识主线：}自发的氧化还原反应（氧化、还原分区进行）$\to$ 原电池（化学能$\to$电能）$\to$ 化学电源；外加电源（强制发生氧化还原）$\to$ 电解池（电能$\to$化学能）$\to$ 氯碱、电镀、精炼、冶金；两者结合 $\to$ 金属的腐蚀与防护。核心思想只有一条：\textbf{负极（阳极）失电子发生氧化反应，正极（阴极）得电子发生还原反应；电子守恒}。

\textbf{高考考情：}广东卷每年选择题中有\textbf{两道}电化学题：2024年第5题（牺牲阳极）、第16题（离子交换膜电解与计算）；2025年第14题（水系可充电电池）、第16题（燃料电池催化机理）；2026年第6题（海上风电防腐）、第16题（电解还原 \ce{CO2} 与海水脱盐）。大题中也常出现电镀、电解回收（2025年第18题、2026年第17、18题）。

\textbf{学习目标：}①判断电极、电子与离子的流向；②规范书写各种介质中的电极反应式；③会分析离子交换膜的作用；④熟练进行电子守恒计算；⑤理解腐蚀原理与防护方法。
\end{daoxue}

\section{原电池的工作原理}

\subsection{从氧化还原到电流}
把锌片直接放入 \ce{CuSO4} 溶液，\ce{Zn + Cu^2+ = Zn^2+ + Cu}，电子直接转移，化学能变成热能。若让氧化反应和还原反应\textbf{在两个区域分别进行}，电子只能通过外电路导线从锌流向铜，就形成了电流——这就是原电池。

\begin{center}
\begin{tikzpicture}[scale=0.8,font=\small]
\draw[thick] (0,0) rectangle (2.6,2.2); \draw[thick] (4.4,0) rectangle (7,2.2);
\fill[mainblue!10] (0.02,0.02) rectangle (2.58,1.6); \fill[mainblue!18] (4.42,0.02) rectangle (6.98,1.6);
\draw[very thick,gray] (1.0,0.4)--(1.0,3.0); \node at (0.6,2.5) {Zn}; \draw[very thick,orange!70!black] (6.0,0.4)--(6.0,3.0); \node at (6.4,2.5) {Cu};
\draw[thick] (1.0,3.0)--(1.0,3.6)--(6.0,3.6)--(6.0,3.0);
\draw[-Stealth,warnred,thick] (2.6,3.75)--(4.4,3.75) node[midway,above]{\scriptsize 电子 $e^-$};
\draw[thick,deepgreen] (1.9,1.3)--(1.9,2.6)--(5.1,2.6)--(5.1,1.3); \draw[thick,deepgreen] (2.2,1.3)--(2.2,2.3)--(4.8,2.3)--(4.8,1.3);
\node[deepgreen] at (3.5,2.95) {\scriptsize 盐桥（KCl 琼胶）};
\node at (1.3,0.8) {\scriptsize \ce{ZnSO4}}; \node at (5.6,0.8) {\scriptsize \ce{CuSO4}};
\node[align=center] at (1.3,-0.6) {\scriptsize 负极：\ce{Zn - 2e- = Zn^2+}};
\node[align=center] at (5.7,-0.6) {\scriptsize 正极：\ce{Cu^2+ + 2e- = Cu}};
\node[align=left,anchor=west] at (7.4,1.6) {\scriptsize 盐桥中 \ce{K+} 移向正极区，\\[-1pt]\scriptsize \ce{Cl-} 移向负极区，\\[-1pt]\scriptsize 使两溶液保持电中性};
\end{tikzpicture}
\end{center}

\begin{gainian}[原电池]
将化学能转化为电能的装置。构成条件：①有自发进行的氧化还原反应；②两个电极；③电解质溶液（或熔融电解质）；④形成闭合回路。
\textbf{负极}：电子流出，发生\textbf{氧化反应}；\textbf{正极}：电子流入，发生\textbf{还原反应}。外电路电子由负极流向正极；内电路阳离子移向正极，阴离子移向负极。
\end{gainian}

\textbf{盐桥的作用：}①连通内电路；②平衡两池电荷（否则 \ce{ZnSO4} 溶液阳离子过多、\ce{CuSO4} 溶液阴离子过多，反应很快停止）；③使氧化剂与还原剂不直接接触，电流稳定、能量利用率高。用离子交换膜代替盐桥也是同样的道理。

\begin{yicuo}
负极不一定是“更活泼的金属”，要看谁与电解质溶液发生自发反应：Mg、Al 插入 NaOH 溶液，Al 与 NaOH 反应，\textbf{Al 作负极}（\ce{Al - 3e- + 4OH- = [Al(OH)4]-}）；Fe、Cu 插入浓硝酸，Fe 钝化，\textbf{Cu 作负极}。另外，电子只在导线中流动，\textbf{“电子不下水，离子不上岸”}。
\end{yicuo}

\subsection{电极反应式的书写}
\begin{fangfa}[三步写电极反应]
①\textbf{写总反应、找得失}：标出化合价变化，确定负极反应物、正极反应物及转移电子数。②\textbf{按介质补粒子}：酸性介质用 \ce{H+} 和 \ce{H2O} 配平，碱性介质用 \ce{OH-} 和 \ce{H2O}，熔融碳酸盐用 \ce{CO3^2-}，固体氧化物用 \ce{O^2-}；酸性介质中不能出现 \ce{OH-}，碱性介质中不能出现 \ce{H+}，也不能生成 \ce{CO2}（应生成 \ce{CO3^2-}）。③\textbf{查守恒}：原子、电荷、电子守恒；两极反应式相加（电子数相同）应得总反应。
\end{fangfa}

\begin{liti}[不同介质中燃料电池的电极反应]
写出下列燃料电池的电极反应式：(1) 甲醇—空气，质子交换膜（酸性）；(2) 肼（\ce{N2H4}）—空气，KOH 溶液；(3) CO—空气，熔融 \ce{K2CO3}（正极通入空气和 \ce{CO2}）；(4) \ce{CH4}—空气，固体氧化物电解质（能传导 \ce{O^2-}）。
\tcblower
\silu 先写正极（只由介质决定），再用“总反应$-$正极反应”得负极。\\
\jie (1) 正：\ce{O2 + 4e- + 4H+ = 2H2O}；负：\ce{CH3OH - 6e- + H2O = CO2 ^ + 6H+}。\\
(2) 正：\ce{O2 + 4e- + 2H2O = 4OH-}；负：\ce{N2H4 - 4e- + 4OH- = N2 ^ + 4H2O}（总反应 \ce{N2H4 + O2 = N2 + 2H2O}）。\\
(3) 正：\ce{O2 + 4e- + 2CO2 = 2CO3^2-}；负：\ce{CO - 2e- + CO3^2- = 2CO2}。负极生成的 \ce{CO2} 可循环到正极使用。\\
(4) 正：\ce{O2 + 4e- = 2O^2-}；负：\ce{CH4 - 8e- + 4O^2- = CO2 + 2H2O}。\\
\xianjing 碱性介质中 C 元素最终以 \ce{CO3^2-} 存在；熔融碳酸盐中不能出现 \ce{OH-}、\ce{H+}；固体氧化物中不能出现 \ce{H2O} 参与传导。
\end{liti}

\textbf{原电池原理的应用：}加快反应速率（锌与稀硫酸反应时滴入少量 \ce{CuSO4} 溶液，Zn 置换出 Cu 形成微电池）；比较金属活动性（作负极者一般更活泼）；设计原电池（选一个自发的氧化还原反应，还原剂作负极或负极反应物，氧化剂溶液作电解质）。

\section{化学电源}

\begin{center}
\small
\begin{tabularx}{\linewidth}{p{2.2cm}p{4.6cm}X}
\toprule
电池 & 总反应 & 电极反应 \\
\midrule
碱性锌锰电池（一次） & \ce{Zn + 2MnO2 + 2H2O = 2MnO(OH) + Zn(OH)2} & 负：\ce{Zn - 2e- + 2OH- = Zn(OH)2}；正：\ce{MnO2 + e- + H2O = MnO(OH) + OH-} \\
铅酸蓄电池（二次） & \ce{Pb + PbO2 + 2H2SO4 <=>[\text{放电}][\text{充电}] 2PbSO4 + 2H2O} & 放电负：\ce{Pb - 2e- + SO4^2- = PbSO4}；放电正：\ce{PbO2 + 2e- + 4H+ + SO4^2- = PbSO4 + 2H2O}；放电时两极质量都增加，\ce{H2SO4} 浓度减小 \\
锂离子电池（二次） & \ce{Li_xC6 + Li_{1-$x$}CoO2 <=>[\text{放电}][\text{充电}] 6C + LiCoO2} & 放电负：\ce{Li_xC6 - $x$e- = $x$Li+ + 6C}；放电正：\ce{Li_{1-$x$}CoO2 + $x$Li+ + $x$e- = LiCoO2}；\ce{Li+} 在两极间“摇摆” \\
氢氧燃料电池 & \ce{2H2 + O2 = 2H2O} & 酸性：负 \ce{H2 - 2e- = 2H+}，正 \ce{O2 + 4e- + 4H+ = 2H2O}；碱性：负 \ce{H2 - 2e- + 2OH- = 2H2O}，正 \ce{O2 + 4e- + 2H2O = 4OH-} \\
甲烷燃料电池（KOH） & \ce{CH4 + 2O2 + 2OH- = CO3^2- + 3H2O} & 负：\ce{CH4 - 8e- + 10OH- = CO3^2- + 7H2O}；正：\ce{2O2 + 8e- + 4H2O = 8OH-}；溶液碱性减弱 \\
\bottomrule
\end{tabularx}
\end{center}

\begin{benzhi}[燃料电池中燃料一定在负极，为什么？]
燃料（\ce{H2}、\ce{CH4}、\ce{CH3OH}、\ce{N2H4} 等）是还原剂，在负极失电子；\ce{O2} 是氧化剂，在正极得电子。正极反应只取决于介质：酸性 \ce{O2 + 4e- + 4H+ = 2H2O}；碱性或中性 \ce{O2 + 4e- + 2H2O = 4OH-}；熔融碳酸盐 \ce{O2 + 4e- + 2CO2 = 2CO3^2-}；固体氧化物 \ce{O2 + 4e- = 2O^2-}。写负极反应时，常用“总反应 $-$ 正极反应”得到，最不易出错。
\end{benzhi}

\textbf{二次电池的充电：}充电是放电的逆过程。\textbf{“负接负、正接正”}——电池负极接外电源负极（充电时作\textbf{阴极}），正极接外电源正极（作\textbf{阳极}）。充电时阴极反应是放电负极反应的逆反应，阳极反应是放电正极反应的逆反应。

\begin{fangfa}[新型电池题的解题步骤]
①\textbf{找变价}：由元素化合价变化（或“放电时某物质减少”）确定正、负极；②\textbf{定流向}：电子由负极经导线到正极，阳离子向正极迁移（充电时向阴极迁移）；③\textbf{写反应}：注意介质和膜允许通过的离子；④\textbf{算数量}：电子守恒，1 mol 电子对应电量 96500 C，即 26.8 A·h。
\end{fangfa}

\begin{liti}[全钒液流电池]
全钒液流电池的两极电解液分别含钒的不同价态离子，用质子交换膜隔开，总反应为
\[\ce{VO2+ + V^2+ + 2H+ <=>[\text{放电}][\text{充电}] VO^2+ + V^3+ + H2O}\]
（颜色：\ce{V^2+} 紫色，\ce{V^3+} 绿色，\ce{VO^2+} 蓝色，\ce{VO2+} 黄色）。(1) 写出放电时正、负极的电极反应式。(2) 放电时 \ce{H+} 通过质子交换膜向哪一极迁移？(3) 充电时阳极室溶液颜色如何变化？(4) 放电过程中每转移 1 mol 电子，正极室溶液中 $n(\ce{H+})$ 净变化多少？
\tcblower
\jie (1) V 元素：\ce{VO2+}（$+5$）$\to$\ce{VO^2+}（$+4$），被还原，在正极：\ce{VO2+ + 2H+ + e- = VO^2+ + H2O}；\ce{V^2+}（$+2$）$\to$\ce{V^3+}（$+3$），在负极：\ce{V^2+ - e- = V^3+}。(2) 移向正极。(3) 充电时阳极（放电正极）发生 \ce{VO^2+ - e- + H2O = VO2+ + 2H+}，溶液由蓝色变为黄色。(4) 每转移 1 mol 电子，正极反应消耗 2 mol \ce{H+}，同时有 1 mol \ce{H+} 经膜从负极室迁入，正极室 $n(\ce{H+})$ 净减少 1 mol。\\
\dianbo 分析离子交换膜问题时，“电极反应消耗或生成的离子”与“经膜迁移的离子”要分开算，再求净变化。
\end{liti}

\begin{zhenti}{2025年广东卷第14题\quad 来源：仓库所附真题文档}
一种高容量水系电池：电极Ⅰ为 S/MnS，电极Ⅱ为 \ce{MnO2}，两极室均为 \ce{MnSO4} 溶液，中间为离子交换膜。已知放电时电极Ⅱ上 \ce{MnO2} 减少；电极材料每转移 1 mol 电子，对应的理论容量为 26.8 A·h。下列说法错误的是
\xxb{充电时Ⅱ为阳极}{放电时Ⅱ极室中溶液的 pH 降低}{放电时负极反应为 \ce{MnS - 2e- = S + Mn^2+}}{充电时 16 g S 能提供的理论容量为 26.8 A·h}
\textbf{解析：}放电时 \ce{MnO2} 减少，被还原为 \ce{Mn^2+}，Ⅱ为正极：\ce{MnO2 + 2e- + 4H+ = Mn^2+ + 2H2O}，消耗 \ce{H+}，pH 升高，B 错误；充电时Ⅱ为阳极，A 对；Ⅰ为负极，C 对；充电时阴极 \ce{S + Mn^2+ + 2e- = MnS}，16 g S 为 0.5 mol，转移 1 mol 电子，D 对。答案：B。
\end{zhenti}

\section{电解池的工作原理}

\begin{gainian}[电解池]
借助电流引起氧化还原反应的装置（电能$\to$化学能）。与电源\textbf{负极}相连的是\textbf{阴极}，得电子发生还原反应；与电源\textbf{正极}相连的是\textbf{阳极}，失电子发生氧化反应。溶液中阳离子移向阴极，阴离子移向阳极。
\end{gainian}

\textbf{离子放电顺序（惰性电极，水溶液中）：}\\
阴极（阳离子得电子）：\ce{Ag+}$>$\allowbreak{}\ce{Fe^3+}$>$\allowbreak{}\ce{Cu^2+}$>$\allowbreak{}\ce{H+}（酸）$>$\allowbreak{}\ce{Pb^2+}$>$\allowbreak{}\ce{Sn^2+}$>$\allowbreak{}\ce{Fe^2+}$>$\allowbreak{}\ce{Zn^2+}$>$\allowbreak{}\ce{H+}（水）$>$\allowbreak{}\ce{Al^3+}$>$\allowbreak{}\ce{Mg^2+}$>$\allowbreak{}\ce{Na+}$>$\allowbreak{}\ce{Ca^2+}$>$\allowbreak{}\ce{K+}。\ce{Al^3+} 及其后的离子在水溶液中不放电，只能电解熔融物得到其单质。\\
阳极：\textbf{活性电极}（Pt、Au、石墨以外的金属）\textbf{优先自身失电子溶解}；惰性电极时阴离子放电：\ce{S^2-}$>$\allowbreak{}\ce{I-}$>$\allowbreak{}\ce{Br-}$>$\allowbreak{}\ce{Cl-}$>$\allowbreak{}\ce{OH-}$>$\allowbreak{}含氧酸根$>$\allowbreak{}\ce{F-}。

\begin{center}
\small
\begin{tabularx}{\linewidth}{p{2.2cm}p{2.6cm}Xp{2.6cm}}
\toprule
类型 & 实例 & 电极产物与溶液变化 & 复原方法 \\
\midrule
电解水型 & \ce{H2SO4}、NaOH、\ce{Na2SO4} & 阴极 \ce{H2}、阳极 \ce{O2}；溶质浓度增大 & 加水 \\
电解电解质型 & HCl、\ce{CuCl2} & 阴极 \ce{H2} 或 Cu，阳极 \ce{Cl2} & 通 HCl、加 \ce{CuCl2} \\
放 \ce{H2} 生碱型 & NaCl & 阴极 \ce{H2}，阳极 \ce{Cl2}；生成 NaOH & 通 HCl \\
放 \ce{O2} 生酸型 & \ce{CuSO4}、\ce{AgNO3} & 阴极 Cu（Ag），阳极 \ce{O2}；生成酸 & 加 CuO（\ce{Ag2O}） \\
\bottomrule
\end{tabularx}
\end{center}

\begin{yicuo}
\begin{enumerate}
\item 复原遵循“\textbf{出什么加什么}”：电解 \ce{CuSO4} 溶液析出 Cu 和 \ce{O2}，应加 CuO（或 \ce{CuCO3}），不能加 \ce{Cu(OH)2}（会多加水）——除非 \ce{Cu^2+} 放电完后又电解了水。
\item 水中 \ce{H+}、\ce{OH-} 放电时，电极反应可写作 \ce{2H2O + 2e- = H2 ^ + 2OH-}（阴极）、\ce{2H2O - 4e- = O2 ^ + 4H+}（阳极）；碱性溶液中阳极也可写 \ce{4OH- - 4e- = O2 ^ + 2H2O}，酸性溶液中阴极可写 \ce{2H+ + 2e- = H2 ^}，以溶液中大量存在的粒子为准。
\end{enumerate}
\end{yicuo}

\begin{liti}[混合溶液的分阶段电解]
用惰性电极电解 1 L 含 0.1 mol \ce{CuSO4} 和 0.1 mol NaCl 的混合溶液。当电路中通过 0.3 mol 电子时：(1) 分析电解的三个阶段；(2) 求阴、阳两极产生的气体的物质的量；(3) 电解后溶液呈酸性还是碱性？
\tcblower
\silu 按放电顺序“排队”：阴极 \ce{Cu^2+}$>$\ce{H+}；阳极 \ce{Cl-}$>$\ce{OH-}。分别算出每种离子放电所需电子数，再按电子数分段。\\
\jie (1) \ce{Cl-} 全部放电需 0.1 mol 电子，\ce{Cu^2+} 全部放电需 0.2 mol 电子。第一阶段（0～0.1 mol 电子）：阴极析 Cu，阳极放 \ce{Cl2}，相当于电解 \ce{CuCl2}；第二阶段（0.1～0.2 mol）：阴极析 Cu，阳极放 \ce{O2}，相当于电解 \ce{CuSO4}，生成 \ce{H2SO4}；第三阶段（0.2 mol 以后）：阴极放 \ce{H2}，阳极放 \ce{O2}，相当于电解水。\\
(2) 阳极：\ce{Cl2} 0.05 mol，\ce{O2} $\frac{0.3-0.1}{4}=0.05$ mol，共 0.1 mol；阴极：Cu 0.1 mol，\ce{H2} $\frac{0.3-0.2}{2}=0.05$ mol。\\
(3) 第二阶段生成 0.1 mol \ce{H+}（$2\ce{Cu^2+} + 2\ce{H2O} = 2\ce{Cu} + \ce{O2} + 4\ce{H+}$，0.1 mol 电子对应 0.1 mol \ce{H+}），第三阶段电解水不改变 \ce{H+} 的量，溶液呈酸性。
\end{liti}

\section{电解原理的应用}

\subsection{氯碱工业、电镀、精炼与电冶金}
\textbf{氯碱工业：}\ce{2NaCl + 2H2O ->[\text{电解}] 2NaOH + H2 ^ + Cl2 ^}。阳极（钛网）\ce{2Cl- - 2e- = Cl2 ^}；阴极（碳钢网）\ce{2H2O + 2e- = H2 ^ + 2OH-}。中间用\textbf{阳离子交换膜}，只允许 \ce{Na+} 由阳极室进入阴极室，既能得到较纯的 NaOH，又能防止 \ce{Cl2} 与 NaOH 反应、防止 \ce{H2} 与 \ce{Cl2} 混合爆炸。

\begin{center}
\small
\begin{tabularx}{\linewidth}{p{1.8cm}XXX}
\toprule
& 电镀（铁件镀铜） & 电解精炼铜 & 电冶金 \\
\midrule
阳极 & 镀层金属 Cu：\ce{Cu - 2e- = Cu^2+} & 粗铜（含 Zn、Fe、Ni、Ag、Au 等） & 惰性电极 \\
阴极 & 镀件（铁件） & 纯铜 & 析出金属 \\
电解质 & 含镀层金属离子的溶液（\ce{CuSO4}） & \ce{CuSO4} 溶液 & 熔融 NaCl、\ce{MgCl2}、\ce{Al2O3}（加冰晶石） \\
特点 & 电解质浓度基本不变 & Zn、Fe、Ni 先溶解，Ag、Au 形成阳极泥；$c(\ce{Cu^2+})$ 略减小 & \ce{2Al2O3 ->[\text{电解}][\text{冰晶石}] 4Al + 3O2 ^}（熔融）；不用 \ce{AlCl3}（共价化合物，熔融不导电） \\
\bottomrule
\end{tabularx}
\end{center}

\subsection{离子交换膜与多室电解}
\begin{fangfa}[判断离子交换膜种类的思路]
先根据电极反应确定各室“多出了什么电荷”（阳极室消耗阴离子或生成 \ce{H+}，阴极室消耗阳离子或生成 \ce{OH-}），再根据\textbf{产品要求}和\textbf{要阻止的副反应}判断哪种离子需要通过：阳离子交换膜只让阳离子通过，阴离子交换膜只让阴离子通过，质子交换膜只让 \ce{H+} 通过，双极膜在中间解离出 \ce{H+} 和 \ce{OH-} 分别移向阴、阳极。
\end{fangfa}

\begin{zhenti}{2024年广东卷第16题\quad 来源：仓库所附真题文档}
一种基于氯碱工艺的新型电解池可用于湿法冶铁：左侧电极室为 NaOH 溶液，电极上附着 \ce{Fe2O3}（电解后得到 Fe），中间为阳离子交换膜，右侧为饱和食盐水，电极上产生气体。下列说法\textbf{不正确}的是
\xxb{阳极反应：\ce{2Cl- - 2e- = Cl2 ^}}{阴极区溶液中 \ce{OH-} 浓度逐渐升高}{理论上每消耗 1 mol \ce{Fe2O3}，阳极室溶液减少 213 g}{理论上每消耗 1 mol \ce{Fe2O3}，阴极室物质最多增加 138 g}
\textbf{解析：}左侧为阴极：\ce{Fe2O3 + 6e- + 3H2O = 2Fe + 6OH-}，$c(\ce{OH-})$ 升高，B 对；右侧为阳极：\ce{2Cl- - 2e- = Cl2 ^}，A 对。每消耗 1 mol \ce{Fe2O3} 转移 6 mol 电子：阳极室放出 3 mol \ce{Cl2}（213 g），同时有 6 mol \ce{Na+}（138 g）迁移到阴极室，阳极室共减少 $213+138=351$ g，C 错；阴极室增加的是迁入的 \ce{Na+}，最多 138 g（\ce{Fe2O3} 转化为 Fe 和 \ce{OH-} 仍在该室），D 对。答案：C。
\end{zhenti}

\begin{liti}[双极膜电渗析]
双极膜由阳离子交换层和阴离子交换层复合而成，在直流电场作用下，膜间的水解离为 \ce{H+} 和 \ce{OH-}，分别向两侧迁移。用“阳极$|$双极膜$|$酸室$|$阴离子交换膜$|$盐室$|$阳离子交换膜$|$碱室$|$双极膜$|$阴极”结构处理 NaCl 溶液，可在酸室得到 HCl、在碱室得到 NaOH。(1) 双极膜中解离出的 \ce{H+} 移向哪一极？(2) 盐室中 \ce{Na+}、\ce{Cl-} 分别向哪个方向迁移？在何处与 \ce{OH-}、\ce{H+} 结合？(3) 与直接电解 NaCl 溶液相比，这种方法的优点是什么？
\tcblower
\jie (1) \ce{H+} 是阳离子，移向阴极方向；\ce{OH-} 移向阳极方向。(2) \ce{Cl-} 透过阴离子交换膜向阳极方向迁移，与双极膜产生的 \ce{H+} 结合生成 HCl；\ce{Na+} 透过阳离子交换膜向阴极方向迁移，与双极膜产生的 \ce{OH-} 结合生成 NaOH。(3) 阳极不产生 \ce{Cl2}（两电极上只发生电解水，阳极放出 \ce{O2}、阴极放出 \ce{H2}），可直接得到较纯的酸和碱；工业上多组膜单元叠放只需一对电极，水在双极膜中解离所需电压低于电解水，能耗低，副产物少，更环保。
\end{liti}

\subsection{电化学计算}
核心是\textbf{电子守恒}：同一电路中各电极转移的电子数相等。常用关系：$4e^-\sim\ce{O2}\sim2\ce{H2}\sim2\ce{Cl2}\sim2\ce{Cu}\sim4\ce{Ag}\sim4\ce{H+}\sim4\ce{OH-}$。电量 $Q=It$，$n(e^-)=\dfrac{Q}{F}$（$F=96500$ C·mol$^{-1}$）。

\begin{liti}[电子守恒计算]
用惰性电极电解 500 mL \ce{CuSO4} 溶液，通电一段时间后阴极增重 0.64 g（无气体产生），设溶液体积不变。(1) 阳极产生气体的体积（标准状况）为多少？(2) 电解后溶液的 $c(\ce{H+})$ 为多少？(3) 若用 0.5 A 的恒定电流，需通电多少秒？
\tcblower
\jie (1) $n(\ce{Cu})=0.01$ mol，转移 0.02 mol 电子；阳极 \ce{2H2O - 4e- = O2 ^ + 4H+}，$n(\ce{O2})=0.005$ mol，$V=0.112$ L。(2) 生成 \ce{H+} 0.02 mol，$c(\ce{H+})=0.04$ mol·L$^{-1}$。(3) $t=\dfrac{nF}{I}=\dfrac{0.02\times96500}{0.5}=3860$ s。
\end{liti}

\section{金属的腐蚀与防护}

\textbf{化学腐蚀}：金属与接触的物质直接反应（如高温下铁与氯气）；\textbf{电化学腐蚀}：不纯的金属与电解质溶液接触形成原电池，较活泼的金属被氧化，更普遍、危害更大。

\begin{center}
\small
\begin{tabular}{lll}
\toprule
& 析氢腐蚀 & 吸氧腐蚀 \\
\midrule
条件 & 水膜酸性较强 & 水膜呈中性、弱酸性或碱性（最普遍） \\
负极 & \ce{Fe - 2e- = Fe^2+} & \ce{Fe - 2e- = Fe^2+} \\
正极 & \ce{2H+ + 2e- = H2 ^} & \ce{O2 + 4e- + 2H2O = 4OH-} \\
后续 & & \ce{Fe^2+ + 2OH- = Fe(OH)2}，\ce{4Fe(OH)2 + O2 + 2H2O = 4Fe(OH)3}，最终形成铁锈 \ce{Fe2O3*$x$H2O} \\
\bottomrule
\end{tabular}
\end{center}

\textbf{防护方法：}①改变金属组成（不锈钢）；②覆盖保护层（涂漆、电镀、搪瓷、钝化）；③电化学保护：\textbf{牺牲阳极法}（原电池原理，连接更活泼的 Zn、Mg 块作负极，被保护金属作正极；锌块会逐渐消耗，需定期更换）和\textbf{外加电流法}（电解原理，被保护金属接电源\textbf{负极}作阴极，辅助阳极接正极）。腐蚀快慢：电解池阳极$>$原电池负极$>$化学腐蚀$>$有防护措施的金属。

\begin{liti}[析氢腐蚀与吸氧腐蚀的实验判断]
将两枚相同的铁钉分别放入两支具支试管：甲中为食盐水浸过的铁钉，乙中为稀醋酸浸过的铁钉，支管均连接插入红墨水中的导管。一段时间后，甲导管中液面上升，乙导管中液面下降。解释原因，并写出甲中正极反应式。
\tcblower
\jie 甲中食盐水呈中性，铁发生吸氧腐蚀，消耗试管内的 \ce{O2}，气体压强减小，导管液面上升；正极：\ce{O2 + 4e- + 2H2O = 4OH-}。乙中稀醋酸呈酸性，发生析氢腐蚀，生成 \ce{H2} 使压强增大，导管液面下降。
\end{liti}

\begin{zhenti}{2026年广东卷第6题\quad 来源：仓库所附真题文档（网络收集版）}
做好防腐是发展海上风电的基础。下列措施\textbf{不能}达到防腐目的的是：A．钢制管桩连接铜块；B．钢制管桩外加富铝涂层；C．钢制管桩接外电源，另一端连接惰性电极；D．钢制管桩连接锌块。

\textbf{解析：}A 中 Cu 不如 Fe 活泼，Fe 作负极被加速腐蚀，不能防腐；B 富铝涂层既隔绝海水，铝又比铁活泼可作牺牲阳极；C 为外加电流法（管桩应接电源负极）；D 为牺牲阳极法。答案：A。
\end{zhenti}

% >>>>>>>> 开始：chapters/ch05ex.tex
% ---------------- 第5章 章末 ----------------
\zhangmo{本章小结}
\subsection*{一、知识总结}
\begin{center}
\small
\begin{tabularx}{\linewidth}{p{1.6cm}XX}
\toprule
& 原电池 & 电解池 \\
\midrule
能量转化 & 化学能$\to$电能（自发） & 电能$\to$化学能（非自发也能进行） \\
电极名称 & 负极（氧化）、正极（还原） & 阳极（氧化，接电源正极）、阴极（还原，接电源负极） \\
离子迁移 & 阳离子$\to$正极，阴离子$\to$负极 & 阳离子$\to$阴极，阴离子$\to$阳极 \\
典型应用 & 化学电源、加快反应、牺牲阳极法 & 氯碱、电镀、精炼、电冶金、外加电流法 \\
\bottomrule
\end{tabularx}
\end{center}

\subsection*{二、方法提升}
\begin{enumerate}
\item \textbf{判电极}：看化合价变化（升高为负极/阳极），或看外接电源，或看现象（溶解、析出、放气）。
\item \textbf{写反应}：先总后分、按介质补粒子；燃料电池“负极$=$总反应$-$正极”。
\item \textbf{判膜}：从电荷平衡和产品需求出发，确定哪种离子必须通过。
\item \textbf{算数量}：抓住电子守恒，建立 $e^-$ 与各物质的关系式。
\end{enumerate}

\subsection*{三、易错点归纳}
\begin{enumerate}
\item 认为活泼金属一定作负极（Mg–Al–NaOH、Fe–Cu–浓硝酸是反例）。
\item 碱性介质中电极反应出现 \ce{H+} 或 \ce{CO2}；酸性介质中出现 \ce{OH-}。
\item 认为电子可以通过电解质溶液；把离子迁移方向弄反。
\item 电解时忽略活性阳极自身溶解；电解后复原“加错物质”。
\item 外加电流法中把被保护金属接电源正极（应接负极）。
\end{enumerate}

\zhangmo{分层练习}
\zu{A组\quad 基础巩固}
\begin{xiti}[start=1]
\item 关于 Zn–Cu 双液原电池（盐桥中为 KCl 琼胶），下列说法正确的是\xz
\xxa{电子由锌片经盐桥流向铜片}{盐桥中的 \ce{K+} 移向 \ce{ZnSO4} 溶液}{铜片上发生的反应为 \ce{Cu^2+ + 2e- = Cu}}{取出盐桥，电流计指针仍发生偏转}

\item 将镁片和铝片用导线相连插入 NaOH 溶液中，下列说法正确的是\xz
\xxa{镁比铝活泼，镁作负极}{负极反应为 \ce{Al - 3e- + 4OH- = [Al(OH)4]-}}{正极上产生 \ce{O2}}{电子由镁片经导线流向铝片}

\item 以 KOH 溶液为电解质的甲烷燃料电池中，负极反应式为\xz
\xxa{\ce{CH4 - 8e- + 2H2O = CO2 + 8H+}}{\ce{CH4 - 8e- + 10OH- = CO3^2- + 7H2O}}{\ce{O2 + 4e- + 2H2O = 4OH-}}{\ce{CH4 + 8e- + 10OH- = CO3^2- + 7H2O}}

\item 关于铅酸蓄电池，下列说法正确的是\xz
\xxa{放电时，负极质量减小}{放电时，电解质溶液的密度增大}{充电时，铅酸蓄电池的负极应与外接电源的正极相连}{充电时，阳极反应为 \ce{PbSO4 - 2e- + 2H2O = PbO2 + 4H+ + SO4^2-}}

\item 用惰性电极电解足量 \ce{CuSO4} 溶液一段时间后，要使溶液恢复到原来的浓度和体积，应加入适量的\xz
\xx{Cu}{CuO}{\ce{Cu(OH)2}}{\ce{CuSO4}}

\item \laiyuan{2024年广东卷第5题，来源：仓库所附真题文档}\ 我国自主设计建造的浮式生产储卸油装置“海葵一号”将在珠江口盆地海域使用，其钢铁外壳镶嵌了锌块，以利用电化学原理延缓外壳的腐蚀。下列有关说法正确的是\xz
\xxb{钢铁外壳为负极}{镶嵌的锌块可永久使用}{该法为外加电流法}{锌发生反应：\ce{Zn - 2e- = Zn^2+}}

\item 用铜作电极电解稀硫酸，下列说法正确的是\xz
\xxa{阳极上放出 \ce{O2}}{阴极上析出 Cu}{溶液逐渐变蓝，阴极上放出 \ce{H2}}{电解一段时间后溶液 pH 减小}
\end{xiti}

\zu{B组\quad 高考典型}
\begin{xiti}
\item \laiyuan{2025年广东卷第16题，来源：仓库所附真题文档}\ 某理论研究认为：氢氧燃料电池（图b，两电极间为质子交换膜）的电极Ⅰ和Ⅱ上所发生反应的催化机理分别如图a和图c。图a：\ce{O2} 在电极Ⅰ表面的催化剂上依次经过 i$\to$ii$\to$iii$\to$iv 四步，每步结合 1 个 \ce{H+} 和 1 个 $e^-$，最终生成 \ce{H2O}；图c：\ce{H2} 在电极Ⅱ表面经两步各失去 1 个电子生成 \ce{H+}。其中 \ce{O2} 获得第一个电子的过程（i$\to$ii）最慢。由此可知，理论上\xz
\xxb{负极反应的催化剂是 i}{图a中，i 到 ii 过程的活化能一定最低}{电池工作过程中，负极室的溶液质量保持不变}{相同时间内，电极Ⅰ和电极Ⅱ上的催化循环完成次数相同}

\item \laiyuan{原创}\ 磷酸铁锂电池的总反应为 \ce{Li_xC6 + Li_{1-$x$}FePO4 <=>[\text{放电}][\text{充电}] 6C + LiFePO4}，电解质为能传导 \ce{Li+} 的有机溶液。(1) 写出放电时正极的电极反应式。(2) 充电时，\ce{Li+} 向哪一极迁移？阳极反应式是什么？(3) 放电时每转移 1 mol 电子，负极材料质量变化多少？(4) 为什么该电池不能用水溶液作电解质？

\item 下列关于电镀和电解精炼的说法正确的是\xz
\xxa{在铁件上镀锌，铁件作阳极，锌片作阴极}{电镀过程中，电解质溶液中镀层金属离子的浓度明显减小}{电解精炼铜时，粗铜作阳极，阳极减少的质量等于阴极增加的质量}{电解精炼铜时，Ag、Au 等不如铜活泼的金属沉积在阳极底部形成阳极泥}

\item \laiyuan{原创}\ 用如下三室装置电解 \ce{Na2SO4} 溶液制备 NaOH 和 \ce{H2SO4}：左侧为阳极室（稀硫酸），中间室为 \ce{Na2SO4} 溶液，右侧为阴极室（稀 NaOH 溶液），两张膜分别为 X 和 Y（X 在阳极室与中间室之间）。(1) X、Y 分别是阴离子交换膜还是阳离子交换膜？(2) 写出阳极、阴极的电极反应式。(3) 以 2.0 A 电流电解 965 s，理论上阴极室生成 NaOH 的物质的量为多少？中间室 \ce{Na2SO4} 减少多少克？

\item \laiyuan{原创}\ 锌—空气电池以 KOH 溶液为电解质，总反应为 \ce{2Zn + O2 + 4OH- + 2H2O = 2[Zn(OH)4]^2-}。(1) 写出正、负极的电极反应式。(2) 放电时，\ce{OH-} 向哪一极移动？电解质溶液中 $c(\ce{OH-})$ 如何变化？(3) 每消耗 6.5 g Zn，理论上消耗标准状况下空气多少升（空气中 \ce{O2} 体积分数按 20\% 计）？

\item \laiyuan{原创}\ 将两个电解池串联：甲池为惰性电极电解足量 \ce{AgNO3} 溶液，乙池为电解精炼铜（粗铜含少量 Zn、Ag）。通电一段时间后甲池阴极增重 10.8 g。(1) 乙池阴极增重多少克？(2) 乙池阳极减少的质量是否等于阴极增加的质量？为什么？(3) 甲池阳极产生气体在标准状况下的体积为多少？
\end{xiti}

\zu{C组\quad 冲刺高分}
\begin{xiti}
\item \laiyuan{2026年广东卷第16题，来源：仓库所附真题文档（网络收集版）}\ 一种利用电解还原 \ce{CO2} 同时实现海水脱盐的装置：从左到右依次为电极a、Ⅰ室、膜①、Ⅱ室（海水）、膜②、Ⅲ室（海水）、膜③、Ⅳ室、电极b；Ⅰ室中通入 \ce{CO2} 并放出 CO，Ⅳ室放出 \ce{O2}；电解液在Ⅰ和Ⅳ室间循环，初始电解液为 \ce{NaHCO3} 溶液；膜①～③为阳离子或阴离子交换膜。下列说法错误的是\xz
\xxb{①为阳离子交换膜，②为阴离子交换膜}{阳极反应为 \ce{4OH- - 4e- = 2H2O + O2 ^}}{电解一段时间后，Ⅲ室海水中盐的浓度降低}{电解过程中，Ⅰ室电解液的 pH 在一段时间内可保持稳定}

\item \laiyuan{原创}\ 以质子交换膜甲醇燃料电池为电源，用惰性电极电解 200 mL 0.50 mol·L$^{-1}$ \ce{CuSO4} 溶液。(1) 写出燃料电池负极、正极的电极反应式。(2) 当燃料电池消耗 0.32 g 甲醇时，电解池阴极析出铜多少克？阳极产生气体在标准状况下的体积为多少？(3) 此时电解液中 $c(\ce{H+})$ 为多少（忽略体积变化）？要使电解液复原，应加入什么物质、多少克？(4) 若继续电解至阴极开始产生气体，此后电解的实质是什么？

\item \laiyuan{原创}\ 用电解法将 \ce{Na2CrO4} 转化为 \ce{Na2Cr2O7}：电解槽用阳离子交换膜隔开，阳极室加入 \ce{Na2CrO4} 溶液，阴极室加入稀 NaOH 溶液，两极均为惰性电极。已知 \ce{2CrO4^2- + 2H+ <=> Cr2O7^2- + H2O}。(1) 写出阳极、阴极的电极反应式。(2) 解释阳极室中 \ce{Na2CrO4} 转化为 \ce{Na2Cr2O7} 的原理。(3) 通过阳离子交换膜迁移的离子是什么？向哪个方向迁移？(4) 若阳极室溶液中 Na 与 Cr 的物质的量之比由开始的 2:1 变为 $d$:1，则此时 \ce{Na2CrO4} 的转化率为\tk（用 $d$ 表示）。
\end{xiti}

\zhangmo{练习解析}
\begin{jiexilist}
\item \ans{C}电子只经导线流动，A 错；盐桥中阳离子 \ce{K+} 移向正极区（\ce{CuSO4} 溶液），B 错；铜为正极，\ce{Cu^2+} 得电子，C 对；取出盐桥，回路断开，无电流，D 错。
\item \ans{B}Mg 不与 NaOH 反应，Al 与 NaOH 反应，Al 作负极，A 错、B 对；正极上 \ce{H2O} 得电子生成 \ce{H2}，C 错；电子由负极（铝）流向正极（镁），D 错。
\item \ans{B}碱性介质中 C 元素以 \ce{CO3^2-} 形式存在，不能写 \ce{H+} 和 \ce{CO2}；C 是正极反应；D 中负极应失电子。
\item \ans{D}放电时负极 Pb 生成 \ce{PbSO4}，质量增大，A 错；消耗 \ce{H2SO4}、生成水，溶液密度减小，B 错；充电时“负接负”，C 错；充电阳极是放电正极反应的逆反应，D 对。
\item \ans{B}电解时阴极析出 Cu，阳极放出 \ce{O2}，“出什么加什么”，相当于减少了 CuO，应加 CuO；加 \ce{Cu(OH)2} 会多加水。
\item \ans{D}Zn 比 Fe 活泼，锌块为负极，钢铁外壳为正极被保护，A 错；锌不断溶解，需定期更换，B 错；该法为牺牲阳极法，C 错；D 对。
\item \ans{C}阳极 Cu 为活性电极，自身溶解：\ce{Cu - 2e- = Cu^2+}，溶液变蓝；阴极开始时 \ce{H+} 放电产生 \ce{H2}（\ce{H+} 浓度大，生成的 \ce{Cu^2+} 较少），总反应 \ce{Cu + H2SO4 ->[\text{电解}] CuSO4 + H2 ^}，消耗 \ce{H+}，pH 增大。
\item \ans{C}电极Ⅰ上 \ce{O2} 被还原，为正极，i 是正极反应的催化剂，A 错；i$\to$ii 最慢，活化能最高，B 错；负极 \ce{H2 - 2e- = 2H+}，每失去 1 个电子就有 1 个 \ce{H+} 通过质子交换膜进入正极室，负极室溶液质量不变，C 对；\ce{O2} 完成一次循环需 4 个电子，\ce{H2} 完成一次循环需 2 个电子，同一时间转移电子数相同，两者循环次数之比为 1:2，D 错。
\item (1) \ce{Li_{1-$x$}FePO4 + $x$Li+ + $x$e- = LiFePO4}。(2) 充电时阳离子移向阴极（即放电时的负极，石墨一侧）；阳极：\ce{LiFePO4 - $x$e- = Li_{1-$x$}FePO4 + $x$Li+}。(3) 负极 \ce{Li_xC6 - $x$e- = $x$Li+ + 6C}，每转移 1 mol 电子，负极减少 1 mol Li，即质量减少 7 g。(4) Li 及嵌锂石墨很活泼，会与水剧烈反应。
\item \ans{D}镀锌时铁件作阴极、锌片作阳极，A 错；电镀时阳极溶解的金属离子与阴极析出的相等，浓度基本不变，B 错；精炼时阳极上 Zn、Fe、Ni 等也溶解且有阳极泥，两极质量变化不相等，C 错；D 对。
\item (1) 阳极室需要 \ce{SO4^2-} 进入（与生成的 \ce{H+} 组成 \ce{H2SO4}），X 为阴离子交换膜；阴极室需要 \ce{Na+} 进入（与生成的 \ce{OH-} 组成 NaOH），Y 为阳离子交换膜。(2) 阳极：\ce{2H2O - 4e- = O2 ^ + 4H+}；阴极：\ce{2H2O + 2e- = H2 ^ + 2OH-}。(3) $n(e^-)=\dfrac{2.0\times965}{96500}=0.020$ mol；生成 NaOH 0.020 mol；中间室迁出 0.020 mol \ce{Na+} 和 0.010 mol \ce{SO4^2-}，即减少 0.010 mol \ce{Na2SO4}，1.42 g。
\item (1) 负极：\ce{Zn - 2e- + 4OH- = [Zn(OH)4]^2-}；正极：\ce{O2 + 4e- + 2H2O = 4OH-}。(2) 移向负极；总反应消耗 \ce{OH-}，$c(\ce{OH-})$ 减小。(3) 0.1 mol Zn 失 0.2 mol 电子，需 0.05 mol \ce{O2}，即 1.12 L，空气 $1.12\div20\%=5.6$ L。
\item (1) 甲池析出 Ag 0.1 mol，转移 0.1 mol 电子；乙池阴极析出 Cu 0.05 mol，增重 3.2 g。(2) 不相等。粗铜中的 Zn 也失电子溶解（摩尔质量与 Cu 不同），Ag 不溶解而落入阳极泥，所以阳极减少的质量与阴极增加的质量不相等，但两极转移的电子数相等。(3) 甲池阳极：\ce{2H2O - 4e- = O2 ^ + 4H+}，0.1 mol 电子生成 0.025 mol \ce{O2}，即 0.56 L。
\item \ans{C}a 极 \ce{CO2} 得电子生成 CO，为阴极：\ce{CO2 + H2O + 2e- = CO + 2OH-}，生成的 \ce{OH-} 与 \ce{HCO3-} 反应生成 \ce{CO3^2-}；b 极为阳极，水电离的 \ce{OH-} 放电生成 \ce{O2}，B 对。若膜③为阴离子交换膜，Ⅲ室 \ce{Cl-} 会进入Ⅳ室在阳极放电生成 \ce{Cl2}，故③应为阳离子交换膜（Ⅳ室 \ce{Na+} 进入Ⅲ室）；要使某室脱盐，该室的 \ce{Na+}、\ce{Cl-} 必须都能移出：Ⅱ室 \ce{Na+} 经膜①（阳膜）移入Ⅰ室，\ce{Cl-} 经膜②（阴膜）移入Ⅲ室，A 对。于是Ⅱ室盐浓度降低（脱盐），Ⅲ室得到 \ce{Cl-}（来自Ⅱ）和 \ce{Na+}（来自Ⅳ），盐浓度升高，C 错。Ⅰ室中 \ce{HCO3-}/\ce{CO3^2-} 构成缓冲体系，且与Ⅳ室电解液循环，pH 可在一段时间内保持稳定，D 对。
\item (1) 负极：\ce{CH3OH - 6e- + H2O = CO2 ^ + 6H+}；正极：\ce{O2 + 4e- + 4H+ = 2H2O}。(2) $n(\ce{CH3OH})=0.01$ mol，转移 0.06 mol 电子。$n(\ce{Cu^2+})=0.10$ mol，足量，析出 Cu 0.03 mol，即 1.92 g；阳极 \ce{O2} 0.015 mol，0.336 L。(3) 生成 \ce{H+} 0.06 mol，$c(\ce{H+})=0.30$ mol·L$^{-1}$；溶液中减少的是 0.03 mol Cu 和 0.015 mol \ce{O2}，相当于 0.03 mol CuO，加入 CuO 2.4 g（或 \ce{CuCO3} 3.72 g）。(4) \ce{Cu^2+} 完全放电后，阴极上 \ce{H+}（来自水和生成的 \ce{H2SO4}）放电，阳极仍放出 \ce{O2}，实质是电解水。
\item (1) 阳极：\ce{2H2O - 4e- = O2 ^ + 4H+}；阴极：\ce{2H2O + 2e- = H2 ^ + 2OH-}。(2) 阳极生成 \ce{H+}，使 \ce{2CrO4^2- + 2H+ <=> Cr2O7^2- + H2O} 正向移动。(3) \ce{Na+}，由阳极室迁移到阴极室（阴极室得到 NaOH）。(4) $2-d$。设起始 \ce{Na2CrO4} 1 mol（Na 2 mol，Cr 1 mol），转化率为 $\alpha$：生成 $\alpha$ mol \ce{H+} 需转移 $\alpha$ mol 电子，同时有 $\alpha$ mol \ce{Na+} 迁出阳极室，剩余 Na $(2-\alpha)$ mol，Cr 仍为 1 mol。$\dfrac{2-\alpha}{1}=d$，$\alpha=2-d$。
\end{jiexilist}
% <<<<<<<< 结束：chapters/ch05ex.tex
% <<<<<<<< 结束：chapters/ch05.tex

% >>>>>>>> 开始：chapters/ch06.tex
% ============================================================
% 第6章 原子结构与元素性质（选必2第一章）
% ============================================================
\chapter{原子结构与元素性质}

\begin{daoxue}
\textbf{知识主线：}电子在原子核外怎样排布（能层、能级、原子轨道、构造原理、三条排布规则）$\to$ 排布决定元素在周期表中的位置（分区）$\to$ 位置决定性质的周期性变化（原子半径、电离能、电负性）$\to$ 用“结构—位置—性质”推断元素。

\textbf{高考考情：}广东卷选择题每年有一道元素推断（2024年第12题由结构式推断、2025年第9题由第一电离能图推断、2026年第11题由轨道“半充满”推断）；大题中常考价层电子排布式或轨道表示式（2025年第19题(1) Ti、2026年第19题(1) S）。

\textbf{学习目标：}①正确书写 1～36 号元素的电子排布式、价层电子排布式和轨道表示式；②理解电离能、电负性的递变规律及特例；③熟练运用“位—构—性”推断元素。
\end{daoxue}

\section{原子核外电子排布}

\subsection{能层、能级与原子轨道}
核外电子按能量不同分成\textbf{能层}（$n=1,2,3,\cdots$，符号 K、L、M、N…），每个能层又分为若干\textbf{能级}（s、p、d、f），能层序数 $n$ 就是该能层的能级数。电子在核外空间出现的概率分布用\textbf{电子云}描述，人们把电子出现概率约 90\% 的空间圈出来，得到\textbf{电子云轮廓图}；量子力学把电子在原子核外的一个空间运动状态称为一个\textbf{原子轨道}：s 轨道呈球形（每个 s 能级 1 个轨道），p 轨道呈哑铃形（每个 p 能级 3 个相互垂直的轨道 $p_x$、$p_y$、$p_z$），d 能级有 5 个轨道，f 能级有 7 个轨道。每个原子轨道最多容纳 2 个自旋相反的电子，所以 s、p、d、f 能级最多容纳 2、6、10、14 个电子，第 $n$ 能层最多容纳 $2n^2$ 个电子。

\subsection{构造原理与排布规则}
\begin{gainian}[构造原理与三条规则]
\textbf{构造原理}（能级填充顺序）：1s 2s 2p 3s 3p \textbf{4s 3d} 4p 5s 4d 5p 6s 4f 5d 6p…（出现“能级交错”：4s 能量低于 3d，先填 4s）。\\
\textbf{能量最低原理}：基态原子电子优先占据能量低的轨道，使整个原子能量最低。\\
\textbf{泡利原理}：一个原子轨道最多容纳 2 个电子，且自旋相反。\\
\textbf{洪特规则}：电子填入同一能级的不同轨道时，优先\textbf{单独}分占不同轨道，且自旋平行。特例：能级处于全满（$p^6$、$d^{10}$）、半满（$p^3$、$d^5$）或全空时较稳定，如 Cr：[Ar]$3d^54s^1$，Cu：[Ar]$3d^{10}4s^1$。
\end{gainian}

\begin{center}
\small
\begin{tabular}{lll}
\toprule
表示方法 & 以 Fe（26 号）为例 & 说明 \\
\midrule
电子排布式 & $1s^22s^22p^63s^23p^63d^64s^2$ & 书写时按能层顺序，3d 写在 4s 前 \\
简化排布式 & [Ar]$3d^64s^2$ & 用稀有气体结构代替内层 \\
价层电子排布式 & $3d^64s^2$ & 主族元素为最外层，过渡元素含次外层 d \\
价层轨道表示式 & 3d：\fbox{$\uparrow\downarrow$}\fbox{$\uparrow$}\fbox{$\uparrow$}\fbox{$\uparrow$}\fbox{$\uparrow$}\ 4s：\fbox{$\uparrow\downarrow$} & 4 个未成对电子 \\
离子 & \ce{Fe^2+}：[Ar]$3d^6$；\ce{Fe^3+}：[Ar]$3d^5$ & \textbf{先失最外层（4s）电子} \\
\bottomrule
\end{tabular}
\end{center}

\begin{benzhi}[为什么 \ce{Fe^3+} 比 \ce{Fe^2+} 稳定，\ce{Cu+} 在固态下比 \ce{Cu^2+} 稳定]
\ce{Fe^3+} 的 3d 能级为 $3d^5$ 半充满，较稳定，所以 \ce{Fe^2+}（$3d^6$）易被氧化为 \ce{Fe^3+}。\ce{Cu+} 为 $3d^{10}$ 全充满，所以高温下 CuO 分解生成 \ce{Cu2O}。但在水溶液中 \ce{Cu+} 易歧化（\ce{2Cu+ = Cu + Cu^2+}），因为 \ce{Cu^2+} 的水合放出更多能量——判断稳定性还要看具体环境。
\end{benzhi}

\begin{yicuo}
\begin{enumerate}
\item 轨道表示式中违反泡利原理（同一方框内两电子自旋相同）或违反洪特规则（未分占、或自旋不平行）都是错误的。
\item 书写过渡元素阳离子的电子排布时，不能先失去 3d 电子。
\item 基态原子“未成对电子数”要画轨道表示式数：第四周期 Cr 有 6 个（$3d^54s^1$），是第四周期最多的。
\end{enumerate}
\end{yicuo}

\begin{liti}[电子排布的书写]
(1) 写出基态 Ga、Br 原子的简化电子排布式；(2) 写出 \ce{Cu^2+}、\ce{Mn^2+} 的价层电子排布式；(3) 基态 Co、Ni 原子的未成对电子数分别是多少？(4) 第四周期中，基态原子未成对电子数与 Fe 相同的元素有哪些？
\tcblower
\jie (1) Ga：[Ar]$3d^{10}4s^24p^1$（不要漏掉 $3d^{10}$）；Br：[Ar]$3d^{10}4s^24p^5$。(2) \ce{Cu^2+}：$3d^9$（Cu 为 $3d^{10}4s^1$，先失 4s 上 1 个电子，再失 3d 上 1 个）；\ce{Mn^2+}：$3d^5$。(3) Co：$3d^74s^2$，3d 上 7 个电子，5 个轨道中 2 个成对、3 个单电子，未成对电子数 3；Ni：$3d^84s^2$，未成对电子数 2。(4) 没有。Fe（$3d^64s^2$）有 4 个未成对电子；未成对电子数为 4 只可能是 $3d^4$ 或 $3d^6$ 构型，而 24 号 Cr 因半满稳定为 $3d^54s^1$（6 个未成对电子），所以第四周期中只有 Fe 有 4 个未成对电子。
\end{liti}

\subsection{原子光谱}
基态原子吸收能量，电子跃迁到较高能级成为激发态；激发态电子跃迁回较低能级时以光的形式释放能量。不同元素的原子发生跃迁时吸收或发射的光的频率不同，形成特征的\textbf{原子光谱}，可用于鉴定元素（光谱分析）。\textbf{焰色}就是电子由激发态跃迁到低能级时放出的光，属于发射光谱，是物理变化。

\section{原子结构与元素周期表}
\textbf{周期}：周期序数$=$能层数；每一周期从 $ns^1$ 开始到 $ns^2np^6$ 结束（第一周期到 $1s^2$）。\textbf{族}：主族序数$=$价层电子数；第ⅢB～ⅦB族序数$=$价层电子（$(n-1)d+ns$）数。\textbf{分区}：

\begin{center}
\small
\begin{tabular}{llll}
\toprule
分区 & 范围 & 价层电子排布 & 特点 \\
\midrule
s 区 & ⅠA、ⅡA（含 H） & $ns^{1\sim2}$ & 活泼金属（H 除外） \\
p 区 & ⅢA～ⅦA、0 族（含 He） & $ns^2np^{1\sim6}$（He 为 $1s^2$） & 包括全部非金属（H 除外） \\
d 区 & ⅢB～ⅦB、Ⅷ族 & $(n-1)d^{1\sim9}ns^{1\sim2}$（Pd 例外） & 过渡金属，多种化合价 \\
ds 区 & ⅠB、ⅡB & $(n-1)d^{10}ns^{1\sim2}$ & Cu、Zn 等 \\
f 区 & 镧系、锕系 & $(n-2)f\,(n-1)d\,ns$ & 性质相似 \\
\bottomrule
\end{tabular}
\end{center}

\section{元素周期律}

\subsection{原子半径}
\textbf{同周期}从左到右，核电荷数增大而能层数不变，核对外层电子的吸引增强，原子半径减小；\textbf{同主族}从上到下，能层数增加，原子半径增大。\textbf{离子半径}：电子层结构相同的离子，核电荷数越大，半径越小，如 $r(\ce{O^2-})>r(\ce{F-})>r(\ce{Na+})>r(\ce{Mg^2+})>r(\ce{Al^3+})$；同种元素 $r(\text{阴离子})>r(\text{原子})>r(\text{阳离子})$，$r(\ce{Fe^2+})>r(\ce{Fe^3+})$。

\subsection{电离能}
\begin{gainian}[第一电离能 $I_1$]
气态电中性基态原子失去一个电子转化为气态基态正离子所需的最低能量。$I_1$ 越小，越容易失电子，金属性越强。同一原子逐级电离能 $I_1<I_2<I_3<\cdots$。
\end{gainian}

\textbf{递变规律：}同周期从左到右 $I_1$ 呈增大趋势，同主族从上到下减小。\textbf{特例}：ⅡA 族（$ns^2$ 全满）和 ⅤA 族（$np^3$ 半满）的 $I_1$ 大于同周期相邻元素。如第二周期：Li$<$B$<$Be$<$C$<$O$<$N$<$F$<$Ne；第三周期：Na$<$Al$<$Mg$<$Si$<$S$<$P$<$Cl$<$Ar。

\begin{center}
\begin{tikzpicture}
\begin{axis}[width=10cm,height=4.6cm,xlabel={原子序数},ylabel={$I_1$/(kJ·mol$^{-1}$)},xmin=0,xmax=21,ymin=0,ymax=2600,
 xtick={1,2,...,20},tick label style={font=\scriptsize},label style={font=\small},grid=major,grid style={gray!20}]
\addplot[thick,mainblue,mark=*,mark size=1.4pt] coordinates {(1,1312)(2,2372)(3,520)(4,900)(5,801)(6,1086)(7,1402)(8,1314)(9,1681)(10,2081)(11,496)(12,738)(13,578)(14,787)(15,1012)(16,1000)(17,1251)(18,1521)(19,419)(20,590)};
\node[font=\scriptsize,above] at (axis cs:2,2372) {He}; \node[font=\scriptsize,above] at (axis cs:10,2081) {Ne}; \node[font=\scriptsize,above] at (axis cs:18,1521) {Ar};
\node[font=\scriptsize,below] at (axis cs:3,520) {Li}; \node[font=\scriptsize,below] at (axis cs:11,496) {Na}; \node[font=\scriptsize,below] at (axis cs:19,419) {K};
\node[font=\scriptsize,above] at (axis cs:4,900) {Be}; \node[font=\scriptsize,above] at (axis cs:7,1402) {N}; \node[font=\scriptsize,above] at (axis cs:12,738) {Mg}; \node[font=\scriptsize,above] at (axis cs:15,1012) {P};
\end{axis}
\end{tikzpicture}
\end{center}
\noindent\small 图：1～20 号元素的第一电离能。每一周期从碱金属（谷）到稀有气体（峰）总体增大；ⅡA（Be、Mg）、ⅤA（N、P）出现“凸起”。2025年广东卷第9题正是据此图形推断元素。\normalsize

\textbf{逐级电离能的应用}：若 $I_{n+1}\gg I_n$（出现突跃），说明失去 $n$ 个电子后达到稳定结构，该元素最外层有 $n$ 个电子，常见化合价为 $+n$。如某元素 $I_1\sim I_4$（kJ·mol$^{-1}$）为 578、1817、2745、11577，$I_4$ 突跃，为 ⅢA 族的 Al。

\begin{liti}[由电离能推断元素]
短周期元素 X、Y 的逐级电离能（kJ·mol$^{-1}$）如下：X：577、1817、2745、11577；Y：496、4562、6912、9543。(1) X、Y 分别是什么元素？(2) X 与 Y 的最高价氧化物对应水化物能否反应？写出离子方程式。
\tcblower
\jie (1) X 的 $I_4\gg I_3$，最外层 3 个电子，为 Al；Y 的 $I_2\gg I_1$，最外层 1 个电子，$I_1$ 很小，为 Na。(2) 能，\ce{Al(OH)3} 是两性氢氧化物：\ce{Al(OH)3 + OH- = [Al(OH)4]-}。
\end{liti}

\subsection{电负性}
\begin{gainian}[电负性]
衡量元素原子在化合物中吸引键合电子能力的标度（F 最大，为 4.0）。同周期从左到右增大，同主族从上到下减小。
\end{gainian}
\textbf{应用：}①判断金属性、非金属性（金属一般 $<1.8$，非金属一般 $>1.8$）；②判断化合物中元素的正负化合价（电负性大的显负价，如 \ce{NaH} 中 H 为 $-1$ 价，\ce{OF2} 中 O 为 $+2$ 价）；③判断化学键类型（电负性差值大于 1.7 一般为离子键，小于 1.7 一般为共价键；差值越大，键的离子性越强）；④对角线规则（Li 与 Mg、Be 与 Al、B 与 Si 性质相似）。

\begin{liti}[电负性与化合价]
已知电负性：H 2.1，C 2.5，N 3.0，F 4.0，Si 1.8。判断下列化合物中加点元素的化合价：\ce{SiH4} 中的 H；\ce{NF3} 中的 N；\ce{CH2Cl2} 中的 C（Cl 为 3.0）。
\tcblower
\jie 电负性大的元素吸引电子能力强，显负价。\ce{SiH4}：H（2.1）$>$Si（1.8），H 为 $-1$ 价（\ce{SiH4} 中 Si 为 $+4$，有强还原性，遇水放出 \ce{H2}）。\ce{NF3}：F 为 $-1$，N 为 $+3$。\ce{CH2Cl2}：C 与 H 成键时 C 吸引电子（H 为 $+1$），与 Cl 成键时 Cl 吸引电子（Cl 为 $-1$），C 为 $0$ 价。
\end{liti}

\begin{yicuo}
第一电离能与电负性的大小顺序\textbf{不完全一致}：电负性 O$>$N，但第一电离能 N$>$O（N 的 $2p^3$ 半满）。比较金属性用电离能，比较非金属性用电负性。
\end{yicuo}

\section{元素推断的方法}

\begin{fangfa}[“位—构—性”三角推断]
\textbf{①由结构式推断}：根据原子形成的共价键数确定族（H、F、Cl 形成 1 个键；O、S 形成 2 个；N、P 形成 3 个，P 还可形成 5 个；C、Si 形成 4 个；B 形成 3 个，带负电时形成 4 个）。\textbf{②由电子排布推断}：如“3p 轨道半充满”为 P，“未成对电子数最多的第四周期元素”为 Cr，“$s$ 电子数与 $p$ 电子数相等”为 O 或 Mg。\textbf{③由电离能图推断}：找 ⅡA、ⅤA 的“凸起”和 0 族的“峰”、ⅠA 的“谷”。\textbf{④由性质推断}：地壳中含量最多的元素为 O，形成化合物种类最多的为 C，单质为淡黄色固体的为 S 等。
\end{fangfa}

\begin{zhenti}{2025年广东卷第9题\quad 来源：仓库所附真题文档}
元素 a～i 为短周期元素，原子序数依次为 $x\sim x+8$，其第一电离能与原子序数的关系为：从 a 到 f 总体增大，但 c 高于 d；f 为最高峰；g 骤降为最低；g、h、i 依次为低—稍高—稍低。下列说法正确的是
\xxb{a 和 g 同主族}{金属性：g$>$h$>$i}{原子半径：e$>$d$>$c}{最简单氢化物沸点：b$>$c}
\textbf{解析：}f 为峰、g 骤降，f 为稀有气体、g 为碱金属；c 比后一个元素 d 高，说明 c 为 ⅤA 族。连续 9 种元素中 f 为 Ne，则 a～i 依次为 B、C、N、O、F、Ne、Na、Mg、Al。验证：h（Mg）高于 i（Al），符合 ⅡA 特例。A：B 与 Na 不同族；B：Na、Mg、Al 金属性依次减弱，正确；C：同周期半径 c（N）$>$d（O）$>$e（F）；D：\ce{NH3} 分子间有氢键，沸点 \ce{NH3}$>$\ce{CH4}（c$>$b）。答案：B。
\end{zhenti}

\begin{zhenti}{2024年广东卷第12题\quad 来源：仓库所附真题文档}
一种可为运动员补充能量的物质，其分子结构式中，R、W、Z、X、Y 为原子序数依次增大的短周期主族元素，Z 和 Y 同族；结构式中 R 均形成 1 个键，W 形成 4 个键，Z 形成 3 个键，X 形成 2 个键，Y 形成 5 个键（含 1 个双键 X=Y）。则
\xxb{沸点：\ce{ZR3}$<$\ce{YR3}}{最高价氧化物的水化物的酸性：Z$<$W}{第一电离能：Z$<$X$<$W}{\ce{ZX3-} 和 \ce{WX3^2-} 空间结构均为平面三角形}
\textbf{解析：}Y 形成 5 个键且与 Z 同族，Z 形成 3 个键，Y 为 P、Z 为 N；W 形成 4 个键且序数小于 N，为 C；R 形成 1 个键、序数最小，为 H；X 形成 2 个键，序数在 N 与 P 之间，为 O。A：\ce{NH3} 分子间有氢键，沸点高于 \ce{PH3}，错；B：酸性 \ce{HNO3}$>$\ce{H2CO3}，错；C：第一电离能 C$<$O$<$N，错；D：\ce{NO3-}、\ce{CO3^2-} 中心原子价层电子对数均为 3，无孤电子对，平面三角形，正确（判断方法见第7章）。答案：D。
\end{zhenti}

% ---------------- 章末 ----------------
\zhangmo{本章小结}
\subsection*{一、知识总结}
\begin{center}
\small
\begin{tabularx}{\linewidth}{p{2.4cm}X}
\toprule
排布规则 & 构造原理（4s 先于 3d）、能量最低、泡利、洪特（含半满全满特例 Cr、Cu） \\
表示方法 & 电子排布式、简化式、价层排布式、轨道表示式；离子先失最外层电子 \\
周期表 & 周期$=$能层数；主族序数$=$价电子数；s、p、d、ds、f 五区 \\
周期律 & 半径：同周期减小、同主族增大；$I_1$：总体增大，ⅡA、ⅤA 特例；电负性：同周期增大、同主族减小 \\
\bottomrule
\end{tabularx}
\end{center}
\subsection*{二、方法提升}
\begin{enumerate}
\item 写排布式先按构造原理填，再按能层顺序整理；遇到 Cr、Cu 想到半满、全满。
\item 逐级电离能找“突跃”定最外层电子数；第一电离能比较注意 ⅡA、ⅤA。
\item 元素推断先抓最有特征的信息（成键数、半满、电离能峰谷），再用其他条件验证。
\end{enumerate}
\subsection*{三、易错点归纳}
\begin{enumerate}
\item \ce{Fe^2+} 写成 [Ar]$3d^44s^2$（错，应先失 4s）。
\item 认为第一电离能 O$>$N、Mg$<$Al（均错）。
\item 认为电负性大的元素第一电离能一定大。
\item 价层电子排布式与最外层电子排布式混淆（过渡元素价层包括 $(n-1)d$）。
\end{enumerate}

\zhangmo{分层练习}
\zu{A组\quad 基础巩固}
\begin{xiti}[start=1]
\item 下列基态原子或离子的电子排布式（或轨道表示式）正确的是\xz
\xxa{Cr：[Ar]$3d^44s^2$}{\ce{Fe^2+}：[Ar]$3d^44s^2$}{Cu：[Ar]$3d^{10}4s^1$}{N 的 2p 轨道：\fbox{$\uparrow\downarrow$}\fbox{$\uparrow$}\fbox{\ }}

\item 下列各组元素第一电离能大小顺序正确的是\xz
\xx{Mg$<$Al$<$Si}{N$<$O$<$F}{Na$<$Mg$<$Al}{Al$<$Mg$<$P}

\item 下列说法正确的是\xz
\xxa{同一原子中，2p、3p、4p 能级的轨道数依次增多}{电子由 3s 能级跃迁到 3p 能级时释放能量}{s 区元素全部是金属元素}{焰色试验是电子跃迁时放出光能，属于物理变化}

\item 某元素的逐级电离能（kJ·mol$^{-1}$）依次为 738、1451、7733、10540、13630……，该元素最可能是\tk，其基态原子的价层电子排布式为\tk。

\item 下列说法正确的是\xz
\xxa{3d 能级最多容纳 6 个电子}{基态原子中 4s 能级的能量高于 3d 能级，所以先填 3d}{p 能级的 3 个原子轨道能量相同、相互垂直}{同一原子中，1s、2s、3s 电子的能量相同}

\item 写出下列粒子的价层电子排布式：\ce{_{24}Cr}\tk；\ce{_{29}Cu+}\tk；\ce{_{26}Fe^3+}\tk；\ce{_{35}Br}\tk。
\end{xiti}

\zu{B组\quad 高考典型}
\begin{xiti}
\item \laiyuan{2026年广东卷第11题，来源：仓库所附真题文档（网络收集版）}\ 化合物 \ce{L2MQR4} 与 \ce{TXR4} 可用作电解质饮料的配料。短周期元素 M、R、L、T、Q、X 原子序数依次增加，R、X 同族，Q 与 X 相邻，Q 基态原子的 3p 轨道为半充满，则下列说法正确的是\xz
\xxb{原子半径：R$>$X$>$M}{金属性：T$>$L$>$M}{第一电离能：Q$>$T$>$L}{最高价氧化物对应水化物的酸性：Q$>$X}

\item \laiyuan{原创}\ 前四周期元素 A、B、C、D、E 原子序数依次增大。A 的基态原子 2p 能级有 2 个未成对电子且 s 电子数与 p 电子数相等；B 是短周期中原子半径最大的主族元素；C 的单质是淡黄色固体；D 与 C 同周期，D 的电负性在同周期主族元素中最大；E 的基态原子 3d 能级有 6 个电子。(1) 写出 A～E 的元素符号。(2) 写出 E 的基态原子电子排布式和 \ce{E^3+} 的价层电子排布式，说明 \ce{E^3+} 较稳定的原因。(3) 比较 A、C、D 的第一电离能和电负性。(4) B 与 A 形成的两种化合物中，含非极性键的是\tk（写化学式）。

\item \laiyuan{原创}\ 短周期主族元素 X、Y、Z、W 原子序数依次增大。X 的基态原子 2p 能级上有 3 个未成对电子；Y 是地壳中含量最多的元素；Z 的第一电离能在同周期主族元素中最小；W 与 Y 同主族。下列说法正确的是\xz
\xxa{第一电离能：Y$>$X$>$W}{电负性：Y$>$W$>$X}{简单离子半径：W$>$Y$>$Z}{Z 与 Y 形成的化合物中只含离子键}

\item \laiyuan{原创}\ 解释下列事实：(1) 第一电离能 Mg$>$Al；(2) 第一电离能 P$>$S；(3) 离子半径 $r(\ce{S^2-})>r(\ce{Cl-})>r(\ce{K+})>r(\ce{Ca^2+})$；(4) Na 的第二电离能远大于其第一电离能。
\end{xiti}

\zu{C组\quad 冲刺高分}
\begin{xiti}
\item \laiyuan{原创}\ (1) 基态 \ce{Cr} 原子的价层轨道表示式为\tk，其未成对电子数为\tk。(2) 比较 Mn 和 Fe 的第三电离能 $I_3$ 的大小并解释原因。(3) 第二周期元素中，第一电离能介于 B 和 N 之间的元素有\tk{}种。(4) 已知 Cu 的 $I_2=1958$ kJ·mol$^{-1}$，Zn 的 $I_2=1733$ kJ·mol$^{-1}$，解释为何 $I_2(\ce{Cu})>I_2(\ce{Zn})$。

\item \laiyuan{原创}\ 锂离子电池正极材料常含 Fe、Co、Ni、Mn 等过渡元素。(1) 写出基态 \ce{Co^3+} 和 \ce{Ni^2+} 的价层电子排布式，并指出各自的未成对电子数。(2) \ce{Mn^2+} 很难被氧化为 \ce{Mn^3+}，而 \ce{Fe^2+} 较易被氧化为 \ce{Fe^3+}，从电子排布角度解释。(3) Co 位于周期表第\tk{}周期第\tk{}族，属于\tk{}区。
\end{xiti}

\zhangmo{练习解析}
\begin{jiexilist}
\item \ans{C}A：Cr 为 [Ar]$3d^54s^1$（半满稳定）；B：\ce{Fe^2+} 为 [Ar]$3d^6$（先失 4s）；C：Cu 为 $3d^{10}4s^1$（全满稳定），正确；D：违反洪特规则，2p 的 3 个电子应分占 3 个轨道且自旋平行。
\item \ans{D}A：Mg（$3s^2$ 全满）$>$Al，应为 Al$<$Mg$<$Si；B：N（$2p^3$ 半满）$>$O，应为 O$<$N$<$F；C：应为 Na$<$Al$<$Mg；D：Al$<$Mg$<$P，正确。
\item \ans{D}A：p 能级都有 3 个轨道；B：3s$\to$3p 是向高能级跃迁，吸收能量；C：H 在 s 区，是非金属；D 正确。
\item Mg；$3s^2$。$I_3\gg I_2$，最外层有 2 个电子，为 ⅡA 族；$I_1=738$ 较小，为第三周期的 Mg（Be 的 $I_1$ 约 900）。
\item \ans{C}A：d 能级有 5 个轨道，最多 10 个电子；B：构造原理中 4s 能量低于 3d，先填 4s；C 正确；D：能层越高能量越高，$E(1s)<E(2s)<E(3s)$。
\item $3d^54s^1$；$3d^{10}$；$3d^5$；$4s^24p^5$。Cr 半满稳定；\ce{Cu+} 由 $3d^{10}4s^1$ 失去 4s 电子；\ce{Fe^3+} 先失 2 个 4s 电子再失 1 个 3d 电子；主族元素的价层电子即最外层电子。
\item \ans{C}Q 的 3p 半充满，为 P；X 与 Q 相邻且在其后，为 S；R 与 X 同族且序数最小，为 O；化合物 \ce{L2MQR4}（如 \ce{K2HPO4} 类）和 \ce{TXR4}（如 \ce{MgSO4}）中，M 序数最小可形成 \ce{L2MPO4}，M 为 H（\ce{Na2HPO4}），L 为 Na，T 为 Mg。验证：原子序数 H$<$O$<$Na$<$Mg$<$P$<$S。A：半径 S$>$O$>$H，即 X$>$R$>$M，错；B：金属性 Na$>$Mg，即 L$>$T，错；C：第一电离能 P$>$Mg$>$Na，即 Q$>$T$>$L，正确；D：酸性 \ce{H2SO4}$>$\ce{H3PO4}，即 X$>$Q，错。答案：C。
\item (1) A 为 O（$1s^22s^22p^4$，s、p 电子均为 4，2p 有 2 个未成对电子）；B 为 Na；C 为 S；D 为 Cl；E 为 Fe（$3d^64s^2$）。(2) Fe：$1s^22s^22p^63s^23p^63d^64s^2$；\ce{Fe^3+}：$3d^5$，3d 能级半充满，能量较低，较稳定。(3) 第一电离能：O$>$Cl$>$S（同周期 S$<$Cl，同主族 O$>$S；O 与 Cl 比较：O 为 1314、Cl 为 1251 kJ·mol$^{-1}$，O$>$Cl）——高考一般只考查同周期或同主族的比较；电负性：O$>$Cl$>$S。(4) \ce{Na2O2}（含 O—O 非极性键）。
\item \ans{C}X 为 N（$2p^3$），Y 为 O，Z 为 Na（同周期 $I_1$ 最小），W 为 S。A：第一电离能 N$>$O$>$S，即 X$>$Y$>$W，错；B：电负性 O$>$N$>$S，即 Y$>$X$>$W，错；C：\ce{S^2-} 有 3 个电子层，半径最大；\ce{O^2-} 与 \ce{Na+} 电子层结构相同，核电荷数小的半径大，\ce{O^2-}$>$\ce{Na+}，正确；D：\ce{Na2O2} 中含 O—O 共价键，错。
\item (1) Mg 的价层电子为 $3s^2$（全满，较稳定），Al 为 $3s^23p^1$，3p 电子能量较高且易失去，所以 Al 的 $I_1$ 较小。(2) P 的 3p 为 $3p^3$ 半满，较稳定；S 为 $3p^4$，失去一个电子后变为半满，较容易。(3) 四种离子电子层结构相同（均为 Ar 型），核电荷数依次增大，核对电子的吸引依次增强，半径依次减小。(4) Na 失去 1 个电子后变为 \ce{Na+}（$2s^22p^6$，稀有气体稳定结构），再失去的电子位于内层，离核近、受吸引强，所以 $I_2\gg I_1$。
\item (1) 3d：\fbox{$\uparrow$}\fbox{$\uparrow$}\fbox{$\uparrow$}\fbox{$\uparrow$}\fbox{$\uparrow$}\ 4s：\fbox{$\uparrow$}；6。(2) $I_3(\ce{Mn})>I_3(\ce{Fe})$。\ce{Mn^2+} 为 $3d^5$ 半满，较稳定，再失去一个电子较难；\ce{Fe^2+} 为 $3d^6$，失去一个电子后变为较稳定的 $3d^5$，较容易。(3) 3 种：Be、C、O（第一电离能 B$<$Be$<$C$<$O$<$N）。(4) Cu 失去 1 个电子后 \ce{Cu+} 为 $3d^{10}$ 全满稳定结构，再失电子较难；Zn 失去 1 个电子后 \ce{Zn+} 为 $3d^{10}4s^1$，再失去 4s 电子较容易。
\item (1) \ce{Co^3+}：$3d^6$，4 个未成对电子；\ce{Ni^2+}：$3d^8$，2 个未成对电子。(2) \ce{Mn^2+} 为 $3d^5$ 半满稳定结构，失电子后破坏半满，较难；\ce{Fe^2+} 为 $3d^6$，失去 1 个电子后变为 $3d^5$ 半满稳定结构，较容易。(3) 四；Ⅷ；d。
\end{jiexilist}
% <<<<<<<< 结束：chapters/ch06.tex

% >>>>>>>> 开始：chapters/ch07.tex
% ============================================================
% 第7章 分子结构与性质（选必2第二章、第三章第四节）
% ============================================================
\chapter{分子结构与性质}

\begin{daoxue}
\textbf{知识主线：}原子怎样结合成分子（共价键：$\sigma$ 键、$\pi$ 键、配位键）$\to$ 分子是什么形状（价层电子对互斥模型、杂化轨道）$\to$ 形状和键决定分子的性质（极性、分子间作用力、氢键、溶解性、手性、酸性）$\to$ 分子与离子的组装（配合物、超分子）。

\textbf{高考考情：}广东卷选择题年年有“结构决定性质”类题（2024年第13题、2025年第13题、2026年第8题）；工艺流程和有机大题中常考杂化方式、空间结构、配位作用（2025年第18题(4)、2026年第18题(6)、第20题(2)）。

\textbf{学习目标：}①会判断 $\sigma$、$\pi$ 键并计数；②会用价层电子对互斥模型判断空间结构，用杂化理论解释；③会比较键角、沸点、溶解性、酸性并说明理由；④理解配合物的组成和超分子的识别作用。
\end{daoxue}

\section{共价键}

\subsection{$\sigma$ 键与 $\pi$ 键}
共价键是原子间通过共用电子对（电子云重叠）形成的化学键，具有\textbf{饱和性}和\textbf{方向性}。按电子云重叠方式分：

\begin{center}
\small
\begin{tabularx}{\linewidth}{lXX}
\toprule
& $\sigma$ 键 & $\pi$ 键 \\
\midrule
重叠方式 & “头碰头”，沿键轴方向重叠，重叠程度大 & “肩并肩”，侧面重叠，重叠程度小 \\
对称性 & 轴对称，可绕键轴旋转 & 镜面对称，不能旋转 \\
牢固程度 & 较牢固 & 一般较弱，易断裂（烯烃易加成） \\
存在 & 共价单键都是 $\sigma$ 键 & 双键含 1 个 $\sigma$、1 个 $\pi$；三键含 1 个 $\sigma$、2 个 $\pi$ \\
\bottomrule
\end{tabularx}
\end{center}

计数举例：\ce{HC#N} 含 2 个 $\sigma$ 键、2 个 $\pi$ 键；\ce{CH2=CH-C#N}（丙烯腈）含 6 个 $\sigma$ 键、3 个 $\pi$ 键；1 mol \ce{[Cu(NH3)4]^2+} 含 $\sigma$ 键 $4\times3+4=16$ mol（配位键也是 $\sigma$ 键）。

\subsection{键参数与配位键}
\textbf{键能}越大、\textbf{键长}越短，共价键越牢固，分子越稳定：如 \ce{N#N} 键能很大，\ce{N2} 化学性质稳定；HF、HCl、HBr、HI 的键长依次增大、键能依次减小，热稳定性依次减弱。\textbf{键角}决定分子的空间结构。

\textbf{配位键}：一方提供孤电子对、另一方提供空轨道形成的共价键，用“A$\to$B”表示（A 提供孤电子对）。如 \ce{NH4+} 中 \ce{H+} 提供空轨道、N 提供孤电子对；形成后 4 个 N—H 键完全相同。

\section{分子的空间结构}

\subsection{价层电子对互斥模型（VSEPR）}
\begin{gainian}[中心原子价层电子对数]
价层电子对数 $=$ $\sigma$ 键电子对数 $+$ 孤电子对数。\\
$\sigma$ 键电子对数 $=$ 与中心原子结合的原子数；\quad 孤电子对数 $=\dfrac12(a-xb)$。\\
$a$：中心原子价电子数（阳离子减去电荷数，阴离子加上电荷数）；$x$：与中心原子结合的原子数；$b$：与中心原子结合的原子最多能接受的电子数（H 为 1，其他原子为“8$-$该原子价电子数”，如 O 为 2、Cl 为 1）。
\end{gainian}

价层电子对相互排斥，尽可能远离，形成 VSEPR 模型；略去孤电子对，就是分子（离子）的实际空间结构：

\begin{center}
\small
\begin{tabular}{cccl}
\toprule
价层电子对数 & VSEPR 模型 & 孤电子对数 & 空间结构及实例 \\
\midrule
2 & 直线形 & 0 & 直线形：\ce{CO2}、\ce{BeCl2}、HCN \\
3 & 平面三角形 & 0 & 平面三角形：\ce{BF3}、\ce{SO3}、\ce{NO3-}、\ce{CO3^2-} \\
3 & 平面三角形 & 1 & V 形：\ce{SO2}（键角约 120°，略小） \\
4 & 四面体形 & 0 & 正四面体形：\ce{CH4}、\ce{NH4+}、\ce{SO4^2-}、\ce{CCl4} \\
4 & 四面体形 & 1 & 三角锥形：\ce{NH3}、\ce{H3O+}、\ce{PCl3}、\ce{SO3^2-}、\ce{ClO3-} \\
4 & 四面体形 & 2 & V 形：\ce{H2O}、\ce{H2S}、\ce{OF2} \\
\bottomrule
\end{tabular}
\end{center}

计算示例：\ce{SO2}：孤电子对数 $=\frac12(6-2\times2)=1$，价层电子对数 $2+1=3$，V 形；\ce{SO3^2-}：$\frac12(6+2-3\times2)=1$，$3+1=4$，三角锥形；\ce{NH4+}：$\frac12(5-1-4\times1)=0$，正四面体形。

\begin{liti}[价层电子对互斥模型的应用]
用价层电子对互斥模型判断下列粒子的空间结构：\ce{PCl3}、\ce{ClO3-}、\ce{NO2-}、\ce{BF4-}、\ce{SO2Cl2}（S 为中心原子）、\ce{H3O+}。
\tcblower
\jie \begin{tabular}[t]{lcccl}
粒子 & $\sigma$ 键对 & 孤电子对 $\frac12(a-xb)$ & 价层电子对 & 空间结构 \\
\ce{PCl3} & 3 & $\frac12(5-3\times1)=1$ & 4 & 三角锥形 \\
\ce{ClO3-} & 3 & $\frac12(7+1-3\times2)=1$ & 4 & 三角锥形 \\
\ce{NO2-} & 2 & $\frac12(5+1-2\times2)=1$ & 3 & V 形 \\
\ce{BF4-} & 4 & $\frac12(3+1-4\times1)=0$ & 4 & 正四面体形 \\
\ce{SO2Cl2} & 4 & $\frac12(6-2\times2-2\times1)=0$ & 4 & 四面体形（非正四面体） \\
\ce{H3O+} & 3 & $\frac12(6-1-3\times1)=1$ & 4 & 三角锥形
\end{tabular}\\
\xianjing \ce{SO2Cl2} 中两种配位原子不同，键长、键角不全相等，只能说“四面体形”，不能说“正四面体形”。
\end{liti}

\subsection{杂化轨道理论}
价层电子对互斥模型能\textbf{预测}形状，杂化轨道理论\textbf{解释}形状：中心原子能量相近的原子轨道“混合”成数目相同、能量相同的杂化轨道，用于形成 $\sigma$ 键或容纳孤电子对。\textbf{杂化轨道数 $=$ 价层电子对数}：

\begin{center}
\small
\begin{tabular}{lllll}
\toprule
杂化类型 & 杂化轨道数 & 夹角 & 实例 & 有机物中的碳原子 \\
\midrule
sp & 2 & 180° & \ce{CO2}、\ce{C2H2} & 三键碳、\ce{CO2} 中的碳 \\
$\mathrm{sp}^2$ & 3 & 120° & \ce{BF3}、\ce{C2H4}、苯 & 双键碳、苯环碳、羰基碳 \\
$\mathrm{sp}^3$ & 4 & 109°28$'$ & \ce{CH4}、\ce{NH3}、\ce{H2O} & 饱和碳 \\
\bottomrule
\end{tabular}
\end{center}

未参与杂化的 p 轨道可“肩并肩”形成 $\pi$ 键。苯分子中 6 个碳原子均为 $\mathrm{sp}^2$ 杂化，各有 1 个未杂化的 p 轨道垂直于环平面，相互重叠形成\textbf{大 $\pi$ 键}（6 中心 6 电子，记作 $\Pi_6^6$），所以苯环中没有单双键交替，6 个碳碳键完全相同（2025年广东卷第2题即考查“具有芳香性的配合物中存在共轭大 $\pi$ 键、不存在碳碳双键”）。

\begin{tuozhan}[大 $\pi$ 键的表示]
大 $\pi$ 键用 $\Pi_n^m$ 表示（$n$ 为参与的原子数，$m$ 为电子数）。常见：苯 $\Pi_6^6$；\ce{CO3^2-}、\ce{NO3-} $\Pi_4^6$；\ce{SO2}、\ce{O3} $\Pi_3^4$。形成大 $\pi$ 键的原子通常共平面且有相互平行的 p 轨道。高考对此只作了解，常结合“共平面”“键长平均化”考查。
\end{tuozhan}

\subsection{键角的比较}
\begin{fangfa}[键角比较三步]
①先看\textbf{杂化方式}：sp（180°）$>\mathrm{sp}^2$（120°）$>\mathrm{sp}^3$（109°28$'$）。②杂化相同时看\textbf{孤电子对数}：孤电子对对成键电子对的排斥大于成键电子对之间的排斥，孤电子对越多，键角越小：\ce{CH4}（109°28$'$）$>$\ce{NH3}（107°）$>$\ce{H2O}（105°）。③孤电子对数也相同时看\textbf{电负性}：中心原子电负性越大，成键电子对越靠近中心原子，排斥越大，键角越大：\ce{NH3}$>$\ce{PH3}，\ce{H2O}$>$\ce{H2S}；配位原子电负性越大，成键电子对越远离中心原子，键角越小：\ce{NF3}$<$\ce{NH3}。
\end{fangfa}

\section{分子的性质}

\subsection{分子的极性}
正电中心与负电中心不重合的分子是极性分子，重合的是非极性分子。双原子分子：键有极性则分子有极性。\textbf{\ce{AB_n} 型分子}：中心原子无孤电子对且所连原子相同（空间结构对称）为非极性分子，如 \ce{CO2}、\ce{BF3}、\ce{CH4}、\ce{CCl4}、\ce{SO3}；否则多为极性分子，如 \ce{H2O}、\ce{NH3}、\ce{SO2}、\ce{CH3Cl}。注意：\ce{O3} 是极性分子，\ce{H2O2} 是极性分子。

\subsection{范德华力与氢键}
\textbf{范德华力}：分子间普遍存在的作用力，比化学键弱得多，影响熔沸点、溶解性。组成和结构相似的分子，相对分子质量越大，范德华力越大；相对分子质量相近时，极性越大，范德华力越大。

\begin{gainian}[氢键]
已与电负性很大的原子（N、O、F）形成共价键的氢原子，与另一个电负性很大、有孤电子对的原子之间的作用力，表示为 X—H$\cdots$Y。氢键比化学键弱、比范德华力强，有方向性和饱和性。
\end{gainian}

\textbf{氢键对性质的影响：}①\textbf{分子间}氢键使熔沸点\textbf{升高}：\ce{H2O}、HF、\ce{NH3} 的沸点在同族氢化物中反常地高；乙醇沸点高于二甲醚。②\textbf{分子内}氢键使熔沸点\textbf{降低}：邻羟基苯甲醛（形成分子内氢键）的沸点低于对羟基苯甲醛（形成分子间氢键）。③与水分子形成氢键使溶解度增大：\ce{NH3} 极易溶于水，乙醇与水互溶。④冰中每个水分子与周围 4 个水分子形成氢键，排列疏松，密度比水小。

\begin{zhenti}{2026年广东卷第8题\quad 来源：仓库所附真题文档（网络收集版）}
结构决定性质。下列事实对应的解释\textbf{错误}的是：A．热稳定性 HCl$>$HI——键能 H—Cl$>$H—I；B．酸性：乙酸$>$丙酸——推电子效应 \ce{-CH3}$<$\ce{-CH2CH3}；C．沸点：邻羟基苯甲醛$>$水——邻羟基苯甲醛可形成分子内氢键；D．在 \ce{CCl4} 中的溶解度：苯$>$\ce{NH3}——苯和 \ce{CCl4} 为非极性分子，\ce{NH3} 为极性分子。

\textbf{解析：}A、B、D 的解释均正确（乙基推电子能力更强，使羧基中 O—H 键极性减弱，丙酸酸性更弱）。C：分子内氢键会使沸点\textbf{降低}，不能用来解释沸点高；邻羟基苯甲醛沸点（约 197 ℃）高于水，是因为其相对分子质量大得多、范德华力大。答案：C。
\end{zhenti}

\begin{liti}[用分子间作用力解释物理性质]
解释下列事实：(1) 沸点：\ce{H2O}（100 ℃）$>$\ce{H2S}（$-60$ ℃）；(2) 沸点：\ce{CH4}$<$\ce{SiH4}$<$\ce{GeH4}；(3) 乙醇与水互溶，而乙酸乙酯在水中溶解度很小；(4) 对羟基苯甲酸的熔点（214 ℃）高于邻羟基苯甲酸（158 ℃）。
\tcblower
\jie (1) 水分子间存在氢键，\ce{H2S} 分子间只有范德华力。(2) 三者组成和结构相似、分子间都无氢键，相对分子质量依次增大，范德华力依次增强。(3) 乙醇分子中的羟基能与水分子形成氢键，且烃基较小；乙酸乙酯不能与水形成足够的氢键，烃基部分疏水。(4) 邻羟基苯甲酸形成分子内氢键，分子间作用减弱；对羟基苯甲酸只能形成分子间氢键，熔点较高。\\
\dianbo 答题时要点明“哪种作用力”“谁比谁强”，不能只写“有氢键”。
\end{liti}

\subsection{溶解性、手性与酸性}
\textbf{溶解性}：“相似相溶”——非极性溶质易溶于非极性溶剂（\ce{I2}、苯易溶于 \ce{CCl4}），极性溶质易溶于极性溶剂；溶质与水形成氢键或与水反应，溶解度增大。

\textbf{手性}：连有 4 个\textbf{不同}原子或原子团的碳原子叫\textbf{手性碳原子}；含手性碳原子的分子存在互为镜像、不能重叠的一对\textbf{对映异构体}（手性异构体）。如乳酸 \ce{CH3-}$\overset{*}{\mathrm{C}}$\ce{H(OH)-COOH} 中带“$*$”的碳。

\textbf{酸性比较}：
\begin{enumerate}
\item \textbf{无机含氧酸}：写成 \ce{(HO)_mRO_n}，\textbf{非羟基氧原子数 $n$ 越多}，R 的正电性越高，O—H 键越易断裂，酸性越强：\ce{HClO}（$n=0$）$<$\ce{HClO2}$<$\ce{HClO3}$<$\ce{HClO4}（$n=3$）；\ce{H2SO3}$<$\ce{H2SO4}。
\item \textbf{羧酸}：烃基上连\textbf{吸电子基团}（F、Cl 等，电负性越大吸电子越强、个数越多越强）使羧基中 O—H 键极性增大，酸性增强：\ce{CF3COOH}$>$\ce{CCl3COOH}$>$\ce{CH2ClCOOH}$>$\ce{CH3COOH}；烃基是\textbf{推电子基团}，烃基越长推电子越强，酸性越弱：\ce{HCOOH}$>$\ce{CH3COOH}$>$\ce{CH3CH2COOH}。
\end{enumerate}

\begin{zhenti}{2025年广东卷第13题\quad 来源：仓库所附真题文档}
由结构\textbf{不能}推测出对应性质的是：A．\ce{SO3} 的 VSEPR 模型为平面三角形——\ce{SO3} 具有氧化性；B．钾和钠的原子结构不同，电子跃迁时能量变化不同——钾和钠的焰色不同；C．乙烯和乙炔分子均含有 $\pi$ 键——两者均可发生加成反应；D．石墨层中未参与杂化的 p 轨道中的电子可在整个碳原子平面中运动——石墨具有类似金属的导电性。

\textbf{解析：}\ce{SO3} 的氧化性源于 S 处于 +6 价（最高价），与其空间结构无关；B、C、D 均为结构与性质的合理对应。答案：A。
\end{zhenti}

\section{配合物与超分子}

\subsection{配合物}
\begin{gainian}[配合物的组成]
由金属离子或原子（\textbf{中心离子}，提供空轨道）与分子或离子（\textbf{配体}，其中提供孤电子对的原子叫\textbf{配位原子}）以配位键结合形成的化合物。以 \ce{[Cu(NH3)4]SO4} 为例：内界 \ce{[Cu(NH3)4]^2+}，外界 \ce{SO4^2-}；中心离子 \ce{Cu^2+}，配体 \ce{NH3}，配位原子 N，配位数 4。内界与外界之间是离子键，在水中完全电离；内界一般难电离。
\end{gainian}

\textbf{常见配合物实验：}
\begin{enumerate}
\item 向 \ce{CuSO4} 溶液中滴加氨水，先生成蓝色沉淀 \ce{Cu(OH)2}，氨水过量后沉淀溶解得深蓝色溶液：\ce{Cu(OH)2 + 4NH3 = [Cu(NH3)4]^2+ + 2OH-}；再加入乙醇（降低溶剂极性，减小配合物溶解度），析出深蓝色晶体 \ce{[Cu(NH3)4]SO4*H2O}。
\item \ce{Fe^3+} 与 \ce{SCN-} 形成红色配合物，用于检验 \ce{Fe^3+}。
\item \ce{AgNO3} 溶液中滴加氨水至沉淀恰好溶解得银氨溶液 \ce{[Ag(NH3)2]OH}；AgCl 溶于氨水。
\item \ce{CuSO4} 溶液呈蓝色是因为形成了 \ce{[Cu(H2O)4]^2+}。
\end{enumerate}

\begin{yicuo}
判断配合物内外界：外界离子在水中能电离，可被检验。如 1 mol \ce{[Co(NH3)5Cl]Cl2} 溶于水，只能电离出 2 mol \ce{Cl-}，加入足量 \ce{AgNO3} 只生成 2 mol AgCl，内界的 1 个 Cl 不沉淀。
\end{yicuo}

\begin{liti}[配合物的结构分析]
铁氰化钾 \ce{K3[Fe(CN)6]} 可用于检验 \ce{Fe^2+}。(1) 指出中心离子、配体、配位原子、配位数和内外界。(2) 1 mol \ce{[Fe(CN)6]^3-} 中含 $\sigma$ 键和 $\pi$ 键各多少？(3) CN$^-$ 中 C 原子作配位原子而不是 N 原子，可能的原因是什么？
\tcblower
\jie (1) 中心离子 \ce{Fe^3+}，配体 \ce{CN-}，配位原子 C，配位数 6；内界 \ce{[Fe(CN)6]^3-}，外界 \ce{K+}。(2) 6 个配位键（$\sigma$）$+$ 6 个 C$\equiv$N（各含 1 个 $\sigma$、2 个 $\pi$）：$\sigma$ 键 12 mol，$\pi$ 键 12 mol。(3) C 的电负性比 N 小，对孤电子对的吸引较弱，更容易给出孤电子对形成配位键。
\end{liti}

\subsection{超分子}
由两种或两种以上分子（或离子）通过\textbf{分子间相互作用}（非共价键，如氢键、静电作用、配位作用）形成的分子聚集体。两个重要特征：\textbf{分子识别}和\textbf{自组装}。例如冠醚的空腔大小不同，能识别大小匹配的碱金属离子：18-冠-6 的空腔与 \ce{K+} 匹配，12-冠-4 与 \ce{Li+} 匹配（2024年广东卷第13题B即以“某冠醚与 \ce{Li+} 能形成超分子，与 \ce{K+} 则不能”考查离子半径与空腔匹配）。2026年广东卷第18题(6)中，杂冠醚选择性识别 \ce{Ag+} 的原因是：与 \ce{Ag+} 有配位作用（S、O、N 原子提供孤电子对），且\textbf{空腔大小与 \ce{Ag+} 的直径相匹配}。

% ---------------- 章末 ----------------
\zhangmo{本章小结}
\subsection*{一、知识总结}
\begin{center}
\small
\begin{tabularx}{\linewidth}{p{2.4cm}X}
\toprule
共价键 & $\sigma$ 键（头碰头）、$\pi$ 键（肩并肩）；单键 1$\sigma$，双键 1$\sigma$1$\pi$，三键 1$\sigma$2$\pi$；配位键属 $\sigma$ 键 \\
空间结构 & 价层电子对数$=\sigma$ 键对数$+\frac12(a-xb)$；杂化轨道数$=$价层电子对数 \\
键角 & 杂化方式$\to$孤电子对数$\to$电负性 \\
分子间作用 & 范德华力（相对分子质量、极性）；氢键（N、O、F；分子间升高、分子内降低沸点） \\
性质 & 极性、相似相溶、手性碳、含氧酸非羟基氧数、羧酸取代基电子效应 \\
配合物与超分子 & 中心离子提供空轨道、配体提供孤电子对；分子识别、自组装 \\
\bottomrule
\end{tabularx}
\end{center}
\subsection*{二、方法提升}
\begin{enumerate}
\item 判断空间结构：算价层电子对数$\to$得 VSEPR 模型$\to$去掉孤电子对得实际形状$\to$得杂化方式。
\item 解释沸点差异的答题模板：“××分子间存在氢键（或××形成分子内氢键）/××相对分子质量更大，范德华力更强，所以……”。
\item 解释酸性差异：抓住“O—H 键极性”这一中间环节。
\end{enumerate}
\subsection*{三、易错点归纳}
\begin{enumerate}
\item 把 VSEPR 模型当作实际空间结构（\ce{H2O} 的 VSEPR 模型是四面体形，实际为 V 形）。
\item 计算孤电子对数时，离子忘记加减电荷数。
\item 认为氢键是化学键；认为分子内氢键使沸点升高。
\item 认为含极性键的分子一定是极性分子（\ce{CO2}、\ce{CH4} 是非极性分子）。
\item 统计 $\sigma$ 键时漏掉配位键和 C—H 键。
\end{enumerate}

\zhangmo{分层练习}
\zu{A组\quad 基础巩固}
\begin{xiti}[start=1]
\item 下列分子或离子中，中心原子的杂化方式与空间结构均正确的是\xz
\xxb{\ce{SO2}：sp，直线形}{\ce{NH3}：$\mathrm{sp}^2$，平面三角形}{\ce{CO3^2-}：$\mathrm{sp}^2$，平面三角形}{\ce{H3O+}：$\mathrm{sp}^3$，正四面体形}

\item 1 mol 丙烯腈（\ce{CH2=CH-C#N}）中含 $\sigma$ 键\tk mol、$\pi$ 键\tk mol；其中碳原子的杂化方式有\tk。

\item 下列说法正确的是\xz
\xxa{\ce{H2O} 比 \ce{H2S} 稳定，是因为水分子间存在氢键}{\ce{CH4}、\ce{CCl4}、\ce{CO2} 都是含极性键的非极性分子}{键角：\ce{H2O}$>$\ce{NH3}$>$\ce{CH4}}{邻羟基苯甲醛的沸点高于对羟基苯甲醛}

\item 下列各组酸中，酸性强弱比较正确的是\xz
\xxa{\ce{HClO4}$<$\ce{HClO3}}{\ce{CH3COOH}$>$\ce{CH2ClCOOH}}{\ce{CF3COOH}$>$\ce{CCl3COOH}}{\ce{CH3CH2COOH}$>$\ce{CH3COOH}}

\item 下列分子中，属于含极性键的非极性分子的是\xz
\xx{\ce{H2O}}{\ce{NH3}}{\ce{BF3}}{\ce{CH3Cl}}

\item 指出下列物质中加点原子的杂化方式：\ce{NH4+} 中的 N\tk；\ce{CO2} 中的 C\tk；甲醛（\ce{HCHO}）中的 C\tk；乙醇中与羟基相连的 C\tk。
\end{xiti}

\zu{B组\quad 高考典型}
\begin{xiti}
\item \laiyuan{原创}\ 比较下列键角大小并说明理由：(1) \ce{NH3} 与 \ce{PH3}；(2) \ce{NH3} 与 \ce{NF3}；(3) \ce{SO4^2-} 与 \ce{SO3^2-}；(4) \ce{BF3} 与 \ce{NF3}。

\item \laiyuan{原创}\ 向 \ce{CuSO4} 溶液中逐滴加入氨水至过量，再加入适量乙醇，析出深蓝色晶体。(1) 描述加氨水过程中的现象并写出相关离子方程式。(2) 加入乙醇的作用是什么？(3) \ce{[Cu(NH3)4]^2+} 中配位原子是\tk，1 mol 该离子中含 $\sigma$ 键\tk mol。(4) \ce{NH3} 中 H—N—H 键角小于 \ce{[Cu(NH3)4]^2+} 中 H—N—H 键角，原因是\tk。

\item \laiyuan{原创}\ 有两种组成均为 \ce{CoCl3*5NH3*H2O} 的配合物 X、Y（配位数均为 6）。各取 1 mol 溶于水，加入足量 \ce{AgNO3} 溶液，X 生成 3 mol AgCl，Y 生成 2 mol AgCl。写出 X、Y 的化学式（用 [ ] 标明内界）。

\item \laiyuan{2024年广东卷第13题，来源：仓库所附真题文档}\ 下列陈述Ⅰ与陈述Ⅱ均正确，且具有因果关系的是\xz
\xxa{酸性：\ce{CF3COOH}$<$\ce{CCl3COOH}——电负性：F$>$Cl}{某冠醚与 \ce{Li+} 能形成超分子，与 \ce{K+} 则不能——\ce{Li+} 与 \ce{K+} 的离子半径不同}{由氨制硝酸：\ce{NH3 -> NO -> NO2 -> HNO3}——\ce{NH3} 和 \ce{NO2} 均具有氧化性}{苯酚与甲醛反应，可合成酚醛树脂——合成酚醛树脂的反应是加聚反应}

\item \laiyuan{原创}\ 硼酸（\ce{H3BO3}）是一元弱酸，其电离方式为 \ce{H3BO3 + H2O <=> [B(OH)4]- + H+}。(1) 硼酸分子中 B 的杂化方式为\tk，\ce{[B(OH)4]-} 中 B 的杂化方式为\tk。(2) 硼酸为什么是一元酸？(3) 硼酸晶体呈片层状，层内分子间以\tk{}相连，层间为\tk，所以硼酸晶体有滑腻感。
\end{xiti}

\zu{C组\quad 冲刺高分}
\begin{xiti}
\item \laiyuan{原创}\ 尿素 \ce{CO(NH2)2} 是重要的氮肥，分子中 C 与 O 以双键相连，C 与两个 N 以单键相连。(1) 尿素中 C 原子的杂化方式为\tk，1 mol 尿素含 $\sigma$ 键\tk mol。(2) 实验测得尿素分子中所有原子共平面，N 原子的杂化方式应为\tk，这说明 N 原子上的孤电子对\tk。(3) 尿素易溶于水，熔点（133 ℃）比相对分子质量相近的乙酸（16.6 ℃）高得多，从分子间作用力角度解释原因。(4) 尿素与 \ce{H2O2} 可通过氢键形成加合物 \ce{CO(NH2)2*H2O2}，该加合物属于\tk（填“配合物”或“超分子”）。

\item \laiyuan{原创}\ 比较并解释酸性强弱：(1) \ce{ClCH2COOH}、\ce{Cl2CHCOOH}、\ce{FCH2COOH}；(2) \ce{H3PO4} 与 \ce{HClO4}；(3) 甲酸、乙酸、丙酸；(4) 为什么 HF 的水溶液是弱酸，而 HCl 是强酸？
\end{xiti}

\zhangmo{练习解析}
\begin{jiexilist}
\item \ans{C}\ce{SO2}：价层电子对数 3，$\mathrm{sp}^2$，V 形；\ce{NH3}：4 对，$\mathrm{sp}^3$，三角锥形；\ce{CO3^2-}：$\frac12(4+2-3\times2)=0$，3 对，$\mathrm{sp}^2$，平面三角形，正确；\ce{H3O+}：4 对含 1 对孤电子对，三角锥形。
\item 6；3；sp 和 $\mathrm{sp}^2$。C—H 3 个、C=C 1$\sigma$1$\pi$、C—C 1 个、C$\equiv$N 1$\sigma$2$\pi$：$\sigma$ 键 $3+1+1+1=6$，$\pi$ 键 3；双键碳 $\mathrm{sp}^2$，三键碳 sp。
\item \ans{B}A：稳定性取决于键能（O—H$>$S—H），与氢键无关；B 正确；C：键角 \ce{CH4}$>$\ce{NH3}$>$\ce{H2O}；D：邻羟基苯甲醛形成分子内氢键，沸点较低。
\item \ans{C}A：\ce{HClO4} 非羟基氧多，酸性更强；B：Cl 吸电子，\ce{CH2ClCOOH} 酸性更强；C：F 电负性大于 Cl，正确；D：乙基推电子能力强于甲基，丙酸酸性更弱。
\item \ans{C}\ce{BF3} 为平面三角形，B—F 键有极性但分子对称，正负电荷中心重合；\ce{H2O}（V 形）、\ce{NH3}（三角锥形）、\ce{CH3Cl}（四面体但不对称）均为极性分子。
\item $\mathrm{sp}^3$；sp；$\mathrm{sp}^2$；$\mathrm{sp}^3$。
\item (1) \ce{NH3}$>$\ce{PH3}：N 的电负性大于 P，成键电子对更靠近 N，彼此排斥更大，键角更大。(2) \ce{NH3}$>$\ce{NF3}：F 的电负性大于 N，N—F 键的成键电子对偏向 F，远离 N，排斥减小，键角变小；而 N—H 键电子对偏向 N。(3) \ce{SO4^2-}$>$\ce{SO3^2-}：均为 $\mathrm{sp}^3$ 杂化，\ce{SO3^2-} 有 1 对孤电子对，排斥使键角减小。(4) \ce{BF3}$>$\ce{NF3}：\ce{BF3} 中 B 为 $\mathrm{sp}^2$ 杂化（120°），\ce{NF3} 中 N 为 $\mathrm{sp}^3$ 杂化且有孤电子对（约 102°）。
\item (1) 先生成蓝色沉淀，后沉淀溶解，得到深蓝色溶液：\ce{Cu^2+ + 2NH3*H2O = Cu(OH)2 v + 2NH4+}；\ce{Cu(OH)2 + 4NH3 = [Cu(NH3)4]^2+ + 2OH-}。(2) 降低溶剂极性，减小 \ce{[Cu(NH3)4]SO4} 的溶解度，使其结晶析出。(3) N；16。(4) \ce{NH3} 中 N 有 1 对孤电子对，对成键电子对排斥较大；形成配合物后，孤电子对与 \ce{Cu^2+} 形成配位键，变为成键电子对，排斥减小，键角增大。
\item X：\ce{[Co(NH3)5(H2O)]Cl3}；Y：\ce{[Co(NH3)5Cl]Cl2*H2O}。X 中 3 个 \ce{Cl-} 都在外界；Y 中有 1 个 Cl 在内界作配体，故只有 2 个 \ce{Cl-} 能沉淀，此时配位数 6 由 5 个 \ce{NH3} 和 1 个 \ce{Cl-} 构成，水在外界。
\item \ans{B}A：F 的电负性大，吸电子能力强，\ce{CF3COOH} 酸性\textbf{强于} \ce{CCl3COOH}，陈述Ⅰ错；B：冠醚空腔与离子大小匹配才能形成超分子，正确且有因果关系；C：氨被氧化体现还原性，陈述Ⅱ错；D：酚醛树脂由缩聚反应生成，陈述Ⅱ错。
\item (1) $\mathrm{sp}^2$；$\mathrm{sp}^3$。(2) B 有空的 p 轨道，硼酸接受水电离出的 \ce{OH-} 的孤电子对形成配位键，生成 \ce{[B(OH)4]-}，同时释放出 1 个 \ce{H+}；它本身的羟基不电离，所以是一元酸（路易斯酸）。(3) 氢键；范德华力。
\item (1) $\mathrm{sp}^2$；7（C=O 1 个、C—N 2 个、N—H 4 个）。(2) $\mathrm{sp}^2$；未参与杂化，处于未杂化的 p 轨道中，与 C=O 的 $\pi$ 键形成离域的大 $\pi$ 键（共轭），使分子共平面。(3) 尿素分子中有 4 个 N—H 键，N、O 原子都有孤电子对，分子间能形成的氢键数目多于乙酸，分子间作用力更强，熔点更高；尿素与水分子间也能形成大量氢键，所以易溶于水。(4) 超分子（分子间通过氢键形成的聚集体，不含配位键）。
\item (1) \ce{Cl2CHCOOH}$>$\ce{FCH2COOH}$>$\ce{ClCH2COOH}。吸电子基团越多、电负性越大，羧基中 O—H 键极性越强，越易电离。（两个 Cl 的吸电子总效应强于一个 F，这是实验事实：p$K_\text{a}$ 分别约为 1.3、2.6、2.9。）(2) \ce{HClO4}（非羟基氧数 3）$>$\ce{H3PO4}（非羟基氧数 1）。(3) 甲酸$>$乙酸$>$丙酸，烷基是推电子基团，烷基越大推电子越强，O—H 键极性越弱。(4) H—F 键键能大（键很牢固），且 F$^-$ 与 \ce{H2O}、HF 形成强氢键，HF 难以完全电离；H—Cl 键较弱，在水中完全电离。
\end{jiexilist}
% <<<<<<<< 结束：chapters/ch07.tex

% >>>>>>>> 开始：chapters/ch08.tex
% ============================================================
% 第8章 晶体结构与性质（选必2第三章）
% ============================================================
\chapter{晶体结构与性质}

% 斜二测立方体与原子绘制（边长 2，纵深方向偏移 (0.8,0.55)）
\newcommand{\cubeat}[1][black]{%
 \draw[#1] (0,0)--(2,0)--(2,2)--(0,2)--cycle;
 \draw[#1] (2,0)--(2.8,0.55)--(2.8,2.55)--(2,2); \draw[#1] (0,2)--(0.8,2.55)--(2.8,2.55);
 \draw[#1,dashed] (0,0)--(0.8,0.55)--(2.8,0.55); \draw[#1,dashed] (0.8,0.55)--(0.8,2.55);}
\newcommand{\at}[5]{\shade[ball color=#4] ({2*#1+0.8*#3},{2*#2+0.55*#3}) circle (#5);}

\begin{daoxue}
\textbf{知识主线：}晶体与非晶体的区别 $\to$ 晶胞与均摊法 $\to$ 四类晶体（构成粒子与作用力决定性质）$\to$ 典型晶胞结构 $\to$ 晶胞计算（化学式、配位数、距离、密度、原子坐标）。

\textbf{高考考情：}广东卷第18题（工艺流程）的最后一问连续三年考晶胞：2024年考 Co、Ga、Ni 合金的粒子个数比和晶胞体积；2025年考 \ce{Ni_xCu_yN_z} 的组成和配位数；2026年考 AlSb 的离子键成分和原子占位。另外选择题常考晶体类型与熔沸点比较（2026年第4题D、2025年第5题D）。

\textbf{学习目标：}①熟练运用均摊法；②掌握四类晶体的性质及熔沸点比较；③记住 NaCl、CsCl、\ce{CaF2}、ZnS（金刚石）、干冰、金属堆积等典型结构；④会计算晶胞密度、原子间距，读懂原子分数坐标。
\end{daoxue}

\section{晶体与晶胞}

\textbf{晶体}内部粒子在三维空间里呈周期性有序排列，具有\textbf{自范性}（能自发呈现多面体外形）、各向异性和固定熔点；\textbf{非晶体}（玻璃、石蜡）则没有。区分晶体与非晶体最可靠的方法是 \textbf{X 射线衍射}：晶体会产生明锐的衍射峰（分立的斑点），非晶体只有弥散的宽峰（2026年广东卷第17题(3)考查了这一点）。

\begin{gainian}[晶胞与均摊法]
晶胞是描述晶体结构的基本单元，晶体可看作数量巨大的晶胞“无隙并置”而成。立方晶胞中，位于\textbf{顶点}的粒子为 8 个晶胞共有，每个晶胞占 $\frac18$；\textbf{棱上}占 $\frac14$；\textbf{面上}占 $\frac12$；\textbf{体内}占 1。（六棱柱晶胞：顶点 $\frac16$，上下底面棱 $\frac14$、侧棱 $\frac13$，面心 $\frac12$。）
\end{gainian}

\section{四类晶体}

\begin{center}
\small
\begin{tabularx}{\linewidth}{p{1.4cm}XXXX}
\toprule
& 分子晶体 & 共价晶体 & 离子晶体 & 金属晶体 \\
\midrule
构成粒子 & 分子 & 原子 & 阴、阳离子 & 金属阳离子、自由电子 \\
作用力 & 范德华力（有的有氢键） & 共价键 & 离子键 & 金属键 \\
熔沸点 & 低 & 很高 & 较高 & 差别大 \\
硬度 & 小 & 很大 & 较大，较脆 & 差别大，有延展性 \\
导电性 & 固态、熔融不导电 & 一般不导电（Si 为半导体） & 熔融或溶于水导电 & 固态导电 \\
实例 & 干冰、冰、\ce{I2}、大多数有机物 & 金刚石、晶体硅、SiC、\ce{SiO2} & NaCl、CsCl、MgO & Na、Cu、Fe \\
\bottomrule
\end{tabularx}
\end{center}

\textbf{金属键的“电子气”理论}：金属原子脱落下来的价电子形成遍布整块晶体的“电子气”，被所有金属阳离子共用，从而把它们维系在一起。这一模型能解释金属的通性：①导电性——在外加电场作用下自由电子定向移动；②导热性——自由电子在运动中与金属阳离子碰撞传递能量；③延展性——金属阳离子层之间发生相对滑动时，“电子气”仍能把它们结合在一起，金属键不被破坏。

\textbf{混合型晶体——石墨}：层内 C 原子 $\mathrm{sp}^2$ 杂化，以共价键形成平面六元环网状结构，未杂化的 p 电子形成遍及整层的大 $\pi$ 键，可在层内移动，所以能导电；层与层之间是范德华力，易滑动，质软。1 mol 石墨含 1.5 mol C—C 键。

\begin{fangfa}[熔沸点比较的思路]
①\textbf{先看晶体类型}：一般共价晶体$>$离子晶体$>$分子晶体，金属晶体差别很大。②\textbf{同类型再比较作用力}：共价晶体——键长越短、键能越大，熔点越高：金刚石$>$SiC$>$晶体硅；离子晶体——离子电荷越高、半径越小，离子键越强：MgO$>$NaCl$>$KCl；分子晶体——有分子间氢键的反常高，组成结构相似的相对分子质量越大越高，相对分子质量相近的极性越大越高；金属晶体——金属阳离子半径越小、价电子越多，金属键越强：Al$>$Mg$>$Na，Li$>$Na$>$K。
\end{fangfa}

\begin{tuozhan}[离子液体、液晶、纳米材料]
离子液体由体积很大的阴、阳离子构成，离子键较弱，室温下呈液态，几乎不挥发，可作溶剂或电解质；液晶介于液态与晶态之间，既有流动性又有类似晶体的各向异性；纳米晶体的表面原子比例大，性质与普通晶体不同。2026年广东卷第19题(5)中的“低共熔溶剂”也属于这类新型溶剂。
\end{tuozhan}

\section{典型晶体结构}

\begin{center}
\begin{tikzpicture}[scale=0.85]
\begin{scope}
\cubeat
\foreach \x/\y/\z in {0/0/0,1/0/0,0/1/0,1/1/0,0/0/1,1/0/1,0/1/1,1/1/1,0.5/0.5/0,0.5/0.5/1,0.5/0/0.5,0.5/1/0.5,0/0.5/0.5,1/0.5/0.5}{\at{\x}{\y}{\z}{mainblue!70}{0.13}}
\foreach \x/\y/\z in {0.5/0/0,0.5/1/0,0/0.5/0,1/0.5/0,0.5/0/1,0.5/1/1,0/0.5/1,1/0.5/1,0/0/0.5,1/0/0.5,0/1/0.5,1/1/0.5,0.5/0.5/0.5}{\at{\x}{\y}{\z}{warnred!80}{0.08}}
\node[font=\scriptsize,align=center] at (1.4,-0.6) {NaCl 型\\蓝 \ce{Cl-}、红 \ce{Na+}，配位数均为 6};
\end{scope}
\begin{scope}[xshift=4.3cm]
\cubeat
\foreach \x/\y/\z in {0/0/0,1/0/0,0/1/0,1/1/0,0/0/1,1/0/1,0/1/1,1/1/1}{\at{\x}{\y}{\z}{mainblue!70}{0.13}}
\at{0.5}{0.5}{0.5}{warnred!80}{0.15}
\node[font=\scriptsize,align=center] at (1.4,-0.6) {CsCl 型\\配位数均为 8};
\end{scope}
\begin{scope}[xshift=8.6cm]
\cubeat
\foreach \x/\y/\z in {0/0/0,1/0/0,0/1/0,1/1/0,0/0/1,1/0/1,0/1/1,1/1/1,0.5/0.5/0,0.5/0.5/1,0.5/0/0.5,0.5/1/0.5,0/0.5/0.5,1/0.5/0.5}{\at{\x}{\y}{\z}{mainblue!70}{0.12}}
\foreach \x/\y/\z in {0.25/0.25/0.25,0.75/0.75/0.25,0.75/0.25/0.75,0.25/0.75/0.75}{\at{\x}{\y}{\z}{warnred!80}{0.12}}
\node[font=\scriptsize,align=center] at (1.4,-0.6) {ZnS 型（闪锌矿）\\4 个红球在体内四面体空隙};
\end{scope}
\end{tikzpicture}
\end{center}

\begin{center}
\small
\begin{tabularx}{\linewidth}{p{2.2cm}Xp{3.6cm}}
\toprule
晶胞 & 结构特点 & 每个晶胞含粒子数、配位数 \\
\midrule
NaCl & \ce{Cl-} 面心立方，\ce{Na+} 位于棱心和体心 & 4 个 NaCl；配位数 6:6 \\
CsCl & \ce{Cl-} 位于顶点，\ce{Cs+} 位于体心 & 1 个 CsCl；配位数 8:8 \\
\ce{CaF2}（萤石） & \ce{Ca^2+} 面心立方，8 个 \ce{F-} 位于 8 个四面体空隙 & 4 个 \ce{CaF2}；\ce{Ca^2+} 配位数 8，\ce{F-} 配位数 4 \\
ZnS（闪锌矿） & \ce{S^2-} 面心立方，\ce{Zn^2+} 占据一半（4 个）四面体空隙 & 4 个 ZnS；配位数 4:4 \\
金刚石（Si、SiC 类似） & 在 ZnS 结构中全部换成 C 原子 & 8 个 C；每个 C 连 4 个 C \\
干冰 & \ce{CO2} 分子面心立方 & 4 个 \ce{CO2}；每个分子周围紧邻 12 个 \\
金属体心立方 & Na、K、Fe 等 & 2 个原子；配位数 8；空间利用率 68\% \\
金属面心立方 & Cu、Ag、Au、Al 等（与六方最密堆积 Mg、Zn 一样为最密堆积） & 4 个原子；配位数 12；空间利用率 74\% \\
\bottomrule
\end{tabularx}
\end{center}

\begin{center}
\begin{tikzpicture}[scale=0.85]
\begin{scope}
\cubeat
\foreach \x/\y/\z in {0/0/0,1/0/0,0/1/0,1/1/0,0/0/1,1/0/1,0/1/1,1/1/1}{\at{\x}{\y}{\z}{orange!80}{0.15}}
\at{0.5}{0.5}{0.5}{orange!80}{0.15}
\node[font=\scriptsize,align=center] at (1.4,-0.6) {体心立方（Na、K、Fe）\\2 个原子，配位数 8};
\end{scope}
\begin{scope}[xshift=4.6cm]
\cubeat
\foreach \x/\y/\z in {0/0/0,1/0/0,0/1/0,1/1/0,0/0/1,1/0/1,0/1/1,1/1/1,0.5/0.5/0,0.5/0.5/1,0.5/0/0.5,0.5/1/0.5,0/0.5/0.5,1/0.5/0.5}{\at{\x}{\y}{\z}{orange!80}{0.14}}
\node[font=\scriptsize,align=center] at (1.4,-0.6) {面心立方（Cu、Ag、Au、Al）\\4 个原子，配位数 12};
\end{scope}
\begin{scope}[xshift=9.2cm]
\cubeat
\foreach \x/\y/\z in {0/0/0,1/0/0,0/1/0,1/1/0,0/0/1,1/0/1,0/1/1,1/1/1}{\at{\x}{\y}{\z}{deepgreen!70}{0.14}}
\at{0.5}{0.5}{0.5}{mainblue!70}{0.15}
\foreach \x/\y/\z in {0.5/0.5/0,0.5/0.5/1,0.5/0/0.5,0.5/1/0.5,0/0.5/0.5,1/0.5/0.5}{\at{\x}{\y}{\z}{warnred!80}{0.10}}
\node[font=\scriptsize,align=center] at (1.4,-0.6) {钙钛矿型 \ce{CaTiO3}\\绿 Ca、蓝 Ti、红 O};
\end{scope}
\end{tikzpicture}
\end{center}
\noindent\small 钙钛矿 \ce{CaTiO3}：Ca 位于顶点（$8\times\frac18=1$），Ti 位于体心（1），O 位于面心（$6\times\frac12=3$）。Ti 周围有 6 个 O（八面体配位）；每个 Ca 周围最近的 O 有 12 个（共用该顶点的 8 个晶胞中，每个晶胞各有 3 个面心 O 与它最近，每个面心 O 为 2 个晶胞共用：$8\times3\div2=12$）。\normalsize

\section{晶胞的计算}

\begin{fangfa}[晶胞计算公式与几何关系]
\textbf{①密度：}$\rho=\dfrac{N\cdot M}{N_\text{A}\cdot V}=\dfrac{N\cdot M}{N_\text{A}\cdot a^3}$（$N$ 为晶胞所含“化学式单元”数）。单位换算：1 nm $=10^{-7}$ cm，1 pm $=10^{-10}$ cm，$a^3$ 的单位 cm$^3$。\\
\textbf{②原子间距（边长 $a$）：}面对角线 $\sqrt2a$，体对角线 $\sqrt3a$。面心立方中最近两原子相距 $\frac{\sqrt2}{2}a$；体心立方中顶点与体心相距 $\frac{\sqrt3}{2}a$；ZnS（金刚石）中最近的 Zn 与 S 相距 $\frac{\sqrt3}{4}a$（体对角线的 $\frac14$）；NaCl 中最近的 \ce{Na+} 与 \ce{Cl-} 相距 $\frac a2$。\\
\textbf{③原子分数坐标：}以晶胞一个顶点为原点，用边长的分数表示位置，如顶点 (0,0,0)，面心 $(\frac12,\frac12,0)$，体心 $(\frac12,\frac12,\frac12)$，ZnS 中体内原子 $(\frac14,\frac14,\frac14)$、$(\frac34,\frac34,\frac14)$ 等。
\end{fangfa}

\begin{liti}[原子分数坐标与投影图]
ZnS（闪锌矿）晶胞中，\ce{S^2-} 位于顶点和面心，\ce{Zn^2+} 位于体内四个互不相邻的四面体空隙中。(1) 若顶点 \ce{S^2-} 的坐标为 (0,0,0)，写出四个 \ce{Zn^2+} 的坐标；(2) 画出晶胞沿 $z$ 轴方向的投影图。
\tcblower
\jie (1) $(\frac14,\frac14,\frac14)$、$(\frac34,\frac34,\frac14)$、$(\frac34,\frac14,\frac34)$、$(\frac14,\frac34,\frac34)$（体对角线的 $\frac14$ 处，四个点互不相邻）。\\
(2) 沿 $z$ 轴投影时只看 $x$、$y$ 坐标：\ce{S^2-} 投影在正方形的四个顶点、四条边的中点和中心；\ce{Zn^2+} 投影在 $(\frac14,\frac14)$、$(\frac34,\frac34)$、$(\frac34,\frac14)$、$(\frac14,\frac34)$ 四点：
\begin{center}
\begin{tikzpicture}[scale=1.3]
\draw (0,0) rectangle (2,2);
\foreach \x/\y in {0/0,2/0,0/2,2/2,1/0,1/2,0/1,2/1,1/1}{\fill[mainblue] (\x,\y) circle (2.2pt);}
\foreach \x/\y in {0.5/0.5,1.5/1.5,1.5/0.5,0.5/1.5}{\fill[warnred] (\x,\y) circle (1.8pt);}
\node[font=\scriptsize,right] at (2.2,1.4) {\textcolor{mainblue}{$\bullet$} \ce{S^2-}};
\node[font=\scriptsize,right] at (2.2,0.9) {\textcolor{warnred}{$\bullet$} \ce{Zn^2+}};
\end{tikzpicture}
\end{center}
\dianbo 投影图中“同一投影点可能对应上下多个原子”，所以不能直接用投影图数原子个数。
\end{liti}

\begin{liti}[密度计算]
铜晶体为面心立方结构，晶胞边长 $a=361$ pm，$M(\ce{Cu})=64$ g·mol$^{-1}$，$N_\text{A}=6.02\times10^{23}$ mol$^{-1}$。求铜的密度和铜原子半径。
\tcblower
\jie 每个晶胞含 Cu：$8\times\frac18+6\times\frac12=4$。$\rho=\dfrac{4\times64}{6.02\times10^{23}\times(361\times10^{-10})^3}\approx9.0$ g·cm$^{-3}$。面心立方中原子沿面对角线相切：$4r=\sqrt2a$，$r=\frac{\sqrt2}{4}\times361\approx128$ pm。\\
\xianjing 忘记把 pm 换算成 cm，或把 $N$ 取成 14（未用均摊法）。
\end{liti}

\begin{liti}[空间利用率]
求面心立方最密堆积的空间利用率（原子视为刚性球）。
\tcblower
\jie 每个晶胞含 4 个原子；沿面对角线相切，$4r=\sqrt2a$，$r=\frac{\sqrt2}{4}a$。空间利用率 $=\dfrac{4\times\frac43\pi r^3}{a^3}=\dfrac{\frac{16}{3}\pi\left(\frac{\sqrt2}{4}\right)^3a^3}{a^3}=\dfrac{\sqrt2\pi}{6}\approx74\%$。同理可得体心立方（沿体对角线相切，$4r=\sqrt3a$，2 个原子）为 $\dfrac{\sqrt3\pi}{8}\approx68\%$。
\end{liti}

\begin{zhenti}{2025年广东卷第18题(5)\quad 来源：仓库所附真题文档}
\ce{Ni_xCu_yN_z} 晶体的立方晶胞中，原子分别位于顶点、面心和体心。已知同种位置原子相同，相邻原子间的最近距离之比 $d_{\text{Ni—Cu}}:d_{\text{Ni—N}}=\sqrt2:1$，则 $x:y:z=$\tk；晶体中与 Cu 原子最近且等距离的原子数目为\tk。

\textbf{解析：}设边长为 $a$：面心到顶点距离 $\frac{\sqrt2}{2}a$，面心到体心距离 $\frac12a$，二者之比为 $\sqrt2:1$。可知 Ni 在面心（同时与 Cu、N 相邻），Cu 在顶点，N 在体心。Ni：$6\times\frac12=3$；Cu：$8\times\frac18=1$；N：1。$x:y:z=3:1:1$。每个顶点 Cu 周围最近的是相邻面心上的 Ni，共 $3\times8\times\frac12=12$ 个。
\end{zhenti}

\begin{zhenti}{2026年广东卷第18题(7)\quad 来源：仓库所附真题文档（网络收集版）}
①Sb 与 N 同族，AlSb 与 GaN 中离子键成分的百分数较大的是\tk。②AlSb 晶体的立方晶胞结构中，原子占据 A（顶点）、B（面心）、C（体内）三种位置，每种位置上只有一种原子，则 Al 占据\tk{}位置，Sb 占据\tk{}位置。

\textbf{解析：}①离子键成分随两元素电负性差增大而增大；N 的电负性远大于 Sb，Ga 与 Al 电负性相近，所以 GaN 的电负性差更大，离子键成分更大。②AlSb 为 1:1 型，A$+$B 共 $8\times\frac18+6\times\frac12=4$ 个，C（体内）4 个，故顶点与面心是同一种原子，体内是另一种原子（与 ZnS 结构相同）：Al 占 A、B 则 Sb 占 C，或 Al 占 C 则 Sb 占 A、B，二者可互换。
\end{zhenti}

% ---------------- 章末 ----------------
\zhangmo{本章小结}
\subsection*{一、知识总结}
\begin{center}
\small
\begin{tabularx}{\linewidth}{p{2.4cm}X}
\toprule
晶体特征 & 自范性、各向异性、固定熔点；X 射线衍射鉴别 \\
均摊法 & 立方晶胞：顶点 $\frac18$、棱 $\frac14$、面 $\frac12$、体内 1 \\
晶体类型 & 分子、共价、离子、金属；混合型石墨；熔沸点“先类型、后作用力” \\
典型结构 & NaCl（6）、CsCl（8）、\ce{CaF2}（8:4）、ZnS/金刚石（4）、干冰（12）、金属 bcc（8）fcc（12） \\
计算 & $\rho=\frac{NM}{N_\text{A}a^3}$；面对角线 $\sqrt2a$、体对角线 $\sqrt3a$；分数坐标 \\
\bottomrule
\end{tabularx}
\end{center}
\subsection*{二、方法提升}
\begin{enumerate}
\item 求化学式：均摊法数出各粒子个数，化为最简比。
\item 求配位数：选一个粒子，找“最近且等距”的异种（或同种）粒子，注意相邻晶胞中的粒子。
\item 求距离：先确定两粒子在晶胞中的几何位置（同一面、体对角线等），再用勾股定理。
\item 密度题的答案常要求“列出计算式”，分子写 $N\cdot M$，分母写 $N_\text{A}\cdot a^3$，单位换算放进式子里。
\end{enumerate}
\subsection*{三、易错点归纳}
\begin{enumerate}
\item 把分子晶体中的共价键强弱当作熔沸点高低的原因（应看分子间作用力）。
\item NaCl 晶胞中 \ce{Na+} 的个数：棱心 $12\times\frac14+$ 体心 $1=4$，不是 13。
\item 金刚石晶胞含 8 个 C 原子，1 mol 金刚石含 2 mol C—C 键；1 mol \ce{SiO2} 含 4 mol Si—O 键。
\item 单位换算错误：$(a\ \mathrm{pm})^3=a^3\times10^{-30}$ cm$^3$。
\end{enumerate}

\zhangmo{分层练习}
\zu{A组\quad 基础巩固}
\begin{xiti}[start=1]
\item 下列说法正确的是\xz
\xxa{具有规则几何外形的固体一定是晶体}{区分晶体与非晶体最可靠的方法是测定熔点}{晶体硅、SiC、金刚石都是共价晶体，熔点：金刚石$>$SiC$>$晶体硅}{分子晶体中分子内共价键越强，熔点越高}

\item 下列各组物质熔点由高到低排列正确的是\xz
\xxa{\ce{CO2}$>$\ce{SiO2}$>$NaCl}{MgO$>$NaCl$>$KCl}{Na$>$Mg$>$Al}{HI$>$HBr$>$HCl$>$HF}

\item 某晶体的立方晶胞中，A 原子位于顶点，B 原子位于面心，C 原子位于体心，则该晶体的化学式为\tk；与 C 原子最近的 B 原子有\tk{}个。

\item 某物质熔点 3550 ℃，硬度很大，不溶于水，固态和熔融态都不导电，该物质最可能是\tk{}晶体；另一物质熔点 801 ℃，易溶于水，固态不导电而熔融态导电，它属于\tk{}晶体。

\item 下列各组晶体中，化学键类型相同、晶体类型也相同的是\xz
\xx{\ce{SO2} 与 \ce{SiO2}}{\ce{CO2} 与 \ce{H2O}}{NaCl 与 HCl}{\ce{CCl4} 与 KCl}
\end{xiti}

\zu{B组\quad 高考典型}
\begin{xiti}
\item \laiyuan{2024年广东卷第18题(6)，来源：仓库所附真题文档}\ 一种含 Ga、Ni、Co 元素的记忆合金的晶体结构可描述为 Ga 与 Ni 交替填充在 Co 构成的立方体体心，形成如下结构单元：边长为 $a$ nm 的小立方体，8 个顶点为 Co，体心为 Ga 或 Ni；相邻两个小立方体的体心分别为 Ga 和 Ni。该合金的晶胞（由 8 个小立方体构成的大立方体）中，粒子个数最简比 Co:Ga:Ni$=$\tk，其立方晶胞的体积为\tk nm$^3$。

\item \laiyuan{原创}\ 金刚石晶胞（与 ZnS 结构类似，全为 C 原子）边长为 $a$ pm。(1) 每个晶胞含 C 原子\tk{}个；(2) 最近两个 C 原子的距离为\tk pm；(3) 若原点处 C 原子坐标为 (0,0,0)，写出另外两个位于体内的 C 原子坐标；(4) 列式表示金刚石的密度（g·cm$^{-3}$，阿伏加德罗常数为 $N_\text{A}$）。

\item \laiyuan{原创}\ \ce{CaF2} 晶胞中 \ce{Ca^2+} 位于顶点和面心，\ce{F-} 位于体内 8 个四面体空隙中，边长为 $a$ nm。(1) \ce{Ca^2+} 和 \ce{F-} 的配位数分别为\tk、\tk；(2) 最近的 \ce{Ca^2+} 与 \ce{F-} 的距离为\tk nm；(3) 晶体密度为\tk g·cm$^{-3}$（列式，$M(\ce{CaF2})=78$ g·mol$^{-1}$）。

\item \laiyuan{原创}\ CsCl 晶胞中 \ce{Cl-} 位于顶点、\ce{Cs+} 位于体心，晶胞边长为 $a$ pm。(1) 每个 \ce{Cs+} 周围距离最近的 \ce{Cs+} 有\tk{}个；(2) \ce{Cs+} 与 \ce{Cl-} 的最近距离为\tk pm；(3) 列式表示晶体密度（$M(\ce{CsCl})=168.5$ g·mol$^{-1}$，阿伏加德罗常数为 $N_\text{A}$）。

\item \laiyuan{原创}\ 钙钛矿 \ce{CaTiO3} 的晶胞如本章图示（Ca 在顶点、Ti 在体心、O 在面心）。(1) 若改为以 Ti 为顶点重新取晶胞，则 Ca 位于\tk，O 位于\tk。(2) 若晶胞边长为 $a$ nm，Ti 与 O 的最近距离为\tk nm，Ca 与 O 的最近距离为\tk nm。
\end{xiti}

\zu{C组\quad 冲刺高分}
\begin{xiti}
\item \laiyuan{原创}\ 氮化硼（BN）有类似石墨的六方结构（h-BN）和类似金刚石的立方结构（c-BN）。(1) h-BN 中 B 原子的杂化方式为\tk，h-BN 不导电而石墨导电，从结构上解释原因。(2) c-BN 的晶胞与 ZnS 相同，B 原子位于体内四面体空隙。c-BN 中 B 原子的杂化方式为\tk，B 的配位数为\tk。(3) c-BN 的硬度仅次于金刚石，熔点却低于金刚石，解释原因。(4) 已知 c-BN 晶胞边长为 361.5 pm，求其密度（保留两位小数，$N_\text{A}=6.02\times10^{23}$ mol$^{-1}$，B 为 10.8，N 为 14.0）。

\item \laiyuan{原创}\ 某铜金合金具有面心立方结构，Au 原子位于顶点、Cu 原子位于面心。(1) 该合金的化学式为\tk。(2) 该合金可以储氢：H 原子进入由 Cu、Au 原子构成的四面体空隙中。若把 Cu、Au 原子等同看待，储氢后晶胞结构与 \ce{CaF2} 相似（H 占据全部四面体空隙），则储氢后的化学式为\tk。(3) 若晶胞边长为 $a$ cm，储氢后的晶体中氢的密度（单位体积内 H 的质量）为多少 g·cm$^{-3}$？
\end{xiti}

\zhangmo{练习解析}
\begin{jiexilist}
\item \ans{C}A：玻璃可制成规则外形，但不是晶体；B：最可靠的是 X 射线衍射；C：键长 C—C$<$Si—C$<$Si—Si，键能依次减小，熔点依次降低，正确；D：分子晶体熔点由分子间作用力决定。
\item \ans{B}A：\ce{SiO2}（共价晶体）$>$NaCl（离子晶体）$>$\ce{CO2}（分子晶体）；B：离子电荷 \ce{Mg^2+}、\ce{O^2-} 高且半径小，MgO 最高，\ce{Na+} 半径小于 \ce{K+}，正确；C：金属键 Al$>$Mg$>$Na；D：HF 分子间有氢键，熔点反常高。
\item \ce{AB3C}；6。A：$8\times\frac18=1$；B：$6\times\frac12=3$；C：1。体心 C 与 6 个面心 B 距离均为 $\frac a2$，最近。
\item 共价；离子。
\item \ans{B}A：\ce{SO2} 为分子晶体，\ce{SiO2} 为共价晶体；B：均为含极性共价键的分子晶体，正确；C：NaCl 含离子键、为离子晶体，HCl 含共价键、为分子晶体；D：\ce{CCl4} 为分子晶体，KCl 为离子晶体。
\item 2:1:1；$8a^3$。大立方体（晶胞）边长为 $2a$ nm，由 8 个小立方体构成；Co 位于大立方体的顶点、面心、棱心和体心：$8\times\frac18+6\times\frac12+12\times\frac14+1=8$ 个；8 个小立方体体心中 Ga、Ni 各 4 个。Co:Ga:Ni$=8:4:4=2:1:1$；晶胞体积 $(2a)^3=8a^3$ nm$^3$。
\item (1) 8（$8\times\frac18+6\times\frac12+4$）。(2) $\frac{\sqrt3}{4}a$（体内原子与最近顶点原子沿体对角线相距 $\frac14$）。(3) 如 $(\frac14,\frac14,\frac14)$、$(\frac34,\frac34,\frac14)$（其余为 $(\frac34,\frac14,\frac34)$、$(\frac14,\frac34,\frac34)$）。(4) $\rho=\dfrac{8\times12}{N_\text{A}\times(a\times10^{-10})^3}$ g·cm$^{-3}$。
\item (1) 8；4。每个 \ce{F-} 位于由 4 个 \ce{Ca^2+} 构成的四面体中心，\ce{F-} 配位数 4；由化学式 \ce{CaF2}，\ce{Ca^2+} 配位数为 $4\times2=8$。(2) $\frac{\sqrt3}{4}a$。(3) $\dfrac{4\times78}{N_\text{A}\times(a\times10^{-7})^3}$。
\item (1) 6（相邻晶胞体心的 \ce{Cs+} 位于上下前后左右，距离为 $a$）。(2) $\frac{\sqrt3}{2}a$（体心到顶点）。(3) $\rho=\dfrac{168.5}{N_\text{A}\times(a\times10^{-10})^3}$ g·cm$^{-3}$（每个晶胞含 1 个 CsCl）。
\item (1) 体心；棱心（以 Ti 为顶点时，原体心的 Ti 成为新顶点，原顶点的 Ca 落在新晶胞体心，原面心的 O 落在新晶胞棱心，组成仍为 \ce{CaTiO3}）。(2) $\frac a2$（体心到面心）；$\frac{\sqrt2}{2}a$（顶点到相邻面心）。
\item (1) $\mathrm{sp}^2$；h-BN 层内 B、N 电负性差较大，N 原子上的 p 电子（孤电子对）被 N 强烈吸引，形成的 $\pi$ 电子基本定域在 N 原子附近，不能像石墨中的大 $\pi$ 电子那样在层内自由移动，所以不导电。(2) $\mathrm{sp}^3$；4。(3) 二者都是共价晶体，B—N 键长略大于 C—C 键长，键能较小，熔点较低；但 B—N 键有一定极性且结构同样是立体网状，硬度仍很大。(4) 每个晶胞含 4 个 BN：$\rho=\dfrac{4\times24.8}{6.02\times10^{23}\times(361.5\times10^{-10})^3}\approx3.49$ g·cm$^{-3}$。
\item (1) \ce{Cu3Au}（Au：$8\times\frac18=1$；Cu：$6\times\frac12=3$）。(2) \ce{Cu3AuH8}：面心立方中共有 8 个四面体空隙，全部被 H 占据，每个晶胞含 8 个 H。(3) 每个晶胞含 8 个 H，质量 $\dfrac{8\times1}{N_\text{A}}$ g，氢的密度为 $\dfrac{8}{N_\text{A}\cdot a^3}$ g·cm$^{-3}$。
\end{jiexilist}
% <<<<<<<< 结束：chapters/ch08.tex

% >>>>>>>> 开始：chapters/ch09.tex
% ============================================================
% 第9章 有机化合物的结构特点与烃（必修二第七章、选必3第一、二章）
% ============================================================
\chapter{有机物的结构与烃}

\begin{daoxue}
\textbf{知识主线：}有机物为什么种类繁多（碳的成键特点、同分异构）$\to$ 怎样分类和命名（官能团、系统命名）$\to$ 怎样确定结构（分离提纯、元素分析、质谱、红外、核磁）$\to$ 最简单的有机物——烃的性质（烷：取代；烯、炔：加成、氧化、加聚；芳香烃：取代为主）。

\textbf{高考考情：}广东卷第20题（有机大题，14分）每年必考官能团名称、分子式、同分异构体数目或结构、核磁共振氢谱、碳原子杂化与共面（2024年第20题(2)、2025年第20题(2)(4)、2026年第20题(2)(3)）；选择题第9题左右考陌生有机物的结构与性质（2024年第9题柳珊瑚酸）。

\textbf{学习目标：}①会判断官能团、碳原子杂化方式和共线共面；②会有序书写同分异构体并计数；③会用谱图信息推断结构；④掌握各类烃的典型反应。
\end{daoxue}

\section{有机物的结构特点}

\subsection{碳原子的成键特点与官能团}
碳原子最外层 4 个电子，一般形成 4 个共价键；碳原子之间可形成单键、双键、三键，可以成链也可以成环——这是有机物种类繁多的原因。决定有机物特殊性质的原子或原子团叫\textbf{官能团}：

\begin{center}
\small
\begin{tabular}{llll}
\toprule
类别 & 官能团 & 类别 & 官能团 \\
\midrule
烯烃 & 碳碳双键 \ce{C=C} & 醛 & 醛基 \ce{-CHO} \\
炔烃 & 碳碳三键 \ce{C#C} & 酮 & 酮羰基 \ce{C=O} \\
卤代烃 & 碳卤键 \ce{-X} & 羧酸 & 羧基 \ce{-COOH} \\
醇 & （醇）羟基 \ce{-OH}（连在饱和碳上） & 酯 & 酯基 \ce{-COOR} \\
酚 & （酚）羟基 \ce{-OH}（直接连苯环） & 胺 & 氨基 \ce{-NH2} \\
醚 & 醚键 \ce{-O-} & 酰胺 & 酰胺基 \ce{-CONH-} \\
\bottomrule
\end{tabular}
\end{center}

\subsection{碳原子的杂化与分子的空间结构}
\textbf{饱和碳}（4 个单键）为 $\mathrm{sp}^3$ 杂化，与相连原子呈四面体；\textbf{双键碳}（含苯环碳、羰基碳）为 $\mathrm{sp}^2$ 杂化，与相连原子共平面；\textbf{三键碳}为 sp 杂化，与相连原子共直线。

\begin{fangfa}[共线共面的判断]
以三种基本模型为“积木”：\textbf{甲烷}（四面体，最多 3 个原子共面）；\textbf{乙烯}（6 个原子共面）；\textbf{乙炔}（4 个原子共线）；\textbf{苯}（12 个原子共面，对位上的 4 个原子共线）。单键可以旋转，所以“可能共面”与“一定共面”要区分：只要分子中有饱和碳原子（如 \ce{-CH3}、\ce{-CH2-}），所有原子就\textbf{不可能}共面。2024年广东卷第9题中，柳珊瑚酸的环上有一个碳同时与 4 个碳相连，为 $\mathrm{sp}^3$ 杂化（四面体），因此“环系结构中 3 个五元环共平面”是错误的。
\end{fangfa}

\subsection{同分异构现象}
分子式相同而结构不同的化合物互为同分异构体。类型：①\textbf{碳架异构}（如正丁烷与异丁烷）；②\textbf{位置异构}（如 1-丁烯与 2-丁烯）；③\textbf{官能团异构}（如乙醇与二甲醚，丙醛与丙酮，羧酸与酯）；④\textbf{顺反异构}（双键两端碳原子各连两个不同基团时，如 2-丁烯有顺、反两种）；⑤\textbf{对映异构}（含手性碳原子）。

\begin{fangfa}[同分异构体的有序书写与计数]
\textbf{①等效氢法求一元取代物数目：}同一碳上的 H 等效；同一碳上所连甲基上的 H 等效；处于对称位置的 H 等效。一元取代物种数 $=$ 等效氢种数。如 \ce{(CH3)3CCH2CH3} 有 3 种等效氢。\\
\textbf{②基元法：}丙基有 2 种、丁基有 4 种、戊基有 8 种。所以 \ce{C4H9Cl} 有 4 种，\ce{C4H9OH}（醇）有 4 种，\ce{C4H9CHO}（醛）、\ce{C4H9COOH}（羧酸）各有 4 种，\ce{C5H11OH}（醇）有 8 种。\\
\textbf{③定一移一法求二元取代物：}先固定一个取代基，再移动另一个。\\
\textbf{④苯环上的取代：}两个取代基有邻、间、对 3 种；三个取代基\textbf{都相同}有 3 种，\textbf{两同一异}有 6 种，\textbf{三个都不同}有 10 种。\\
\textbf{⑤酯类：}按“酸的碳数$+$醇的碳数”拆分，如 \ce{C4H8O2} 的酯：甲酸丙酯（2 种）、乙酸乙酯（1 种）、丙酸甲酯（1 种），共 4 种。
\end{fangfa}

\begin{zhenti}{2025年广东卷第20题(2)②\quad 来源：仓库所附真题文档}
化合物 2a（苯乙烯，\ce{C8H8}）可与 \ce{H2O} 发生加成反应生成化合物Ⅰ。在Ⅰ的同分异构体中，同时含有苯环和醇羟基结构的共\tk{}种（含化合物Ⅰ，不考虑立体异构）。

\textbf{解析：}Ⅰ的分子式为 \ce{C8H10O}，除苯环外剩 2 个 C，均为饱和结构。醇羟基不能直接连在苯环上：①一个侧链：\ce{C6H5CH2CH2OH}、\ce{C6H5CH(OH)CH3}，2 种；②两个侧链 \ce{-CH3} 和 \ce{-CH2OH}：邻、间、对 3 种。共 5 种。\\
\dianbo “醇羟基”排除了酚（\ce{-OH} 直接连苯环）；同时要排除醚（如 \ce{C6H5CH2OCH3}）。
\end{zhenti}

\begin{liti}[限定条件的同分异构体]
分子式为 \ce{C8H8O2}、含苯环且属于酯类的同分异构体共有几种？其中核磁共振氢谱有 4 组峰的是哪些？
\tcblower
\silu 不饱和度 $\Omega=\frac{2\times8+2-8}{2}=5$：苯环 4，酯基 1，其余均饱和。按“酸 $+$ 醇（酚）”拆分。\\
\jie ①甲酸酯：\ce{HCOOCH2C6H5}（1 种）、\ce{HCOOC6H4CH3}（甲基在邻、间、对，3 种）；②乙酸苯酯 \ce{CH3COOC6H5}（1 种）；③苯甲酸甲酯 \ce{C6H5COOCH3}（1 种）。共 6 种。\\
4 组峰：\ce{HCOOCH2C6H5}（苯环上 3 种 H $+$ \ce{CH2} $+$ HCOO 中的 H，共 5 组，不符）；甲酸对甲基苯酯 \ce{HCOO-C6H4-CH3}（对位）有 HCOO、苯环 2 种、\ce{CH3}，共 4 组，符合；\ce{CH3COOC6H5} 有 \ce{CH3} 和苯环 3 种，共 4 组，符合；\ce{C6H5COOCH3} 同理 4 组，符合。所以有 3 种。\\
\xianjing 一取代苯环上有邻、间、对 3 种 H；对位二取代（两基团不同）苯环上只有 2 种 H。
\end{liti}

\begin{liti}[共线共面]
\ce{CH3-CH=CH-C#C-CH3} 分子中，最多有几个碳原子共面？一定共线的碳原子有几个？
\tcblower
\jie 以 C=C 为中心：C1、C2、C3、C4 一定共面（乙烯模型）；C3—C4$\equiv$C5—C6 中，C4、C5 为 sp 杂化，C3、C4、C5、C6 一定共线，这条直线上的 C3、C4 已在平面内，所以 C5、C6 也在该平面内。6 个碳原子都共面；一定共线的是 C3、C4、C5、C6 共 4 个（C2—C3 不在这条直线上，因为 C3 是 $\mathrm{sp}^2$ 杂化，键角 120°）。
\end{liti}

\subsection{有机物的命名}
\textbf{烷烃系统命名}：选最长碳链为主链（等长时选支链多的），从离支链最近的一端编号（位次和最小），写成“位次—取代基名称—主链名称”，相同取代基合并并用“二、三”表示，不同取代基简单的写在前面。如 \ce{(CH3)2CHCH2C(CH3)3} 为 2,2,4-三甲基戊烷。\textbf{含官能团的有机物}：选含官能团的最长碳链，从离官能团最近的一端编号，标出官能团位置，如 \ce{CH2=CHCH(CH3)CH3} 为 3-甲基-1-丁烯，\ce{CH3CH(OH)CH3} 为 2-丙醇。苯的同系物可用邻、间、对或 1,2-二甲苯等表示。

\section{研究有机物的一般方法}
\textbf{分离提纯}：蒸馏（沸点相差较大的液体混合物）；重结晶（溶解度随温度变化差异大的固体，如粗苯甲酸：加热溶解$\to$趁热过滤$\to$冷却结晶$\to$过滤洗涤，2024年广东卷第6题、2025年第17题(2)均考查）；萃取分液；色谱法。

\textbf{确定实验式和分子式}：元素分析（燃烧法测 C、H，差值得 O）求实验式；\textbf{质谱}中最大质荷比即相对分子质量。

\textbf{确定结构}：\textbf{红外光谱}判断化学键和官能团；\textbf{核磁共振氢谱}中，吸收峰的\textbf{组数}$=$等效氢的种类数，峰面积之比$=$各类氢原子数之比。如乙醇 \ce{CH3CH2OH} 有 3 组峰，面积比 3:2:1；二甲醚 \ce{CH3OCH3} 只有 1 组峰。2025年广东卷第20题(4)：丙酮只有 1 组峰，溴代后生成 \ce{CH3COCH2Br}，单峰消失，出现两组新峰，面积比 3:2。

\section{烃的结构与性质}

\subsection{烷烃}
结构：碳原子均为 $\mathrm{sp}^3$ 杂化，链状烷烃通式 \ce{C_nH_{2n+2}}。性质稳定，不能使酸性 \ce{KMnO4} 溶液、溴水褪色；特征反应是光照下与卤素单质的\textbf{取代反应}（连锁反应，产物为多种卤代物的混合物）：\ce{CH4 + Cl2 ->[\text{光}] CH3Cl + HCl}。

\subsection{烯烃和炔烃}
\textbf{(1) 加成反应}：与 \ce{Br2}（溴水或溴的 \ce{CCl4} 溶液褪色）、\ce{H2}、HX、\ce{H2O} 等加成。不对称烯烃与 HX 加成时，主要产物中 H 加到含 H 较多的双键碳上（马氏规则），如 \ce{CH3CH=CH2 + HBr -> CH3CHBrCH3}（主要产物）。共轭二烯烃有 1,2-加成和 1,4-加成：\ce{CH2=CH-CH=CH2 + Br2 -> CH2Br-CH=CH-CH2Br}（1,4-加成）。\\
\textbf{(2) 氧化反应}：使酸性 \ce{KMnO4} 溶液褪色。\ce{RCH=CH2} 被氧化生成 \ce{RCOOH} 和 \ce{CO2}；\ce{R2C=CHR$'$} 生成酮 \ce{R2C=O} 和 \ce{R$'$COOH}——这一规律可用于由氧化产物反推双键位置。\\
\textbf{(3) 加聚反应}：\ce{$n$CH2=CH2 ->[\text{催化剂}] -[CH2-CH2]_{$n$}-}；\ce{$n$CH2=CHCl -> -[CH2-CHCl]_{$n$}-}（聚氯乙烯，单体是氯乙烯 \ce{CH2=CHCl}，不是氯乙烷——2025年广东卷第5题A）。\\
\textbf{(4) 乙炔的实验室制法}：\ce{CaC2 + 2H2O -> Ca(OH)2 + CH#CH ^}；用饱和食盐水代替水以减缓反应速率；用 \ce{CuSO4} 溶液除去 \ce{H2S} 等杂质。

\subsection{苯和苯的同系物}
\textbf{苯}：6 个碳原子均为 $\mathrm{sp}^2$ 杂化，形成大 $\pi$ 键，6 个碳碳键完全相同。苯不能使溴水（化学反应）和酸性 \ce{KMnO4} 溶液褪色；易取代、难加成、难氧化：
\begin{itemize}
\item 溴代：\ce{C6H6 + Br2 ->[\ce{FeBr3}] C6H5Br + HBr}（用液溴，不用溴水；产物溴苯因溶有 \ce{Br2} 呈褐色，用 NaOH 溶液洗涤后分液）。
\item 硝化：\ce{C6H6 + HNO3 ->[\text{浓硫酸}][\text{50～60 ℃}] C6H5NO2 + H2O}（水浴加热；浓硫酸作催化剂和吸水剂）。
\item 加成：\ce{C6H6 + 3H2 ->[\ce{Ni}][\text{△}] C6H12}（环己烷）。
\end{itemize}

\textbf{苯的同系物}（如甲苯）——\textbf{基团间相互影响}：①苯环使侧链活化：与苯环直接相连的碳上有 H 时，侧链能被酸性 \ce{KMnO4} 溶液氧化为 \ce{-COOH}（如甲苯$\to$苯甲酸，2025年广东卷第17题(1)现象为“溶液由紫色变为无色”）；②侧链使苯环活化：甲苯硝化在邻、对位发生，可生成 2,4,6-三硝基甲苯（TNT）。

\begin{yicuo}
\begin{enumerate}
\item 苯与溴水混合振荡，溴水褪色是\textbf{萃取}（物理变化），不是加成。
\item 苯的同系物侧链被 \ce{KMnO4} 氧化的前提是“与苯环直接相连的碳原子上有氢原子”，如叔丁基苯不能被氧化。
\item 光照下甲苯与 \ce{Cl2} 取代侧链上的 H；\ce{FeCl3} 催化下取代苯环上的 H——条件不同，位置不同。
\end{enumerate}
\end{yicuo}

\begin{center}
\small
\begin{tabularx}{\linewidth}{p{1.8cm}p{1.9cm}p{2.0cm}X}
\toprule
烃 & 溴水（或 \ce{Br2}/\ce{CCl4}） & 酸性 \ce{KMnO4} & 特征反应 \\
\midrule
烷烃 & 不褪色 & 不褪色 & 光照取代 \\
烯烃、炔烃 & 褪色（加成） & 褪色（氧化） & 加成、氧化、加聚 \\
苯 & 不反应（萃取溴水中的溴，水层褪色；与 \ce{Br2}/\ce{CCl4} 互溶，不褪色） & 不褪色 & \ce{FeBr3} 催化溴代、硝化、加氢 \\
甲苯等 & 不反应（同苯，只发生萃取） & 褪色（侧链氧化） & 侧链氧化；苯环邻、对位取代；光照侧链取代 \\
\bottomrule
\end{tabularx}
\end{center}

\begin{fangfa}[烃的燃烧规律]
\ce{C_xH_y + ($x$ + $\frac y4$)O2 -> $x$CO2 + $\frac y2$H2O}。①等物质的量的烃完全燃烧，耗氧量看 $x+\frac y4$；②等质量的烃完全燃烧，H 的质量分数越大，耗氧越多（\ce{CH4} 最多）；③120 ℃ 以上（水为气体）燃烧前后气体体积不变时，$y=4$（如 \ce{CH4}、\ce{C2H4}、\ce{C3H4}）。
\end{fangfa}

% ---------------- 章末 ----------------
\zhangmo{本章小结}
\subsection*{一、知识总结}
\begin{center}
\small
\begin{tabularx}{\linewidth}{p{2.2cm}X}
\toprule
结构 & 饱和碳 $\mathrm{sp}^3$、双键碳 $\mathrm{sp}^2$、三键碳 sp；甲烷、乙烯、乙炔、苯四种空间模型 \\
同分异构 & 碳架、位置、官能团、顺反、对映；等效氢、基元法、定一移一、苯环取代 3/6/10 \\
研究方法 & 蒸馏、重结晶、萃取、色谱；质谱（相对分子质量）、红外（官能团）、核磁（氢的种类与数目比） \\
烷烃 & 光照取代 \\
烯、炔 & 加成（\ce{Br2}、\ce{H2}、HX、\ce{H2O}）、氧化（\ce{KMnO4} 褪色）、加聚 \\
芳香烃 & 苯：溴代（\ce{FeBr3}）、硝化、加成；甲苯：侧链氧化、邻对位取代 \\
\bottomrule
\end{tabularx}
\end{center}
\subsection*{二、方法提升}
\begin{enumerate}
\item 同分异构体书写“三步走”：先算不饱和度、定类别（官能团），再定碳架，最后移位置；用对称性去重。
\item 用核磁共振氢谱“峰组数$+$峰面积比”验证结构；限定条件题先把“必须有的结构”拿出来，剩余部分再组合。
\end{enumerate}
\subsection*{三、易错点归纳}
\begin{enumerate}
\item 认为含 \ce{-CH3} 的分子中所有原子可能共面（错）。
\item 苯环上二元取代计数遗漏或重复；“两同一异”记成 10 种。
\item 写加聚产物时把双键碳以外的原子放进主链。
\item 苯与溴水“褪色”误认为加成反应。
\end{enumerate}

\zhangmo{分层练习}
\zu{A组\quad 基础巩固}
\begin{xiti}[start=1]
\item 下列有机物分子中所有原子一定共平面的是\xz
\xxb{\ce{CH3CH=CH2}}{苯乙烯}{\ce{CH2=CH-CH=CH2}}{\ce{C6H5Cl}（氯苯）}

\item 分子式为 \ce{C4H8O2} 且属于酯的同分异构体有\tk{}种；分子式为 \ce{C5H10O2} 且能与 \ce{NaHCO3} 反应放出气体的有\tk{}种。

\item 某烃的相对分子质量为 72，其核磁共振氢谱只有 1 组峰，则其结构简式为\tk，系统命名为\tk。

\item 下列关于苯及其同系物的说法正确的是\xz
\xxa{苯能使溴水褪色，发生了加成反应}{甲苯能使酸性 \ce{KMnO4} 溶液褪色，说明苯环对甲基有影响}{苯与液溴在光照条件下生成溴苯}{苯的硝化反应中，浓硫酸只作吸水剂}

\item 用系统命名法命名：(1) \ce{CH2=C(CH3)-CH=CH2}\tk；(2) \ce{CH3CH(C2H5)CH2CH3}\tk；(3) \ce{CH#C-CH(CH3)2}\tk。

\item 等质量的 \ce{CH4}、\ce{C2H4}、\ce{C2H2}、\ce{C6H6} 完全燃烧，消耗 \ce{O2} 最多的是\tk；等物质的量的上述四种烃完全燃烧，消耗 \ce{O2} 最多的是\tk。
\end{xiti}

\zu{B组\quad 高考典型}
\begin{xiti}
\item \laiyuan{2024年广东卷第20题(2)，来源：仓库所附真题文档}\ 化合物Ⅱ为苯乙酮（\ce{C6H5COCH3}），其中含氧官能团的名称是\tk。化合物Ⅱ的某同分异构体含有苯环，在核磁共振氢谱上只有 4 组峰，且能够发生银镜反应，其结构简式为\tk。

\item \laiyuan{原创}\ 某烯烃 A（\ce{C6H12}）被酸性 \ce{KMnO4} 溶液氧化，只得到一种产物——丙酮（\ce{CH3COCH3}）。(1) 写出 A 的结构简式并命名。(2) A 与 HBr 加成的产物有几种？(3) A 是否存在顺反异构？

\item \laiyuan{原创}\ 分子式为 \ce{C8H10} 的芳香烃共有\tk{}种；其中苯环上的一氯代物只有一种的是\tk（写结构简式）；被酸性 \ce{KMnO4} 溶液氧化后生成苯甲酸的是\tk（写名称）。

\item \laiyuan{原创}\ 分子式为 \ce{C5H10} 的链状烯烃，不考虑顺反异构时有\tk{}种；考虑顺反异构时有\tk{}种。其中与 \ce{H2} 加成后得到 2-甲基丁烷的有\tk{}种。

\item \laiyuan{原创}\ 分子式为 \ce{C9H12} 的芳香烃共有\tk{}种；其中苯环上的一溴代物只有一种的结构简式为\tk。
\end{xiti}

\zu{C组\quad 冲刺高分}
\begin{xiti}
\item \laiyuan{原创}\ 某有机物 X 只含 C、H、O 三种元素。取 1.48 g X 完全燃烧，生成 3.52 g \ce{CO2} 和 1.80 g \ce{H2O}；质谱显示最大质荷比为 74；红外光谱显示有 O—H 键；核磁共振氢谱有 2 组峰，面积比为 9:1。(1) 求 X 的分子式。(2) 写出 X 的结构简式并命名。(3) 写出与 X 同类别（醇）的同分异构体中，能被催化氧化为醛的结构简式。

\item \laiyuan{原创}\ 0.1 mol 烃 A 完全燃烧生成 0.6 mol \ce{CO2} 和 0.5 mol \ce{H2O}；A 能使溴的 \ce{CCl4} 溶液褪色，1 mol A 最多与 1 mol \ce{H2} 加成；A 被酸性 \ce{KMnO4} 溶液氧化只得到己二酸 \ce{HOOC(CH2)4COOH}。(1) 求 A 的分子式并计算不饱和度。(2) 写出 A 的结构简式（名称）。(3) 写出 A 与 \ce{Br2} 反应的化学方程式。(4) A 的同分异构体中，属于链状二烯烃且核磁共振氢谱只有 2 组峰的结构简式为\tk。
\end{xiti}

\zhangmo{练习解析}
\begin{jiexilist}
\item \ans{D}A 含 \ce{-CH3}，不可能全部共面；B、C 中苯环与乙烯基、两个乙烯基之间以单键相连，可旋转，“可能”共面但不一定；D 中 Cl 取代苯环上的 H，仍在苯环平面内，一定共面。
\item 4；4。酯：甲酸正丙酯、甲酸异丙酯、乙酸乙酯、丙酸甲酯。羧酸 \ce{C4H9COOH}：丁基 4 种。
\item \ce{C(CH3)4}；2,2-二甲基丙烷。$72=14\times5+2$，符合 \ce{C_nH_{2n+2}}（$n=5$），为 \ce{C5H12}；只有一种等效氢，为新戊烷。
\item \ans{B}A：萃取；B：甲基本身不能被 \ce{KMnO4} 氧化，在苯环影响下才被氧化为羧基，正确；C：需 \ce{FeBr3} 催化，不是光照；D：浓硫酸作催化剂和吸水剂。
\item (1) 2-甲基-1,3-丁二烯（异戊二烯）；(2) 3-甲基戊烷（主链应选含乙基的 5 个碳，不能命名为“2-乙基丁烷”）；(3) 3-甲基-1-丁炔。
\item \ce{CH4}；\ce{C6H6}。等质量时含 H 质量分数越大耗氧越多（\ce{CH4} 最大）；等物质的量时看 $x+\frac y4$：\ce{CH4} 2、\ce{C2H4} 3、\ce{C2H2} 2.5、\ce{C6H6} 7.5。
\item 酮羰基；对甲基苯甲醛（\ce{CH3C6H4CHO}，对位）。能发生银镜反应说明含 \ce{-CHO}；\ce{C8H8O} 除苯环和醛基外剩 1 个 C，为甲基；4 组峰要求对称性高，甲基与醛基处于对位（氢的种类：\ce{-CH3}、苯环上两种、\ce{-CHO}）。
\item (1) \ce{(CH3)2C=C(CH3)2}，2,3-二甲基-2-丁烯。氧化只得丙酮，说明双键两端各连两个甲基。(2) 1 种。分子结构对称，与 HBr 加成只得一种产物 \ce{(CH3)2CBr-CH(CH3)2}。(3) 不存在，双键碳上各连两个相同的甲基。
\item 4；\ce{CH3-C6H4-CH3}（对二甲苯）；乙苯。\ce{C8H10} 芳香烃有乙苯和邻、间、对二甲苯 4 种。对二甲苯苯环上的 4 个 H 等效，一氯代物只有 1 种。侧链被酸性 \ce{KMnO4} 氧化时，无论侧链多长，与苯环相连的碳都变为 \ce{-COOH}：乙苯（1 个侧链）生成苯甲酸，二甲苯（2 个侧链）生成苯二甲酸。
\item 5；6；3。碳架 C—C—C—C—C：1-戊烯、2-戊烯（有顺反异构）；碳架 C—C(C)—C—C：2-甲基-1-丁烯、3-甲基-1-丁烯、2-甲基-2-丁烯；新戊烷碳架不能形成双键。加氢得 2-甲基丁烷的是带支链的 3 种。
\item 8；1,3,5-三甲苯（均三甲苯），结构简式如下：\par{\centering\setchemfig{atom sep=1.5em}\chemfig{*6(-(-CH_3)=-(-CH_3)=-(-CH_3)=)}\par}一个侧链：丙苯、异丙苯（2 种）；两个侧链（甲基、乙基）：邻、间、对（3 种）；三个甲基：连三甲苯、偏三甲苯、均三甲苯（3 种）。均三甲苯高度对称，苯环上 3 个 H 等效。
\item (1) $n(\ce{C})=3.52/44=0.08$ mol，$n(\ce{H})=2\times1.80/18=0.20$ mol，$m(\ce{O})=1.48-0.96-0.20=0.32$ g，$n(\ce{O})=0.02$ mol；C:H:O$=4:10:1$，实验式 \ce{C4H10O}，式量 74，分子式 \ce{C4H10O}。(2) 有 O—H 键，为醇；2 组峰面积比 9:1，为 \ce{(CH3)3COH}，2-甲基-2-丙醇。(3) 能氧化为醛的醇含 \ce{-CH2OH}：\ce{CH3CH2CH2CH2OH}、\ce{(CH3)2CHCH2OH}。
\item (1) \ce{C6H10}，$\Omega=\frac{2\times6+2-10}{2}=2$。(2) 只与 1 mol \ce{H2} 加成，说明只有 1 个碳碳双键，另 1 个不饱和度为环；氧化只得到直链二元酸，双键在六元环中：环己烯。(3) 加成反应，生成 1,2-二溴环己烷：\par{\centering\setchemfig{atom sep=1.5em}\chemfig{*6(-=----)}\quad\ce{+ Br2 ->}\quad\chemfig{*6(-(-Br)-(-Br)----)}\par}(4) \ce{CH2=C(CH3)-C(CH3)=CH2}（2,3-二甲基-1,3-丁二烯，\ce{CH2} 与 \ce{CH3} 两种 H，面积比 4:6）。
\end{jiexilist}
% <<<<<<<< 结束：chapters/ch09.tex

% >>>>>>>> 开始：chapters/ch10.tex
% ============================================================
% 第10章 烃的衍生物（必修二第七章、选必3第三章）
% ============================================================
\chapter{烃的衍生物}

\begin{daoxue}
\textbf{知识主线：}官能团决定性质。卤代烃（水解、消去）$\leftrightarrow$ 醇（取代、消去、氧化、酯化）$\to$ 醛（氧化、还原、加成）$\to$ 羧酸（酸性、酯化）$\to$ 酯（水解）；酚是“羟基$+$苯环”相互影响的典型；胺与酰胺是含氮衍生物。掌握这条“转化链”，有机推断就有了骨架。

\textbf{高考考情：}广东卷每年第20题的反应类型、方程式书写、合成路线都建立在本章官能团转化之上；选择题考陌生物质的官能团性质（2024年第9题）、方程式正误（2026年第14题D 乙醛被氢氧化铜氧化）、实验（2025年第11题乙醇催化氧化与银镜反应、2026年第12题卤代烃的消去与水解）。

\textbf{学习目标：}①掌握各类官能团的特征反应及条件；②理解基团间相互影响；③会进行官能团的定量计算（消耗 \ce{H2}、NaOH、\ce{Br2} 等）。
\end{daoxue}

\section{卤代烃}
C—X 键极性较强，易断裂。卤代烃一般难溶于水，密度多比水大（一氯代烷例外）。两个核心反应条件不同、产物不同：

\begin{center}
\small
\begin{tabularx}{\linewidth}{p{1.6cm}XX}
\toprule
& 水解（取代）反应 & 消去反应 \\
\midrule
条件 & NaOH \textbf{水溶液}、加热 & NaOH \textbf{醇溶液}、加热 \\
实例 & \ce{CH3CH2Br + NaOH ->[\ce{H2O}][\text{△}] CH3CH2OH + NaBr} & \ce{CH3CH2Br + NaOH ->[\text{乙醇}][\text{△}] CH2=CH2 ^ + NaBr + H2O} \\
结构要求 & 所有卤代烃都能水解 & 与 X 相连碳的\textbf{相邻碳上必须有 H}（β-H），如 \ce{CH3Br}、\ce{(CH3)3CCH2Br} 不能消去 \\
\bottomrule
\end{tabularx}
\end{center}

\begin{yicuo}
检验卤代烃中的卤素：水解（或消去）后，必须\textbf{先加稀硝酸酸化}中和过量的 NaOH，再加 \ce{AgNO3} 溶液观察沉淀颜色；否则 \ce{Ag+} 与 \ce{OH-} 反应生成褐色 \ce{Ag2O} 沉淀干扰判断。卤代烃本身不能电离出 \ce{X-}，直接加 \ce{AgNO3} 不产生沉淀。
\end{yicuo}

\begin{zhenti}{2026年广东卷第12题\quad 来源：仓库所附真题文档（网络收集版）}
用如图装置探究卤代烃的性质：烧瓶Ⅰ中加入物质 a 和 NaOH 溶液（溶剂为 b）及碎瓷片加热，产生的气体依次通过Ⅱ（盛水，可注入酸性 \ce{K2Cr2O7} 溶液）、Ⅲ（用酸性 \ce{KMnO4} 溶液浸湿的石英棉）、Ⅳ（溴的 \ce{CCl4} 溶液）。对应叙述正确的是：A．a 为 2-溴丁烷、b 为乙醇：Ⅰ中发生消去反应，Ⅳ中溶液褪色，发生取代反应；B．a 为 2-溴丁烷、b 为水：Ⅰ中生成 2-丁醇，Ⅲ中石英棉和Ⅳ中溶液均褪色；C．a 为溴乙烷、b 为水：Ⅱ中注入少量 \ce{K2Cr2O7} 溶液后，溶液颜色无明显变化；D．a 为溴乙烷、b 为乙醇：冷却后向烧瓶中加过量稀硝酸，再加 \ce{AgNO3} 溶液，生成沉淀。

\textbf{解析：}A：消去生成丁烯，使溴的 \ce{CCl4} 溶液褪色的是\textbf{加成}反应，错。B：水解生成 2-丁醇，Ⅳ中溴的 \ce{CCl4} 溶液不会因醇而褪色，错。C：水解生成乙醇，乙醇蒸气溶于Ⅱ的水中，能被酸性 \ce{K2Cr2O7} 氧化，溶液由橙色变为绿色，错。D：无论发生消去还是水解，都生成 \ce{Br-}，酸化后加 \ce{AgNO3} 生成淡黄色 AgBr 沉淀，正确。答案：D。
\end{zhenti}

\section{醇和酚}

\subsection{醇}
醇分子间能形成氢键，沸点比相对分子质量相近的烃高得多；低级醇与水互溶。以乙醇为例，按断键位置理解反应：

\begin{center}
\small
\begin{tabularx}{\linewidth}{p{2.4cm}Xp{3.2cm}}
\toprule
反应 & 方程式（条件） & 断键与结构要求 \\
\midrule
与活泼金属 & \ce{2CH3CH2OH + 2Na -> 2CH3CH2ONa + H2 ^} & O—H 键；比水与 Na 反应缓和 \\
与氢卤酸取代 & \ce{CH3CH2OH + HBr ->[\text{△}] CH3CH2Br + H2O} & C—O 键 \\
分子内脱水（消去） & \ce{CH3CH2OH ->[\text{浓硫酸}][170\ \text{℃}] CH2=CH2 ^ + H2O} & 需 β-H；温度迅速升至 170 ℃ \\
分子间脱水（取代） & \ce{2CH3CH2OH ->[\text{浓硫酸}][140\ \text{℃}] CH3CH2OCH2CH3 + H2O} & 生成醚 \\
催化氧化 & \ce{2CH3CH2OH + O2 ->[\ce{Cu}][\text{△}] 2CH3CHO + 2H2O} & 需 α-H（羟基所连碳上有 H） \\
酯化 & 见10.4节 & O—H 键 \\
\bottomrule
\end{tabularx}
\end{center}

\begin{fangfa}[醇的催化氧化规律]
羟基所连碳上有 \textbf{2～3 个 H}（\ce{-CH2OH}）$\to$ 醛；有 \textbf{1 个 H}（\ce{-CH(OH)-}）$\to$ 酮；\textbf{没有 H}（叔醇）$\to$ 不能催化氧化。用铜丝做实验时，红色铜丝在火焰上变黑（\ce{2Cu + O2 ->[\text{△}] 2CuO}），插入乙醇后又变红（\ce{CuO + CH3CH2OH ->[\text{△}] Cu + CH3CHO + H2O}），说明铜是催化剂、乙醇被氧化（2025年广东卷第11题C）。醇还能被酸性 \ce{KMnO4}、酸性 \ce{K2Cr2O7} 溶液氧化（后者由橙色变绿色，可用于检测酒驾）。
\end{fangfa}

\subsection{苯酚}
羟基直接连在苯环上的化合物叫酚。苯酚无色晶体，露置空气中因被氧化而呈粉红色；常温下在水中溶解度不大，65 ℃ 以上与水互溶；有毒，沾到皮肤上用乙醇洗。

\begin{benzhi}[苯环与羟基的相互影响]
\textbf{苯环使羟基上的氢更活泼}——苯酚有弱酸性（俗称石炭酸），能与 NaOH 反应，而乙醇不能。但苯酚酸性很弱，不能使石蕊变色，酸性：\ce{H2CO3}$>$苯酚$>$\ce{HCO3-}，所以：
\[\ce{C6H5ONa + CO2 + H2O -> C6H5OH + NaHCO3}\quad(\text{无论 \ce{CO2} 多少都生成 \ce{NaHCO3}})\]
苯酚能与 \ce{Na2CO3} 反应生成 \ce{NaHCO3}，但不能与 \ce{NaHCO3} 反应。\\
\textbf{羟基使苯环上邻、对位的氢更活泼}——苯酚与\textbf{浓溴水}在常温下即发生取代，生成 2,4,6-三溴苯酚白色沉淀（苯与溴反应却需液溴和催化剂）：
\begin{center}
\chemfig{*6(-=-(-OH)=-=)} \ \ce{+ 3Br2 ->} \ \chemfig{*6((-Br)-=(-Br)-(-OH)=(-Br)-=)} \ce{v + 3HBr}
\end{center}
\end{benzhi}

其他性质：\textbf{显色反应}——遇 \ce{FeCl3} 溶液显紫色（检验酚羟基）；\textbf{缩聚反应}——与甲醛在酸或碱催化下生成酚醛树脂（2024年广东卷第13题D：合成酚醛树脂是\textbf{缩聚}反应，不是加聚反应）。

\section{醛和酮}

\textbf{醛基（\ce{-CHO}）既有氧化性又有还原性}：
\begin{enumerate}
\item \textbf{银镜反应}（水浴加热）：\ce{CH3CHO + 2[Ag(NH3)2]OH ->[\text{△}] CH3COONH4 + 2Ag v + 3NH3 + H2O}。配制银氨溶液：向 \ce{AgNO3} 溶液中逐滴加入稀氨水至最初产生的沉淀\textbf{恰好溶解}。
\item \textbf{与新制 \ce{Cu(OH)2} 反应}（碱过量、加热至沸腾）：\ce{CH3CHO + 2Cu(OH)2 + NaOH ->[\text{△}] CH3COONa + Cu2O v + 3H2O}，生成砖红色沉淀；离子方程式为 \ce{CH3CHO + 2Cu(OH)2 + OH- ->[\text{△}] CH3COO- + Cu2O v + 3H2O}（2026年广东卷第14题D）。
\item \textbf{还原（加氢）}：\ce{CH3CHO + H2 ->[\ce{Ni}][\text{△}] CH3CH2OH}；有机反应中“加氢”或“去氧”都叫还原反应。
\item \textbf{与 HCN 加成}（增长碳链）：\ce{CH3CHO + HCN -> CH3CH(OH)CN}，所得\textbf{羟基腈}水解可得 α-羟基酸（2026年广东卷第20题(5)即以“酮 $\to$ 羟基腈”为合成步骤）。
\end{enumerate}
\textbf{甲醛}（\ce{HCHO}）相当于含两个醛基：1 mol HCHO 与足量银氨溶液反应生成 4 mol Ag。\textbf{酮}（如丙酮）中羰基两端都连烃基，不能发生银镜反应，但能加氢还原为醇、能与 HCN 加成。醛和酮羰基邻位的 C—H（α-H）较活泼，可被卤素取代（2025年广东卷第20题(4)：丙酮与 \ce{Br2} 取代生成 \ce{CH3COCH2Br}）。

\begin{yicuo}
醛基能被溴水、酸性 \ce{KMnO4} 溶液氧化而使其褪色。所以检验 \ce{CH2=CHCHO} 中的碳碳双键时，要先用银氨溶液或新制 \ce{Cu(OH)2} 把醛基氧化，酸化后再加溴水。
\end{yicuo}

\begin{liti}[官能团检验的顺序]
设计实验检验丙烯醛（\ce{CH2=CHCHO}）中的两种官能团，写出操作要点和化学方程式。
\tcblower
\jie ①检验醛基：取少量样品，加入足量新制 \ce{Cu(OH)2} 悬浊液（碱性），加热至沸腾，产生砖红色沉淀：\ce{CH2=CHCHO + 2Cu(OH)2 + NaOH ->[\text{△}] CH2=CHCOONa + Cu2O v + 3H2O}。②检验碳碳双键：取①反应后的上层清液，\textbf{先加稀硫酸酸化}（中和 NaOH），再加溴水，溴水褪色：\ce{CH2=CHCOOH + Br2 -> CH2BrCHBrCOOH}。\\
\xianjing 不能直接用溴水或酸性 \ce{KMnO4} 检验双键，因为醛基也会使它们褪色；不酸化时溴水会与 NaOH 反应而褪色，造成误判。
\end{liti}

\section{羧酸及其衍生物}

\subsection{羧酸与酯化反应}
羧酸具有酸的通性，乙酸酸性强于碳酸：\ce{2CH3COOH + Na2CO3 -> 2CH3COONa + CO2 ^ + H2O}。\textbf{只有羧基能与 \ce{NaHCO3} 反应放出 \ce{CO2}}，这是区分羧基与酚羟基的常用方法。

\begin{gainian}[酯化反应]
\ce{CH3COOH + CH3CH2{}^{18}OH <=>[\text{浓硫酸}][\text{△}] CH3CO{}^{18}OCH2CH3 + H2O}\\
机理：\textbf{酸脱羟基、醇脱氢}（用 \ce{^{18}O} 标记乙醇，\ce{^{18}O} 出现在酯中而不在水中）。反应可逆；浓硫酸作催化剂和吸水剂。
\end{gainian}

\textbf{实验要点}：加试剂顺序为乙醇$\to$浓硫酸$\to$乙酸；加碎瓷片防暴沸；用\textbf{饱和 \ce{Na2CO3} 溶液}收集——中和乙酸、溶解乙醇、降低乙酸乙酯的溶解度便于分层；导管口不能伸入液面以下，防倒吸；产品在上层，分液得到。

\subsection{酯、油脂、胺与酰胺}
\textbf{酯的水解}：酸性条件下可逆：\ce{CH3COOC2H5 + H2O <=>[\ce{H+}][\text{△}] CH3COOH + C2H5OH}；碱性条件下彻底（生成羧酸盐）：\ce{CH3COOC2H5 + NaOH ->[\text{△}] CH3COONa + C2H5OH}。注意\textbf{酚酯}水解生成的酚还能继续消耗 NaOH：1 mol \ce{CH3COOC6H5} 最多消耗 2 mol NaOH。

\textbf{油脂}是高级脂肪酸的甘油酯。油脂在碱性条件下的水解叫\textbf{皂化反应}（生成高级脂肪酸钠，即肥皂的主要成分，和甘油）；液态油含较多不饱和键，经催化加氢（油脂的氢化、硬化）变为固态脂肪。油脂不是高分子，也不是芳香烃（2024年广东卷第4题B）。

\textbf{胺}（如 \ce{CH3NH2}、苯胺 \ce{C6H5NH2}）显碱性，能与酸反应生成盐：\ce{C6H5NH2 + HCl -> C6H5NH3Cl}。\textbf{酰胺}（含 \ce{-CO-NH-}）在酸或碱存在并加热时水解：\ce{RCONH2 + H2O + HCl ->[\text{△}] RCOOH + NH4Cl}，\ce{RCONH2 + NaOH ->[\text{△}] RCOONa + NH3 ^}。

\section{官能团转化网络}
\begin{center}
\begin{tikzpicture}[font=\small,node distance=1.0cm and 1.6cm,
  box/.style={draw,rounded corners=2pt,fill=lightblue,minimum width=1.4cm,minimum height=0.65cm},
  lab/.style={font=\scriptsize,fill=white,inner sep=1pt}]
\node[box] (alkene) at (0,2.2) {烯烃};
\node[box] (rx) at (0,0) {卤代烃};
\node[box] (alc) at (3.6,0) {醇};
\node[box] (ald) at (7.0,0) {醛};
\node[box] (acid) at (10.4,0) {羧酸};
\node[box] (ester) at (7.0,-2.0) {酯};
\node[box] (alkane) at (-2.9,0) {烷烃};
\draw[-Stealth,thick] (alkane)--node[lab,above]{\ce{X2}，光} (rx);
\draw[-Stealth,thick] ([xshift=-2mm]alkene.south)--node[lab,left]{HX 加成} ([xshift=-2mm]rx.north);
\draw[-Stealth,thick] ([xshift=2mm]rx.north)--node[lab,right]{NaOH/醇，△} ([xshift=2mm]alkene.south);
\draw[-Stealth,thick] ([yshift=1.5mm]rx.east)--node[lab,above]{NaOH/\ce{H2O}，△} ([yshift=1.5mm]alc.west);
\draw[-Stealth,thick] ([yshift=-1.5mm]alc.west)--node[lab,below]{HX，△} ([yshift=-1.5mm]rx.east);
\draw[-Stealth,thick] (alkene.east) to[bend left=12] node[lab,above]{\ce{H2O} 加成} (alc.north);
\draw[-Stealth,thick] ([yshift=1.5mm]alc.east)--node[lab,above]{\ce{O2}/Cu，△} ([yshift=1.5mm]ald.west);
\draw[-Stealth,thick] ([yshift=-1.5mm]ald.west)--node[lab,below]{\ce{H2}/Ni，△} ([yshift=-1.5mm]alc.east);
\draw[-Stealth,thick] (ald)--node[lab,above]{\ce{O2} 或银氨} (acid);
\draw[-Stealth,thick] (acid.south) |- node[lab,pos=0.75,below]{醇，浓硫酸，△（酯化）} (ester.east);
\draw[-Stealth,thick] (ester.west) -| node[lab,pos=0.25,below]{稀硫酸或 NaOH，△（水解）} (alc.south);
\end{tikzpicture}
\end{center}
这张图是有机推断和合成路线设计的“地图”：已知起点和终点，沿箭头找路径；箭头上的条件就是写合成路线时必须标注的条件。

\section{官能团性质的定量关系}
\begin{center}
\small
\begin{tabularx}{\linewidth}{p{1.6cm}X}
\toprule
试剂 & 能反应的官能团（1 mol 官能团消耗试剂的物质的量） \\
\midrule
Na & 醇羟基、酚羟基、羧基（各 1 mol 生成 0.5 mol \ce{H2}） \\
NaOH & 酚羟基 1；羧基 1；普通酯基 1；\textbf{酚酯基 2}；卤代烃 1（若生成酚还需再加 1） \\
\ce{NaHCO3} & 只有羧基 1（放出 1 mol \ce{CO2}） \\
\ce{Na2CO3} & 羧基；酚羟基（生成 \ce{NaHCO3}） \\
\ce{H2} & 碳碳双键 1；碳碳三键 2；苯环 3；醛基 1；酮羰基 1；\textbf{羧基、酯基一般不与 \ce{H2} 加成} \\
\ce{Br2} & 碳碳双键 1；碳碳三键 2；酚羟基邻、对位上的每个 H 1；醛基被溴水氧化 1 \\
\bottomrule
\end{tabularx}
\end{center}

\begin{liti}[多官能团物质的定量计算]
有机物 M 的结构如下。1 mol M 分别与足量的 NaOH 溶液、\ce{H2}（Ni，加热）、浓溴水反应，最多消耗各试剂多少摩尔？
\begin{center}
\chemfig{HO-[:0]*6(-(-CHO)=-=(-O-[:-30]C(=[:-90]O)-[:30]CH=[:-30]CH-[:30]COOH)-=)}
\end{center}
\tcblower
\silu 先圈出官能团：酚羟基、醛基、酚酯基（酯基的 O 直接连在苯环上）、碳碳双键、羧基。\\
\jie NaOH：酚羟基 1 $+$ 羧基 1 $+$ 酚酯基 2（水解生成的酚羟基再中和 1）$=4$ mol。\ce{H2}：苯环 3 $+$ 碳碳双键 1 $+$ 醛基 1 $=5$ mol（羧基、酯基不加氢）。\ce{Br2}：酚羟基的另一个邻位和对位上的 2 个 H 被取代，消耗 2 $+$ 碳碳双键加成 1 $+$ 醛基被溴水氧化 1 $=4$ mol。\\
\xianjing 酚酯基消耗 2 mol NaOH；醛基也能被溴水氧化——这两点最常漏。
\end{liti}

% ---------------- 章末 ----------------
\zhangmo{本章小结}
\subsection*{一、知识总结}
\begin{center}
\small
\begin{tabularx}{\linewidth}{p{1.8cm}X}
\toprule
卤代烃 & NaOH 水溶液水解成醇；NaOH 醇溶液消去成烯（需 β-H）；检验 X 先酸化 \\
醇 & 与 Na、HX；浓硫酸 170 ℃ 消去、140 ℃ 成醚；催化氧化（看 α-H）；酯化 \\
酚 & 弱酸性（\ce{H2CO3}$>$苯酚$>$\ce{HCO3-}）；浓溴水白色沉淀；\ce{FeCl3} 显紫色；与 HCHO 缩聚 \\
醛、酮 & 醛：银镜、新制 \ce{Cu(OH)2}、加氢、加 HCN；HCHO 相当于二元醛；酮不发生银镜 \\
羧酸及衍生物 & 酸性强于碳酸；酯化（酸脱羟基醇脱氢）；酯水解（碱性彻底，酚酯耗 2 NaOH）；胺显碱性，酰胺可水解 \\
\bottomrule
\end{tabularx}
\end{center}
\subsection*{二、方法提升}
\begin{enumerate}
\item \textbf{由条件推反应}：NaOH/水$\to$水解；NaOH/醇$\to$消去；浓硫酸/△$\to$消去、酯化、成醚；Cu 或 Ag/\ce{O2}$\to$醇氧化；银氨或新制 \ce{Cu(OH)2}$\to$醛氧化；\ce{H2}/Ni$\to$加成（还原）；稀硫酸/△$\to$酯或酰胺水解。
\item \textbf{官能团的检验顺序}：互相干扰时先检验（或消除）易被氧化的醛基，再检验碳碳双键。
\end{enumerate}
\subsection*{三、易错点归纳}
\begin{enumerate}
\item 卤代烃水解、消去条件混淆；检验卤素忘记酸化。
\item 苯酚钠溶液通 \ce{CO2} 写成生成 \ce{Na2CO3}。
\item 银镜反应方程式中漏写 \ce{NH3}，或系数配错；配制银氨溶液时把氨水加过量太多。
\item 酯化反应中水的 O 来源弄错；忘记写“\ce{<=>}”。
\item 计算 NaOH 用量时把酚酯基按 1 mol 算。
\end{enumerate}

\zhangmo{分层练习}
\zu{A组\quad 基础巩固}
\begin{xiti}[start=1]
\item 下列卤代烃中，既能发生水解反应，又能发生消去反应且消去产物只有一种的是\xz
\xx{\ce{CH3Cl}}{\ce{(CH3)3CCH2Br}}{\ce{CH3CHBrCH3}}{\ce{CH3CH2CHBrCH3}}

\item 下列醇中，能被催化氧化生成醛的是\xz
\xxb{\ce{(CH3)3COH}}{\ce{CH3CH(OH)CH3}}{\ce{(CH3)2CHCH2OH}}{\ce{(CH3)2C(OH)CH2CH3}}

\item 下列关于苯酚的说法正确的是\xz
\xxa{苯酚的酸性比碳酸强}{向苯酚钠溶液中通入少量 \ce{CO2}，生成苯酚和 \ce{Na2CO3}}{苯酚与浓溴水反应生成白色沉淀，可用于苯酚的定性检验和定量测定}{苯酚和乙醇都能与 NaOH 溶液反应}

\item 实验室制备乙酸乙酯时，下列操作或说法正确的是\xz
\xxa{先加浓硫酸，再加乙醇和乙酸}{用 NaOH 溶液代替饱和 \ce{Na2CO3} 溶液收集产品效果更好}{导气管末端伸入饱和 \ce{Na2CO3} 溶液液面以下，使产品充分吸收}{饱和 \ce{Na2CO3} 溶液的作用之一是降低乙酸乙酯的溶解度}

\item 0.1 mol HCHO 与 \ce{CH3CHO} 的混合物与足量银氨溶液反应，生成 0.3 mol Ag，则混合物中 $n(\ce{HCHO}):n(\ce{CH3CHO})=$\tk。

\item 只用一种试剂即可鉴别乙醇、乙醛、乙酸、甲酸四种无色液体，该试剂是\tk，写出各物质对应的现象。
\end{xiti}

\zu{B组\quad 高考典型}
\begin{xiti}
\item \laiyuan{2025年广东卷第11题，来源：仓库所附真题文档}\ 利用如下装置实验：b 中为 \ce{MnO2}，分液漏斗中为 10\% \ce{H2O2} 溶液；a 中为酸化的淀粉-KI 溶液；c 为加热的铜丝（右端可滴入无水乙醇）；d 中为银氨溶液。打开 \ce{K1}、\ce{K2}，一段时间后，a 中溶液变蓝；关闭 \ce{K1}，打开 \ce{K3}，点燃酒精灯加热数分钟后，滴入无水乙醇。下列说法\textbf{错误}的是\xz
\xxb{a 中现象体现了 \ce{I-} 的还原性}{b 中 \ce{H2O2} 既作氧化剂也作还原剂}{乙醇滴加过程中，c 中的铜丝由黑变红，说明乙醇被氧化}{d 中有银镜反应发生，说明 c 中产物有乙酸}

\item \laiyuan{原创}\ 有机物 N 的结构简式为 \ce{HO-C6H4-CH2-CH=CH-COOCH3}（\ce{-OH} 与侧链在苯环对位）。(1) N 中含有的官能团名称为\tk。(2) 1 mol N 最多与\tk mol NaOH、\tk mol \ce{H2}、\tk mol \ce{Br2}（浓溴水）反应。(3) 写出 N 与足量 NaOH 溶液共热的化学方程式。

\item \laiyuan{2024年广东卷第20题(3)(4)，来源：仓库所附真题文档}\ 将 3D 打印制备的固载铜离子陶瓷催化材料用于化学催化和生物催化一体化技术，实现化合物Ⅲ的绿色合成：苯乙烯（Ⅰ）在固载铜离子催化体系中与 \ce{H2O}、\ce{O2} 转化为苯乙酮（Ⅱ），Ⅱ在有机酶催化体系中被葡萄糖还原为手性醇Ⅲ（\ce{C6H5CH(OH)CH3}，同时葡萄糖转化为葡萄糖酸内酯）。(3) 关于上述物质及转化，下列说法正确的有\tk。A．由化合物Ⅰ到Ⅱ的转化中，有 $\pi$ 键的断裂与形成\quad B．由葡萄糖到葡萄糖酸内酯的转化中，葡萄糖被还原\quad C．葡萄糖易溶于水，是因为其分子中有多个羟基，能与水分子形成氢键\quad D．由化合物Ⅱ到Ⅲ的转化中，存在 C、O 原子杂化方式的改变，有手性碳原子形成。(4) 对化合物Ⅲ分析预测其可能的化学性质：①若苯环转化为饱和六元环，所需反应试剂、条件为\tk，反应类型为\tk；②若发生取代反应生成溴代物，所需反应试剂、条件为\tk，生成的新结构为\tk。

\item \laiyuan{原创}\ 酯 X 的分子式为 \ce{C4H8O2}，在酸性条件下水解得到 A 和 B。A 能发生银镜反应；B 能被催化氧化，但产物不能发生银镜反应。写出 X、A、B 的结构简式及 X 在 NaOH 溶液中水解的化学方程式。
\end{xiti}

\zu{C组\quad 冲刺高分}
\begin{xiti}
\item \laiyuan{原创}\ 以乙烯为原料可合成乙二酸二乙酯：\ce{CH2=CH2 ->[\text{①}] A ->[\text{②}] B ->[\text{③}] C ->[\text{④}] D}，再由 D 与乙醇反应得到产品。已知 A 为 1,2-二溴乙烷，B 为乙二醇，D 为乙二酸。(1) 写出①～④的试剂与条件及反应类型。(2) 写出 D 与乙醇反应生成乙二酸二乙酯的化学方程式。(3) 若用 B 与 D 在一定条件下反应，可能得到哪些产物？（写出环状酯和高分子的结构简式）

\item \laiyuan{原创}\ 咖啡酸存在于咖啡、蔬菜中，结构简式为
\begin{center}\chemfig{HO-[:0]*6(-(-OH)=-(-CH=CH-COOH)=-=)}\end{center}
(1) 咖啡酸的分子式为\tk，含氧官能团名称为\tk。(2) 1 mol 咖啡酸分别与足量 Na、NaOH、\ce{NaHCO3}、浓溴水、\ce{H2}（一定条件）反应，最多消耗的物质的量分别为多少？(3) 咖啡酸与乙醇发生酯化反应的产物分子式为\tk。
\end{xiti}

\zhangmo{练习解析}
\begin{jiexilist}
\item \ans{C}A：\ce{CH3Cl} 无 β-C，不能消去；B：与 Br 相连碳的邻位碳上无 H，不能消去；C：两侧 β-H 等效，消去只得丙烯；D：消去可得 1-丁烯和 2-丁烯两种。
\item \ans{C}A、D 为叔醇，不能催化氧化；B 为仲醇，氧化得丙酮；C 含 \ce{-CH2OH}，氧化得醛。
\item \ans{C}A：苯酚酸性弱于碳酸；B：生成 \ce{NaHCO3}；C 正确；D：乙醇不与 NaOH 反应。
\item \ans{D}A：应先加乙醇，再缓慢加浓硫酸（类似稀释浓硫酸），最后加乙酸；B：NaOH 会使乙酸乙酯水解；C：会倒吸；D 正确。
\item 1:1。设 HCHO $x$ mol、\ce{CH3CHO} $y$ mol：$x+y=0.1$，$4x+2y=0.3$，解得 $x=y=0.05$。
\item 新制 \ce{Cu(OH)2} 悬浊液。乙醇：无明显现象（加热也无砖红色沉淀）；乙醛：常温不溶解沉淀，加热生成砖红色沉淀；乙酸：沉淀溶解得蓝色溶液，加热无砖红色沉淀；甲酸：沉淀溶解得蓝色溶液，加热（碱过量时）生成砖红色沉淀（甲酸含醛基结构）。
\item \ans{D}A：\ce{O2} 将 \ce{I-} 氧化为 \ce{I2}，\ce{I-} 体现还原性，对；B：\ce{2H2O2 ->[\ce{MnO2}] 2H2O + O2 ^}，O 元素既升又降，对；C：铜丝先被 \ce{O2} 氧化为黑色 CuO，CuO 被乙醇还原为红色 Cu，乙醇被氧化，对；D：银镜反应说明有\textbf{乙醛}，乙酸不能发生银镜反应，错。
\item (1) 羟基（酚羟基）、碳碳双键、酯基。(2) 2；4；3。NaOH：酚羟基 1 $+$ 酯基 1（醇酯，生成甲醇）$=2$；\ce{H2}：苯环 3 $+$ 碳碳双键 1 $=4$（酯基不加氢）；\ce{Br2}：酚羟基的 2 个邻位 H（对位已被占据）取代 2 $+$ 碳碳双键加成 1 $=3$。(3) \ce{HO-C6H4-CH2CH=CHCOOCH3 + 2NaOH ->[\text{△}] NaO-C6H4-CH2CH=CHCOONa + CH3OH + H2O}。
\item (3) ACD。A：Ⅰ中碳碳双键（$\pi$ 键）断裂，Ⅱ中形成碳氧双键（$\pi$ 键），对；B：葡萄糖被\textbf{氧化}（醛基转化），错；C 对；D：Ⅱ中羰基 C、O 为 $\mathrm{sp}^2$ 杂化，Ⅲ中对应 C、O 为 $\mathrm{sp}^3$ 杂化，且与羟基相连的碳为手性碳，对。(4) ①\ce{H2}，催化剂、加热；加成反应（还原反应）。②HBr，加热；\ce{C6H5CHBrCH3}（羟基被溴原子取代）。
\item A 能发生银镜反应，A 为甲酸 HCOOH；B 为 \ce{C3H7OH}，催化氧化产物不能发生银镜反应（为酮），B 为 2-丙醇 \ce{CH3CH(OH)CH3}。X 为甲酸异丙酯 \ce{HCOOCH(CH3)2}。\ce{HCOOCH(CH3)2 + NaOH ->[\text{△}] HCOONa + (CH3)2CHOH}。
\item (1) ①\ce{Br2} 的 \ce{CCl4} 溶液（或溴水），加成反应；②NaOH 水溶液、加热，水解（取代）反应；③\ce{O2}、Cu（或 Ag）、加热，氧化反应，C 为乙二醛 \ce{OHC-CHO}；④\ce{O2}、催化剂、加热（或银氨溶液/新制 \ce{Cu(OH)2} 氧化后酸化），氧化反应。(2) \ce{HOOC-COOH + 2CH3CH2OH <=>[\text{浓硫酸}][\text{△}] CH3CH2OOC-COOCH2CH3 + 2H2O}。(3) 环状酯：乙二酸乙二酯（六元环，\ce{-OOC-COO-CH2CH2-} 首尾相连成环）；高分子：\ce{H-[OCH2CH2OOC-CO]_{$n$}-OH}（缩聚，同时生成 $(2n-1)$ \ce{H2O}）；还可能生成链状的单酯 \ce{HOCH2CH2OOC-COOH}。
\item (1) \ce{C9H8O4}；（酚）羟基、羧基。(2) Na：2 个酚羟基 $+$ 1 个羧基，3 mol（生成 1.5 mol \ce{H2}）；NaOH：3 mol；\ce{NaHCO3}：只与羧基反应，1 mol；浓溴水：苯环上与酚羟基处于邻、对位的 3 个 H 都可被取代（3 mol）$+$ 碳碳双键加成（1 mol），共 4 mol；\ce{H2}：苯环 3 mol $+$ 碳碳双键 1 mol，共 4 mol（羧基不加氢）。(3) \ce{C11H12O4}（\ce{C9H8O4 + C2H6O - H2O}）。
\end{jiexilist}
% <<<<<<<< 结束：chapters/ch10.tex

% >>>>>>>> 开始：chapters/ch11.tex
% ============================================================
% 第11章 生物大分子、合成高分子与有机合成（选必3第四、五章，第三章第五节）
% ============================================================
\chapter{生物大分子、合成高分子与有机合成}

\begin{daoxue}
\textbf{知识主线：}生命中的大分子（糖类、蛋白质、核酸）$\to$ 人工合成的大分子（加聚、缩聚）$\to$ 把前面所有官能团转化串起来——有机合成（碳骨架构建、官能团引入与保护、逆合成分析）$\to$ 高考有机大题的解题策略。

\textbf{高考考情：}选择题常以生活、文物情境考糖类、油脂、蛋白质（2024年第1、4题，2025年第3、4题，2026年第10题），以材料情境考高分子（2025年第5题、2026年第14题C）；第20题后两问几乎每年都是“设计合成路线”或“逆推合成步骤”（2024年第20题(5)(6)、2025年第20题(5)(6)、2026年第20题(5)）。

\textbf{学习目标：}①掌握糖类、蛋白质的组成与性质；②会写加聚、缩聚反应方程式并由高分子推单体；③掌握碳链增长、官能团保护等合成策略；④能用逆合成分析法设计 3～4 步合成路线。
\end{daoxue}

\section{生物大分子}

\subsection{糖类}
\begin{center}
\small
\begin{tabularx}{\linewidth}{p{1.4cm}p{2.6cm}X}
\toprule
类别 & 代表物 & 结构与性质 \\
\midrule
单糖 & 葡萄糖、果糖（\ce{C6H12O6}，互为同分异构体） & 葡萄糖为多羟基醛，能发生银镜反应、与新制 \ce{Cu(OH)2} 反应；果糖为多羟基酮。不能水解 \\
二糖 & 蔗糖、麦芽糖（\ce{C12H22O11}） & 1 mol 蔗糖水解得 1 mol 葡萄糖和 1 mol 果糖；麦芽糖水解得 2 mol 葡萄糖。麦芽糖有还原性，蔗糖没有 \\
多糖 & 淀粉、纤维素（\ce{(C6H10O5)_{$n$}}） & 天然高分子；$n$ 值不同，\textbf{不是}同分异构体；水解最终产物都是葡萄糖；淀粉遇碘变蓝 \\
\bottomrule
\end{tabularx}
\end{center}

\begin{fangfa}[检验淀粉水解程度]
取水解液分两份：一份加碘水，检验是否还有淀粉；另一份\textbf{先加 NaOH 溶液中和作催化剂的稀硫酸至碱性}，再加银氨溶液（或新制 \ce{Cu(OH)2}）并加热，检验是否生成葡萄糖。“不变蓝、有银镜”为完全水解；“变蓝、有银镜”为部分水解；“变蓝、无银镜”为未水解。
\end{fangfa}

\subsection{氨基酸、蛋白质与核酸}
\textbf{氨基酸}分子中含 \ce{-NH2}（碱性）和 \ce{-COOH}（酸性），是两性化合物。两个氨基酸分子脱水缩合形成肽键（\ce{-CO-NH-}）。蛋白质是由多种氨基酸缩聚形成的高分子，含 C、H、O、N 等元素（2026年广东卷第10题B：牛肉中的蛋白质含 N、O 元素）。

\begin{center}
\small
\begin{tabularx}{\linewidth}{p{1.6cm}XX}
\toprule
& 盐析 & 变性 \\
\midrule
条件 & 加入浓的轻金属盐或铵盐溶液，如 \ce{(NH4)2SO4}、\ce{Na2SO4}、NaCl & 加热、紫外线、重金属盐、强酸强碱、乙醇、甲醛等 \\
变化 & 溶解度降低而析出，物理变化，\textbf{可逆} & 结构改变、失去生理活性，化学变化，\textbf{不可逆} \\
应用 & 分离、提纯蛋白质；用 \ce{MgCl2} 等点卤制豆腐（2025年广东卷第4题C） & 消毒杀菌；重金属盐中毒可服用牛奶急救 \\
\bottomrule
\end{tabularx}
\end{center}

其他性质：含苯环的蛋白质遇\textbf{浓硝酸}显黄色（显色反应）；灼烧有烧焦羽毛的气味（区分蚕丝与棉、合成纤维）；\textbf{酶}是对生物体内化学反应具有催化作用的有机物，其中绝大多数是蛋白质，具有高效性、专一性，需要温和条件。\textbf{核酸}（DNA、RNA）是由核苷酸（磷酸、戊糖、碱基）聚合成的生物大分子，DNA 双螺旋两条链上的碱基通过氢键配对。

\section{合成高分子}

\subsection{加聚反应与缩聚反应}
\begin{center}
\small
\begin{tabularx}{\linewidth}{p{1.6cm}XX}
\toprule
& 加成聚合（加聚） & 缩合聚合（缩聚） \\
\midrule
单体结构 & 含碳碳双键等不饱和键 & 含两个（或以上）能相互反应的官能团，如 \ce{-OH} 与 \ce{-COOH}、\ce{-NH2} 与 \ce{-COOH} \\
产物 & 只生成高分子，链节组成与单体相同 & 生成高分子的同时有小分子（\ce{H2O} 等）生成 \\
实例 & 聚乙烯、聚氯乙烯、聚苯乙烯、聚甲基丙烯酸甲酯（有机玻璃）、顺丁橡胶 & 聚酯（涤纶）、聚酰胺（尼龙）、酚醛树脂、聚乳酸 \\
\bottomrule
\end{tabularx}
\end{center}

\textbf{缩聚方程式的写法}：①高分子写端基；②小分子的计量数：由一种单体（如羟基酸）缩聚生成 $(n-1)$ 个 \ce{H2O}；由两种单体（如二元酸与二元醇）缩聚生成 $(2n-1)$ 个 \ce{H2O}。例如：
\[\ce{$n$HOOC(CH2)5OH ->[\text{催化剂}][\text{△}] H-[O(CH2)5CO]_{$n$}-OH + ($n$-1)H2O}\]
\[\ce{$n$HOOC-C6H4-COOH + $n$HOCH2CH2OH -> HO-[OC-C6H4-COOCH2CH2O]_{$n$}-H + (2$n$-1)H2O}\]

\begin{zhenti}{2026年广东卷第14题\quad 来源：仓库所附真题文档（网络收集版）}
下列物质性质实验对应的反应方程式书写\textbf{错误}的是：A．\ce{Al2O3} 溶于 NaOH 溶液：\ce{Al2O3 + 2OH- + 3H2O = 2[Al(OH)4]-}；B．铜丝插入热的浓硫酸中：\ce{Cu + 2H2SO4}（浓）\ce{->[\text{△}] CuSO4 + 2H2O + SO2 ^}；C．6-羟基己酸的缩聚反应：\ce{$n$HOOC(CH2)5OH -> HO-[CO(CH2)5O]_{$n$}-H + $n$H2O}；D．乙醛被 \ce{Cu(OH)2} 氧化：\ce{CH3CHO + 2Cu(OH)2 + OH- ->[\text{△}] CH3COO- + Cu2O v + 3H2O}。

\textbf{解析：}一种单体缩聚生成 $n$ 个链节，形成 $(n-1)$ 个酯基，只脱去 $(n-1)$ 个水分子，C 中写“$n$\ce{H2O}”错误。答案：C。
\end{zhenti}

\subsection{由高分子推断单体}
\begin{fangfa}[推单体的方法]
\textbf{加聚产物}：主链全是碳原子时，去掉方括号和 $n$，从链节一端开始“每两个碳原子断开”，单键变双键；若主链中有双键，则以双键为中心的 4 个碳是二烯烃单体（如顺丁橡胶 \ce{-[CH2-CH=CH-CH2]_{$n$}-} 的单体是 1,3-丁二烯）。\textbf{缩聚产物}：找到酯基或酰胺基，从 C—O 键或 C—N 键处断开，羰基一侧补 \ce{-OH}，O、N 一侧补 \ce{-H}。如尼龙-66 \ce{H-[NH(CH2)6NHCO(CH2)4CO]_{$n$}-OH} 的单体是己二胺和己二酸。
\end{fangfa}

\textbf{高分子材料}：塑料分热塑性（线型结构，可反复加热成型，如聚乙烯）和热固性（网状结构，如酚醛树脂）；合成纤维如涤纶、锦纶、腈纶；合成橡胶硫化后形成二硫键交联的网状结构，强度和弹性提高（2024年广东卷第7题C）。功能高分子如高吸水性树脂（含大量亲水基团的网状结构）、离子交换树脂、导电高分子等。

\section{有机合成}

\subsection{碳骨架的构建}
\begin{center}
\small
\begin{tabularx}{\linewidth}{p{2.0cm}X}
\toprule
目的 & 常用反应 \\
\midrule
增长碳链 & ①卤代烃与 NaCN：\ce{CH3CH2Br + NaCN -> CH3CH2CN + NaBr}，再水解：\ce{CH3CH2CN ->[\ce{H+}/\ce{H2O}] CH3CH2COOH}；②醛、酮与 HCN 加成生成羟基腈；③羟醛缩合（含 α-H 的醛在稀碱下）：\ce{2CH3CHO ->[\text{稀 NaOH}] CH3CH(OH)CH2CHO}，加热脱水得 \ce{CH3CH=CHCHO}；④加聚、缩聚 \\
缩短碳链 & 烯烃、炔烃被酸性 \ce{KMnO4} 氧化断裂；苯的同系物侧链氧化为 \ce{-COOH}；羧酸盐脱羧 \\
成环 & 二元醇脱水成环醚；羟基酸、二元酸与二元醇分子内（间）酯化成环酯；二烯烃与烯烃的环加成（第尔斯—阿尔德反应，题目会给出信息） \\
\bottomrule
\end{tabularx}
\end{center}

\subsection{官能团的引入、转化与保护}
\textbf{引入碳碳双键}：醇消去、卤代烃消去、炔烃部分加氢。\textbf{引入卤原子}：烷烃或苯环取代、烯烃加成 HX 或 \ce{X2}、醇与 HX 取代。\textbf{引入羟基}：烯烃与水加成、卤代烃水解、醛酮还原、酯水解。\textbf{引入醛基或羰基}：醇的催化氧化、烯烃氧化。\textbf{引入羧基}：醛氧化、苯的同系物侧链氧化、腈或酯水解。

\begin{benzhi}[为什么要“保护”官能团]
如果分子中有两个官能团都能与某试剂反应，而我们只想让其中一个反应，就要先把另一个“保护”起来，反应后再“脱保护”。例如：氧化甲基为羧基时，酚羟基也会被氧化，可先用 \ce{CH3I} 使其转化为 \ce{-OCH3}（醚），氧化后再用 HI 恢复；需要氧化醛基而保留碳碳双键时，用弱氧化剂（银氨溶液、新制 \ce{Cu(OH)2}）而不用 \ce{KMnO4}；需要氧化羟基而保留碳碳双键时，可先用 HBr 与双键加成，氧化后再用 NaOH 醇溶液消去恢复双键。
\end{benzhi}

\subsection{逆合成分析法}
从目标分子出发，找出最后一步可能由哪些“前体”生成（关注目标分子中官能团和新形成的化学键），再对前体重复这一过程，直到得到题目给定的原料。书写时：箭头上方写试剂，下方写条件；每步都要合理、可行，步数尽量少。

\begin{zhenti}{2026年广东卷第20题(5)\quad 来源：仓库所附真题文档（网络收集版）}
某兴趣小组设计以环己烷为原料，经四步反应合成化合物Ⅲa（分子中不含手性碳原子）的路线，Ⅲa 可经羟基迁移路径合成化合物Ⅶa（环己烷环上相邻两个碳分别连有 \ce{-CN} 和 \ce{-OH}）。①Ⅲa 的结构简式为\tk。②从Ⅲa出发逆推合成步骤：羟基腈 $\Rightarrow$\tk$\Rightarrow$ 醇 $\Rightarrow$\tk$\Rightarrow$ 环烷烃（填有机物类别）。③合成Ⅲa的步骤中，第一步为取代反应，其化学方程式为\tk（注明条件）；第四步所需无机试剂为\tk。④1 mol Ⅶa 可与 2 mol \ce{H2} 发生反应，其产物的结构简式为\tk。

\textbf{解析：}①“羟基迁移”前，\ce{-CN} 与 \ce{-OH} 在同一个碳上，且无手性碳（该碳连有两个相同的环上 \ce{-CH2-}），Ⅲa 为 1-羟基环己基甲腈（环己烷中一个碳同时连 \ce{-OH} 和 \ce{-CN}）。②羟基腈由酮与 HCN 加成得到，酮由醇氧化得到，醇由卤代烃水解得到，卤代烃由环烷烃取代得到：羟基腈 $\Rightarrow$ \textbf{酮} $\Rightarrow$ 醇 $\Rightarrow$ \textbf{卤代烃} $\Rightarrow$ 环烷烃。③\ce{C6H12 + Cl2 ->[\text{光照}] C6H11Cl + HCl}；HCN。④\ce{-C#N} 加 2 mol \ce{H2} 生成 \ce{-CH2NH2}，产物为 2-（氨甲基）环己醇（环上相邻碳分别连 \ce{-CH2NH2} 和 \ce{-OH}）。
\end{zhenti}

\subsection{有机反应类型总览}
\begin{center}
\small
\begin{tabularx}{\linewidth}{p{1.8cm}Xp{4.2cm}}
\toprule
反应类型 & 特点与常见实例 & 典型条件 \\
\midrule
取代反应 & 原子或原子团被替代：烷烃卤代、苯的溴代与硝化、卤代烃水解、醇与 HX、酯化、酯与酰胺水解、α-H 卤代 & 光照；\ce{FeBr3}；NaOH/水，△；浓硫酸，△ \\
加成反应 & 不饱和键断开结合其他原子：烯、炔与 \ce{X2}、HX、\ce{H2O}、\ce{H2}；苯加氢；醛酮加 \ce{H2}、HCN & 催化剂；Ni，△ \\
消去反应 & 生成不饱和键并脱去小分子：醇、卤代烃的消去 & 浓硫酸 170 ℃；NaOH/醇，△ \\
氧化反应 & 加氧或去氢：燃烧；醇催化氧化；醛的银镜反应；烯、苯同系物被 \ce{KMnO4} 氧化 & \ce{O2}/Cu，△；银氨溶液 \\
还原反应 & 加氢或去氧：醛、酮、烯与 \ce{H2} 加成；硝基苯还原为苯胺 & \ce{H2}/Ni；Fe/HCl \\
加聚反应 & 含双键的单体聚合，无小分子生成 & 催化剂 \\
缩聚反应 & 双官能团单体聚合，生成小分子 & 催化剂，△ \\
\bottomrule
\end{tabularx}
\end{center}

\subsection{有机推断题的突破口}
\begin{center}
\small
\begin{tabularx}{\linewidth}{p{4.2cm}X}
\toprule
信息 & 推断 \\
\midrule
能与 \ce{NaHCO3} 反应放气 & 含 \ce{-COOH} \\
遇 \ce{FeCl3} 显紫色 & 含酚羟基 \\
能发生银镜反应 & 含 \ce{-CHO}（醛、甲酸、甲酸酯、葡萄糖等） \\
NaOH/醇、加热；浓硫酸、170 ℃ & 消去反应 \\
NaOH/水、加热；稀硫酸、加热 & 卤代烃或酯水解 \\
连续两步氧化（\ce{A ->[O2] B ->[O2] C}） & 醇$\to$醛$\to$酸，说明含 \ce{-CH2OH} \\
分子式中 H 减少 2、O 不变 & 醇的催化氧化（去氢，生成醛或酮） \\
分子式中减少 1 个 \ce{H2O} & 醇的消去（或成醚、分子内酯化等脱水） \\
不饱和度 $\Omega=\dfrac{2C+2+N-H-X}{2}$ & 苯环 4，双键 1，三键 2，环 1 \\
\bottomrule
\end{tabularx}
\end{center}

\begin{liti}[有机推断完整示例]
以丙烯为原料合成聚乳酸，路线如下：
\[\ce{CH2=CHCH3 ->[\text{溴的 CCl4 溶液}] B ->[\text{NaOH 水溶液}][\text{△}] C ->[\text{O2, Cu}][\text{△}] D}\,(\ce{C3H4O2})\]
\[\ce{D ->[\text{①新制 Cu(OH)2}][\text{②酸化}] E ->[\text{H2, Ni}][\text{△}] F ->[\text{催化剂}] \text{聚乳酸}}\]
(1) 写出 B～F 的结构简式。(2) 指出 B$\to$C、E$\to$F 的反应类型。(3) F 在浓硫酸、加热条件下两分子间脱水可生成六元环状化合物 G，写出 G 的结构简式。
\tcblower
\jie (1) B：\ce{CH2BrCHBrCH3}（加成）；C：\ce{HOCH2CH(OH)CH3}（1,2-丙二醇）；D：\ce{OHC-CO-CH3}（伯醇氧化为醛基、仲醇氧化为酮羰基，分子式 \ce{C3H4O2} 吻合）；E：\ce{CH3COCOOH}（醛基被氧化为羧基，酮羰基不被氧化）；F：\ce{CH3CH(OH)COOH}（酮羰基加氢还原为羟基，即乳酸）。(2) 取代（水解）反应；加成（还原）反应。(3) G 为两分子乳酸的羧基与羟基相互酯化形成的六元环二酯（丙交酯），分子式 \ce{C6H8O4}：环上依次为 —O—CH(CH$_3$)—CO—O—CH(CH$_3$)—CO—。\\
\dianbo 推断时抓住“分子式变化”：C$\to$D 少了 4 个 H，说明两个羟基都被氧化；E$\to$F 多了 2 个 H 且条件为 \ce{H2}/Ni，说明 C=O 被还原。
\end{liti}

\begin{fangfa}[广东卷第20题的作答策略]
①\textbf{先整体后局部}：比较相邻物质的分子式和结构，找出“变了哪里”，判断反应类型；②\textbf{善用题给新反应}：陌生反应（如 2025 年的三组分反应、2026 年的羟基迁移）往往在后面的合成设计中要“照猫画虎”；③\textbf{同分异构限定条件}：先计算不饱和度，再把“一定含有”的结构拿出来；④\textbf{合成路线}：优先用题给反应和课本中最熟悉的转化，条件写全（如“NaOH 水溶液，△”“浓硫酸，△”）。
\end{fangfa}

% ---------------- 章末 ----------------
\zhangmo{本章小结}
\subsection*{一、知识总结}
\begin{center}
\small
\begin{tabularx}{\linewidth}{p{2.2cm}X}
\toprule
糖类 & 葡萄糖（多羟基醛，还原性）、蔗糖（无还原性）、淀粉与纤维素（不是同分异构体） \\
蛋白质 & 氨基酸两性、肽键；盐析可逆、变性不可逆；浓硝酸显黄色 \\
高分子 & 加聚（不饱和单体）；缩聚（双官能团，生成小分子，$(n-1)$ 或 $(2n-1)$）；推单体 \\
有机合成 & 增碳（\ce{-CN}、HCN、羟醛缩合）、减碳（氧化断裂）、成环；官能团引入、转化、保护；逆合成分析 \\
\bottomrule
\end{tabularx}
\end{center}
\subsection*{二、方法提升}
\begin{enumerate}
\item 合成路线设计：先比较原料与目标的\textbf{碳骨架}（碳数、环、支链），再比较\textbf{官能团}的种类和位置，找出需要的转化，最后排定顺序（注意保护）。
\item 写缩聚方程式：先写链节，再补端基，最后数小分子。
\end{enumerate}
\subsection*{三、易错点归纳}
\begin{enumerate}
\item 检验淀粉水解产物时，未先用 NaOH 中和稀硫酸。
\item 把盐析与变性混淆（重金属盐使蛋白质变性）。
\item 缩聚方程式中小分子计量数写成 $n$。
\item 合成路线中漏写条件，或使用了会破坏其他官能团的试剂（如用 \ce{KMnO4} 氧化醛基而破坏碳碳双键）。
\end{enumerate}

\zhangmo{分层练习}
\zu{A组\quad 基础巩固}
\begin{xiti}[start=1]
\item 下列说法正确的是\xz
\xxa{淀粉和纤维素的化学式均为 \ce{(C6H10O5)_{$n$}}，二者互为同分异构体}{蔗糖和麦芽糖水解的最终产物都只有葡萄糖}{向鸡蛋清溶液中加入饱和 \ce{(NH4)2SO4} 溶液，有沉淀析出，加水后沉淀又溶解}{油脂是高分子化合物，在碱性条件下水解称为皂化反应}

\item 下列高分子化合物是通过缩聚反应得到的是\xz
\xx{聚氯乙烯}{有机玻璃}{涤纶}{顺丁橡胶}

\item 聚乳酸（\ce{-[O-CH(CH3)-CO]_{$n$}-}）是一种可降解塑料，其单体的结构简式为\tk，由单体合成聚乳酸的反应类型为\tk。

\item 下列关于葡萄糖的说法\textbf{错误}的是\xz
\xxa{葡萄糖是多羟基醛，能发生银镜反应}{葡萄糖能发生水解反应生成两种单糖}{葡萄糖与果糖互为同分异构体}{葡萄糖在人体内被氧化生成 \ce{CO2} 和 \ce{H2O} 并放出能量}

\item 下列关于蛋白质的说法正确的是\xz
\xxa{向蛋白质溶液中加入 \ce{CuSO4} 溶液产生沉淀，加水后沉淀溶解}{用灼烧的方法可以区分蚕丝和人造纤维}{蛋白质水解的最终产物是葡萄糖}{酶是蛋白质，温度越高催化效果越好}
\end{xiti}

\zu{B组\quad 高考典型}
\begin{xiti}
\item \laiyuan{2024年广东卷第20题(5)(6)，来源：仓库所附真题文档}\ (5) 在一定条件下，以原子利用率 100\% 的反应制备 \ce{HOCH(CH3)2}。该反应中：①若反应物之一为非极性分子，则另一反应物为\tk（写结构简式）；②若反应物之一为 V 形结构分子，则另一反应物为\tk（写结构简式）。(6) 以 2-溴丙烷为唯一有机原料，合成 \ce{CH3COOCH(CH3)2}。基于你设计的合成路线，回答：①最后一步反应的化学方程式为\tk（注明反应条件）；②第一步反应的化学方程式为\tk（写一个即可，注明反应条件）。

\item \laiyuan{原创}\ 写出下列高分子的单体：(1) 涤纶 \ce{HO-[OC-C6H4-COOCH2CH2O]_{$n$}-H}（\ce{C6H4} 为对位亚苯基）；(2) \ce{-[CH2-C(CH3)(COOCH3)]_{$n$}-}；(3) \ce{-[CH2-CH=CH-CH2-CH2-CH(C6H5)]_{$n$}-}（丁苯橡胶）。

\item \laiyuan{2025年广东卷第20题(5)，来源：仓库所附真题文档}\ 已知：羧酸在一定条件下，可发生类似于丙酮的 α-H 取代反应（如丙酮与 \ce{Br2} 生成 \ce{CH3COCH2Br}）。某小组以乙酸为原料，分三步合成一种含酯基的目标化合物Ⅱ（第三步为题给的光酶催化三组分反应，需要 \ce{BrCH2COOCH3} 作反应物）。①第一步，引入溴：其反应的化学方程式为\tk。②第二步，进行\tk（填具体反应类型）：其反应的化学方程式为\tk（注明反应条件）。

\item \laiyuan{原创}\ 阿司匹林（乙酰水杨酸）的结构简式为
\begin{center}\chemfig{*6(=-=(-O-[:30]C(=[:90]O)-[:-30]CH_3)-(-COOH)=-)}\end{center}
(1) 写出其含氧官能团名称。(2) 1 mol 阿司匹林与足量 NaOH 溶液共热，最多消耗 NaOH 多少摩尔？写出化学方程式。(3) 阿司匹林可由水杨酸（邻羟基苯甲酸）与乙酸酐 [\ce{(CH3CO)2O}] 反应制得，写出化学方程式并指出反应类型。
\end{xiti}

\zu{C组\quad 冲刺高分}
\begin{xiti}
\item \laiyuan{原创}\ 肉桂酸乙酯（\ce{C6H5CH=CHCOOC2H5}）是一种香料，可按下列路线合成：
\[\text{A}(\ce{C7H8})\xrightarrow[\text{光照}]{\ce{Cl2}}\text{B}\xrightarrow[\text{△}]{\text{NaOH 水溶液}}\text{C}\xrightarrow[\text{△}]{\ce{O2},\ \ce{Cu}}\text{D}\xrightarrow[\text{②△}]{\text{①}\ce{CH3CHO},\ \text{稀 NaOH}}\text{E}(\ce{C9H8O})\xrightarrow[\text{②酸化}]{\text{①新制}\ce{Cu(OH)2},\ \text{△}}\text{F}\xrightarrow[\text{浓硫酸，△}]{\ce{C2H5OH}}\text{产品}\]
已知：\ce{RCHO + CH3CHO ->[\text{稀 NaOH}] RCH(OH)CH2CHO ->[\text{△}] RCH=CHCHO + H2O}。
(1) A 的名称为\tk，C 中官能团名称为\tk。(2) B$\to$C、D$\to$E 的反应类型分别为\tk、\tk。(3) 写出 C$\to$D 的化学方程式。(4) E 中的碳碳双键能否用酸性 \ce{KMnO4} 溶液检验？说明理由。(5) F 的同分异构体中，同时满足下列条件的有\tk{}种：①含苯环且苯环上只有两个取代基；②能与 \ce{NaHCO3} 反应放出 \ce{CO2}；③含碳碳双键（不考虑立体异构）。(6) 参照上述路线，以乙醇为唯一有机原料（无机试剂任选），设计合成 2-丁烯醛（\ce{CH3CH=CHCHO}）的路线。

\item \laiyuan{原创}\ 已知（第尔斯—阿尔德反应）：共轭二烯烃与含碳碳双键的化合物可发生环加成，生成六元环，如 \ce{CH2=CH-CH=CH2 + CH2=CH2 ->[\text{△}]} 环己烯。以 1,3-丁二烯和丙烯酸甲酯（\ce{CH2=CHCOOCH3}）为原料：1,3-丁二烯 $+$ 丙烯酸甲酯 $\xrightarrow{\text{△}}$ M $\xrightarrow{\text{NaOH 溶液，△，再酸化}}$ N $\xrightarrow{\ce{H2},\ \text{Ni，△}}$ P。(1) M 中的官能团名称为\tk，M 中有\tk{}个手性碳原子。(2) N 的分子式为\tk。(3) P 的结构简式（名称）为\tk。(4) 若用 1,3-丁二烯与乙炔发生上述反应，产物的结构简式为\tk。
\end{xiti}

\zhangmo{练习解析}
\begin{jiexilist}
\item \ans{C}A：$n$ 值不同，不是同分异构体；B：蔗糖水解得葡萄糖和果糖；C：盐析可逆，正确；D：油脂不是高分子。
\item \ans{C}涤纶由对苯二甲酸与乙二醇缩聚得到；其余三种为加聚产物。
\item \ce{CH3CH(OH)COOH}（乳酸）；缩聚反应。
\item \ans{B}葡萄糖是单糖，不能水解。
\item \ans{B}A：重金属盐使蛋白质变性，不可逆；B：蚕丝是蛋白质，灼烧有烧焦羽毛气味，正确；C：最终产物是氨基酸；D：温度过高酶会变性失活。
\item (5) ①\ce{CH3COCH3}（丙酮与 \ce{H2} 加成，\ce{H2} 为非极性分子）；②\ce{CH3CH=CH2}（丙烯与 \ce{H2O} 加成，\ce{H2O} 为 V 形分子）。(6) 思路：乙酸异丙酯需要乙酸和 2-丙醇。2-丙醇由 2-溴丙烷水解得到；乙酸可由 2-溴丙烷消去得丙烯，再用酸性 \ce{KMnO4} 溶液氧化断裂双键（\ce{CH3CH=CH2 -> CH3COOH + CO2}）得到。①\ce{CH3COOH + (CH3)2CHOH <=>[\text{浓硫酸}][\text{△}] CH3COOCH(CH3)2 + H2O}；②\ce{(CH3)2CHBr + NaOH ->[\ce{H2O}][\text{△}] (CH3)2CHOH + NaBr}（或 \ce{(CH3)2CHBr + NaOH ->[\text{醇}][\text{△}] CH3CH=CH2 ^ + NaBr + H2O}）。
\item (1) 对苯二甲酸（\ce{HOOC-C6H4-COOH}）和乙二醇（\ce{HOCH2CH2OH}）。(2) 甲基丙烯酸甲酯 \ce{CH2=C(CH3)COOCH3}。(3) 1,3-丁二烯 \ce{CH2=CH-CH=CH2} 和苯乙烯 \ce{C6H5CH=CH2}（主链中有双键，以双键为中心取 4 个碳为二烯烃单体，剩余两个碳为乙烯基单体）。
\item ①\ce{CH3COOH + Br2 ->[\text{一定条件}] BrCH2COOH + HBr}；②酯化反应（取代反应）：\ce{BrCH2COOH + CH3OH <=>[\text{浓硫酸}][\text{△}] BrCH2COOCH3 + H2O}。
\item (1) 羧基、酯基。(2) 3 mol：羧基 1 mol；酚酯基水解生成乙酸钠和酚，酚羟基再消耗 1 mol，共 2 mol。\ce{CH3COO-C6H4-COOH + 3NaOH ->[\text{△}] NaO-C6H4-COONa + CH3COONa + 2H2O}（\ce{-C6H4-} 为邻位亚苯基）。(3) 水杨酸 $+$ \ce{(CH3CO)2O} $\to$ 乙酰水杨酸 $+$ \ce{CH3COOH}；取代反应（酯化反应，酚羟基被乙酰化）。
\item (1) 甲苯；羟基。(2) 取代反应（水解反应）；加成反应和消去反应（先羟醛加成，再加热消去）。(3) \ce{2C6H5CH2OH + O2 ->[\ce{Cu}][\text{△}] 2C6H5CHO + 2H2O}。(4) 不能。E 为 \ce{C6H5CH=CHCHO}，醛基也能使酸性 \ce{KMnO4} 溶液褪色，会干扰双键的检验；应先用新制 \ce{Cu(OH)2}（或银氨溶液）氧化醛基，酸化后再用溴水检验双键。(5) F 为肉桂酸 \ce{C6H5CH=CHCOOH}（\ce{C9H8O2}）。苯环上两个取代基且含 \ce{-COOH} 和碳碳双键：取代基组合为 \ce{-COOH} 与 \ce{-CH=CH2}，邻、间、对 3 种。共 3 种。(6) \ce{CH3CH2OH ->[\ce{O2},\ \ce{Cu}][\text{△}] CH3CHO ->[\text{稀 NaOH}] CH3CH(OH)CH2CHO ->[\text{△}] CH3CH=CHCHO}。
\item (1) 碳碳双键、酯基；1（环上连接 \ce{-COOCH3} 的碳原子，所连 4 个基团各不相同）。M 为 3-环己烯-1-甲酸甲酯（3-环己烯甲酸甲酯）。(2) N 为 3-环己烯-1-甲酸，\ce{C7H10O2}。(3) P 为环己烷甲酸（环己基甲酸，\ce{C6H11COOH}）。(4) 1,4-环己二烯（六元环中两个双键处于 1,4 位）。
\end{jiexilist}
% <<<<<<<< 结束：chapters/ch11.tex

% >>>>>>>> 开始：chapters/ch12.tex
% ============================================================
% 第12章 化学实验综合（必修实验基础 + 高三综合实验，对应广东卷第17题）
% ============================================================
\chapter{化学实验综合}

\begin{daoxue}
\textbf{知识主线：}会用仪器（仪器选择、基本操作、安全）$\to$ 会分离与检验（分离提纯方法、离子与气体检验）$\to$ 会制备（气体制备装置的连接）$\to$ 会测定（滴定等定量实验、误差分析）$\to$ 会探究（控制变量、提出假设、设计与评价方案）。

\textbf{高考考情：}广东卷选择题中有 2～3 道实验题（2024年第6、8、11题；2025年第6、11、15题；2026年第9、12题），第17题为 14 分的综合实验探究题，三年的共同特点是“基础操作$+$定量滴定$+$探究与设计”：2024年探究沉淀转化的速率与平衡建立，2025年测苯甲酸的浓度与 $K_\text{a}$、设计测 HClO 的 $K_\text{a}$，2026年探究铜氨配合物与酸的反应。

\textbf{学习目标：}①熟悉常用仪器和基本操作；②掌握常见离子、气体的检验；③会连接气体制备装置；④熟练进行滴定计算与误差分析；⑤会用规范语言描述操作、现象和结论，能设计和评价实验方案。
\end{daoxue}

\section{常用仪器与基本操作}

\begin{center}
\small
\begin{tabularx}{\linewidth}{p{2.6cm}X}
\toprule
类别 & 仪器与要点 \\
\midrule
可直接加热 & 试管、蒸发皿、坩埚（灼烧固体用坩埚，配泥三角、坩埚钳） \\
垫陶土网加热 & 烧杯、烧瓶、锥形瓶 \\
计量仪器 & 量筒（粗量，0.1 mL，无“0”刻度）；滴定管（0.01 mL，“0”刻度在上；酸式盛酸性、氧化性溶液，碱式盛碱性溶液，聚四氟乙烯活塞的滴定管酸碱通用）；容量瓶（配制一定物质的量浓度的溶液，使用前检漏，不能加热、不能长期贮存溶液）；托盘天平（0.1 g）；pH 计 \\
分离仪器 & 普通漏斗（过滤）、分液漏斗（萃取分液、滴加液体）、蒸馏烧瓶与直形冷凝管（蒸馏，冷却水下进上出，温度计水银球在支管口处）、球形冷凝管（冷凝回流） \\
\bottomrule
\end{tabularx}
\end{center}

\textbf{基本操作要点：}
\begin{enumerate}
\item \textbf{检查气密性}：微热法（导管末端插入水中，微热容器，导管口有气泡冒出，停止加热后导管内形成一段稳定水柱）；液差法（关闭止水夹，从长颈漏斗加水，漏斗中液面高于容器内液面，一段时间内液面差不变）。
\item \textbf{测 pH}：用洁净干燥的玻璃棒蘸取待测液点在 pH 试纸中部，与标准比色卡对照；试纸不能润湿，不能测氯水等漂白性溶液的 pH。
\item \textbf{沉淀洗涤}：向过滤器中加蒸馏水至浸没沉淀，待水自然流下，重复 2～3 次。\textbf{检验是否洗净}：取最后一次洗涤液少许于试管中，加入××试剂，若无××现象，则已洗净。
\item \textbf{配制溶液}：计算$\to$称量（量取）$\to$溶解（稀释）并冷却$\to$转移$\to$洗涤 2～3 次并转移$\to$定容（加水至刻度线下 1～2 cm 时改用胶头滴管）$\to$摇匀。定容时俯视，体积偏小，浓度偏高；仰视则偏低。
\item \textbf{安全}：稀释浓硫酸时将浓硫酸沿器壁缓慢注入水中并不断搅拌；加热液体加碎瓷片防暴沸；有毒气体（\ce{Cl2}、\ce{SO2}、\ce{NO2}、\ce{H2S}）要有尾气处理；极易溶于水的气体（\ce{NH3}、HCl）吸收时要防倒吸（倒扣漏斗、干燥管或 \ce{CCl4} 层）。
\end{enumerate}

\begin{zhenti}{2024年广东卷第8题\quad 来源：仓库所附真题文档}
1810 年，戴维首次确认“氯气”是一种新元素组成的单质。兴趣小组利用以下装置进行实验，其中\textbf{难以}达到预期目的的是：A．制备 \ce{Cl2}（分液漏斗盛盐酸，烧瓶中为 \ce{MnO2}，无加热装置）；B．净化、干燥 \ce{Cl2}（饱和食盐水洗气后接无水 \ce{CaCl2} 干燥管）；C．收集 \ce{Cl2}（集气瓶长进短出）；D．验证 \ce{Cl2} 的氧化性（\ce{H2} 在盛 \ce{Cl2} 的集气瓶中燃烧）。

\textbf{解析：}\ce{MnO2} 与浓盐酸反应需要加热，图中缺少加热装置，A 不能达到目的；B 饱和食盐水除 HCl，无水 \ce{CaCl2} 可干燥 \ce{Cl2}；C \ce{Cl2} 密度比空气大，长进短出（向上排空气法）；D \ce{H2} 在 \ce{Cl2} 中安静燃烧，发出苍白色火焰，体现 \ce{Cl2} 的氧化性。答案：A。
\end{zhenti}

\section{物质的分离、提纯与检验}

\subsection{分离与提纯}
\begin{center}
\small
\begin{tabularx}{\linewidth}{p{2.4cm}Xp{4.4cm}}
\toprule
方法 & 适用对象 & 要点 \\
\midrule
过滤 & 固体与液体 & 一贴二低三靠；趁热过滤防止溶质析出 \\
蒸发结晶 & 溶解度受温度影响小的溶质（NaCl） & 出现大量晶体时停止加热，用余热蒸干 \\
冷却结晶 & 溶解度受温度影响大、或带结晶水、或受热易分解的物质 & 蒸发浓缩$\to$冷却结晶$\to$过滤$\to$洗涤$\to$干燥 \\
重结晶 & 提纯固体（如粗苯甲酸含 NaCl、泥沙） & 加热溶解$\to$趁热过滤$\to$冷却结晶$\to$过滤洗涤 \\
萃取分液 & 溶质在两种互不相溶溶剂中溶解度差异大 & 萃取剂与原溶剂不互溶、不反应；下层液体从下口放出，上层液体从上口倒出 \\
蒸馏 & 沸点相差较大的液体混合物 & 加碎瓷片；温度计位置；冷凝水下进上出 \\
洗气 & 气体除杂 & 长进短出 \\
\bottomrule
\end{tabularx}
\end{center}

\textbf{除杂四原则}：不增（不引入新杂质）、不减（不减少被提纯物质）、易分离、易复原。常见除杂：\ce{CO2}（HCl）——饱和 \ce{NaHCO3} 溶液；\ce{Cl2}（HCl）——饱和食盐水；\ce{SO2}（HCl）——饱和 \ce{NaHSO3} 溶液；\ce{FeCl2}（\ce{FeCl3}）——加过量铁粉后过滤；\ce{FeCl3}（\ce{FeCl2}）——通 \ce{Cl2} 或加 \ce{H2O2}；\ce{NaHCO3} 溶液（\ce{Na2CO3}）——通过量 \ce{CO2}；乙酸乙酯（乙酸）——饱和 \ce{Na2CO3} 溶液后分液。

\subsection{常见离子与气体的检验}
\begin{center}
\small
\begin{tabularx}{\linewidth}{p{1.4cm}X}
\toprule
离子 & 方法与现象 \\
\midrule
\ce{Fe^3+} & 加 KSCN 溶液，溶液变红；或加 NaOH 溶液生成红褐色沉淀 \\
\ce{Fe^2+} & 加 \ce{K3[Fe(CN)6]} 溶液生成蓝色沉淀（\ce{Fe^3+} 存在时也适用，2025年广东卷第8题C）；或先加 KSCN 无现象，再加氯水变红（无 \ce{Fe^3+} 时） \\
\ce{NH4+} & 加浓 NaOH 溶液并加热，产生使湿润的红色石蕊试纸变蓝的气体 \\
\ce{SO4^2-} & 先加稀盐酸无现象（排除 \ce{CO3^2-}、\ce{SO3^2-}、\ce{Ag+} 干扰），再加 \ce{BaCl2} 溶液生成白色沉淀 \\
\ce{Cl-} & 加稀硝酸酸化的 \ce{AgNO3} 溶液，生成白色沉淀 \\
\ce{I-} & 加氯水后再加淀粉溶液变蓝；或加氯水后用 \ce{CCl4} 萃取，下层呈紫红色 \\
\ce{Na+}、\ce{K+} & 焰色试验：Na 黄色；K 紫色（透过蓝色钴玻璃观察） \\
\bottomrule
\end{tabularx}
\end{center}

气体检验：\ce{SO2} 使品红溶液褪色、加热又恢复红色；\ce{CO2} 使澄清石灰水变浑浊（需先除 \ce{SO2}）；\ce{NH3} 使湿润的红色石蕊试纸变蓝、遇浓盐酸产生白烟；\ce{Cl2} 使湿润的淀粉-KI 试纸变蓝。

\section{气体的制备与性质实验}

\begin{benzhi}[物质制备实验的一般思路]
无论制备气体、无机物还是有机物，都遵循“\textbf{原理}$\to$\textbf{装置}$\to$\textbf{操作}$\to$\textbf{分离提纯}$\to$\textbf{产率计算}”的主线：①根据目标产物选反应原理（原料易得、反应条件温和、副反应少、绿色）；②根据反应物状态和条件选发生装置，根据产物和杂质的性质选除杂、收集装置；③根据产物的特殊性质确定操作要点（易水解要隔绝水汽、易氧化要通保护气、易分解要控温、易挥发要冷凝回流）；④选择合适方法分离提纯；⑤以不足量的反应物计算理论产量，求产率。
\end{benzhi}

\begin{fangfa}[气体制备装置的连接顺序]
\textbf{发生装置}（固$+$固加热：\ce{NH3}、\ce{O2}；固$+$液不加热：\ce{H2}、\ce{CO2}、\ce{SO2}、\ce{KMnO4} 与浓盐酸制 \ce{Cl2}；固$+$液加热：\ce{MnO2} 与浓盐酸制 \ce{Cl2}）$\to$ \textbf{除杂}（先除杂后干燥）$\to$ \textbf{干燥}（浓硫酸不能干燥 \ce{NH3}、\ce{H2S}；碱石灰不能干燥酸性气体；无水 \ce{CaCl2} 不能干燥 \ce{NH3}）$\to$ \textbf{收集或性质实验} $\to$ \textbf{尾气处理}。
\end{fangfa}

\begin{center}
\small
\begin{tabularx}{\linewidth}{p{1.0cm}p{5.4cm}X}
\toprule
气体 & 实验室制法 & 净化、收集、尾气 \\
\midrule
\ce{Cl2} & \ce{MnO2 + 4HCl}（浓）\ce{->[\text{△}] MnCl2 + Cl2 ^ + 2H2O}；或 \ce{2KMnO4 + 16HCl}（浓）\ce{= 2KCl + 2MnCl2 + 5Cl2 ^ + 8H2O} & 饱和食盐水除 HCl，浓硫酸干燥；向上排空气；NaOH 溶液吸收 \\
\ce{NH3} & \ce{2NH4Cl + Ca(OH)2 ->[\text{△}] CaCl2 + 2NH3 ^ + 2H2O}；或浓氨水滴到 NaOH 固体或 CaO 上 & 碱石灰干燥；向下排空气；水或稀硫酸吸收（防倒吸） \\
\ce{SO2} & \ce{Na2SO3 + H2SO4}（浓）\ce{= Na2SO4 + SO2 ^ + H2O}（常用 70\% 左右的硫酸） & 浓硫酸干燥；向上排空气；NaOH 溶液吸收 \\
\bottomrule
\end{tabularx}
\end{center}

\begin{zhenti}{2026年广东卷第9题\quad 来源：仓库所附真题文档（网络收集版）}
用如图装置进行实验：\ce{Ca(ClO)2} 中滴加盐酸产生 \ce{Cl2}，依次通过 KI 溶液、饱和食盐水、浓硫酸，进入装有铁丝的 U 形管（加热），最后用 NaOH 溶液吸收尾气。反应一段时间后，关闭活塞和止水夹，熄灭酒精灯，冷却后提起铁丝，注入稀硫酸。下列说法\textbf{错误}的是：A．KI 溶液变棕黄色，体现 \ce{I-} 的还原性；B．浓硫酸的作用是吸收水分，干燥 \ce{Cl2}；C．U 形管中溶液遇 KSCN 溶液变红，说明有 \ce{Fe^3+} 存在；D．烧杯溶液中最终存在 $c(\ce{Na+})+c(\ce{H+})=c(\ce{Cl-})+c(\ce{ClO-})$。

\textbf{解析：}A：\ce{Cl2 + 2I- = 2Cl- + I2}，对；B 对；C：Fe 在 \ce{Cl2} 中燃烧生成 \ce{FeCl3}，溶于稀硫酸后遇 KSCN 变红，对；D：NaOH 吸收 \ce{Cl2} 后溶液中还有 \ce{OH-}，电荷守恒应为 $c(\ce{Na+})+c(\ce{H+})=c(\ce{Cl-})+c(\ce{ClO-})+c(\ce{OH-})$，错。答案：D。
\end{zhenti}

\begin{liti}[易潮解物质的制备]
实验室用 \ce{Cl2} 与铁粉加热制备无水 \ce{FeCl3}（\ce{FeCl3} 易潮解，约 300 ℃ 以上升华）。装置依次为：A（\ce{MnO2} 与浓盐酸加热制 \ce{Cl2}）$\to$ B $\to$ C $\to$ D（硬质玻璃管中加热铁粉）$\to$ E（广口瓶，冷水浴收集产品）$\to$ F $\to$ G（NaOH 溶液）。(1) B、C 中分别盛放什么试剂？(2) 为什么要先通一段时间 \ce{Cl2} 再点燃 D 处酒精灯？(3) 产品为什么在 E 中而不是 D 中得到？(4) F 中应盛放什么？作用是什么？
\tcblower
\jie (1) B：饱和食盐水（除 HCl）；C：浓硫酸（干燥）。(2) 排尽装置中的空气，防止铁被 \ce{O2} 氧化、产品混有杂质。(3) 生成的 \ce{FeCl3} 受热升华，随气流进入 E 中冷凝为固体。(4) 浓硫酸（或装有无水 \ce{CaCl2} 的干燥管），防止 G 中的水蒸气进入 E 使 \ce{FeCl3} 潮解。\\
\dianbo 制备“易水解/易潮解”物质的实验，产品收集装置前后都要“防水”：前面干燥气体，后面加干燥装置隔绝尾气处理装置中的水蒸气。
\end{liti}

\section{定量实验}

定量实验的核心是“\textbf{准确测量}”与“\textbf{守恒计算}”：

\begin{enumerate}
\item \textbf{滴定法}：酸碱滴定、氧化还原滴定（\ce{KMnO4} 法：\ce{MnO4^- + 5Fe^2+ + 8H+ = Mn^2+ + 5Fe^3+ + 4H2O}；碘量法：\ce{I2 + 2S2O3^2- = 2I- + S4O6^2-}）、沉淀滴定（\ce{AgNO3} 测 \ce{Cl-}，以 \ce{K2CrO4} 为指示剂）、配位滴定。计算时用“关系式法”：找出待测物与滴定剂之间的物质的量关系（常借助电子守恒）。
\item \textbf{重量法}：将待测成分转化为沉淀（或气体被吸收），称量质量。要保证沉淀完全、洗净、干燥（灼烧）至恒重（两次称量相差不超过 0.1 g）。
\item \textbf{气体体积法}：读数前要恢复到室温、调节量气管两侧液面相平、视线与凹液面最低处相平。
\end{enumerate}

\begin{liti}[氧化还原滴定计算]
称取 2.000 g 某 \ce{FeSO4*7H2O}（$M=278$ g·mol$^{-1}$）样品，溶于稀硫酸配成 100.00 mL 溶液。取 25.00 mL，用 0.02000 mol·L$^{-1}$ \ce{KMnO4} 标准溶液滴定，消耗 16.00 mL。(1) 判断滴定终点的现象。(2) 求样品中 \ce{FeSO4*7H2O} 的质量分数。(3) 若滴定前滴定管尖嘴有气泡、滴定后消失，结果如何变化？
\tcblower
\jie (1) 滴入最后半滴 \ce{KMnO4} 溶液时，溶液由（近）无色变为浅紫红色，且半分钟内不褪色。\\
(2) \ce{MnO4^-} $\sim$ 5\ce{Fe^2+}：$n(\ce{Fe^2+})=5\times0.02000\times0.01600=1.600\times10^{-3}$ mol（25.00 mL 中），总量 $6.400\times10^{-3}$ mol，质量 $6.400\times10^{-3}\times278=1.779$ g，质量分数 $1.779/2.000\approx88.96\%$。\\
(3) 消耗标准液体积读数偏大，结果偏高。
\end{liti}

\begin{liti}[气体体积法测定纯度]
用如图所示的量气装置（锥形瓶中放样品，分液漏斗中为足量稀盐酸，产生的气体进入量气管，量气管与水准管用橡胶管相连）测定某镁带样品中 Mg 的质量分数（杂质不与酸反应）。取样品 0.120 g，实验结束后测得气体体积为 100.8 mL（已换算为标准状况）。(1) 读数时应注意哪些问题？(2) 求 Mg 的质量分数。(3) 若读数时水准管液面高于量气管液面，测得结果偏大还是偏小？
\tcblower
\jie (1) ①冷却至室温再读数；②上下移动水准管，使水准管与量气管内液面相平（使管内气体压强等于大气压）；③视线与凹液面最低处相平。(2) $n(\ce{H2})=\dfrac{0.1008}{22.4}=4.5\times10^{-3}$ mol，$n(\ce{Mg})=4.5\times10^{-3}$ mol，$m=0.108$ g，质量分数 90.0\%。(3) 水准管液面高，量气管内气体被压缩，压强大于大气压，读出的体积偏小，结果偏小。
\end{liti}

\section{探究性实验与方案设计}

\subsection{探究实验的一般思路}
\textbf{提出问题}$\to$\textbf{提出假设}（往往有两种或多种可能）$\to$\textbf{设计实验}（控制变量；设置对照实验或空白实验）$\to$\textbf{实验并记录}$\to$\textbf{分析得出结论}$\to$\textbf{反思评价}。2024年第17题中，乙同学质疑“pH 不变就是平衡”，提出“硫酸钙已耗尽”和“被碳酸钙包裹”两种假设并逐一验证；2026年第17题中，实验结果与理论预测（$V_\text{a}:V_\text{b}=1$）不符，进而提出“沉淀中含 \ce{Cu2(OH)2SO4}”的假设并定量验证——这正是高考所倡导的科学探究过程。

\begin{benzhi}[为什么要做“空白实验”]
空白实验是在不加待测物、其他条件完全相同的情况下进行的实验，用来扣除试剂、溶剂本身带来的影响。2026年广东卷第17题中，(iv) 用同浓度硫酸滴定等体积的稀氨水，测出氨水自身消耗的酸 $V_1$，才能由总体积 $V_2$ 求出配合物反应消耗的酸 $V_2-V_1$。
\end{benzhi}

\subsection{答题规范}
\begin{fangfa}[规范表述模板]
\textbf{①现象描述}（“海陆空”）：溶液颜色由\tk 变为\tk；产生\tk 色沉淀；有\tk 色气体产生、试纸变\tk。\\
\textbf{②检验操作}：取少量待测液于试管中，滴加\tk（试剂），若观察到\tk（现象），则说明\tk（结论）。\\
\textbf{③目的与原因}：说清“这样做是为了……，否则会……”。如“趁热过滤的目的是防止苯甲酸冷却结晶析出而损失”。\\
\textbf{④方案设计}：写清仪器、试剂、操作步骤、需测量的数据及数据处理方法（参见2025年广东卷第17题(5)用 pH 计测 NaClO 溶液 pH 求 HClO 的 $K_\text{a}$）。
\end{fangfa}

\begin{zhenti}{2025年广东卷第17题(1)(2)(4)\quad 来源：仓库所附真题文档}
(1) 向盛有 2 mL 甲苯的试管中加入几滴酸性 \ce{KMnO4} 溶液，振荡，观察到体系颜色\tk。(2) 某苯甲酸粗品含少量泥沙和氯化钠，用重结晶法提纯，需要的操作及顺序为：加热溶解、\tk（从“a．冷却结晶，b．常温过滤，c．蒸馏，d．趁热过滤”中选择并排序）。(4) 甲同学根据“卤原子电负性越大，羧酸 $K_\text{a}$ 越大”推测邻氟苯甲酸（Ⅰ）$>$邻溴苯甲酸（Ⅱ）$>$苯甲酸（Ⅲ）的酸性，并测得常温下三种酸的\textbf{饱和溶液} pH 大小顺序为Ⅲ$>$Ⅱ$>$Ⅰ，据此推断假设成立。乙同学认为该推断依据不足，因为\tk；乙同学用 pH 滴定法测得三种酸的 $K_\text{a}$ 顺序为Ⅱ$>$Ⅰ$>$Ⅲ。

\textbf{答案：}(1) 由紫色变为无色（甲苯被氧化为苯甲酸）。(2) dab（趁热过滤除泥沙，冷却结晶，常温过滤得晶体，NaCl 留在母液中）。(4) 常温下三种酸的饱和溶液浓度不同，不能直接用 pH 比较 $K_\text{a}$。\\
\dianbo 这道题告诉我们：用 pH 比较酸性强弱，必须在\textbf{浓度相同}的前提下；“规律”的外推要谨慎（邻位取代基的影响很复杂）。
\end{zhenti}

% ---------------- 章末 ----------------
\zhangmo{本章小结}
\subsection*{一、知识总结}
\begin{center}
\small
\begin{tabularx}{\linewidth}{p{2.2cm}X}
\toprule
仪器操作 & 精度（量筒 0.1、滴定管 0.01、天平 0.1）；气密性；pH 试纸；洗涤与检验洗净；配制溶液 \\
分离检验 & 过滤、结晶（蒸发/冷却/重结晶）、萃取、蒸馏、洗气；离子检验关注干扰与顺序 \\
气体制备 & 发生$\to$除杂$\to$干燥$\to$收集/性质$\to$尾气 \\
定量实验 & 滴定（关系式、电子守恒）、重量、气体体积；误差归结到测量量 \\
探究设计 & 假设$\to$控制变量与对照$\to$规范表述 \\
\bottomrule
\end{tabularx}
\end{center}
\subsection*{二、方法提升}
\begin{enumerate}
\item 判断装置能否达到目的：逐一检查“原理是否正确、试剂是否合适、装置是否合理（加热、长短管、防倒吸）、操作顺序”。
\item 定量计算三步：写反应（找关系式）$\to$代数据（注意取用比例，如 25.00 mL/100.00 mL）$\to$算结果（有效数字）。
\end{enumerate}
\subsection*{三、易错点归纳}
\begin{enumerate}
\item 用润湿的 pH 试纸测 pH；用 pH 试纸测氯水的 pH。
\item 检验 \ce{SO4^2-} 时用稀硝酸酸化或先加 \ce{BaCl2}（可能有 \ce{SO3^2-}、\ce{Ag+} 干扰）。
\item 气体实验中先干燥后除杂（除杂会带入水蒸气）。
\item 滴定计算忘记换算取样比例；氧化还原滴定找错计量关系。
\item 描述现象只写结论（如“生成 \ce{Fe(OH)3}”）而不写现象（“生成红褐色沉淀”）。
\end{enumerate}

\zhangmo{分层练习}
\zu{A组\quad 基础巩固}
\begin{xiti}[start=1]
\item 下列实验操作正确的是\xz
\xxa{用 10 mL 量筒量取 8.58 mL 稀盐酸}{用碱式滴定管盛装 \ce{KMnO4} 标准溶液}{蒸馏时，温度计水银球置于蒸馏烧瓶的支管口处}{分液时，上层液体也从下口放出}

\item 检验某溶液中是否含有 \ce{Fe^2+}（溶液中同时含有 \ce{Fe^3+}），最合适的试剂是\xz
\xx{KSCN 溶液}{氯水和 KSCN 溶液}{\ce{K3[Fe(CN)6]} 溶液}{NaOH 溶液}

\item \laiyuan{2024年广东卷第6题，来源：仓库所附真题文档}\ 提纯 2.0 g 苯甲酸粗品（含少量 NaCl 和泥沙）的过程：苯甲酸粗品 $\xrightarrow{\text{加水加热}}$ 浊液 $\xrightarrow{\text{趁热过滤}}$ 滤液 $\xrightarrow{\text{X}}$ $\xrightarrow{\text{过滤洗涤}}$ 苯甲酸。其中，操作 X 为\xz
\xx{加热蒸馏}{加水稀释}{冷却结晶}{萃取分液}

\item 选择合适的分离方法（填序号：a.过滤\ b.蒸发结晶\ c.冷却结晶\ d.萃取分液\ e.蒸馏）：(1) 从食盐水中获得 NaCl\tk；(2) 从溴水中提取溴\tk；(3) 分离乙醇与水\tk；(4) 除去 \ce{KNO3} 中少量的 NaCl\tk。

\item 下列试剂的保存方法\textbf{错误}的是\xz
\xxa{浓硝酸保存在棕色试剂瓶中，置于阴凉处}{NaOH 溶液保存在带磨口玻璃塞的试剂瓶中}{少量金属钠保存在煤油中}{氢氟酸保存在塑料瓶中}
\end{xiti}

\zu{B组\quad 高考典型}
\begin{xiti}
\item \laiyuan{2025年广东卷第15题，来源：仓库所附真题文档}\ 按如下装置进行实验：试管 a 中为稀硝酸，铜丝可上下抽动；a 连接烧瓶 b（空瓶），b 再连接试管 c（含酚酞的 \ce{NaHCO3} 溶液）。将铜丝插入溶液中，当 c 中红色褪去时，将铜丝拔离液面。下列叙述\textbf{错误}的是\xz
\xxa{a 中有化合反应发生，并有颜色变化}{b 中气体变红棕色时，所含氮氧化物至少有两种}{c 中溶液红色刚好褪去时，\ce{HCO3-} 恰好完全反应}{若将 a 中稀硝酸换为浓硫酸并加热，则 c 中溶液颜色会褪去}

\item \laiyuan{原创}\ 用如下方法测定某漂白粉中“有效氯”（以 \ce{Cl2} 计的质量分数）：称取 1.000 g 样品，加水研磨后配成 250.0 mL 溶液；取 25.00 mL，加入过量 KI 溶液和稀硫酸（\ce{ClO- + 2I- + 2H+ = Cl- + I2 + H2O}），用 0.1000 mol·L$^{-1}$ \ce{Na2S2O3} 标准溶液滴定生成的 \ce{I2}，以淀粉为指示剂，消耗 9.00 mL。(1) 终点现象为\tk。(2) 求有效氯的质量分数。(3) 若滴定时间过长，空气中的 \ce{O2} 氧化了部分 \ce{I-}，结果偏\tk。

\item \laiyuan{2026年广东卷第17题(1)～(3)，来源：仓库所附真题文档（网络收集版）}\ 某兴趣小组对铜的配合物进行探究实验。(1) 氨水的制备：从下列装置中选择需要的装置\tk。A．加热 NaOH 溶液与 \ce{NH4Cl} 固体的发生装置；B．烧杯中盛水、倒扣漏斗的吸收装置；C．盛 \ce{P2O5} 的干燥管；D．分液漏斗中盛甲苯与 \ce{NH4Cl} 水溶液的萃取装置。(2) 制备铜氨溶液，以其为电解液在铁件表面镀铜，铁件连接直流电源\tk{}极。(3) 铜氨配合物的制备与表征：向 \ce{CuSO4} 溶液中滴加氨水至过量，得深蓝色溶液，加入乙醇，析出固体，过滤，\tk，晾干得产物 Q。用 X 射线衍射仪测试 Q，观察到\tk，说明 Q 为晶体。

\item \laiyuan{原创，与第1章联系}\ 用 50 mL 0.50 mol·L$^{-1}$ 盐酸与 50 mL 0.55 mol·L$^{-1}$ NaOH 溶液测定中和反应反应热，三次实验的温度差分别为 3.30 ℃、3.40 ℃、3.35 ℃（溶液密度按 1 g·mL$^{-1}$、比热容按 4.18 J·g$^{-1}$·℃$^{-1}$ 计算）。(1) 求测得的中和反应反应热 $\Delta H$（保留一位小数）。(2) 与理论值 $-57.3$ kJ·mol$^{-1}$ 相比，指出两条可能的原因。
\end{xiti}

\zu{C组\quad 冲刺高分}
\begin{xiti}
\item \laiyuan{原创}\ 某小组探究 \ce{FeCl3} 溶液与 \ce{Na2SO3} 溶液的反应。实验：向 2 mL 1 mol·L$^{-1}$ \ce{FeCl3} 溶液中滴加 1 mol·L$^{-1}$ \ce{Na2SO3} 溶液，溶液由黄色变为红褐色，无气体和沉淀产生；放置一段时间后，溶液变为浅绿色。
(1) 甲同学认为红褐色是因为发生了相互促进的水解反应生成了 \ce{Fe(OH)3} 胶体，写出该反应的离子方程式。
(2) 乙同学认为发生了氧化还原反应，写出离子方程式。
(3) 设计实验证明最终溶液中存在 \ce{Fe^2+} 和 \ce{SO4^2-}（写出操作、现象和结论）。
(4) 结合实验现象，说明两种反应的速率和限度有何不同。

\item \laiyuan{原创}\ 铜与浓硝酸反应后的溶液呈绿色，而铜与稀硝酸反应后的溶液呈蓝色。甲同学认为是 \ce{Cu(NO3)2} 浓度较大所致；乙同学认为是溶液中溶解了 \ce{NO2}（\ce{NO2} 溶于 \ce{Cu(NO3)2} 溶液呈黄色，与蓝色混合呈绿色）。(1) 设计实验验证甲的假设（写出操作和预期现象）。(2) 设计实验验证乙的假设。(3) 若实验证实乙的假设正确，如何使绿色溶液转变为蓝色？
\end{xiti}

\zhangmo{练习解析}
\begin{jiexilist}
\item \ans{C}A：量筒精度 0.1 mL；B：\ce{KMnO4} 有强氧化性，会腐蚀碱式滴定管的橡胶管，应用酸式滴定管；C 正确；D：上层液体从上口倒出。
\item \ans{C}\ce{Fe^3+} 存在时，KSCN 法无法判断 \ce{Fe^2+}；\ce{K3[Fe(CN)6]} 与 \ce{Fe^2+} 生成蓝色沉淀，不受 \ce{Fe^3+} 干扰。
\item \ans{C}苯甲酸溶解度随温度变化大，NaCl 变化小；趁热过滤除泥沙后，冷却结晶使苯甲酸析出，NaCl 留在溶液中。
\item (1) b；(2) d（用 \ce{CCl4} 或苯萃取）；(3) e（可先加 CaO 吸水后蒸馏）；(4) c（\ce{KNO3} 溶解度随温度变化大）。
\item \ans{B}NaOH 与玻璃中的 \ce{SiO2} 反应生成具有黏性的 \ce{Na2SiO3}，会使瓶塞与瓶口粘住，应用橡胶塞。
\item \ans{C}A：Cu 与稀硝酸生成 NO，NO 与上方空气中 \ce{O2} 化合生成红棕色 \ce{NO2}，溶液也变蓝，对。B：b 中出现 \ce{NO2} 说明有 \ce{O2} 与 NO 反应，且存在 \ce{2NO2 <=> N2O4}，至少有两种氮氧化物，对。C：酚酞变色范围 pH 约 8.2～10，红色褪去时溶液可能仍呈弱碱性，不能说明 \ce{HCO3-} 恰好完全反应，错。D：浓硫酸与铜加热生成 \ce{SO2}，通入 \ce{NaHCO3} 溶液生成 \ce{NaHSO3}，溶液呈酸性，红色褪去，对。答案：C。
\item (1) 滴入最后半滴 \ce{Na2S2O3} 溶液时，溶液蓝色褪去，且半分钟内不恢复。(2) 关系式：\ce{Cl2}（有效氯）$\sim$ \ce{ClO-} $\sim$ \ce{I2} $\sim$ 2\ce{S2O3^2-}。$n(\ce{S2O3^2-})=9.000\times10^{-4}$ mol，$n(\ce{Cl2})=4.500\times10^{-4}$ mol（25.00 mL 中），250.0 mL 中共 $4.500\times10^{-3}$ mol，$m(\ce{Cl2})=0.3195$ g，有效氯约为 32.0\%。（有效氯的定义：样品与盐酸作用放出的 \ce{Cl2} 的质量分数，\ce{ClO- + Cl- + 2H+ = Cl2 ^ + H2O}，故 1 mol \ce{ClO-} 相当于 1 mol \ce{Cl2}。）(3) 高（\ce{O2} 氧化 \ce{I-} 生成额外的 \ce{I2}，消耗 \ce{Na2S2O3} 偏多）。
\item (1) AB。用 NaOH 与 \ce{NH4Cl} 加热制 \ce{NH3}（A），用倒扣漏斗将 \ce{NH3} 溶于水防倒吸（B）；\ce{P2O5} 是酸性干燥剂，会与 \ce{NH3} 反应；D 是萃取装置，与制氨水无关。(2) 负（镀件作阴极）。(3) 洗涤；明锐的衍射峰（或分立的斑点）。
\item (1) 平均温差 $\Delta t=\frac{3.30+3.40+3.35}{3}=3.35$ ℃；$Q=100\times4.18\times3.35=1400$ J$=1.40$ kJ；$\Delta H=-\dfrac{1.40}{0.025}\approx-56.0$ kJ·mol$^{-1}$。(2) 测得值偏小，可能原因：装置保温隔热效果不好，有热量散失；NaOH 溶液分多次倒入；测完酸的温度后未冲洗温度计就测碱的温度；溶液比热容、密度近似取值带来误差等（任答两条）。
\item (1) \ce{2Fe^3+ + 3SO3^2- + 6H2O = 2Fe(OH)3}（胶体）\ce{+ 3H2SO3}。(2) \ce{2Fe^3+ + SO3^2- + H2O = 2Fe^2+ + SO4^2- + 2H+}。(3) 取少量最终溶液，滴加 \ce{K3[Fe(CN)6]} 溶液，产生蓝色沉淀，说明含 \ce{Fe^2+}；另取少量溶液，先加足量稀盐酸，再加 \ce{BaCl2} 溶液，产生白色沉淀，说明含 \ce{SO4^2-}（先加盐酸可除去 \ce{SO3^2-} 的干扰）。(4) 开始时溶液迅速变为红褐色，说明相互促进的水解反应速率较快；放置后变为浅绿色，说明氧化还原反应虽然较慢，但进行的程度（限度）更大，最终占主导——体现了“速率”与“限度”是两个不同的问题。
\item (1) 取适量 \ce{Cu(NO3)2} 晶体，配制成与反应后溶液 \ce{Cu^2+} 浓度相同的溶液，观察颜色：若呈蓝色，则甲的假设不成立；若呈绿色，则成立（控制变量：只保留“浓度”这一因素）。(2) 向绿色溶液中通入空气（或 \ce{N2}）并微热，赶出溶解的 \ce{NO2}，若溶液由绿色变为蓝色，说明乙的假设成立；也可向蓝色的 \ce{Cu(NO3)2} 溶液中通入 \ce{NO2}，若变为绿色，同样支持乙的假设。(3) 通入空气（或 \ce{O2}、\ce{N2}）并加热，或加水稀释，使 \ce{NO2} 逸出或反应（\ce{4NO2 + O2 + 2H2O = 4HNO3}）。
\end{jiexilist}
% <<<<<<<< 结束：chapters/ch12.tex

% >>>>>>>> 开始：chapters/ch13.tex
% ============================================================
% 第13章 元素化合物与化工流程（必修元素化学 + 高三工艺流程，对应广东卷第18题）
% ============================================================
\chapter{元素化合物与化工流程}

\begin{daoxue}
\textbf{知识主线：}元素化合物不是孤立的方程式，而是一张“价—类”网络：同一元素按\textbf{化合价}（氧化还原）和\textbf{物质类别}（酸碱性、复分解）相互转化。化工流程题就是在这张网上，用反应原理（速率、平衡、$K_\text{sp}$、电化学）和实验方法（分离提纯）解决“从原料到产品”的实际问题。

\textbf{高考考情：}选择题有价类二维图（2024年第14题、2026年第13题）、物质转化（2025年第7题）、陈述因果（2025年第8题）；第18题为 14 分的工艺流程题，三年分别是从炭渣中提镓（2024）、从多金属精矿提取 Fe、Cu、Ni（2025）、从金精矿提金（2026），综合考查预处理、氧化还原方程式、$K_\text{sp}$ 计算、配合物与晶胞。

\textbf{学习目标：}①构建 Na、Al、Fe、Cu、Cl、S、N、Si 的价类网络；②掌握流程题的“预处理—转化—分离—计算”框架；③熟练书写陌生氧化还原（离子）方程式。
\end{daoxue}

\section{元素化合物的核心网络}

\subsection{金属元素}
\begin{center}
\small
\begin{tabularx}{\linewidth}{p{1.0cm}X}
\toprule
元素 & 核心性质与方程式 \\
\midrule
Na & \ce{2Na + 2H2O = 2NaOH + H2 ^}；\ce{2Na2O2 + 2H2O = 4NaOH + O2 ^}，\ce{2Na2O2 + 2CO2 = 2Na2CO3 + O2}（\ce{Na2O2} 既是氧化剂又是还原剂，每生成 1 mol \ce{O2} 转移 2 mol 电子）；\ce{2NaHCO3 ->[\text{△}] Na2CO3 + CO2 ^ + H2O}；侯氏制碱：\ce{NaCl + NH3 + CO2 + H2O = NaHCO3 v + NH4Cl} \\
Al & 两性：\ce{2Al + 2NaOH + 6H2O = 2Na[Al(OH)4] + 3H2 ^}；\ce{Al2O3 + 2OH- + 3H2O = 2[Al(OH)4]-}；\ce{Al(OH)3 + OH- = [Al(OH)4]-}；\ce{[Al(OH)4]- + CO2}（过量）\ce{= Al(OH)3 v + HCO3-}；常温下遇浓硫酸、浓硝酸钝化 \\
Fe & \ce{3Fe + 4H2O(g) ->[\text{高温}] Fe3O4 + 4H2}；\ce{2Fe^3+ + Fe = 3Fe^2+}；\ce{2Fe^3+ + Cu = 2Fe^2+ + Cu^2+}；\ce{2Fe^2+ + Cl2 = 2Fe^3+ + 2Cl-}；\ce{2Fe^2+ + H2O2 + 2H+ = 2Fe^3+ + 2H2O}；\ce{4Fe(OH)2 + O2 + 2H2O = 4Fe(OH)3}（白色$\to$灰绿色$\to$红褐色） \\
Cu & \ce{Cu + 2H2SO4}（浓）\ce{->[\text{△}] CuSO4 + SO2 ^ + 2H2O}；\ce{3Cu + 8HNO3}（稀）\ce{= 3Cu(NO3)2 + 2NO ^ + 4H2O}；\ce{Cu + 4HNO3}（浓）\ce{= Cu(NO3)2 + 2NO2 ^ + 2H2O}；\ce{Cu2O} 在酸中歧化：\ce{Cu2O + 2H+ = Cu + Cu^2+ + H2O} \\
\bottomrule
\end{tabularx}
\end{center}

\subsection{非金属元素}
\begin{center}
\small
\begin{tabularx}{\linewidth}{p{1.0cm}X}
\toprule
元素 & 核心性质与方程式 \\
\midrule
Cl & \ce{Cl2 + H2O <=> HCl + HClO}；\ce{Cl2 + 2NaOH = NaCl + NaClO + H2O}；\ce{2Cl2 + 2Ca(OH)2 = CaCl2 + Ca(ClO)2 + 2H2O}（漂白粉）；\ce{HClO} 有强氧化性、漂白性，见光分解；“84”消毒液与洁厕灵（含盐酸）不能混用：\ce{ClO- + Cl- + 2H+ = Cl2 ^ + H2O} \\
S & \ce{SO2}：酸性氧化物；漂白性（与有色物质化合，可逆，不能漂白酸碱指示剂）；还原性 \ce{SO2 + Br2 + 2H2O = H2SO4 + 2HBr}；氧化性 \ce{SO2 + 2H2S = 3S + 2H2O}。浓硫酸：吸水性（作干燥剂）、脱水性（使蔗糖变黑）、强氧化性（与 Cu、C 反应，常温下使 Fe、Al 钝化） \\
N & \ce{N2 + O2 ->[\text{放电}] 2NO}，\ce{2NO + O2 = 2NO2}，\ce{3NO2 + H2O = 2HNO3 + NO}；\ce{4NH3 + 5O2 ->[\text{催化剂}][\text{△}] 4NO + 6H2O}；硝酸见光分解（棕色瓶保存）、强氧化性；铵盐受热分解、与碱反应放 \ce{NH3} \\
Si & Si 为半导体材料；\ce{SiO2} 制光导纤维，\ce{SiO2 + 4HF = SiF4 ^ + 2H2O}（刻蚀玻璃），\ce{SiO2 + 2NaOH = Na2SiO3 + H2O}；传统硅酸盐材料：玻璃、水泥、陶瓷；新型无机非金属材料：SiC、\ce{Si3N4}、碳纳米材料 \\
\bottomrule
\end{tabularx}
\end{center}

\begin{center}
\small
\begin{tabularx}{\linewidth}{p{2.2cm}X}
\toprule
推断“突破口” & 常见物质 \\
\midrule
特征颜色 & 淡黄色固体：S、\ce{Na2O2}、AgBr；红棕色气体：\ce{NO2}、溴蒸气；黄绿色气体：\ce{Cl2}；红褐色沉淀：\ce{Fe(OH)3}；蓝色沉淀：\ce{Cu(OH)2}；砖红色沉淀：\ce{Cu2O}；黑色固体：CuO、FeO、\ce{Fe3O4}、\ce{MnO2}、C；溶液颜色：\ce{Fe^3+} 黄、\ce{Fe^2+} 浅绿、\ce{Cu^2+} 蓝、\ce{MnO4-} 紫红、\ce{[Cu(NH3)4]^2+} 深蓝 \\
特征现象 & 白色沉淀$\to$灰绿色$\to$红褐色：\ce{Fe(OH)2} 被氧化；遇空气变红棕色：NO；使品红褪色、加热复原：\ce{SO2}；遇 KSCN 变红：\ce{Fe^3+}；淀粉变蓝：\ce{I2} \\
特征反应 & 既能与强酸又能与强碱反应：Al、\ce{Al2O3}、\ce{Al(OH)3}、弱酸的铵盐、弱酸的酸式盐；与水反应放气：Na、\ce{Na2O2}、\ce{CaC2}；两种气体相遇生成白烟：\ce{NH3} 与 HCl \\
\bottomrule
\end{tabularx}
\end{center}

\begin{fangfa}[价类二维图的解题方法]
横轴是物质类别（单质、氧化物、酸、碱、盐），纵轴是化合价。①\textbf{同一横线}（化合价不变）上的转化是非氧化还原反应（酸碱、复分解）；②\textbf{纵向}转化需要氧化剂或还原剂；③先由“特征点”锁定元素：如 $-1$ 价的酸是 HCl，$+5$ 价的酸可能是 \ce{HNO3}，同时有 $+2$、$+3$ 价氧化物和碱的金属是 Fe。2026年广东卷第13题中，“$+4$ 价氧化物 c 可由 $+4$ 价的盐 e 受热分解产生”对应 \ce{CaCO3 ->[\text{高温}] CaO + CO2 ^}；“a、c、f 均可用于自来水消毒”错在 f（$-1$ 价的酸 HCl）不能消毒（a 可为 \ce{Cl2}，c 可为 \ce{ClO2}）。
\end{fangfa}

\begin{zhenti}{2025年广东卷第7题\quad 来源：仓库所附真题文档}
能满足下列物质间直接转化关系，且推理成立的是：单质 X $\xrightarrow{\ce{O2}}$ 氧化物1 $\xrightarrow{\ce{O2}}$ 氧化物2 $\xrightarrow{\ce{H2O}}$ 酸（或碱）$\xrightarrow{\text{NaOH（或 HCl）}}$ 盐。A．X 可为铝，盐的水溶液一定显酸性；B．X 可为硫，氧化物1可使品红溶液褪色；C．X 可为钠，氧化物2可与水反应生成 \ce{H2}；D．X 可为碳，盐的热稳定性：\ce{NaHCO3}$>$\ce{Na2CO3}。

\textbf{解析：}A：\ce{Al2O3} 不能再被 \ce{O2} 氧化，也不与水反应，错；B：\ce{S -> SO2 -> SO3 -> H2SO4 -> Na2SO4}，\ce{SO2} 能使品红褪色，对；C：\ce{Na2O2} 与水反应生成 \ce{O2}，错；D：\ce{Na2CO3} 比 \ce{NaHCO3} 稳定，错。答案：B。
\end{zhenti}

\section{化工流程题的解题框架}

\subsection{流程的基本结构}
\begin{center}
\begin{tikzpicture}[node distance=0.55cm,font=\small,box/.style={draw,rounded corners=2pt,fill=lightblue,minimum height=0.7cm,align=center}]
\node[box] (a) {原料};
\node[box,right=of a] (b) {预处理\\\scriptsize 粉碎、焙烧、浸取};
\node[box,right=of b] (c) {核心转化\\\scriptsize 氧化还原、调 pH 沉淀};
\node[box,right=of c] (d) {分离提纯\\\scriptsize 过滤、萃取、结晶};
\node[box,right=of d] (e) {产品};
\draw[-Stealth,thick] (a)--(b); \draw[-Stealth,thick] (b)--(c); \draw[-Stealth,thick] (c)--(d); \draw[-Stealth,thick] (d)--(e);
\draw[-Stealth,thick,deepgreen,dashed] (d.south) -- ++(0,-0.5) -| node[pos=0.25,below]{\scriptsize 循环利用} (b.south);
\end{tikzpicture}
\end{center}

\subsection{各环节的设问与答题要点}
\begin{center}
\small
\begin{tabularx}{\linewidth}{p{2.4cm}X}
\toprule
环节（常见设问） & 答题要点 \\
\midrule
粉碎、研磨 & 增大反应物接触面积，加快反应速率，提高浸取率（原料利用率） \\
焙烧、煅烧 & 除去有机物或硫；使物质结构改变、更易溶解；将低价态物质氧化 \\
酸浸、碱浸 & 使目标元素溶解进入溶液，或溶解除去杂质（如碱浸除 Al、Si） \\
提高浸取率 & 搅拌、适当升温、适当增大酸（碱）浓度、延长浸取时间、多次浸取 \\
加氧化剂 & 将 \ce{Fe^2+} 氧化为 \ce{Fe^3+}，便于调 pH 除去；\ce{H2O2} 为“绿色氧化剂”（还原产物为水，不引入杂质），但温度不宜过高（防分解） \\
调 pH & 使某些离子沉淀完全（$\le10^{-5}$ mol·L$^{-1}$）而目标离子不沉淀；所用试剂应能消耗 \ce{H+} 且不引入杂质（如用 ZnO 调含 \ce{Zn^2+} 的溶液） \\
控制温度 & 加快速率；防止物质（\ce{H2O2}、\ce{NH4HCO3}、浓硝酸）分解或（HCl、氨）挥发；使催化剂活性最大；利用溶解度差异结晶 \\
萃取与反萃取 & 利用目标物在有机相中溶解度大（或形成配合物进入有机相），再反萃取回水相 \\
洗涤 & 冰水洗：降低溶解损失；乙醇洗：减少溶解、乙醇易挥发便于干燥 \\
循环利用 & 找“流程中既是前面的原料、又是后面的生成物”的物质 \\
\bottomrule
\end{tabularx}
\end{center}

\subsection{陌生氧化还原方程式的书写}
\begin{fangfa}[四步配平法（以 2026年广东卷第18题(5)为例）]
\textbf{例：}“沉金”时 \ce{[AuCl4]-} 被 \ce{HSO3-} 在弱酸性条件下还原为 Au。\\
①\textbf{定物质}：氧化剂 \ce{[AuCl4]-}$\to$Au，还原剂 \ce{HSO3-}$\to$\ce{SO4^2-}；Cl 以 \ce{Cl-} 形式存在。\\
②\textbf{配电子}：Au 由 $+3\to0$，每个得 3 个电子；S 由 $+4\to+6$，每个失 2 个电子。最小公倍数 6：2 \ce{[AuCl4]-}、3 \ce{HSO3-}。\\
③\textbf{配电荷}（看介质）：左边 $2\times(-1)+3\times(-1)=-5$；右边 3 \ce{SO4^2-} 和 8 \ce{Cl-} 共 $-14$。弱酸性介质用 \ce{H+} 配平：右边加 9 \ce{H+}，电荷 $-14+9=-5$。\\
④\textbf{配原子}（补 \ce{H2O}）：右边 H 为 9，左边 3 \ce{HSO3-} 含 3 个 H，左边补 3 \ce{H2O}；检查 O：左 $3\times3+3=12$，右 $3\times4=12$。\\
\ce{2[AuCl4]- + 3HSO3- + 3H2O = 2Au v + 3SO4^2- + 8Cl- + 9H+}
\end{fangfa}

\section{真题精析：2025年广东卷第18题}

\begin{zhenti}{2025年广东卷第18题（节选）\quad 来源：仓库所附真题文档}
一种从多金属精矿（主要含 Fe、Al、Cu、Ni、O 等元素）中提取 Fe、Cu、Ni 并探究新型绿色冶铁方法的工艺如下：矿粉通入 \ce{SO2} \textbf{酸浸}（固体残渣分离）$\to$ 浸取液调 pH$=3.0$ $\to$ 通入空气\textbf{高压加热}，得 \ce{Fe2O3} 和滤液1 $\to$ 滤液1 常温\textbf{沉铝}得滤饼和滤液2 $\to$ 滤液2 \textbf{选择萃取}得 Cu、Ni 配合物（最终制得 \ce{Ni_xCu_yN_z}）。\ce{Fe2O3} 一部分与 \ce{NH3}、\ce{H2} 混合在 700 ℃ 加热制 Fe，另一部分在 NaOH 溶液中\textbf{电解}制 Fe。$K_\text{sp}$（298 K）：\ce{Fe(OH)3} $2.8\times10^{-39}$，\ce{Al(OH)3} $1.3\times10^{-33}$，\ce{Cu(OH)2} $2.2\times10^{-20}$，\ce{Ni(OH)2} $5.5\times10^{-16}$。

\begin{center}
\begin{tikzpicture}[font=\scriptsize,node distance=0.4cm and 0.85cm,
 st/.style={draw,rounded corners=1.5pt,fill=lightblue,minimum height=0.55cm,inner sep=2pt},
 pr/.style={inner sep=1pt}]
\node[pr] (ore) {矿粉};
\node[st,right=of ore] (leach) {酸浸};
\node[st,right=of leach] (hp) {高压加热};
\node[st,right=of hp] (al) {沉铝};
\node[st,right=of al] (ext) {选择萃取};
\node[pr,right=of ext,align=left] (cu) {Cu 配合物\\Ni 配合物};
\draw[-Stealth] (ore)--(leach); \draw[-Stealth] (leach)--node[above]{浸取液}node[below]{pH$=3.0$}(hp);
\draw[-Stealth] (hp)--node[above]{滤液1}(al); \draw[-Stealth] (al)--node[above]{滤液2}(ext); \draw[-Stealth] (ext)--(cu);
\node[pr,above=0.3cm of leach] (so2) {\ce{SO2}}; \draw[-Stealth] (so2)--(leach);
\node[pr,above=0.3cm of hp] (air) {空气}; \draw[-Stealth] (air)--(hp);
\node[pr,below=0.3cm of leach] (res) {固体残渣}; \draw[-Stealth] (leach)--(res);
\node[pr,below=0.3cm of al] (cake) {滤饼}; \draw[-Stealth] (al)--(cake);
\node[pr,below=0.3cm of hp] (fe2o3) {\ce{Fe2O3}}; \draw[-Stealth] (hp)--(fe2o3);
\node[st,below left=0.35cm and 0.1cm of fe2o3] (h700) {700 ℃ 加热（\ce{NH3}、\ce{H2}）};
\node[st,below right=0.35cm and 0.1cm of fe2o3] (elec) {电解（NaOH 溶液）};
\draw[-Stealth] (fe2o3)--(h700); \draw[-Stealth] (fe2o3)--(elec);
\node[pr,below=0.25cm of h700] (fe1) {Fe}; \node[pr,below=0.25cm of elec] (fe2) {Fe};
\draw[-Stealth] (h700)--(fe1); \draw[-Stealth] (elec)--(fe2);
\end{tikzpicture}
\end{center}
(1)“酸浸”中，提高浸取速率的措施有\tk（写一条）。
(2)“高压加热”时，生成 \ce{Fe2O3} 的离子方程式为\tk。
(3)“沉铝”时，pH 最高可调至\tk（已知滤液1中 $c(\ce{Cu^2+})=0.022$ mol·L$^{-1}$，$c(\ce{Ni^2+})=0.042$ mol·L$^{-1}$）。
(4)“选择萃取”中，镍形成的配合物中 Ni 与吡啶环上的 N、磺酸根上的 O、水分子中的 O 配位，且配体带有长烷基链。镍易进入有机相的原因有\tk（填字母）。A．镍与 N、O 形成配位键\quad B．配位时 \ce{Ni^2+} 被还原\quad C．配合物与水能形成分子间氢键\quad D．烷基链有疏水性
(6)①“700 ℃ 加热”步骤中，混合气体中仅加少量 \ce{H2}，但借助工业合成氨的逆反应，可使 Fe 不断生成。写出该步骤发生反应的化学方程式。②“电解”时，\ce{Fe2O3} 颗粒分散于 NaOH 溶液中，以 Fe 片、石墨棒为电极，说明电极连接方式。③与传统高炉炼铁相比，两种新型冶铁方法体现“绿色化学”思想的共同点是\tk。

\textbf{解析与答案：}
(1) 搅拌、粉碎矿石、适当提高酸浸温度（任写一条）。
(2) 通入 \ce{SO2} 把 \ce{Fe^3+} 还原为 \ce{Fe^2+}，所以浸取液中铁以 \ce{Fe^2+} 存在；高压加热时被 \ce{O2} 氧化并水解：\ce{4Fe^2+ + O2 + 4H2O ->[\text{高压}][\text{△}] 2Fe2O3 v + 8H+}。
(3) 5（计算见第4章真题链接：\ce{Cu^2+} 开始沉淀时 $c(\ce{OH-})=10^{-9}$ mol·L$^{-1}$）。
(4) AD。配位键使镍与萃取剂结合；长烷基链的疏水性使配合物易溶于有机相。配位时镍仍为 +2 价，B 错；与水形成氢键反而有利于留在水相，C 错。
(6) ①\ce{2NH3 <=>[\text{催化剂}][\text{△}] N2 + 3H2}、\ce{Fe2O3 + 3H2 ->[700\ \text{℃}] 2Fe + 3H2O}（\ce{H2} 消耗使氨分解平衡正移，持续提供 \ce{H2}）。②Fe 片接电源负极作阴极（\ce{Fe2O3} 在阴极被还原为 Fe），石墨棒接正极作阳极。③都没有产生 CO 等污染性气体（不用焦炭作还原剂）。
\end{zhenti}

% ---------------- 章末 ----------------
\zhangmo{本章小结}
\subsection*{一、知识总结}
\begin{center}
\small
\begin{tabularx}{\linewidth}{p{2.2cm}X}
\toprule
金属 & Na（\ce{Na2O2}、\ce{Na2CO3}/\ce{NaHCO3}）、Al（两性）、Fe（\ce{Fe^2+}/\ce{Fe^3+} 转化与检验）、Cu（与浓硫酸、硝酸） \\
非金属 & Cl（\ce{Cl2}、HClO、含氯消毒剂）、S（\ce{SO2} 三性、浓硫酸三性）、N（氮氧化物、硝酸、氨）、Si（\ce{SiO2}、硅酸盐） \\
流程框架 & 预处理（加快速率、提高浸取率）$\to$转化（氧化还原、调 pH）$\to$分离（过滤、萃取、结晶）$\to$计算（$K_\text{sp}$、产率） \\
方程式 & 定物质$\to$配电子$\to$配电荷（看介质）$\to$配原子 \\
\bottomrule
\end{tabularx}
\end{center}
\subsection*{二、方法提升}
\begin{enumerate}
\item 读流程“三看”：看箭头进出（加了什么、出去了什么）、看元素去向（目标元素在哪一相）、看条件（温度、pH、氧化剂）。
\item 不会的小问先跳过，流程题各问相对独立，常有“送分”问（如提高浸取率的措施）。
\end{enumerate}
\subsection*{三、易错点归纳}
\begin{enumerate}
\item 酸性条件下的离子方程式中出现 \ce{OH-}，或碱性条件下出现 \ce{H+}。
\item 调 pH 时选用会引入杂质的试剂（如用 NaOH 调含 \ce{Cu^2+} 溶液的 pH 以制纯净的 \ce{CuSO4}）。
\item 认为 \ce{SO2} 能使紫色石蕊先变红后褪色（只变红不褪色）。
\item 认为 \ce{Na2O2} 是碱性氧化物；\ce{Na2O2} 与 \ce{CO2} 反应时 \ce{CO2} 是氧化剂（实际 \ce{Na2O2} 自身发生歧化）。
\end{enumerate}

\zhangmo{分层练习}
\zu{A组\quad 基础巩固}
\begin{xiti}[start=1]
\item 下列陈述Ⅰ和陈述Ⅱ均正确，且具有因果关系的是\xz
\xxa{浓硝酸保存在棕色试剂瓶中——浓硝酸具有强氧化性}{向蔗糖中加入适量浓硫酸，蔗糖变黑——浓硫酸具有脱水性}{\ce{Fe^2+} 与 \ce{K3[Fe(CN)6]} 可生成蓝色物质——KSCN 溶液可用于检验 \ce{Fe^3+}}{\ce{MnO2} 与浓盐酸共热生成黄绿色气体——\ce{Cl2 + SO2 + 2H2O = H2SO4 + 2HCl}}

\item 下列关于 \ce{SO2} 的说法正确的是\xz
\xxa{\ce{SO2} 能使品红溶液和酸性 \ce{KMnO4} 溶液褪色，原理相同}{\ce{SO2} 通入 \ce{BaCl2} 溶液中产生白色沉淀}{\ce{SO2} 与 \ce{H2S} 反应生成 S，体现 \ce{SO2} 的氧化性}{将 \ce{SO2} 和 \ce{Cl2} 等体积混合后通入品红溶液，漂白效果更好}

\item 写出下列反应的离子方程式：(1) \ce{FeCl3} 溶液腐蚀铜电路板；(2) \ce{Na[Al(OH)4]} 溶液中通入过量 \ce{CO2}；(3) 酸性条件下 \ce{KMnO4} 溶液氧化 \ce{H2O2}。

\item 将 7.8 g \ce{Na2O2} 投入足量水中，生成 \ce{O2} 在标准状况下的体积为\tk L，转移电子\tk mol；若改为通入足量 \ce{CO2}，固体质量增加\tk g。

\item 向 \ce{FeCl3}、\ce{CuCl2} 的混合溶液中加入一定量铁粉，充分反应后过滤，滤渣中只有 Cu。则滤液中一定含有的金属阳离子是\tk，可能含有的是\tk，一定不含的是\tk。
\end{xiti}

\zu{B组\quad 高考典型}
\begin{xiti}
\item \laiyuan{2024年广东卷第14题，来源：仓库所附真题文档}\ 部分含 Mg 或 Al 或 Fe 物质的分类与相应化合价关系：a 为 0 价单质；b、g 分别为 $+3$、$+2$ 价氧化物；c、f 分别为 $+3$、$+2$ 价的盐；d、e 分别为 $+3$、$+2$ 价的碱。下列推断合理的是\xz
\xxa{若 a 在沸水中可生成 e，则 a$\to$f 的反应一定是化合反应}{在 g$\to$f$\to$e$\to$d 转化过程中，一定存在物质颜色的变化}{加热 c 的饱和溶液，一定会形成能产生丁达尔效应的红棕色分散系}{若 b 和 d 均能与同一物质反应生成 c，则组成 a 的元素一定位于周期表 p 区}

\item \laiyuan{2024年广东卷第18题(2)～(5)，来源：仓库所附真题文档}\ 镓（Ga）在半导体等领域有重要应用。从电解铝的副产品炭渣（含 C、Na、Al、F 和少量的 Ga、Fe、K、Ca 等元素）中提取镓及循环利用铝的工艺：炭渣焙烧（尾气处理）$\to$ 用 6 mol·L$^{-1}$ 盐酸浸取 $\to$ 浸取液加铝片“还原”$\to$ 原料液用 LAEM（一种新型阴离子交换膜，允许带负电荷的配离子从高浓度区扩散至低浓度区）提取：Ⅰ室原料液中的 \ce{[GaCl4]-} 进入Ⅱ室（pH$=2$）转化为 \ce{Ga^3+} 和 \ce{Cl-}，Ⅱ室溶液进一步处理得镓；Ⅰ室溶液加入含 \ce{F-} 的废液“调 pH”后结晶，得冰晶石 \ce{Na3[AlF6]} 用于电解铝。已知：①p$K_\text{a}$(HF)$=3.2$；②$K_\text{sp}$(\ce{Na3[AlF6]})$=4.0\times10^{-10}$；③浸取液中，Ga(Ⅲ) 和 Fe(Ⅲ) 以 \ce{[MCl_m]^{(m-3)-}}（$m=0\sim4$）微粒形式存在，\ce{Fe^2+} 最多可与 2 个 \ce{Cl-} 配位，其他金属离子与 \ce{Cl-} 的配位可忽略。
(2)“浸取”中，由 \ce{Ga^3+} 形成 \ce{[GaCl4]-} 的离子方程式为\tk。
(3)“还原”的目的：避免\tk{}元素以\tk（填化学式）微粒的形式通过 LAEM，从而有利于 Ga 的分离。
(4)“LAEM 提取”中，原料液的 \ce{Cl-} 浓度越\tk，越有利于 Ga 的提取；研究表明，原料液酸度过高会降低 Ga 的提取率。因此，在不提高原料液酸度的前提下，可向Ⅰ室中加入\tk（填化学式），以进一步提高 Ga 的提取率。
(5)“调 pH”中，pH 至少应大于\tk，使溶液中 $c(\ce{F-})>c(\ce{HF})$，有利于 \ce{[AlF6]^3-} 配离子及 \ce{Na3[AlF6]} 晶体的生成。若“结晶”后溶液中 $c(\ce{Na+})=0.10$ mol·L$^{-1}$，则 \ce{[AlF6]^3-} 浓度为\tk mol·L$^{-1}$。

\item \laiyuan{2026年广东卷第18题(4)，来源：仓库所附真题文档（网络收集版）}\ 从金精矿提金流程中，含金银粉体需经步骤“X”处理，X 包含多个步骤，需同时用到 HCl 和 \ce{H2O2} 溶液。一定条件下，Au、Ag 浸取率随溶液中 $c(\ce{H2O2})$ 的变化为：$c(\ce{H2O2})=x_1$ 时 Ag 浸取率接近 100\%、Au 浸取率很低；$c(\ce{H2O2})=x_3$（$x_3>x_2>x_1$）时 Au 浸取率接近 100\%。Au 在滤液2中主要以 \ce{[AuCl4]-} 存在。为尽可能提高金的提取率和纯度，“X”包括的步骤为\tk。

\item \laiyuan{2026年广东卷第13题，来源：仓库所附真题文档（网络收集版）}\ 部分含 C 或 N 或 Cl 物质的分类与相应化合价关系：a 为 0 价单质；b、c 分别为 $+2$、$+4$ 价氧化物；d 为 $+5$ 价的酸；f 为 $-1$ 价的酸；e 为 $+4$ 价的盐。下列说法\textbf{错误}的是\xz
\xxb{a 分子中可存在 $\pi$ 键和 $\sigma$ 键}{c 可由 e 受热分解产生}{可存在 a$\to$b$\to$c$\to$d 的转化}{a、c、f 均可用于自来水消毒}

\item \laiyuan{原创}\ 以软锰矿（主要成分 \ce{MnO2}，含少量 \ce{Fe2O3}、\ce{SiO2}）为原料，用 \ce{SO2} 还原浸出制备 \ce{MnSO4}：(1) 写出浸出时 \ce{MnO2} 与 \ce{SO2} 反应的化学方程式。(2) 浸出液中的铁主要以 \ce{Fe^2+} 存在，除铁时先加 \ce{MnO2}（适量）再调 pH，加 \ce{MnO2} 的作用是什么？写出离子方程式。(3) 滤渣的主要成分是什么？
\end{xiti}

\zu{C组\quad 冲刺高分}
\begin{xiti}
\item \laiyuan{原创}\ 从废旧锂离子电池正极材料（主要成分 \ce{LiCoO2}，含少量 Al、Fe 杂质）中回收钴和锂的流程：正极材料 $\xrightarrow{\text{NaOH 溶液}}$ 滤渣 $\xrightarrow{\text{稀硫酸、}\ce{H2O2}}$ 浸出液 $\xrightarrow{\text{调 pH}}$ 除铁 $\xrightarrow{\ce{(NH4)2C2O4}}$ \ce{CoC2O4} 沉淀；滤液 $\xrightarrow{\ce{Na2CO3}}$ \ce{Li2CO3}。已知：$K_\text{sp}[\ce{Fe(OH)3}]=2.8\times10^{-39}$，$K_\text{sp}[\ce{Co(OH)2}]=1.6\times10^{-15}$；\ce{Li2CO3} 的溶解度随温度升高而减小。
(1) 碱浸的目的是什么？写出相关反应的离子方程式。
(2) 酸浸时 \ce{H2O2} 的作用是什么？写出 \ce{LiCoO2} 溶解的离子方程式。为何酸浸温度不宜过高？
(3) 若浸出液中 $c(\ce{Co^2+})=0.16$ mol·L$^{-1}$，除铁时应控制的 pH 范围是多少？（$\lg2.8\approx0.45$，$\lg6.5\approx0.8$）
(4) 沉锂时为什么用热水洗涤 \ce{Li2CO3}？

\item \laiyuan{原创}\ 由铬铁矿（主要成分 \ce{FeCr2O4}）制备重铬酸钾：铬铁矿与纯碱在空气中焙烧，生成 \ce{Na2CrO4}、\ce{Fe2O3} 和 \ce{CO2}；水浸、过滤后，滤液酸化使 \ce{Na2CrO4} 转化为 \ce{Na2Cr2O7}；再加入 KCl，经蒸发浓缩、冷却结晶得 \ce{K2Cr2O7}。已知 \ce{2CrO4^2- + 2H+ <=> Cr2O7^2- + H2O}，25 ℃ 时 $K=1.0\times10^{14}$。(1) 写出焙烧反应的化学方程式。(2) 酸化后溶液中 $c(\ce{Cr2O7^2-})=0.25$ mol·L$^{-1}$，若要求 $c(\ce{CrO4^2-})\le1.0\times10^{-5}$ mol·L$^{-1}$，溶液的 $c(\ce{H+})$ 至少为多少？(3) 加入 KCl 能得到 \ce{K2Cr2O7} 的原因是什么？
\end{xiti}

\zhangmo{练习解析}
\begin{jiexilist}
\item \ans{B}（2025年广东卷第8题）A：浓硝酸见光易分解，因果不对应；B：浓硫酸使蔗糖脱水炭化，正确；C：两陈述均正确但无因果关系；D：两者无因果关系。
\item \ans{C}A：使品红褪色是漂白性，使 \ce{KMnO4} 褪色是还原性；B：\ce{H2SO3} 酸性弱于 HCl，不反应；C：\ce{SO2} 中 S 由 $+4$ 降为 0，体现氧化性，正确；D：\ce{SO2 + Cl2 + 2H2O = H2SO4 + 2HCl}，漂白性消失。
\item (1) \ce{2Fe^3+ + Cu = 2Fe^2+ + Cu^2+}；(2) \ce{[Al(OH)4]- + CO2 = Al(OH)3 v + HCO3-}；(3) \ce{2MnO4^- + 5H2O2 + 6H+ = 2Mn^2+ + 5O2 ^ + 8H2O}。
\item 1.12；0.1；2.8。$n(\ce{Na2O2})=0.1$ mol；\ce{2Na2O2 + 2H2O = 4NaOH + O2 ^}，生成 0.05 mol \ce{O2}，每生成 1 mol \ce{O2} 转移 2 mol 电子，共 0.1 mol；\ce{2Na2O2 + 2CO2 = 2Na2CO3 + O2}，固体增重相当于“CO”的质量：$0.1\times28=2.8$ g。
\item \ce{Fe^2+}；\ce{Cu^2+}；\ce{Fe^3+}。氧化性 \ce{Fe^3+}$>$\ce{Cu^2+}，铁先还原 \ce{Fe^3+}，再置换 \ce{Cu^2+}；有 Cu 析出说明 \ce{Fe^3+} 已反应完（Cu 也能与 \ce{Fe^3+} 反应），滤渣中无 Fe 说明 \ce{Cu^2+} 可能有剩余。
\item \ans{B}A：a 为 Mg，Mg 与盐酸生成 \ce{MgCl2} 是置换反应，不一定是化合反应；B：同时有 $+2$、$+3$ 价且有碱的只能是 Fe，\ce{FeO}（黑）$\to$\ce{Fe^2+}（浅绿）$\to$\ce{Fe(OH)2}（白）$\to$\ce{Fe(OH)3}（红褐），一定有颜色变化，正确；C：加热饱和 \ce{FeCl3} 溶液可能直接得到 \ce{Fe(OH)3} 沉淀，不一定是胶体；D：\ce{Fe2O3}、\ce{Fe(OH)3} 都能与盐酸反应生成 \ce{FeCl3}，Fe 位于 d 区。
\item (2) \ce{Ga^3+ + 4Cl- = [GaCl4]-}。(3) 铁；\ce{[FeCl4]-}。加铝片将 \ce{Fe^3+} 还原为 \ce{Fe^2+}（\ce{Fe^2+} 最多结合 2 个 \ce{Cl-}，不能形成带负电荷的配离子），避免铁随 \ce{[FeCl4]-} 进入Ⅱ室。(4) 高；NaCl（增大 $c(\ce{Cl-})$ 而不增大酸度，使 \ce{Ga^3+ + 4Cl- <=> [GaCl4]-} 正移）。(5) 3.2；$4.0\times10^{-7}$。由 $K_\text{a}=\dfrac{c(\ce{H+})c(\ce{F-})}{c(\ce{HF})}=10^{-3.2}$，$c(\ce{F-})>c(\ce{HF})$ 要求 $c(\ce{H+})<10^{-3.2}$，pH$>3.2$；$c(\ce{[AlF6]^3-})=\dfrac{K_\text{sp}}{c^3(\ce{Na+})}=\dfrac{4.0\times10^{-10}}{0.10^3}=4.0\times10^{-7}$ mol·L$^{-1}$。
\item 先控制 $c(\ce{H2O2})$ 为 $x_1$，与 HCl 溶液浸出 Ag 后过滤除去（Ag 进入溶液，Au 留在滤渣中）；再控制 $c(\ce{H2O2})$ 为 $x_3$，与 HCl 溶液浸出 Au（得到含 \ce{[AuCl4]-} 的滤液2）。这样既除去了银、提高了金的纯度，又使金基本完全浸出。
\item \ans{D}a 为 0 价单质，可为 \ce{N2}（含 $\sigma$ 键和 $\pi$ 键），A 对；c 为 $+4$ 价氧化物、e 为 $+4$ 价的盐，如 \ce{CaCO3 ->[\text{高温}] CaO + CO2 ^}，B 对；\ce{N2 -> NO -> NO2 -> HNO3}，C 对；f 为 $-1$ 价的酸 HCl，不能用于自来水消毒（a 可为 \ce{Cl2}、c 可为 \ce{ClO2}），D 错。
\item (1) \ce{MnO2 + SO2 = MnSO4}。(2) 将 \ce{Fe^2+} 氧化为 \ce{Fe^3+}，便于在较低 pH 下沉淀完全，且不引入杂质（还原产物为 \ce{Mn^2+}）：\ce{MnO2 + 2Fe^2+ + 4H+ = Mn^2+ + 2Fe^3+ + 2H2O}。(3) \ce{SiO2}（浸出时不溶）和 \ce{Fe(OH)3}（调 pH 时生成）。
\item (1) 溶解除去 Al：\ce{2Al + 2OH- + 6H2O = 2[Al(OH)4]- + 3H2 ^}（Al 若以 \ce{Al2O3} 存在：\ce{Al2O3 + 2OH- + 3H2O = 2[Al(OH)4]-}）。(2) 作还原剂，将 +3 价 Co 还原为 \ce{Co^2+}：\ce{2LiCoO2 + H2O2 + 6H+ = 2Li+ + 2Co^2+ + O2 ^ + 4H2O}；温度过高 \ce{H2O2} 分解，还原效率降低。(3) \ce{Fe^3+} 沉淀完全：$c(\ce{OH-})=\sqrt[3]{2.8\times10^{-39}/10^{-5}}\approx6.5\times10^{-12}$，pH$\approx2.8$；\ce{Co^2+} 开始沉淀：$c(\ce{OH-})=\sqrt{1.6\times10^{-15}/0.16}=10^{-7}$，pH$=7$。控制 $2.8\le$pH$<7$。（若浸出液中铁以 \ce{Fe^2+} 存在，需先氧化为 \ce{Fe^3+}。）(4) \ce{Li2CO3} 溶解度随温度升高而减小，用热水洗涤可减少洗涤过程中的溶解损失。
\item (1) Fe 由 $+2\to+3$，Cr 由 $+3\to+6$，每个 \ce{FeCr2O4} 失去 $1+2\times3=7$ 个电子，\ce{O2} 得 4 个电子：\ce{4FeCr2O4 + 8Na2CO3 + 7O2 ->[\text{高温}] 8Na2CrO4 + 2Fe2O3 + 8CO2}。(2) $K=\dfrac{c(\ce{Cr2O7^2-})}{c^2(\ce{CrO4^2-})\cdot c^2(\ce{H+})}$，$c^2(\ce{H+})=\dfrac{0.25}{1.0\times10^{14}\times(1.0\times10^{-5})^2}=2.5\times10^{-5}$，$c(\ce{H+})=5.0\times10^{-3}$ mol·L$^{-1}$（pH$\approx2.3$），即 $c(\ce{H+})\ge5.0\times10^{-3}$ mol·L$^{-1}$。(3) 低温下 \ce{K2Cr2O7} 的溶解度比 \ce{Na2Cr2O7}、NaCl、KCl 都小（且随温度降低明显减小），发生复分解 \ce{Na2Cr2O7 + 2KCl = K2Cr2O7 v + 2NaCl}，冷却结晶析出 \ce{K2Cr2O7}。
\end{jiexilist}
% <<<<<<<< 结束：chapters/ch13.tex

% >>>>>>>> 开始：chapters/ch14.tex
% ============================================================
% 第14章 高考综合大题训练（四道原创大题，按广东卷第17～20题的题型与分值设计）
% ============================================================
\chapter{高考综合大题训练}

\begin{daoxue}
\textbf{使用说明：}本章四道题均为原创，分别对应广东卷第17题（实验探究）、第18题（工艺流程）、第19题（原理综合）、第20题（有机综合），每题 14 分，难度略高于高考平均水平，用于在学完全书后检验“知识迁移与综合应用”能力。建议每题限时 13～15 分钟，独立完成后再对照“参考答案与评分要点”自评；得分低于 11 分的题，请回到对应章节复习。
\end{daoxue}

\section{实验探究题}
\noindent\textbf{第1题}（14 分）某小组探究 \ce{FeCl3} 与 KI 的反应，并测定某 \ce{FeCl3} 溶液的浓度。

(1) 配制 100 mL 0.10 mol·L$^{-1}$ \ce{FeCl3} 溶液，需要的玻璃仪器除烧杯、玻璃棒、量筒外，还有\tk；配制时先将 \ce{FeCl3} 固体溶于浓盐酸再加水稀释，目的是\tk。（2 分）

(2) 探究反应的限度：取 5 mL 0.10 mol·L$^{-1}$ KI 溶液，滴加 1 mL 0.10 mol·L$^{-1}$ \ce{FeCl3} 溶液，充分反应后分成两份。实验①：向第一份中加入 \ce{CCl4}，振荡静置，下层呈紫红色；实验②：向第二份中滴加 KSCN 溶液，溶液变红。
\quad ①写出 \ce{FeCl3} 与 KI 反应的离子方程式。（1 分）
\quad ②实验中为什么要让 KI 过量？实验②的现象说明了什么？（2 分）

(3) 探究浓度对平衡的影响：向实验②的红色溶液中加入少量 KI 固体，红色变浅。用平衡移动原理解释。（2 分）

(4) 测定浓度：取 25.00 mL 待测 \ce{FeCl3} 溶液于锥形瓶中，加入过量 KI 溶液，充分反应后，用 0.1000 mol·L$^{-1}$ \ce{Na2S2O3} 标准溶液滴定（\ce{I2 + 2S2O3^2- = 2I- + S4O6^2-}），消耗 20.00 mL。
\quad ①应选用的指示剂为\tk，滴定终点的现象为\tk。（2 分）
\quad ②该 \ce{FeCl3} 溶液的浓度为\tk mol·L$^{-1}$。（2 分）
\quad ③若加入 KI 后锥形瓶敞口放置过久，测定结果\tk（填“偏高”“偏低”或“无影响”），原因是\tk。（2 分）

(5) 有同学认为，(4) 中“反应可逆”会使测定结果偏低。你是否同意？说明理由。（1 分）

\section{工艺流程题}
\noindent\textbf{第2题}（14 分）以含钴废料（主要含 \ce{Co2O3}、CoO，还含少量 \ce{Fe2O3}、\ce{Al2O3}、\ce{SiO2}）为原料制备草酸钴晶体（\ce{CoC2O4*2H2O}）和 \ce{Co3O4} 的流程如下：
\begin{center}
\begin{tikzpicture}[font=\footnotesize,node distance=0.45cm and 0.8cm,st/.style={draw,rounded corners=1.5pt,fill=lightblue,minimum height=0.6cm,inner sep=3pt},pr/.style={inner sep=1pt,align=center}]
\node[pr] (a) {含钴废料};
\node[st,right=of a] (b) {酸浸};
\node[st,right=of b] (c) {氧化};
\node[st,right=of c] (d) {调 pH};
\node[st,right=of d] (e) {沉钴};
\node[st,right=2.3cm of e] (f) {煅烧};
\node[pr,right=of f] (g) {\ce{Co3O4}};
\draw[-Stealth] (a)--(b); \draw[-Stealth] (b)--(c); \draw[-Stealth] (c)--(d); \draw[-Stealth] (d)--(e); \draw[-Stealth] (e)--node[below=1pt]{\ce{CoC2O4*2H2O}}(f); \draw[-Stealth] (f)--(g);
\node[pr,above=of b] (rb) {稀 \ce{H2SO4}\\\ce{Na2SO3}}; \draw[-Stealth] (rb)--(b);
\node[pr,above=of c] (rc) {\ce{NaClO3}}; \draw[-Stealth] (rc)--(c);
\node[pr,above=of d] (rd) {\ce{Na2CO3}}; \draw[-Stealth] (rd)--(d);
\node[pr,above=of e] (re) {\ce{(NH4)2C2O4}}; \draw[-Stealth] (re)--(e);
\node[pr,below=of b] (r1) {滤渣1}; \draw[-Stealth] (b)--(r1);
\node[pr,below=of d] (r2) {滤渣2}; \draw[-Stealth] (d)--(r2);
\node[pr,below=of e] (r3) {滤液}; \draw[-Stealth] (e)--(r3);
\end{tikzpicture}
\end{center}
已知：$K_\text{sp}[\ce{Fe(OH)3}]=2.8\times10^{-39}$，$K_\text{sp}[\ce{Al(OH)3}]=1.3\times10^{-33}$，$K_\text{sp}[\ce{Co(OH)2}]=1.6\times10^{-15}$；$\lg5.1\approx0.7$，$\lg1.26\approx0.1$。

(1)“酸浸”时，提高浸取速率的措施有\tk（写一条）；滤渣1的主要成分是\tk。（2 分）

(2)“酸浸”时加入 \ce{Na2SO3} 的作用是\tk，写出 \ce{Co2O3} 参与反应的离子方程式。（2 分）

(3)“氧化”时 \ce{NaClO3} 将 \ce{Fe^2+} 氧化，写出离子方程式。已知 $K_\text{sp}[\ce{Fe(OH)2}]=4.9\times10^{-17}$，说明为什么要先把 \ce{Fe^2+} 氧化为 \ce{Fe^3+} 再调 pH。（2 分）

(4) 若“调 pH”前溶液中 $c(\ce{Co^2+})=0.10$ mol·L$^{-1}$，为使 \ce{Fe^3+}、\ce{Al^3+} 沉淀完全（离子浓度 $\le1.0\times10^{-5}$ mol·L$^{-1}$）而 \ce{Co^2+} 不沉淀，应控制 pH 的范围为\tk。（2 分）

(5)“沉钴”的离子方程式为\tk。用乙醇洗涤 \ce{CoC2O4*2H2O} 晶体的目的是\tk。（2 分）

(6) 取 18.3 g \ce{CoC2O4*2H2O}（$M=183$ g·mol$^{-1}$）在空气中加热，固体质量随温度变化：$300$ ℃ 时剩余 14.7 g，$450$ ℃ 以上剩余 8.03 g 并保持不变。
\quad ①300 ℃ 时剩余固体的化学式为\tk。（1 分）
\quad ②写出 300～450 ℃ 之间发生反应的化学方程式。（2 分）

(7) CoO 的晶胞与 NaCl 相同，晶胞边长为 0.426 nm，则 CoO 晶体的密度为\tk g·cm$^{-3}$（列出算式即可）。（1 分）

\section{原理综合题}
\noindent\textbf{第3题}（14 分）丙烷脱氢制丙烯是丙烯生产的重要途径，主反应和副反应如下：\\
主反应 Ⅰ：\ce{C3H8(g) <=> C3H6(g) + H2(g)} \quad $\Delta H_1$\\
副反应 Ⅱ：\ce{C3H8(g) <=> C2H4(g) + CH4(g)} \quad $\Delta H_2=+81$ kJ·mol$^{-1}$

(1) 已知键能（kJ·mol$^{-1}$）：C—C 348，C=C 615，C—H 413，H—H 436。估算 $\Delta H_1=$\tk kJ·mol$^{-1}$。（2 分）

(2) 反应Ⅰ在\tk（填“高温”或“低温”）下能自发进行；下列措施中，既能加快反应Ⅰ的速率、又能提高 \ce{C3H8} 平衡转化率的是\tk（填字母）。a．升高温度\quad b．增大压强\quad c．使用高效催化剂\quad d．恒容时充入 \ce{C3H8}（2 分）

(3) 在恒温恒压条件下进行反应，下列能说明体系达到平衡状态的是\tk（填字母）。A．单位时间内生成 $n$ mol \ce{C3H6} 的同时生成 $n$ mol \ce{H2}\quad B．混合气体的密度不再变化\quad C．$n(\ce{C3H6}):n(\ce{H2})=1:1$\quad D．\ce{C3H8} 的体积分数不再变化（2 分）

(4) 某温度下，在总压恒为 100 kPa 的条件下，向反应器中通入 1 mol \ce{C3H8} 和 4 mol Ar，达到平衡时 \ce{C3H8} 的转化率为 50\%，丙烯的选择性为 80\%（丙烯的选择性 $=\dfrac{\text{转化为丙烯的 }\ce{C3H8}\text{ 的量}}{\text{转化的 }\ce{C3H8}\text{ 总量}}\times100\%$）。
\quad ①平衡时 \ce{H2} 的物质的量为\tk mol，气体总物质的量为\tk mol。（2 分）
\quad ②反应Ⅰ的平衡常数 $K_p=$\tk kPa（保留两位小数）。（2 分）

(5) 实验发现：在总压不变时，增大 Ar 的比例，\ce{C3H8} 的平衡转化率增大；但 Ar 比例过大时，单位时间内 \ce{C3H8} 的消耗量明显下降。分别解释原因。（2 分）

(6) 长时间反应后催化剂表面会积炭而失活。研究发现，在原料气中加入适量 \ce{CO2} 可以减少积炭，用化学方程式表示其原理。（2 分）

\section{有机综合题}
\noindent\textbf{第4题}（14 分）对乙酰氨基酚（扑热息痛，\ce{HO-C6H4-NHCOCH3}，\ce{-OH} 与 \ce{-NHCOCH3} 处于苯环对位）是常用解热镇痛药，一种合成路线如下：
\[\text{苯酚}\xrightarrow[\text{①}]{\text{稀}\ce{HNO3}}\text{A}(\ce{C6H5NO3}\text{，对位产物})\xrightarrow[\text{②}]{\ce{Fe},\ \ce{HCl}}\text{B}(\ce{C6H7NO})\xrightarrow[\text{③}]{\ce{(CH3CO)2O}}\text{对乙酰氨基酚}\]
已知：\ce{-NH2} 与乙酸酐的反应活性远大于酚羟基；\ce{R-NO2 ->[\ce{Fe},\ \ce{HCl}] R-NH2}。

(1) 苯酚与稀硝酸反应生成 A 的反应类型为\tk；A 的名称为\tk。（2 分）

(2) 反应②中有机物发生了\tk（填“氧化”或“还原”）反应；B 的结构简式为\tk。（2 分）

(3) 写出反应③的化学方程式。（2 分）

(4) 对乙酰氨基酚分子中碳原子的杂化方式有\tk；1 mol 对乙酰氨基酚与足量 NaOH 溶液共热，最多消耗 NaOH\tk mol。（2 分）

(5) 对乙酰氨基酚的同分异构体中，同时满足下列条件的有\tk{}种。①含苯环，苯环上只有两个取代基；②其中一个取代基为 \ce{-NH2}；③能发生银镜反应，也能在 NaOH 溶液中水解。（2 分）

(6) 以苯胺为原料在苯环对位引入硝基等基团时，通常先用乙酸酐把 \ce{-NH2} 转化为 \ce{-NHCOCH3}（乙酰苯胺），反应结束后再水解恢复 \ce{-NH2}。说明这样做的目的。（2 分）

(7) 乙酸酐可由乙酸制得。以乙烯为原料（无机试剂任选）制备乙酸，合成路线为：乙烯 $\to$ \tk $\to$ \tk $\to$ 乙酸（写出中间产物的结构简式并注明各步反应条件）。（2 分）

\zhangmo{参考答案与评分要点}
\begin{jiexilist}
\item \textbf{第1题}\\
(1) 100 mL 容量瓶、胶头滴管（1 分）；抑制 \ce{Fe^3+} 水解（1 分）。\\
(2) ①\ce{2Fe^3+ + 2I- <=> 2Fe^2+ + I2}（1 分，写“$=$”不得分——本实验正是要证明其可逆）。②KI 过量，若反应能进行到底，\ce{Fe^3+} 应完全反应（1 分）；加 KSCN 变红说明 \ce{Fe^3+} 仍有剩余，即该反应是可逆反应，存在限度（1 分）。\\
(3) 加入 KI，$c(\ce{I-})$ 增大，\ce{2Fe^3+ + 2I- <=> 2Fe^2+ + I2} 正向移动，$c(\ce{Fe^3+})$ 减小，使 \ce{Fe^3+ + 3SCN- <=> Fe(SCN)3} 逆向移动，红色变浅（2 分，两个平衡各 1 分）。\\
(4) ①淀粉溶液（1 分）；滴入最后半滴 \ce{Na2S2O3} 溶液时，溶液蓝色褪去，且半分钟内不恢复（1 分）。②关系式 \ce{2Fe^3+}～\ce{I2}～\ce{2S2O3^2-}，$n(\ce{Fe^3+})=n(\ce{S2O3^2-})=2.000\times10^{-3}$ mol，$c=\dfrac{2.000\times10^{-3}}{0.02500}=0.08000$ mol·L$^{-1}$（2 分）。③偏高（1 分）；空气中的 \ce{O2} 将 \ce{I-} 氧化为 \ce{I2}（\ce{4I- + O2 + 4H+ = 2I2 + 2H2O}），消耗 \ce{Na2S2O3} 偏多（1 分）。\\
(5) 不同意（1 分，理由合理即可）：KI 大量过量，且滴定过程中 \ce{I2} 不断被 \ce{S2O3^2-} 消耗，使平衡持续正向移动，\ce{Fe^3+} 最终能完全反应。

\item \textbf{第2题}\\
(1) 粉碎废料、适当升温、适当增大硫酸浓度、搅拌（任一条，1 分）；\ce{SiO2}（1 分）。\\
(2) 作还原剂，将 $+3$ 价 Co 还原为 \ce{Co^2+}（1 分）；\ce{Co2O3 + SO3^2- + 4H+ = 2Co^2+ + SO4^2- + 2H2O}（1 分）。（\ce{Fe^3+} 同时被还原为 \ce{Fe^2+}。）\\
(3) \ce{ClO3^- + 6Fe^2+ + 6H+ = Cl- + 6Fe^3+ + 3H2O}（1 分）；\ce{Fe(OH)2} 与 \ce{Co(OH)2} 的 $K_\text{sp}$ 接近，\ce{Fe^2+} 沉淀完全时 \ce{Co^2+} 也已大量沉淀，无法通过调 pH 分离；\ce{Fe(OH)3} 的 $K_\text{sp}$ 极小，\ce{Fe^3+} 在 pH 约 2.8 时即沉淀完全，可与 \ce{Co^2+} 分离（1 分）。\\
(4) $4.7\le$pH$<7.1$（2 分）。\ce{Fe(OH)3} 的 $K_\text{sp}$ 更小，\ce{Al^3+} 完全沉淀时 \ce{Fe^3+} 早已沉淀完全：$c(\ce{OH-})=\sqrt[3]{\dfrac{1.3\times10^{-33}}{1.0\times10^{-5}}}\approx5.1\times10^{-10}$，pOH$\approx9.3$，pH$\approx4.7$；\ce{Co^2+} 开始沉淀时 $c(\ce{OH-})=\sqrt{\dfrac{1.6\times10^{-15}}{0.10}}\approx1.26\times10^{-7}$，pOH$\approx6.9$，pH$\approx7.1$。\\
(5) \ce{Co^2+ + C2O4^2- + 2H2O = CoC2O4*2H2O v}（1 分）；减少晶体的溶解损失，乙醇易挥发，便于晶体干燥（1 分）。\\
(6) ①$n=0.1$ mol，300 ℃ 时固体 $M=147$ g·mol$^{-1}$，为 \ce{CoC2O4}（失去全部结晶水）（1 分）。②8.03 g 中含 Co 0.1 mol（5.9 g），O 为 $8.03-5.9=2.13$ g，约 0.133 mol，$n(\ce{Co}):n(\ce{O})=3:4$，为 \ce{Co3O4}：\ce{3CoC2O4 + 2O2 ->[\text{△}] Co3O4 + 6CO2}（2 分）。\\
(7) $\rho=\dfrac{4\times75}{N_\text{A}\times(0.426\times10^{-7})^3}$（约 6.4 g·cm$^{-3}$）（1 分）。

\item \textbf{第3题}\\
(1) $+123$（2 分）。断键：1 个 C—C 和 2 个 C—H（反应物中多出的部分）；成键：1 个 C=C 和 1 个 H—H。$\Delta H_1=348+2\times413-615-436=+123$ kJ·mol$^{-1}$。（可只计算发生变化的键。）\\
(2) 高温（$\Delta H>0$、$\Delta S>0$）（1 分）；a（1 分）。b 使平衡逆移；c 不改变平衡转化率；d 相当于增大压强，\ce{C3H8} 的平衡转化率降低。\\
(3) BD（2 分，选对一个得 1 分，错选不得分）。两个反应都是气体分子数增大的反应，恒压下未平衡时气体总物质的量在变、容器体积在变，而气体总质量不变，故密度是“变量”，不变时即达平衡，B 正确；D 中 \ce{C3H8} 的体积分数是变量，不变即平衡，D 正确；A 中“生成 \ce{C3H6}”与“生成 \ce{H2}”都是正反应方向，任何时刻都成立，不能说明 $v_\text{正}=v_\text{逆}$；C 中 \ce{C3H6} 与 \ce{H2} 只由反应Ⅰ按 $1:1$ 生成，比值始终为 $1:1$，不能作为标志。\\
(4) ①转化 \ce{C3H8} 0.5 mol，其中 0.4 mol 生成丙烯，0.1 mol 发生副反应。\ce{H2} 0.4 mol（1 分）；平衡时 \ce{C3H8} 0.5、\ce{C3H6} 0.4、\ce{H2} 0.4、\ce{C2H4} 0.1、\ce{CH4} 0.1、Ar 4，总 5.5 mol（1 分）。②$K_p=\dfrac{(\frac{0.4}{5.5}\times100)^2}{\frac{0.5}{5.5}\times100}=\dfrac{0.16\times100}{0.5\times5.5}\approx5.82$ kPa（2 分）。\\
(5) 总压不变时增大 Ar 比例，反应物、生成物分压同比减小，相当于减压，平衡向气体分子数增大的正方向移动，转化率增大（1 分）；Ar 比例过大，\ce{C3H8} 分压（浓度）过低，反应速率下降（1 分）。\\
(6) \ce{C + CO2 ->[\text{高温}] 2CO}（2 分）。

\item \textbf{第4题}\\
(1) 取代反应（硝化反应）（1 分）；对硝基苯酚（4-硝基苯酚）（1 分）。\\
(2) 还原（1 分）；\ce{HO-C6H4-NH2}（对氨基苯酚，对位）（1 分）。\\
(3) \ce{HO-C6H4-NH2 + (CH3CO)2O -> HO-C6H4-NHCOCH3 + CH3COOH}（2 分，产物、配平各 1 分）。\\
(4) $\mathrm{sp}^2$、$\mathrm{sp}^3$（1 分）；2（1 分）。酚羟基与 NaOH 中和消耗 1 mol；酰胺基在碱性条件下水解生成 \ce{CH3COONa}（与对氨基苯酚钠），消耗 1 mol，共 2 mol。\\
(5) 3（2 分）。对乙酰氨基酚的分子式为 \ce{C8H9NO2}。既能发生银镜反应又能水解，说明含甲酸酯基 \ce{HCOO-}。苯环（\ce{C6H4}）上一个取代基为 \ce{-NH2}，则另一个取代基的组成为 \ce{C2H3O2} 且含 \ce{HCOO-}，只能是 \ce{HCOOCH2-}（\ce{-CH2OOCH}）；\ce{-OOCCH3}、\ce{-COOCH3} 不含醛基结构，排除。两个取代基在苯环上有邻、间、对 3 种位置关系，共 3 种。\\
(6) \ce{-NH2} 还原性强，直接硝化时会被浓硝酸氧化（1 分）；乙酰化后 \ce{-NHCOCH3} 不易被氧化，起“保护氨基”的作用，同时它对苯环的活化作用减弱、体积较大，可减少多取代并有利于得到对位产物，反应后水解即可恢复 \ce{-NH2}（1 分）。\\
(7) \ce{CH2=CH2 ->[H2O][\text{催化剂，加热、加压}] CH3CH2OH ->[O2, Cu][\text{△}] CH3CHO ->[O2][\text{催化剂，△}] CH3COOH}（中间产物各 1 分；条件缺失酌情扣分）。
\end{jiexilist}
% <<<<<<<< 结束：chapters/ch14.tex

\appendix
% >>>>>>>> 开始：chapters/appendix.tex
% ============================================================
% 附录
% ============================================================
\ctexset{appendix/name={附录,}, appendix/number={\Alph{chapter}}}
\renewcommand{\chaptermark}[1]{\markboth{附录\,\thechapter\quad #1}{}}

\chapter{课件勘误汇总}
编写本讲义时逐页审读了选择性必修1的七份课件。课件整体质量较好，以下是发现的笔误、不规范表述或与现行教材不一致之处（正文中均在对应位置作了说明，多数以“课件勘误”栏目呈现）：

\begin{center}
\small
\begin{tabularx}{\linewidth}{p{2.6cm}XXp{1.2cm}}
\toprule
课件位置 & 原表述 & 更正 & 所在节 \\
\midrule
1.1 反应热，第47页例1 & \ce{N2(g) + 3H2(aq) = 2NH3(g)} & \ce{H2} 的状态应为 g：\ce{N2(g) + 3H2(g) = 2NH3(g)} & 1.3 节 \\
1.1 反应热，第56页 & 燃烧的指定产物中列有“Cl$\to$\ce{Cl2(g)}” & 教材只规定 C、H、S、N 四种元素的指定产物；含氯物质燃烧时氯一般转化为 HCl，不宜记忆为 \ce{Cl2(g)} & 1.4 节 \\
1.1 反应热，第15页 & “中和热” & 新教材称“中和反应反应热”（含义相同），表述时应与教材一致 & 1.2 节 \\
1.2 反应热的计算，第19页 & “$\Delta H=\sum$反应物的燃烧热$-\sum$生成物的燃烧热” & 只有把燃烧热写成带负号的 $\Delta_\text{c}H$ 时才成立；用正值燃烧热 $Q$ 代入应为 $\Delta H=\sum Q_\text{生}-\sum Q_\text{反}$ & 1.6 节 \\
2.1 化学反应速率，第4页 & “纯液体表示速率可用 L/s 或 L/min” & 规范说法：纯液体、纯固体不能用浓度变化表示速率，可用单位时间内物质的量或质量的变化表示 & 2.1 节 \\
2.2 化学平衡，第38～39页 & 方程式写作“\ce{$m$A(g) + $n$B(g) <=> $e$C(g) + $f$d(g)}” & 应为 \ce{$f$D(g)}（大小写笔误） & 3.2 节 \\
2.3 化学反应方向，第16页 & \ce{Fe(OH)2(s) + O2(g) + H2O(g) = Fe(OH)3(s)}，并称“制备氢氧化亚铁” & 未配平，且是氢氧化亚铁被氧化：\ce{4Fe(OH)2(s) + O2(g) + 2H2O(l) = 4Fe(OH)3(s)} & 3.5 节 \\
\bottomrule
\end{tabularx}
\end{center}

\chapter{常用数据（25 ℃）}
\begin{center}
\small
\begin{tabular}{llll}
\toprule
弱电解质 & 电离常数 & 难溶电解质 & 溶度积 $K_\text{sp}$ \\
\midrule
\ce{CH3COOH} & $K_\text{a}=1.75\times10^{-5}$ & AgCl & $1.8\times10^{-10}$ \\
\ce{NH3*H2O} & $K_\text{b}=1.8\times10^{-5}$ & AgI & $8.5\times10^{-17}$ \\
\ce{H2CO3} & $K_{\text{a}1}=4.5\times10^{-7}$，$K_{\text{a}2}=4.7\times10^{-11}$ & \ce{Ag2CrO4} & $1.1\times10^{-12}$ \\
\ce{H2SO3} & $K_{\text{a}1}=1.4\times10^{-2}$，$K_{\text{a}2}=6.0\times10^{-8}$ & \ce{BaSO4} & $1.1\times10^{-10}$ \\
HClO & $K_\text{a}=4.0\times10^{-8}$ & \ce{BaCO3} & $2.6\times10^{-9}$ \\
HF & p$K_\text{a}=3.2$（2024年广东卷所给） & \ce{CaCO3} & $3.4\times10^{-9}$ \\
HCN & $K_\text{a}=6.2\times10^{-10}$ & \ce{CaSO4} & $4.9\times10^{-5}$ \\
& & \ce{Mg(OH)2} & $5.6\times10^{-12}$ \\
& & \ce{Fe(OH)2} & $4.9\times10^{-17}$ \\
& & \ce{Fe(OH)3} & $2.8\times10^{-39}$ \\
& & \ce{Al(OH)3} & $1.3\times10^{-33}$ \\
& & \ce{Cu(OH)2} & $2.2\times10^{-20}$ \\
\bottomrule
\end{tabular}
\end{center}
\noindent\small 说明：电离常数、溶度积数值随资料来源略有差异，考试时一律以题目所给数据为准。其他常数：$N_\text{A}=6.02\times10^{23}$ mol$^{-1}$，$F=96500$ C·mol$^{-1}$，$R=8.314$ J·mol$^{-1}$·K$^{-1}$，$K_\text{w}=1.0\times10^{-14}$（25 ℃），标准状况下 $V_\text{m}=22.4$ L·mol$^{-1}$。\normalsize

\chapter{高考化学答题规范用语}
\begin{center}
\small
\begin{tabularx}{\linewidth}{p{3.2cm}X}
\toprule
设问 & 规范表述 \\
\midrule
条件选择的原因 & 速率角度：……反应速率较快；平衡角度：……平衡转化率较高；实际角度：催化剂活性最大 / 设备要求 / 成本 / 安全环保 \\
平衡移动解释 & 改变……（条件），使……（平衡）向……方向移动，导致……（结果） \\
沸点或溶解性比较 & ××分子间（与水分子间）能形成氢键 / ××相对分子质量大，范德华力强 / 形成分子内氢键，沸点降低 \\
酸性强弱比较 & ××的电负性大，吸电子能力强，使羧基中 O—H 键极性增强，更易电离出 \ce{H+} \\
滴定终点 & 当滴入最后半滴××溶液时，溶液由×色变为×色，且半分钟内不恢复原色 \\
检验沉淀是否洗净 & 取最后一次洗涤液少许于试管中，滴加××，若无××现象，则已洗净 \\
离子检验 & 取少量待测液于试管中，滴加××，若出现××，则含有××（必要时先排除干扰） \\
提高浸取率 & 粉碎原料、搅拌、适当升温、适当增大浸取液浓度、延长浸取时间 \\
调节 pH 的原因 & 使××完全沉淀而××不沉淀（附 $K_\text{sp}$ 计算） \\
控制温度不宜过高 & 防止××分解（挥发）；防止催化剂失活；平衡逆向移动（放热反应） \\
冰水 / 乙醇洗涤 & 减少晶体溶解损失；乙醇易挥发，便于干燥 \\
\bottomrule
\end{tabularx}
\end{center}
% <<<<<<<< 结束：chapters/appendix.tex

\end{document}
